\input{preamble}

% OK, start here.
%
\begin{document}

\title{Catégories dérivées}


\maketitle

\phantomsection
\label{section-phantom}

\tableofcontents

\section{Introduction}
\label{section-introduction}

\noindent
Nous étudions d'abord les catégories triangulées et la localisation dans les
catégories triangulées. Nous démontrons ensuite que la catégorie d'homotopie
des complexes d'une catégorie additive est une catégorie triangulée. Cela fait,
nous définissons la catégorie dérivée d'une catégorie abélienne comme la
localisation de la catégorie d'homotopie relativement aux quasi-isomorphismes.
La thèse de Verdier \cite{Verdier} constitue une bonne référence.



\section{Catégories triangulées}
\label{section-triangulated-categories}

\noindent
Les catégories triangulées sont un outil commode pour décrire le type de
structure inhérent à la catégorie dérivée d'une catégorie abélienne. On pourra
consulter \cite{Verdier}, \cite{KS} et \cite{Neeman}.




\section{Définition d'une catégorie triangulée}
\label{section-triangulated-definitions}

\noindent
Nous rassemblons dans cette section la plupart des définitions relatives aux
catégories triangulées et prétriangulées.

\begin{definition}
\label{definition-triangle}
Soit $\mathcal{D}$ une catégorie additive. Soit
$[1] : \mathcal{D} \to \mathcal{D}$, $E \mapsto E[1]$,
un foncteur additif qui est une autoéquivalence de $\mathcal{D}$.
\begin{enumerate}
\item Un {\it triangle} est un sextuple
$(X, Y, Z, f, g, h)$ où $X, Y, Z \in \Ob(\mathcal{D})$ et où
$f : X \to Y$, $g : Y \to Z$ et $h : Z \to X[1]$ sont des morphismes
de $\mathcal{D}$.
\item Un {\it morphisme de triangles}
$(X, Y, Z, f, g, h) \to (X', Y', Z', f', g', h')$
est la donnée de morphismes $a : X \to X'$, $b : Y \to Y'$ et $c : Z \to Z'$
de $\mathcal{D}$ tels que
$b \circ f = f' \circ a$, $c  \circ g = g' \circ b$ et
$a[1] \circ h = h' \circ c$.
\end{enumerate}
\end{definition}

\noindent
Un morphisme de triangles se représente par le diagramme commutatif suivant :
$$
\xymatrix{
X \ar[r] \ar[d]^a &
Y \ar[r] \ar[d]^b &
Z \ar[r] \ar[d]^c &
X[1] \ar[d]^{a[1]} \\
X' \ar[r] &
Y' \ar[r] &
Z' \ar[r] &
X'[1]
}
$$
Dans la situation de la définition \ref{definition-triangle}, on écrit
$[0] = \text{id}$ ; pour $n > 0$, on note $[n]$ le composé $n$ fois de
$[1]$ ; on choisit un quasi-inverse $[-1]$ de $[1]$ et l'on pose que $[-n]$
est le composé $n$ fois de $[-1]$. Ainsi, $\{[n]\}_{n \in \mathbf{Z}}$
est une famille d'autoéquivalences additives de $\mathcal{D}$ indexée par
$n \in \mathbf{Z}$, munie d'isomorphismes de foncteurs
$[n] \circ [m] \cong [n + m]$.

\medskip\noindent
Voici la définition d'une catégorie triangulée donnée dans la thèse de Verdier.

\begin{definition}
\label{definition-triangulated-category}
Une {\it catégorie triangulée} est un triplet
$(\mathcal{D}, \{[n]\}_{n\in \mathbf{Z}}, \mathcal{T})$
où
\begin{enumerate}
\item $\mathcal{D}$ est une catégorie additive,
\item $[1] : \mathcal{D} \to \mathcal{D}$, $E \mapsto E[1]$, est une
autoéquivalence additive et $[n]$, pour $n \in \mathbf{Z}$, est défini comme
ci-dessus, et
\item $\mathcal{T}$ est un ensemble de triangles
(définition \ref{definition-triangle}), appelés les {\it triangles distingués},
\end{enumerate}
soumis aux conditions suivantes :
\begin{enumerate}
\item[TR1] Tout triangle isomorphe à un triangle distingué est distingué.
Tout triangle de la forme $(X, X, 0, \text{id}, 0, 0)$ est distingué.
Pour tout morphisme $f : X \to Y$ de $\mathcal{D}$, il existe un triangle
distingué de la forme $(X, Y, Z, f, g, h)$.
\item[TR2] Le triangle $(X, Y, Z, f, g, h)$ est distingué si et seulement si
le triangle $(Y, Z, X[1], g, h, -f[1])$ l'est.
\item[TR3] Étant donné un diagramme aux flèches pleines
$$
\xymatrix{
X \ar[r]^f \ar[d]^a &
Y \ar[r]^g \ar[d]^b &
Z \ar[r]^h \ar@{-->}[d] &
X[1] \ar[d]^{a[1]} \\
X' \ar[r]^{f'} &
Y' \ar[r]^{g'} &
Z' \ar[r]^{h'} &
X'[1]
}
$$
dont les lignes sont des triangles distingués et qui satisfait
$b \circ f = f' \circ a$, il existe un morphisme $c : Z \to Z'$ tel que
$(a, b, c)$ soit un morphisme de triangles.
\item[TR4] Étant donnés des objets $X$, $Y$, $Z$ de $\mathcal{D}$, des
morphismes $f : X \to Y$, $g : Y \to Z$, et des triangles distingués
$(X, Y, Q_1, f, p_1, d_1)$,
$(X, Z, Q_2, g \circ f, p_2, d_2)$
et
$(Y, Z, Q_3, g, p_3, d_3)$,
il existe des morphismes $a : Q_1 \to Q_2$ et $b : Q_2 \to Q_3$ tels que
\begin{enumerate}
\item $(Q_1, Q_2, Q_3, a, b, p_1[1] \circ d_3)$ soit un triangle distingué,
\item le triplet $(\text{id}_X, g, a)$ soit un morphisme de triangles
$(X, Y, Q_1, f, p_1, d_1) \to (X, Z, Q_2, g \circ f, p_2, d_2)$, et
\item le triplet $(f, \text{id}_Z, b)$ soit un morphisme de triangles
$(X, Z, Q_2, g \circ f, p_2, d_2) \to (Y, Z, Q_3, g, p_3, d_3)$.
\end{enumerate}
\end{enumerate}
Nous dirons que $(\mathcal{D}, [\ ], \mathcal{T})$ est une
{\it catégorie prétriangulée} si TR1, TR2 et TR3 sont
vérifiés.\footnote{Nous employons $[\ ]$ comme abréviation de la famille
$\{[n]\}_{n\in \mathbf{Z}}$.}
\end{definition}

\noindent
L'idée sous-jacente à TR4 est la suivante : si l'on considère $Q_1$ comme
$Y/X$, $Q_2$ comme $Z/X$ et $Q_3$ comme $Z/Y$, alors TR4(a) exprime
l'isomorphisme $(Z/X)/(Y/X) \cong Z/Y$, tandis que TR4(b) et TR4(c) expriment
que l'on peut comparer les triangles $X \to Y \to Q_1 \to X[1]$, etc.,
par des morphismes de triangles. Pour une reformulation plus précise de cette
idée, voir la démonstration du lemme \ref{lemma-two-split-injections}.

\medskip\noindent
Le signe dans TR2 signifie que, si $(X, Y, Z, f, g, h)$ est un triangle
distingué, alors chaque suite de quatre termes dans la suite longue
\begin{equation}
\label{equation-rotate}
\ldots \to
Z[-1] \xrightarrow{-h[-1]}
X \xrightarrow{f}
Y \xrightarrow{g}
Z \xrightarrow{h}
X[1] \xrightarrow{-f[1]}
Y[1] \xrightarrow{-g[1]}
Z[1] \to \ldots
\end{equation}
donne un triangle distingué.

\medskip\noindent
Comme d'ordinaire, nous abusons de la notation et parlons simplement d'une
catégorie (pré)triangulée $\mathcal{D}$, sans introduire explicitement de
notation pour les données supplémentaires. La notion de catégorie
prétriangulée est utile pour trouver des énoncés équivalents à TR4.

\medskip\noindent
Nous avons la définition suivante d'un foncteur triangulé.

\begin{definition}
\label{definition-exact-functor-triangulated-categories}
Soient $\mathcal{D}$ et $\mathcal{D}'$ des catégories prétriangulées.
Un {\it foncteur exact}, ou un {\it foncteur triangulé}, de $\mathcal{D}$
dans $\mathcal{D}'$ est un foncteur $F : \mathcal{D} \to \mathcal{D}'$
muni d'isomorphismes fonctoriels $\xi_X : F(X[1]) \to F(X)[1]$ tels que,
pour tout triangle distingué $(X, Y, Z, f, g, h)$ de $\mathcal{D}$, le triangle
$(F(X), F(Y), F(Z), F(f), F(g), \xi_X \circ F(h))$ soit distingué dans
$\mathcal{D}'$.
\end{definition}

\noindent
Un foncteur exact est additif ; voir le
lemme \ref{lemma-exact-functor-additive}.
Lorsque nous disons que deux catégories triangulées sont équivalentes, nous
entendons qu'elles sont équivalentes dans la $2$-catégorie des catégories
triangulées. Un $2$-morphisme $a : (F, \xi) \to (F', \xi')$ dans cette
$2$-catégorie est simplement une transformation de foncteurs $a : F \to F'$
compatible avec $\xi$ et $\xi'$, c'est-à-dire que le diagramme
$$
\xymatrix{
F \circ [1] \ar[r]_\xi \ar[d]_{a \star 1} & [1] \circ F \ar[d]^{1 \star a} \\
F' \circ [1] \ar[r]^{\xi'} & [1] \circ F'
}
$$
est commutatif.

\begin{definition}
\label{definition-triangulated-subcategory}
Soit $(\mathcal{D}, [\ ], \mathcal{T})$ une catégorie prétriangulée.
Une {\it sous-catégorie prétriangulée}\footnote{Cette définition n'est
peut-être pas standard. Si $\mathcal{D}'$ est une sous-catégorie pleine, alors
$\mathcal{T}'$ est l'intersection de l'ensemble des triangles de
$\mathcal{D}'$ avec $\mathcal{T}$ ; voir le
lemme \ref{lemma-triangulated-subcategory}. Dans ce cas, nous omettons
$\mathcal{T}'$ de la notation.}
est un couple $(\mathcal{D}', \mathcal{T}')$ tel que
\begin{enumerate}
\item $\mathcal{D}'$ soit une sous-catégorie additive de $\mathcal{D}$,
stable par $[1]$, et que $[1] : \mathcal{D}' \to \mathcal{D}'$ soit une
autoéquivalence,
\item $\mathcal{T}' \subset \mathcal{T}$ soit un sous-ensemble tel que, pour
tout $(X, Y, Z, f, g, h) \in \mathcal{T}'$, on ait
$X, Y, Z \in \Ob(\mathcal{D}')$ et
$f, g, h \in \text{Flèches}(\mathcal{D}')$, et
\item $(\mathcal{D}', [\ ], \mathcal{T}')$ soit une catégorie prétriangulée.
\end{enumerate}
Si $\mathcal{D}$ est une catégorie triangulée, nous disons que
$(\mathcal{D}', \mathcal{T}')$ est une {\it sous-catégorie triangulée} si
c'est une sous-catégorie prétriangulée et si
$(\mathcal{D}', [\ ], \mathcal{T}')$ est une catégorie triangulée.
\end{definition}

\noindent
Dans cette situation, le foncteur d'inclusion
$\mathcal{D}' \to \mathcal{D}$ est exact, avec
$\xi_X : X[1] \to X[1]$ donné par l'identité de $X[1]$.

\medskip\noindent
Nous verrons au lemme \ref{lemma-composition-zero} que, pour un triangle
distingué $(X, Y, Z, f, g, h)$ d'une catégorie prétriangulée, le composé
$g \circ f : X \to Z$ est nul. La suite (\ref{equation-rotate}) est donc un
complexe. Un foncteur homologique est un foncteur qui transforme ce complexe
en une suite exacte longue.

\begin{definition}
\label{definition-homological}
Soit $\mathcal{D}$ une catégorie prétriangulée et soit $\mathcal{A}$ une
catégorie abélienne. Un foncteur additif $H : \mathcal{D} \to \mathcal{A}$
est dit {\it homologique} si, pour tout triangle distingué
$(X, Y, Z, f, g, h)$, la suite
$$
H(X) \to H(Y) \to H(Z)
$$
est exacte dans la catégorie abélienne $\mathcal{A}$. Un foncteur additif
$H : \mathcal{D}^{opp} \to \mathcal{A}$ est dit {\it cohomologique} si le
foncteur correspondant $\mathcal{D} \to \mathcal{A}^{opp}$ est homologique.
\end{definition}

\noindent
Si $H : \mathcal{D} \to \mathcal{A}$ est un foncteur homologique, nous
écrivons souvent $H^n(X) = H(X[n])$, de sorte que $H(X) = H^0(X)$.
La discussion de TR2 ci-dessus montre qu'un triangle distingué
$(X, Y, Z, f, g, h)$ détermine une suite exacte longue
\begin{equation}
\label{equation-long-exact-cohomology-sequence}
\xymatrix@C=3pc{
H^{-1}(Z) \ar[r]^{H(h[-1])} &
H^0(X) \ar[r]^{H(f)} &
H^0(Y) \ar[r]^{H(g)} &
H^0(Z) \ar[r]^{H(h)} &
H^1(X)
}
\end{equation}
Nous l'appellerons la {\it suite exacte longue} associée au triangle distingué
et au foncteur homologique. Comme l'indique la notation, nous n'introduirons
aucun signe dans les morphismes de cette suite exacte longue. Il en résulte que
les applications de la suite exacte longue associée à la rotation (TR2) d'un
triangle distingué diffèrent parfois, par des signes, des applications de la
suite ci-dessus.

\begin{definition}
\label{definition-delta-functor}
Soient $\mathcal{A}$ une catégorie abélienne et $\mathcal{D}$ une catégorie
triangulée. Un {\it $\delta$-foncteur de $\mathcal{A}$ dans $\mathcal{D}$}
est la donnée d'un foncteur $G : \mathcal{A} \to \mathcal{D}$ et d'une règle
qui associe à toute suite exacte courte
$$
0 \to A \xrightarrow{a} B \xrightarrow{b} C \to 0
$$
un morphisme $\delta = \delta_{A \to B \to C} : G(C) \to G(A)[1]$ tel que
\begin{enumerate}
\item le triangle
$(G(A), G(B), G(C), G(a), G(b), \delta_{A \to B \to C})$
soit un triangle distingué de $\mathcal{D}$ pour toute suite exacte courte
comme ci-dessus, et
\item pour tout morphisme
$(A \to B \to C) \to (A' \to B' \to C')$ de suites exactes courtes, le
diagramme
$$
\xymatrix{
G(C) \ar[d] \ar[rr]_{\delta_{A \to B \to C}} & &
G(A)[1] \ar[d] \\
G(C') \ar[rr]^{\delta_{A' \to B' \to C'}} & &
G(A')[1]
}
$$
soit commutatif.
\end{enumerate}
Dans cette situation, nous appelons
$(G(A), G(B), G(C), G(a), G(b), \delta_{A \to B \to C})$
l'{\it image de la suite exacte courte par le $\delta$-foncteur donné}.
\end{definition}

\noindent
On remarquera qu'un $\delta$-foncteur est muni d'une structure supplémentaire.
À strictement parler, il n'est pas correct de dire qu'un foncteur donné
$\mathcal{A} \to \mathcal{D}$ est un $\delta$-foncteur ; nous le ferons
néanmoins souvent.









\section{Résultats élémentaires sur les catégories triangulées}
\label{section-elementary-results}

\noindent
La plupart des résultats de cette section sont démontrés pour les catégories
prétriangulées et sont donc, a fortiori, valables dans toute catégorie
triangulée.

\begin{lemma}
\label{lemma-composition-zero}
Soit $\mathcal{D}$ une catégorie prétriangulée. Soit
$(X, Y, Z, f, g, h)$ un triangle distingué. Alors
$g \circ f = 0$, $h \circ g = 0$ et $f[1] \circ h = 0$.
\end{lemma}

\begin{proof}
Par TR1, $(X, X, 0, 1, 0, 0)$ est un triangle distingué. Appliquons TR3 à
$$
\xymatrix{
X \ar[r] \ar[d]^1 &
X \ar[r] \ar[d]^f &
0 \ar[r] \ar@{-->}[d] &
X[1] \ar[d]^{1[1]} \\
X \ar[r]^f &
Y \ar[r]^g &
Z \ar[r]^h &
X[1]
}
$$
Bien entendu, la flèche pointillée est le morphisme nul. La commutativité du
diagramme implique donc $g \circ f = 0$. Pour les autres cas, on effectue une
rotation du triangle, c'est-à-dire que l'on applique TR2.
\end{proof}

\begin{lemma}
\label{lemma-representable-homological}
Soit $\mathcal{D}$ une catégorie prétriangulée. Pour tout objet $W$ de
$\mathcal{D}$, le foncteur $\Hom_\mathcal{D}(W, -)$ est homologique et le
foncteur $\Hom_\mathcal{D}(-, W)$ est cohomologique.
\end{lemma}

\begin{proof}
Considérons un triangle distingué $(X, Y, Z, f, g, h)$. Nous avons déjà vu que
$g \circ f = 0$ ; voir le lemme \ref{lemma-composition-zero}. Supposons que
$a : W \to Y$ soit un morphisme tel que $g \circ a = 0$. On obtient alors le
diagramme commutatif
$$
\xymatrix{
W \ar[r]_1 \ar@{..>}[d]^b &
W \ar[r] \ar[d]^a &
0 \ar[r] \ar[d]^0 &
W[1] \ar@{..>}[d]^{b[1]} \\
X \ar[r] & Y \ar[r] & Z \ar[r] & X[1]
}
$$
Les deux lignes sont des triangles distingués (utiliser TR1 pour la ligne
supérieure). On peut donc compléter la flèche pointillée $b$ (on effectue
d'abord une rotation par TR2, puis on applique TR3, et enfin on effectue la
rotation inverse). Cela démontre le lemme.
\end{proof}

\begin{lemma}
\label{lemma-third-isomorphism-triangle}
Soit $\mathcal{D}$ une catégorie prétriangulée. Soit
$$
(a, b, c) : (X, Y, Z, f, g, h) \to (X', Y', Z', f', g', h')
$$
un morphisme de triangles distingués. Si deux des morphismes $a, b, c$ sont des
isomorphismes, alors le troisième l'est également.
\end{lemma}

\begin{proof}
Supposons que $a$ et $c$ soient des isomorphismes. Pour tout objet $W$ de
$\mathcal{D}$, écrivons $H_W( - ) = \Hom_\mathcal{D}(W, -)$. On obtient un
diagramme commutatif de groupes abéliens
$$
\xymatrix{
H_W(Z[-1]) \ar[r] \ar[d] &
H_W(X) \ar[r] \ar[d] &
H_W(Y) \ar[r] \ar[d] &
H_W(Z) \ar[r] \ar[d] &
H_W(X[1]) \ar[d] \\
H_W(Z'[-1]) \ar[r] &
H_W(X') \ar[r] &
H_W(Y') \ar[r] &
H_W(Z') \ar[r] &
H_W(X'[1])
}
$$
Par hypothèse, les deux flèches verticales de droite et les deux flèches
verticales de gauche sont bijectives. Comme $H_W$ est homologique par le
lemme \ref{lemma-representable-homological}, le lemme des cinq
(Homologie, lemme \ref{homology-lemma-five-lemma}) montre que la flèche
verticale du milieu est un isomorphisme. Le lemme de Yoneda, voir
Catégories, lemme \ref{categories-lemma-yoneda}, montre donc que $b$ est un
isomorphisme. On en déduit les autres cas par rotation (en utilisant TR2).
\end{proof}

\begin{remark}
\label{remark-special-triangles}
Soit $\mathcal{D}$ une catégorie additive munie de foncteurs de translation
$[n]$ comme dans la définition \ref{definition-triangle}. Disons qu'un triangle
$(X, Y, Z, f, g, h)$ est {\it spécial}\footnote{Cette notation n'est pas
standard.} si, pour tout objet $W$ de $\mathcal{D}$, la suite longue de groupes
abéliens
$$
\ldots \to
\Hom_\mathcal{D}(W, X) \to
\Hom_\mathcal{D}(W, Y) \to
\Hom_\mathcal{D}(W, Z) \to
\Hom_\mathcal{D}(W, X[1]) \to \ldots
$$
est exacte. La démonstration du lemme
\ref{lemma-third-isomorphism-triangle} montre que, si
$$
(a, b, c) : (X, Y, Z, f, g, h) \to (X', Y', Z', f', g', h')
$$
est un morphisme de triangles spéciaux et si deux des morphismes $a, b, c$ sont
des isomorphismes, alors le troisième l'est aussi. On dispose d'un énoncé dual
pour les triangles {\it cospéciaux}, c'est-à-dire les triangles que le foncteur
$\Hom_\mathcal{D}(-, W)$ transforme en suites exactes longues. Ainsi, les
triangles distingués sont spéciaux et cospéciaux, mais il existe en général
beaucoup plus de triangles (co)spéciaux que de triangles distingués.
\end{remark}

\begin{lemma}
\label{lemma-third-map-square-zero}
Soit $\mathcal{D}$ une catégorie prétriangulée. Soient
$$
(0, b, 0), (0, b', 0) : (X, Y, Z, f, g, h) \to (X, Y, Z, f, g, h)
$$
des endomorphismes d'un triangle distingué. Alors $bb' = 0$.
\end{lemma}

\begin{proof}
Considérons le diagramme
$$
\xymatrix{
X \ar[r] \ar[d]^0 &
Y \ar[r] \ar[d]^{b, b'} \ar@{..>}[ld]^\alpha &
Z \ar[r] \ar[d]^0 \ar@{..>}[ld]^\beta &
X[1] \ar[d]^0 \\
X \ar[r] & Y \ar[r] & Z \ar[r] & X[1]
}
$$
En appliquant le lemme \ref{lemma-representable-homological}, on trouve des
flèches pointillées $\alpha$ et $\beta$ telles que
$b' = f \circ \alpha$ et $b = \beta \circ g$. Ainsi,
$bb' = \beta \circ g \circ f \circ \alpha = 0$, puisque
$g \circ f = 0$ par le lemme \ref{lemma-composition-zero}.
\end{proof}

\begin{lemma}
\label{lemma-third-map-idempotent}
Soit $\mathcal{D}$ une catégorie prétriangulée. Soit
$(X, Y, Z, f, g, h)$ un triangle distingué. Si
$$
\xymatrix{
Z \ar[r]_h \ar[d]_c & X[1] \ar[d]^{a[1]} \\
Z \ar[r]^h & X[1]
}
$$
est commutatif et si $a^2 = a$, $c^2 = c$, alors il existe un morphisme
$b : Y \to Y$ tel que $b^2 = b$ et que $(a, b, c)$ soit un endomorphisme du
triangle $(X, Y, Z, f, g, h)$.
\end{lemma}

\begin{proof}
Par TR3, il existe un morphisme $b'$ tel que $(a, b', c)$ soit un endomorphisme
de $(X, Y, Z, f, g, h)$. Alors $(0, (b')^2 - b', 0)$ est également un
endomorphisme. Le lemme \ref{lemma-third-map-square-zero} montre que
$(b')^2 - b'$ est de carré nul. Posons
$b = b' - (2b' - 1)((b')^2 - b') = 3(b')^2 - 2(b')^3$. Un calcul montre que
$(a, b, c)$ est un endomorphisme et que
$b^2 - b = (4(b')^2 - 4b' - 3)((b')^2 - b')^2 = 0$.
\end{proof}

\begin{lemma}
\label{lemma-cone-triangle-unique-isomorphism}
Soit $\mathcal{D}$ une catégorie prétriangulée. Soit $f : X \to Y$ un
morphisme de $\mathcal{D}$. Il existe un triangle distingué
$(X, Y, Z, f, g, h)$, unique à isomorphisme (non unique) de triangles près.
Plus précisément, étant donné un second triangle distingué de cette forme
$(X, Y, Z', f, g', h')$, il existe un isomorphisme
$$
(1, 1, c) : (X, Y, Z, f, g, h) \longrightarrow (X, Y, Z', f, g', h')
$$
\end{lemma}

\begin{proof}
L'existence résulte de TR1. L'unicité à isomorphisme près résulte de TR3 et du
lemme \ref{lemma-third-isomorphism-triangle}.
\end{proof}

\begin{lemma}
\label{lemma-uniqueness-third-arrow}
Soit $\mathcal{D}$ une catégorie prétriangulée. Soit
$$
(a, b, c) : (X, Y, Z, f, g, h) \to (X', Y', Z', f', g', h')
$$
un morphisme de triangles distingués. Si l'une des conditions suivantes est
vérifiée :
\begin{enumerate}
\item $\Hom(Y, X') = 0$,
\item $\Hom(Z, Y') = 0$,
\item $\Hom(X, X') = \Hom(Z, X') = 0$,
\item $\Hom(Z, X') = \Hom(Z, Z') = 0$, ou
\item $\Hom(X[1], Z') = \Hom(Z, X') = 0$,
\end{enumerate}
alors $b$ est l'unique morphisme $Y \to Y'$ tel que $(a, b, c)$ soit un
morphisme de triangles.
\end{lemma}

\begin{proof}
S'il existe un second morphisme de triangles $(a, b', c)$, alors
$(0, b - b', 0)$ est un morphisme de triangles. Il suffit donc de montrer que
le seul morphisme $b : Y \to Y'$ tel que les composés
$X \to Y \to Y'$ et $Y \to Y' \to Z'$ soient nuls est $0$.
Nous utiliserons sans autre mention le
lemme \ref{lemma-representable-homological}. En particulier, la condition (3)
implique (1). Sous l'hypothèse (1), si le composé
$g' \circ b : Y \to Y' \to Z'$ est nul, alors $b$ se relève en un morphisme
$Y \to X'$, qui est nécessairement nul. Cela démontre (1).

\medskip\noindent
Les démonstrations de (2) et de (4) sont duales de cet argument.

\medskip\noindent
Supposons (5). Considérons le diagramme
$$
\xymatrix{
X \ar[r]_f \ar[d]^0 &
Y \ar[r]_g \ar[d]^b &
Z \ar[r]_h \ar[d]^0 \ar@{..>}[ld]^\epsilon &
X[1] \ar[d]^0 \\
X' \ar[r]^{f'} &
Y' \ar[r]^{g'} &
Z' \ar[r]^{h'} &
X'[1]
}
$$
On peut choisir $\epsilon$ de telle sorte que $b = \epsilon \circ g$.
Alors $g' \circ \epsilon \circ g = 0$, ce qui implique
$g' \circ \epsilon = \delta \circ h$ pour un certain
$\delta \in \Hom(X[1], Z')$. Puisque $\Hom(X[1], Z') = 0$, on en déduit
$g' \circ \epsilon = 0$. Ainsi, $\epsilon = f' \circ \gamma$ pour un
certain $\gamma \in \Hom(Z, X')$. Comme $\Hom(Z, X') = 0$, on obtient
$\epsilon = 0$, et donc $b = 0$, comme voulu.
\end{proof}

\begin{lemma}
\label{lemma-third-object-zero}
Soit $\mathcal{D}$ une catégorie prétriangulée. Soit $f : X \to Y$ un
morphisme de $\mathcal{D}$. Les conditions suivantes sont équivalentes :
\begin{enumerate}
\item $f$ est un isomorphisme,
\item $(X, Y, 0, f, 0, 0)$ est un triangle distingué, et
\item pour tout triangle distingué $(X, Y, Z, f, g, h)$, on a $Z = 0$.
\end{enumerate}
\end{lemma}

\begin{proof}
Par TR1, le triangle $(X, X, 0, 1, 0, 0)$ est distingué. Soit
$(X, Y, Z, f, g, h)$ un triangle distingué. Par TR3, il existe un morphisme de
triangles distingués $(1, f, 0) : (X, X, 0) \to (X, Y, Z)$. Si $f$ est un
isomorphisme, alors $(1, f, 0)$ est un isomorphisme de triangles par le
lemme \ref{lemma-third-isomorphism-triangle}, et $Z = 0$. Réciproquement, si
$Z = 0$, alors $(1, f, 0)$ est encore un isomorphisme de triangles, et $f$ est
donc un isomorphisme.
\end{proof}

\begin{lemma}
\label{lemma-direct-sum-triangles}
Soient $\mathcal{D}$ une catégorie prétriangulée et
$(X, Y, Z, f, g, h)$, $(X', Y', Z', f', g', h')$ des triangles. Les conditions
suivantes sont équivalentes :
\begin{enumerate}
\item $(X \oplus X', Y \oplus Y', Z \oplus Z',
f \oplus f', g \oplus g', h \oplus h')$ est un triangle distingué,
\item les deux triangles $(X, Y, Z, f, g, h)$ et
$(X', Y', Z', f', g', h')$ sont distingués.
\end{enumerate}
\end{lemma}

\begin{proof}
Supposons (2). Par TR1, on peut choisir un triangle distingué
$(X \oplus X', Y \oplus Y', Q, f \oplus f', g'', h'')$. Par TR3, on peut
trouver des morphismes de triangles distingués
$(X, Y, Z, f, g, h) \to
(X \oplus X', Y \oplus Y', Q, f \oplus f', g'', h'')$
et
$(X', Y', Z', f', g', h') \to
(X \oplus X', Y \oplus Y', Q, f \oplus f', g'', h'')$.
En prenant la somme directe de ces morphismes, on obtient un morphisme de
triangles
$$
\xymatrix{
(X \oplus X', Y \oplus Y', Z \oplus Z',
f \oplus f', g \oplus g', h \oplus h')
\ar[d]^{(1, 1, c)} \\
(X \oplus X', Y \oplus Y', Q, f \oplus f', g'', h'').
}
$$
Selon la terminologie de la remarque \ref{remark-special-triangles}, c'est un
morphisme de triangles spéciaux (car une somme directe de triangles spéciaux
est spéciale), et l'on en déduit que $c$ est un isomorphisme. Ainsi, (1) est
vérifiée.

\medskip\noindent
Supposons (1). Montrons que $(X, Y, Z, f, g, h)$ est un triangle distingué.
Remarquons d'abord que $(X, Y, Z, f, g, h)$ est un triangle spécial
(terminologie de la remarque \ref{remark-special-triangles}), puisqu'il est un
facteur direct du triangle distingué, donc spécial,
$(X \oplus X', Y \oplus Y', Z \oplus Z',
f \oplus f', g \oplus g', h \oplus h')$. Par TR1, choisissons un triangle
distingué $(X, Y, Q, f, g'', h'')$. Par TR3, il existe un morphisme de
triangles distingués
$(X \oplus X', Y \oplus Y', Z \oplus Z',
f \oplus f', g \oplus g', h \oplus h') \to (X, Y, Q, f, g'', h'')$.
En le composant avec l'inclusion, on obtient un morphisme de triangles
$$
(1, 1, c) :
(X, Y, Z, f, g, h)
\longrightarrow
(X, Y, Q, f, g'', h'')
$$
La remarque \ref{remark-special-triangles} montre que $c$ est un isomorphisme,
d'où (2).
\end{proof}

\begin{lemma}
\label{lemma-split}
Soit $\mathcal{D}$ une catégorie prétriangulée.
Soit $(X, Y, Z, f, g, h)$ un triangle distingué.
\begin{enumerate}
\item Si $h = 0$, alors il existe un inverse à droite $s : Z \to Y$ de $g$.
\item Pour tout inverse à droite $s : Z \to Y$ de $g$, l'application
$f \oplus s : X \oplus Z \to Y$ est un isomorphisme.
\item Pour tous objets $X', Z'$ de $\mathcal{D}$, le triangle
$(X', X' \oplus Z', Z', (1, 0), (0, 1), 0)$ est distingué.
\end{enumerate}
\end{lemma}

\begin{proof}
Pour (1), on utilise l'exactitude de
$\Hom_\mathcal{D}(Z, Y) \to \Hom_\mathcal{D}(Z, Z) \to
\Hom_\mathcal{D}(Z, X[1])$, donnée par le
lemme \ref{lemma-representable-homological}.
De même, si $s$ est comme dans (2), alors $h = 0$ et la suite
$$
0 \to \Hom_\mathcal{D}(W, X) \to \Hom_\mathcal{D}(W, Y)
\to \Hom_\mathcal{D}(W, Z) \to 0
$$
est exacte scindée (par $s : Z \to Y$). Le lemme de Yoneda montre donc que
$X \oplus Z \to Y$ est un isomorphisme. La dernière assertion résulte de TR1
et du lemme \ref{lemma-direct-sum-triangles}.
\end{proof}

\begin{lemma}
\label{lemma-when-split}
Soit $\mathcal{D}$ une catégorie prétriangulée.
Soit $f : X \to Y$ un morphisme de $\mathcal{D}$.
Les conditions suivantes sont équivalentes :
\begin{enumerate}
\item $f$ possède un noyau,
\item $f$ possède un conoyau,
\item $f$ est isomorphe à une composée
$K \oplus Z \to Z \to Z \oplus Q$ d'une projection et d'une coprojection,
pour certains objets $K, Z, Q$ de $\mathcal{D}$.
\end{enumerate}
\end{lemma}

\begin{proof}
Tout morphisme isomorphe à une application de la forme
$X' \oplus Z \to Z \oplus Y'$ possède à la fois un noyau et un conoyau.
Ainsi (3) $\Rightarrow$ (1), (2).
Montrons ensuite que (1) $\Rightarrow$ (3).
Supposons d'abord que $f : X \to Y$ soit un monomorphisme, c'est-à-dire que son
noyau soit nul. Par TR1, il existe un triangle distingué
$(X, Y, Z, f, g, h)$. D'après le lemme \ref{lemma-composition-zero}, la composée
$f \circ h[-1] = 0$. Comme $f$ est un monomorphisme, on a $h[-1] = 0$,
et donc $h = 0$. Le lemme \ref{lemma-split} entraîne alors
$Y = X \oplus Z$ ; la condition (3) est donc satisfaite.
Supposons maintenant que $f$ possède un noyau $K$. Comme $K \to X$ est un
monomorphisme, on obtient $X = K \oplus X'$ et
$f|_{X'} : X' \to Y$ est un monomorphisme. Ainsi $Y = X' \oplus Y'$,
ce qui conclut la démonstration.
L'implication (2) $\Rightarrow$ (3) est duale.
\end{proof}

\begin{lemma}
\label{lemma-products-sums-shifts-triangles}
Soit $\mathcal{D}$ une catégorie prétriangulée.
Soit $I$ un ensemble.
\begin{enumerate}
\item Soit $X_i$, $i \in I$, une famille d'objets de $\mathcal{D}$.
\begin{enumerate}
\item Si $\prod X_i$ existe, alors $(\prod X_i)[1] = \prod X_i[1]$.
\item Si $\bigoplus X_i$ existe, alors
$(\bigoplus X_i)[1] = \bigoplus X_i[1]$.
\end{enumerate}
\item Soit $X_i \to Y_i \to Z_i \to X_i[1]$ une famille de triangles
distingués de $\mathcal{D}$.
\begin{enumerate}
\item Si $\prod X_i$, $\prod Y_i$, $\prod Z_i$ existent, alors
$\prod X_i \to \prod Y_i \to \prod Z_i \to \prod X_i[1]$
est un triangle distingué.
\item Si $\bigoplus X_i$, $\bigoplus Y_i$,
$\bigoplus Z_i$ existent, alors
$\bigoplus X_i \to \bigoplus Y_i \to \bigoplus Z_i \to \bigoplus X_i[1]$
est un triangle distingué.
\end{enumerate}
\end{enumerate}
\end{lemma}

\begin{proof}
La partie (1) résulte de ce que $[1]$ est une autoéquivalence de $\mathcal{D}$
et de ce que les sommes directes et les produits sont définis par la structure
de catégorie. Démontrons (2)(a). Choisissons un triangle distingué
$\prod X_i \to \prod Y_i \to Z \to \prod X_i[1]$. Pour chaque $j$, TR3 permet
de choisir un morphisme $p_j : Z \to Z_j$ qui, avec les projections
$\prod X_i \to X_j$ et $\prod Y_i \to Y_j$, forme un morphisme de triangles
distingués. Par la définition des produits, on obtient une application
$\prod p_i : Z \to \prod Z_i$ qui complète un morphisme du triangle distingué
vers le triangle formé par les produits. Remarquons que le triangle « produit »
$\prod X_i \to \prod Y_i \to \prod Z_i \to \prod X_i[1]$
est spécial au sens de la remarque \ref{remark-special-triangles}, car les
produits de suites exactes de groupes abéliens sont exacts.
La remarque \ref{remark-special-triangles} montre donc que le morphisme de
triangles est un isomorphisme, et l'on conclut par TR1.
La démonstration de (2)(b) est duale.
\end{proof}

\begin{lemma}
\label{lemma-projectors-have-images-triangulated}
Soit $\mathcal{D}$ une catégorie prétriangulée.
Si $\mathcal{D}$ admet les produits dénombrables, alors $\mathcal{D}$
est karoubienne.
Si $\mathcal{D}$ admet les coproduits dénombrables, alors $\mathcal{D}$
est karoubienne.
\end{lemma}

\begin{proof}
Supposons que $\mathcal{D}$ admette les produits dénombrables. D'après
Algèbre homologique, lemme \ref{homology-lemma-projectors-have-images},
il suffit de vérifier que les morphismes possédant un inverse à droite ont un
noyau. Tout morphisme possédant un inverse à droite est un épimorphisme et
possède donc un noyau d'après le lemme \ref{lemma-when-split}.
La seconde assertion est duale de la première.
\end{proof}

\noindent
Le lemme suivant simplifie légèrement la démonstration du fait qu'une catégorie
prétriangulée est triangulée.

\begin{lemma}
\label{lemma-easier-axiom-four}
Soit $\mathcal{D}$ une catégorie prétriangulée.
Pour démontrer TR4, il suffit de montrer que, pour toute paire de morphismes
composables $f : X \to Y$ et $g : Y \to Z$, il existe
\begin{enumerate}
\item des isomorphismes $i : X' \to X$, $j : Y' \to Y$ et
$k : Z' \to Z$, puis, en posant $f' = j^{-1}fi : X' \to Y'$ et
$g' = k^{-1}gj : Y' \to Z'$, il existe
\item des triangles distingués
$(X', Y', Q_1, f', p_1, d_1)$,
$(X', Z', Q_2, g' \circ f', p_2, d_2)$
et
$(Y', Z', Q_3, g', p_3, d_3)$,
pour lesquels l'assertion de TR4 est satisfaite.
\end{enumerate}
\end{lemma}

\begin{proof}
Le remplacement de $X, Y, Z$ par $X', Y', Z'$ est inoffensif, par notre
définition des triangles distingués et de leurs isomorphismes.
Le lemme résulte du fait que les triangles distingués
$(X', Y', Q_1, f', p_1, d_1)$,
$(X', Z', Q_2, g' \circ f', p_2, d_2)$
et
$(Y', Z', Q_3, g', p_3, d_3)$
sont uniques à isomorphisme près, d'après le
lemme \ref{lemma-cone-triangle-unique-isomorphism}.
\end{proof}

\begin{lemma}
\label{lemma-triangulated-subcategory}
Soit $\mathcal{D}$ une catégorie prétriangulée.
Supposons que $\mathcal{D}'$ soit une sous-catégorie additive pleine de
$\mathcal{D}$. Les conditions suivantes sont équivalentes :
\begin{enumerate}
\item il existe un ensemble de triangles $\mathcal{T}'$ tel que
$(\mathcal{D}', \mathcal{T}')$ soit une sous-catégorie prétriangulée de
$\mathcal{D}$,
\item $\mathcal{D}'$ est stable par $[1]$,
$[1] : \mathcal{D}' \to \mathcal{D}'$ est une autoéquivalence, et pour tout
morphisme $f : X \to Y$ de $\mathcal{D}'$, il existe un triangle distingué
$(X, Y, Z, f, g, h)$ dans $\mathcal{D}$ tel que $Z$ soit isomorphe à un objet
de $\mathcal{D}'$.
\end{enumerate}
Dans ce cas, $\mathcal{T}'$ dans (1) est l'ensemble des triangles distingués
$(X, Y, Z, f, g, h)$ de $\mathcal{D}$ tels que
$X, Y, Z \in \Ob(\mathcal{D}')$. Enfin, si $\mathcal{D}$ est une catégorie
triangulée, alors (1) et (2) sont également équivalentes à
\begin{enumerate}
\item[(3)] $\mathcal{D}'$ est une sous-catégorie triangulée.
\end{enumerate}
\end{lemma}

\begin{proof}
Démonstration omise.
\end{proof}

\begin{lemma}
\label{lemma-exact-functor-additive}
Un foncteur exact entre catégories prétriangulées est additif.
\end{lemma}

\begin{proof}
Soit $F : \mathcal{D} \to \mathcal{D}'$ un foncteur exact entre catégories
prétriangulées. Puisque $(0, 0, 0, 1_0, 1_0, 0)$ est un triangle distingué de
$\mathcal{D}$, le triangle
$$
(F(0), F(0), F(0), 1_{F(0)}, 1_{F(0)}, F(0))
$$
est distingué dans $\mathcal{D}'$.
Il en résulte que $1_{F(0)} \circ 1_{F(0)}$ est nul ; voir le
lemme \ref{lemma-composition-zero}.
Ainsi $F(0)$ est l'objet nul de $\mathcal{D}'$. Cela montre également que
$F$ envoie tout morphisme nul sur zéro (car un morphisme d'une catégorie
additive est nul si et seulement s'il se factorise par l'objet nul).
Ensuite, le lemme \ref{lemma-split}, partie (3), affirme que
$(X, X \oplus Y, Y, (1, 0), (0, 1), 0)$ est un triangle distingué ; par
conséquent,
$(F(X), F(X \oplus Y), F(Y), F(1, 0), F(0, 1), 0)$ en est également un.
Le lemme \ref{lemma-split}, partie (2), montre alors que l'application
$F(X) \oplus F(Y) \to F(X \oplus Y)$ est un isomorphisme. Pour conclure,
on applique Algèbre homologique, lemme \ref{homology-lemma-additive-functor}.
\end{proof}

\begin{lemma}
\label{lemma-exact-equivalence}
Soit $F : \mathcal{D} \to \mathcal{D}'$ un foncteur exact pleinement fidèle
entre catégories prétriangulées. Alors un triangle $(X, Y, Z, f, g, h)$ de
$\mathcal{D}$ est distingué si et seulement si
$(F(X), F(Y), F(Z), F(f), F(g), F(h))$ est distingué dans $\mathcal{D}'$.
\end{lemma}

\begin{proof}
Le sens direct est clair. Supposons $(F(X), F(Y), F(Z))$ distingué dans
$\mathcal{D}'$. Choisissons un triangle distingué
$(X, Y, Z', f, g', h')$ dans $\mathcal{D}$. D'après le
lemme \ref{lemma-cone-triangle-unique-isomorphism}, il existe un isomorphisme de
triangles
$$
(1, 1, c') : (F(X), F(Y), F(Z)) \longrightarrow (F(X), F(Y), F(Z')).
$$
Comme $F$ est pleinement fidèle, il existe un morphisme $c : Z \to Z'$ tel que
$F(c) = c'$. Alors $(1, 1, c)$ est un isomorphisme entre $(X, Y, Z)$ et
$(X, Y, Z')$. Par TR1, $(X, Y, Z)$ est donc distingué.
\end{proof}

\begin{lemma}
\label{lemma-composition-exact}
Soient $\mathcal{D}, \mathcal{D}', \mathcal{D}''$ des catégories
prétriangulées. Soient $F : \mathcal{D} \to \mathcal{D}'$ et
$F' : \mathcal{D}' \to \mathcal{D}''$ des foncteurs exacts.
Alors $F' \circ F$ est un foncteur exact.
\end{lemma}

\begin{proof}
Démonstration omise.
\end{proof}

\begin{lemma}
\label{lemma-exact-compose-homological-functor}
Soit $\mathcal{D}$ une catégorie prétriangulée.
Soit $\mathcal{A}$ une catégorie abélienne.
Soit $H : \mathcal{D} \to \mathcal{A}$ un foncteur homologique.
\begin{enumerate}
\item Soit $\mathcal{D}'$ une catégorie prétriangulée.
Soit $F : \mathcal{D}' \to \mathcal{D}$ un foncteur exact.
Alors la composée $H \circ F$ est aussi un foncteur homologique.
\item Soit $\mathcal{A}'$ une catégorie abélienne. Soit
$G : \mathcal{A} \to \mathcal{A}'$ un foncteur exact.
Alors $G \circ H$ est aussi un foncteur homologique.
\end{enumerate}
\end{lemma}

\begin{proof}
Démonstration omise.
\end{proof}

\begin{lemma}
\label{lemma-exact-compose-delta-functor}
Soit $\mathcal{D}$ une catégorie triangulée.
Soit $\mathcal{A}$ une catégorie abélienne.
Soit $G : \mathcal{A} \to \mathcal{D}$ un $\delta$-foncteur.
\begin{enumerate}
\item Soit $\mathcal{D}'$ une catégorie triangulée.
Soit $F : \mathcal{D} \to \mathcal{D}'$ un foncteur exact.
Alors la composée $F \circ G$ est aussi un $\delta$-foncteur.
\item Soit $\mathcal{A}'$ une catégorie abélienne. Soit
$H : \mathcal{A}' \to \mathcal{A}$ un foncteur exact.
Alors $G \circ H$ est aussi un $\delta$-foncteur.
\end{enumerate}
\end{lemma}

\begin{proof}
Démonstration omise.
\end{proof}

\begin{lemma}
\label{lemma-compose-delta-functor-homological}
Soit $\mathcal{D}$ une catégorie triangulée.
Soient $\mathcal{A}$ et $\mathcal{B}$ des catégories abéliennes.
Soit $G : \mathcal{A} \to \mathcal{D}$ un $\delta$-foncteur.
Soit $H : \mathcal{D} \to \mathcal{B}$ un foncteur homologique.
Supposons que $H^{-1}(G(A)) = 0$ pour tout $A$ dans $\mathcal{A}$.
Alors la famille
$$
\{H^n \circ G, H^n(\delta_{A \to B \to C})\}_{n \geq 0}
$$
est un $\delta$-foncteur de $\mathcal{A} \to \mathcal{B}$ ; voir
Algèbre homologique, définition
\ref{homology-definition-cohomological-delta-functor}.
\end{lemma}

\begin{proof}
La notation signifie ce qui suit. Si
$0 \to A \xrightarrow{a} B \xrightarrow{b} C \to 0$ est une suite exacte
courte dans $\mathcal{A}$, alors
$$
\delta = \delta_{A \to B \to C} : G(C) \to G(A)[1]
$$
est un morphisme de $\mathcal{D}$ tel que
$(G(A), G(B), G(C), a, b, \delta)$ soit un triangle distingué ; voir la
définition \ref{definition-delta-functor}.
Alors $H^n(\delta) : H^n(G(C)) \to H^n(G(A)[1]) = H^{n + 1}(G(A))$ dépend
manifestement fonctoriellement de la suite exacte courte.
Enfin, la longue suite exacte de cohomologie
(\ref{equation-long-exact-cohomology-sequence}), combinée à l'annulation de
$H^{-1}(G(C))$, donne une suite exacte longue
$$
0 \to H^0(G(A)) \to H^0(G(B)) \to H^0(G(C))
\xrightarrow{H^0(\delta)} H^1(G(A)) \to \ldots
$$
dans $\mathcal{B}$, comme voulu.
\end{proof}

\noindent
La démonstration du résultat suivant utilise TR4.

\begin{proposition}
\label{proposition-9}
Soit $\mathcal{D}$ une catégorie triangulée. Tout diagramme commutatif
$$
\xymatrix{
X \ar[r] \ar[d] & Y \ar[d] \\
X' \ar[r] & Y'
}
$$
se complète en un diagramme
$$
\xymatrix{
X \ar[r] \ar[d] & Y \ar[r] \ar[d] & Z \ar[r] \ar[d] & X[1] \ar[d] \\
X' \ar[r] \ar[d] & Y' \ar[r] \ar[d] & Z' \ar[r] \ar[d] & X'[1] \ar[d] \\
X'' \ar[r] \ar[d] & Y'' \ar[r] \ar[d] & Z'' \ar[r] \ar[d] & X''[1] \ar[d] \\
X[1] \ar[r] & Y[1] \ar[r] & Z[1] \ar[r] & X[2]
}
$$
où tous les carrés sont commutatifs, sauf le carré inférieur droit qui est
anticommutatif. De plus, chacune des lignes et des colonnes est un triangle
distingué. Enfin, les morphismes de la ligne inférieure (respectivement de la
colonne de droite) s'obtiennent à partir de ceux de la ligne supérieure
(respectivement de la colonne de gauche) par application de $[1]$.
\end{proposition}

\begin{proof}
Dans cette démonstration, nous n'indiquons pas les flèches afin de préserver la
lisibilité. Choisissons des triangles distingués
$(X, Y, Z)$, $(X', Y', Z')$, $(X, X', X'')$, $(Y, Y', Y'')$ et
$(X, Y', A)$. Remarquons que le morphisme $X \to Y'$ est à la fois égal à la
composée $X \to Y \to Y'$ et à la composée $X \to X' \to Y'$. On peut donc
trouver des morphismes
\begin{enumerate}
\item $a : Z \to A$ et $b : A \to Y''$, et
\item $a' : X'' \to A$ et $b' : A \to Z'$,
\end{enumerate}
comme dans TR4. Notons $c : Y'' \to Z[1]$ la composée
$Y'' \to Y[1] \to Z[1]$, et $c' : Z' \to X''[1]$ la composée
$Z' \to X'[1] \to X''[1]$. Notre application de TR4 donne :
\begin{enumerate}
\item $(Z, A, Y'', a, b, c)$, $(X'', A, Z', a', b', c')$
sont des triangles distingués,
\item $(X, Y, Z) \to (X, Y', A)$,
$(X, Y', A) \to (Y, Y', Y'')$,
$(X, X', X'') \to (X, Y', A)$,
$(X, Y', A) \to (X', Y', Z')$
sont des morphismes de triangles.
\end{enumerate}
En utilisant d'abord que $(X, X', X'') \to (X, Y', A)$ et
$(X, Y', A) \to (Y, Y', Y'')$ sont des morphismes de triangles, on voit que
le premier des diagrammes
$$
\vcenter{
\xymatrix{
X' \ar[r] \ar[d] & Y' \ar[d] \\
X'' \ar[r]^{b \circ a'} \ar[d] & Y'' \ar[d] \\
X[1] \ar[r] & Y[1]
}
}
\quad\text{et}\quad
\vcenter{
\xymatrix{
Y \ar[r] \ar[d] & Z \ar[d]^{b' \circ a} \ar[r] & X[1] \ar[d] \\
Y' \ar[r] & Z' \ar[r] & X'[1]
}
}
$$
est commutatif. Le second l'est également, puisque
$(X, Y, Z) \to (X, Y', A)$ et $(X, Y', A) \to (X', Y', Z')$ sont des
morphismes de triangles. Choisissons à présent un triangle distingué
$(X'', Y'' , Z'')$ commençant par l'application $b \circ a' : X'' \to Y''$.

\medskip\noindent
Appliquons encore une fois TR4 aux morphismes $X'' \to A \to Y''$ et aux
triangles
$(X'', A, Z', a', b', c')$,
$(X'', Y'', Z'')$, et
$(A, Y'', Z[1], b, c , -a[1])$, afin d'obtenir des morphismes
$a'' : Z' \to Z''$ et $b'' : Z'' \to Z[1]$.
Alors $(Z', Z'', Z[1], a'', b'', - b'[1] \circ a[1])$ est un triangle
distingué, donc $(Z, Z', Z'', -b' \circ a, a'', -b'')$ l'est aussi, et par
conséquent $(Z, Z', Z'', b' \circ a, a'', b'')$ également.
De plus, $(X'', A, Z') \to (X'', Y'', Z'')$ et
$(X'', Y'', Z'') \to (A, Y'', Z[1], b, c , -a[1])$
sont des morphismes de triangles.
Nous avons maintenant défini tous les triangles distingués et tous les
morphismes ; il ne reste qu'à vérifier quelques relations de commutativité.

\medskip\noindent
Pour voir que le carré central du diagramme commute, remarquons que la flèche
$Y' \to Z'$ se factorise en $Y' \to A \to Z'$, car
$(X, Y', A) \to (X', Y', Z')$ est un morphisme de triangles.
De même, le morphisme $Y' \to Y''$ se factorise en
$Y' \to A \to Y''$, car $(X, Y', A) \to (Y, Y', Y'')$ est un morphisme de
triangles. Le carré central commute donc : le carré de côtés
$(A, Z', Z'', Y'')$ commute puisque
$(X'', A, Z') \to (X'', Y'', Z'')$ est un morphisme de triangles (par TR4).
Le carré de côtés $(Y'', Z'', Y[1], Z[1])$ commute parce que
$(X'', Y'', Z'') \to (A, Y'', Z[1], b, c , -a[1])$ est un morphisme de
triangles et que $c : Y'' \to Z[1]$ est la composée
$Y'' \to Y[1] \to Z[1]$.
Le carré de côtés $(Z', X'[1], X''[1], Z'')$ est commutatif parce que
$(X'', A, Z') \to (X'', Y'', Z'')$ est un morphisme de triangles et que
$c' : Z' \to X''[1]$ est la composée $Z' \to X'[1] \to X''[1]$.
Enfin, il faut montrer que le carré de côtés
$(Z'', X''[1], Z[1], X[2])$ est anticommutatif. Cela résulte du fait que
$(X'', Y'', Z'') \to (A, Y'', Z[1], b, c , -a[1])$ est un morphisme de
triangles, ce qui conclut.
\end{proof}

\section{Localisation des catégories triangulées}
\label{section-localization}

\noindent
Pour construire la catégorie dérivée à partir de la catégorie homotopique des
complexes, nous utiliserons un procédé de localisation.

\begin{definition}
\label{definition-localization}
Soit $\mathcal{D}$ une catégorie prétriangulée. On dit qu'un système
multiplicatif $S$ est {\it compatible avec la structure triangulée} si les deux
conditions suivantes sont satisfaites :
\begin{enumerate}
\item[MS5] Pour tout morphisme $f$ de $\mathcal{D}$, on a
$f \in S \Leftrightarrow f[1] \in S$\footnote{Voir la remarque
\ref{remark-MS5}.}.
\item[MS6] Étant donné un carré commutatif plein
$$
\xymatrix{
X \ar[r] \ar[d]^s &
Y \ar[r] \ar[d]^{s'} &
Z \ar[r] \ar@{-->}[d] &
X[1] \ar[d]^{s[1]} \\
X' \ar[r] &
Y' \ar[r] &
Z' \ar[r] &
X'[1]
}
$$
dont les lignes sont des triangles distingués et où $s, s' \in S$, il existe
un morphisme $s'' : Z \to Z'$ appartenant à $S$ tel que $(s, s', s'')$ soit un
morphisme de triangles.
\end{enumerate}
\end{definition}

\noindent
Il s'avère que ces axiomes ne sont pas indépendants des axiomes définissant les
systèmes multiplicatifs.

\begin{lemma}
\label{lemma-localization-conditions}
Soit $\mathcal{D}$ une catégorie prétriangulée.
Soit $S \subset \text{Flèches}(\mathcal{D})$.
\begin{enumerate}
\item Si $S$ contient toutes les identités et si MS6 est satisfaite
(définition \ref{definition-localization}), alors tout isomorphisme de
$\mathcal{D}$ appartient à $S$.
\item Si MS1, MS5 (Catégories, définition
\ref{categories-definition-multiplicative-system}) et MS6 sont satisfaites,
alors MS2 l'est aussi.
\end{enumerate}
\end{lemma}

\begin{proof}
Supposons que $S$ contienne toutes les identités et que MS6 soit satisfaite.
Soit $f : X \to Y$ un isomorphisme de $\mathcal{D}$.
Considérons le diagramme
$$
\xymatrix{
0 \ar[r] \ar[d]^1 &
X \ar[r]_1 \ar[d]^1 &
X \ar[r] \ar@{-->}[d] &
0[1] \ar[d]^{1[1]} \\
0 \ar[r] &
X \ar[r]^f &
Y \ar[r] &
0[1]
}
$$
Les lignes sont des triangles distingués d'après le
lemme \ref{lemma-third-object-zero}. Par MS6, la flèche pointillée existe et
appartient à $S$ ; ainsi $f$ appartient à $S$.

\medskip\noindent
Supposons MS1, MS5 et MS6 satisfaites.
Soient $f : X \to Y$ un morphisme de $\mathcal{D}$ et
$t : X \to X'$ un élément de $S$. Choisissons un triangle distingué
$(X, Y, Z, f, g, h)$. Choisissons ensuite un triangle distingué
$(X', Y', Z, f', g', t[1] \circ h)$ (on utilise ici TR1 et TR2).
Par MS5 et MS6 (après une rotation par TR2), on peut trouver la flèche
pointillée du diagramme commutatif
$$
\xymatrix{
X \ar[r] \ar[d]^t &
Y \ar[r] \ar@{..>}[d]^{s'} &
Z \ar[r] \ar[d]^1 &
X[1] \ar[d]^{t[1]} \\
X' \ar[r] &
Y' \ar[r] &
Z \ar[r] &
X'[1]
}
$$
avec, de plus, $s' \in S$. Cela démontre LMS2. La démonstration de RMS2 est
duale.
\end{proof}

\begin{remark}
\label{remark-MS5}
En présence de MS1 et MS6, la condition MS5 équivaut à demander que
$s[n] \in S$ pour tous $s \in S$ et $n \in \mathbf{Z}$. Supposons par exemple
MS5 satisfaite, soit $s \in S$, et montrons que $s[-1] \in S$. Ce n'est pas
immédiat, car $s[-1][1]$ n'est pas égal à $s$, mais est seulement isomorphe à
$s$ comme flèche de $\mathcal{D}$. Cela implique néanmoins que
$s[-1][1] = f \circ s \circ g$ pour certains isomorphismes $f$, $g$.
Le lemme \ref{lemma-localization-conditions}, partie (1), donne
$f, g \in S$, puis MS1 donne $s[-1][1] \in S$, et enfin MS5 donne
$s[-1] \in S$. Nous laissons la démonstration complète en exercice au lecteur.
\end{remark}

\begin{lemma}
\label{lemma-triangle-functor-localize}
Soit $F : \mathcal{D} \to \mathcal{D}'$ un foncteur exact entre catégories
prétriangulées. Posons
$$
S = \{f \in \text{Flèches}(\mathcal{D})
\mid F(f)\text{ est un isomorphisme}\}
$$
Alors $S$ est un système multiplicatif saturé (voir Catégories, définition
\ref{categories-definition-saturated-multiplicative-system}), compatible avec
la structure triangulée de $\mathcal{D}$.
\end{lemma}

\begin{proof}
Il faut démontrer les axiomes MS1 -- MS6 ; voir Catégories, définitions
\ref{categories-definition-multiplicative-system} et
\ref{categories-definition-saturated-multiplicative-system}, ainsi que la
définition \ref{definition-localization}.
MS1, MS4 et MS5 résultent directement des définitions. MS6 résulte de TR3 et du
lemme \ref{lemma-third-isomorphism-triangle}.
Le lemme \ref{lemma-localization-conditions} donne alors MS2. Pour achever la
démonstration, il reste à établir MS3. Soient donc $f, g : X \to Y$ des
morphismes de $\mathcal{D}$ et $t : Z \to X$ un élément de $S$ tel que
$f \circ t = g \circ t$. Comme $\mathcal{D}$ est additive, cela signifie
simplement que $a \circ t = 0$, où $a = f - g$. Par TR1, choisissons un
triangle distingué $(Z, X, Q, t, d, h)$. Puisque $a \circ t = 0$, le
lemme \ref{lemma-representable-homological} fournit un morphisme
$i : Q \to Y$ tel que $i \circ d = a$. Enfin, une nouvelle application de TR1
permet de choisir un triangle $(Q, Y, W, i, j, k)$. Voici la situation :
$$
\xymatrix{
Z \ar[r]_t & X \ar[r]_d \ar[d]^1 & Q \ar[r] \ar[d]^i & Z[1] \\
& X \ar[r]_a & Y \ar[d]^j \\
& & W
}
$$
Appliquons maintenant le foncteur $F$ à ce diagramme.
Comme $t \in S$, on a $F(Q) = 0$ ; voir le
lemme \ref{lemma-third-object-zero}.
Le même lemme montre donc que $F(j)$ est un isomorphisme, c'est-à-dire que
$j \in S$. Enfin, $j \circ a = j \circ i \circ d = 0$, puisque
$j \circ i = 0$. Ainsi $j \circ f = j \circ g$, ce qui établit LMS3.
La démonstration de RMS3 est duale.
\end{proof}

\begin{lemma}
\label{lemma-homological-functor-localize}
Soit $H : \mathcal{D} \to \mathcal{A}$ un foncteur homologique d'une catégorie
prétriangulée vers une catégorie abélienne. Posons
$$
S = \{f \in \text{Flèches}(\mathcal{D})
\mid H^i(f)\text{ est un isomorphisme pour tout }i \in \mathbf{Z}\}
$$
Alors $S$ est un système multiplicatif saturé (voir Catégories, définition
\ref{categories-definition-saturated-multiplicative-system}), compatible avec
la structure triangulée de $\mathcal{D}$.
\end{lemma}

\begin{proof}
Il faut démontrer les axiomes MS1 -- MS6 ; voir Catégories, définitions
\ref{categories-definition-multiplicative-system} et
\ref{categories-definition-saturated-multiplicative-system}, ainsi que la
définition \ref{definition-localization}.
MS1, MS4 et MS5 résultent directement des définitions.
MS6 résulte de TR3 et de la longue suite exacte de cohomologie
(\ref{equation-long-exact-cohomology-sequence}).
Le lemme \ref{lemma-localization-conditions} donne alors MS2. Pour achever la
démonstration, il reste à établir MS3. Soient donc $f, g : X \to Y$ des
morphismes de $\mathcal{D}$ et $t : Z \to X$ un élément de $S$ tel que
$f \circ t = g \circ t$. Comme $\mathcal{D}$ est additive, cela signifie
simplement que $a \circ t = 0$, où $a = f - g$. Par TR1 et TR2, choisissons un
triangle distingué $(Z, X, Q, t, u, v)$. Puisque $a \circ t = 0$, le
lemme \ref{lemma-representable-homological} fournit un morphisme
$i : Q \to Y$ tel que $i \circ u = a$. Enfin, une nouvelle application de TR1
permet de choisir un triangle $(Q, Y, W, i, j, k)$. Voici la situation :
$$
\xymatrix{
Z \ar[r]_t & X \ar[r]_u \ar[d]^1 & Q \ar[r]_v \ar[d]^i & Z[1] \\
& X \ar[r]_a & Y \ar[d]^j \\
& & W
}
$$
Appliquons maintenant les foncteurs $H^i$ à ce diagramme.
Comme $t \in S$, la longue suite exacte de cohomologie
(\ref{equation-long-exact-cohomology-sequence}) donne $H^i(Q) = 0$.
Le même argument montre donc que $H^i(j)$ est un isomorphisme pour tout $i$,
c'est-à-dire que $j \in S$. Enfin, $j \circ a = j \circ i \circ u = 0$,
puisque $j \circ i = 0$. Ainsi $j \circ f = j \circ g$, ce qui établit LMS3.
La démonstration de RMS3 est duale.
\end{proof}

\begin{proposition}
\label{proposition-construct-localization}
Soit $\mathcal{D}$ une catégorie prétriangulée. Soit $S$ un système
multiplicatif compatible avec la structure triangulée.
Il existe alors une unique structure de catégorie prétriangulée sur
$S^{-1}\mathcal{D}$ telle que $[1] \circ Q = Q \circ [1]$ et telle que le
foncteur de localisation $Q : \mathcal{D} \to S^{-1}\mathcal{D}$ soit exact.
De plus, si $\mathcal{D}$ est une catégorie triangulée, alors
$S^{-1}\mathcal{D}$ l'est aussi.
\end{proposition}

\begin{proof}
Nous avons vu dans Algèbre homologique, lemme
\ref{homology-lemma-localization-additive}, que $S^{-1}\mathcal{D}$ est une
catégorie additive et que le foncteur de localisation $Q$ est additif.
Il résulte de MS5 qu'il existe une unique autoéquivalence additive
$[1] : S^{-1}\mathcal{D} \to S^{-1}\mathcal{D}$ telle que
$Q \circ [1] = [1] \circ Q$ (égalité de foncteurs) ; nous omettons les détails.
On dit qu'un triangle de $S^{-1}\mathcal{D}$ est distingué s'il est isomorphe à
l'image par le foncteur de localisation $Q$ d'un triangle distingué.

\medskip\noindent
Démonstration de TR1. Le seul point à établir est que, si
$a : Q(X) \to Q(Y)$ est un morphisme de $S^{-1}\mathcal{D}$, alors $a$ s'insère
dans un triangle distingué. Écrivons $a = Q(s)^{-1} \circ Q(f)$, où
$s : Y \to Y'$ appartient à $S$ et $f : X \to Y'$. Choisissons un triangle
distingué $(X, Y', Z, f, g, h)$ dans $\mathcal{D}$. Alors
$(Q(X), Q(Y), Q(Z), a, Q(g) \circ Q(s), Q(h))$ est un triangle distingué de
$S^{-1}\mathcal{D}$.

\medskip\noindent
Démonstration de TR2. Elle résulte immédiatement des définitions.

\medskip\noindent
Démonstration de TR3. Remarquons que l'existence de la flèche pointillée requise
peut être démontrée après remplacement des deux triangles par des triangles
isomorphes. On peut donc supposer donnés des triangles distingués
$(X, Y, Z, f, g, h)$ et $(X', Y', Z', f', g', h')$ de $\mathcal{D}$, ainsi qu'un
diagramme commutatif
$$
\xymatrix{
Q(X) \ar[r]_{Q(f)} \ar[d]_a & Q(Y) \ar[d]^b \\
Q(X') \ar[r]^{Q(f')} & Q(Y')
}
$$
dans $S^{-1}\mathcal{D}$. Appliquons maintenant Catégories, lemme
\ref{categories-lemma-left-localization-diagram}, pour trouver un morphisme
$f'' : X'' \to Y''$ de $\mathcal{D}$ et un diagramme commutatif
$$
\xymatrix{
X \ar[d]_f \ar[r]_k & X'' \ar[d]^{f''} & X' \ar[d]^{f'} \ar[l]^s \\
Y \ar[r]^l & Y'' & Y' \ar[l]_t
}
$$
dans $\mathcal{D}$, avec $s, t \in S$ et
$a = s^{-1}k$, $b = t^{-1}l$.
On peut alors appliquer TR3 dans $\mathcal{D}$ et MS6 pour obtenir un diagramme
commutatif
$$
\xymatrix{
X \ar[r] \ar[d]^k &
Y \ar[r] \ar[d]^l &
Z \ar[r] \ar[d]^m &
X[1] \ar[d]^{k[1]} \\
X'' \ar[r] &
Y'' \ar[r] &
Z'' \ar[r] &
X''[1] \\
X' \ar[r] \ar[u]_s &
Y' \ar[r] \ar[u]_t &
Z' \ar[r] \ar[u]_r &
X'[1] \ar[u]_{s[1]}
}
$$
avec $r \in S$. En posant $c = Q(r)^{-1}Q(m)$, on obtient donc le morphisme de
triangles recherché
$$
\xymatrix{
(Q(X), Q(Y), Q(Z), Q(f), Q(g), Q(h))
\ar[d]^{(a, b, c)} \\
(Q(X'), Q(Y'), Q(Z'), Q(f'), Q(g'), Q(h'))
}
$$

\medskip\noindent
Cela démontre la première assertion de la proposition. Si $\mathcal{D}$ est en
outre une catégorie triangulée, il reste à démontrer TR4 pour conclure que
$S^{-1}\mathcal{D}$ est aussi triangulée. Le lemme
\ref{lemma-easier-axiom-four} ramène cela à l'énoncé suivant : étant donnés des
morphismes composables $a : Q(X) \to Q(Y)$ et $b : Q(Y) \to Q(Z)$, il faut
construire un octaèdre, quitte à remplacer $Q(X), Q(Y), Q(Z)$ par des objets
isomorphes. On peut d'abord remplacer $Y$ par un objet tel que
$a = Q(f)$ pour un morphisme $f : X \to Y$ de $\mathcal{D}$. (Plus précisément,
écrivons $a = s^{-1}f$, avec $s : Y \to Y'$ dans $S$ et $f : X \to Y'$, puis
remplaçons $Y$ par $Y'$.) De même, remplaçons ensuite $Z$ par un objet tel que
$b = Q(g)$ pour un morphisme $g : Y \to Z$. On peut alors trouver des triangles
distingués
$(X, Y, Q_1, f, p_1, d_1)$,
$(X, Z, Q_2, g \circ f, p_2, d_2)$ et
$(Y, Z, Q_3, g, p_3, d_3)$ dans $\mathcal{D}$ (par TR1), ainsi que des morphismes
$a : Q_1 \to Q_2$ et $b : Q_2 \to Q_3$ comme dans TR4.
Il est alors immédiat que l'application du foncteur $Q$ à toutes ces données
fournit la structure correspondante dans $S^{-1}\mathcal{D}$.
\end{proof}

\noindent
La propriété universelle de la localisation d'une catégorie triangulée est la
suivante (nous la formulons pour les catégories prétriangulées ; elle vaut donc
a fortiori pour les catégories triangulées).

\begin{lemma}
\label{lemma-universal-property-localization}
Soit $\mathcal{D}$ une catégorie prétriangulée. Soit $S$ un système
multiplicatif compatible avec la structure triangulée. Soit
$Q : \mathcal{D} \to S^{-1}\mathcal{D}$ le foncteur de localisation ; voir la
proposition \ref{proposition-construct-localization}.
\begin{enumerate}
\item Si $H : \mathcal{D} \to \mathcal{A}$ est un foncteur homologique à
valeurs dans une catégorie abélienne $\mathcal{A}$ tel que $H(s)$ soit un
isomorphisme pour tout $s \in S$, alors l'unique factorisation
$H' : S^{-1}\mathcal{D} \to \mathcal{A}$ telle que $H = H' \circ Q$ (voir
Catégories, lemme \ref{categories-lemma-properties-left-localization}) est aussi
un foncteur homologique.
\item Si $F : \mathcal{D} \to \mathcal{D}'$ est un foncteur exact à valeurs
dans une catégorie prétriangulée $\mathcal{D}'$ tel que $F(s)$ soit un
isomorphisme pour tout $s \in S$, alors l'unique factorisation
$F' : S^{-1}\mathcal{D} \to \mathcal{D}'$ telle que $F = F' \circ Q$ (voir
Catégories, lemme \ref{categories-lemma-properties-left-localization}) est aussi
un foncteur exact.
\end{enumerate}
\end{lemma}

\begin{proof}
Le lemme est immédiat. Nous omettons les détails.
\end{proof}

\begin{lemma}
\label{lemma-localization-subcategory}
Soit $\mathcal{D}$ une catégorie prétriangulée et soit
$\mathcal{D}' \subset \mathcal{D}$ une sous-catégorie pleine prétriangulée.
Soit $S$ un système multiplicatif saturé de $\mathcal{D}$, compatible avec la
structure triangulée. Supposons que, pour tout $X$ dans $\mathcal{D}$, il existe
un $s : X' \to X$ dans $S$ tel que $X'$ soit un objet de $\mathcal{D}'$.
Alors $S' = S \cap \text{Flèches}(\mathcal{D}')$ est un système multiplicatif
saturé compatible avec la structure triangulée, et le foncteur
$$
(S')^{-1}\mathcal{D}' \longrightarrow S^{-1}\mathcal{D}
$$
est une équivalence de catégories prétriangulées.
\end{lemma}

\begin{proof}
Considérons le foncteur quotient $Q : \mathcal{D} \to S^{-1}\mathcal{D}$ de la
proposition \ref{proposition-construct-localization}.
Comme $S$ est saturé, un morphisme $f : X \to Y$ appartient à $S$ si et
seulement si $Q(f)$ est inversible ; voir Catégories, lemme
\ref{categories-lemma-what-gets-inverted}.
Ainsi $S'$ est la famille des flèches que la composée
$\mathcal{D}' \to \mathcal{D} \to S^{-1}\mathcal{D}$ transforme en
isomorphismes. Le lemme \ref{lemma-triangle-functor-localize} montre donc que
$S'$ est un système multiplicatif saturé compatible avec la structure
triangulée. Le lemme \ref{lemma-universal-property-localization} fournit le
foncteur exact $(S')^{-1}\mathcal{D}' \to S^{-1}\mathcal{D}$ entre catégories
prétriangulées. Par hypothèse, ce foncteur est essentiellement surjectif.
Soient $X', Y'$ des objets de $\mathcal{D}'$. D'après Catégories, remarque
\ref{categories-remark-right-localization-morphisms-colimit}, on a
$$
\Mor_{S^{-1}\mathcal{D}}(X', Y') =
\colim_{s : X \to X'\text{ dans }S} \Mor_\mathcal{D}(X, Y')
$$
Notre hypothèse implique que, pour tout $s : X \to X'$ dans $S$, il existe un
morphisme $s' : X'' \to X$ dans $S$ avec $X''$ dans $\mathcal{D}'$.
Alors $s \circ s' : X'' \to X'$ appartient à $S'$. La colimite ci-dessus est
donc égale à
$$
\colim_{s'' : X'' \to X'\text{ dans }S'} \Mor_{\mathcal{D}'}(X'', Y') =
\Mor_{(S')^{-1}\mathcal{D}'}(X', Y')
$$
Cela montre que notre foncteur est aussi pleinement fidèle et achève la
démonstration.
\end{proof}

\noindent
Le lemme suivant décrit le noyau (voir la
définition \ref{definition-kernel-category}) du foncteur de localisation.

\begin{lemma}
\label{lemma-kernel-localization}
Soit $\mathcal{D}$ une catégorie prétriangulée. Soit $S$ un système
multiplicatif compatible avec la structure triangulée. Soit $Z$ un objet de
$\mathcal{D}$. Les conditions suivantes sont équivalentes :
\begin{enumerate}
\item $Q(Z) = 0$ dans $S^{-1}\mathcal{D}$,
\item il existe $Z' \in \Ob(\mathcal{D})$ tel que
$0 : Z \to Z'$ appartienne à $S$,
\item il existe $Z' \in \Ob(\mathcal{D})$ tel que
$0 : Z' \to Z$ appartienne à $S$, et
\item il existe un objet $Z'$ et un triangle distingué
$(X, Y, Z \oplus Z', f, g, h)$ tels que $f \in S$.
\end{enumerate}
Si $S$ est saturé, ces conditions sont aussi équivalentes à
\begin{enumerate}
\item[(5)] le morphisme $0 \to Z$ appartient à $S$,
\item[(6)] le morphisme $Z \to 0$ appartient à $S$,
\item[(7)] il existe un triangle distingué $(X, Y, Z, f, g, h)$ tel que
$f \in S$.
\end{enumerate}
\end{lemma}

\begin{proof}
L'équivalence de (1), (2) et (3) est donnée par Algèbre homologique, lemme
\ref{homology-lemma-kernel-localization}.
Si (2) est satisfaite, alors
$(Z'[-1], Z'[-1] \oplus Z, Z, (1, 0), (0, 1), 0)$ est un triangle distingué
(voir le lemme \ref{lemma-split}) avec « $0 \in S$ ». Une rotation montre donc
que (4) est satisfaite. Si $(X, Y, Z \oplus Z', f, g, h)$ est un triangle
distingué avec $f \in S$, alors $Q(f)$ est un isomorphisme, donc
$Q(Z \oplus Z') = 0$, et donc $Q(Z) = 0$. Ainsi (1) -- (4) sont toutes
équivalentes.

\medskip\noindent
Supposons maintenant $S$ saturé. Chacune des conditions (5), (6), (7) implique
l'une des conditions équivalentes (1) -- (4). Supposons que $Q(Z) = 0$.
Alors $0 \to Z$ est un morphisme de $\mathcal{D}$ qui devient un isomorphisme
dans $S^{-1}\mathcal{D}$. D'après Catégories, lemme
\ref{categories-lemma-what-gets-inverted}, la saturation de $S$ implique que
$0 \to Z$ appartient à $S$. Ainsi (1) $\Rightarrow$ (5). Par dualité,
(1) $\Rightarrow$ (6). Enfin, si $0 \to Z$ appartient à $S$, alors le triangle
$(0, Z, Z, 0, \text{id}_Z, 0)$ est distingué par TR1 et TR2, et c'est un
triangle comme dans (7).
\end{proof}

\begin{lemma}
\label{lemma-limit-triangles}
Soit $\mathcal{D}$ une catégorie prétriangulée.
Soit $S$ un système multiplicatif saturé de $\mathcal{D}$, compatible avec la
structure triangulée.
Soit $(X, Y, Z, f, g, h)$ un triangle distingué de $\mathcal{D}$.
Considérons la catégorie des morphismes de triangles
$$
\mathcal{I} =
\{(s, s', s'') : (X, Y, Z, f, g, h) \to (X', Y', Z', f', g', h')
\mid s, s', s'' \in S\}
$$
Alors $\mathcal{I}$ est une catégorie filtrante, et les foncteurs
$\mathcal{I} \to X/S$, $\mathcal{I} \to Y/S$ et $\mathcal{I} \to Z/S$ sont
cofinaux.
\end{lemma}

\begin{proof}
Nous conseillons vivement au lecteur de sauter cette démonstration et de la
retrouver plutôt sur un coin de table.

\medskip\noindent
La première observation est que la rotation des triangles distingués (TR2)
donne une équivalence de catégories entre $\mathcal{I}$ et la catégorie
correspondante au triangle distingué $(Y, Z, X[1], g, h, -f[1])$.
On voit ainsi, par exemple, que si le foncteur $\mathcal{I} \to X/S$ est
cofinal, il en va de même des foncteurs $\mathcal{I} \to Y/S$ et
$\mathcal{I} \to Z/S$.

\medskip\noindent
Remarquons que si $s : X \to X'$ est un morphisme de $S$, alors MS2 permet de
trouver $s' : Y \to Y'$ et $f' : X' \to Y'$ tels que
$f' \circ s = s' \circ f$ ; MS6 permet ensuite de compléter ces données en un
objet de $\mathcal{I}$. Le foncteur $\mathcal{I} \to X/S$ est donc surjectif
sur les objets. Par rotation comme ci-dessus, il en va de même des foncteurs
$\mathcal{I} \to Y/S$ et $\mathcal{I} \to Z/S$.

\medskip\noindent
Soient $s_1 : X \to X_1 $ et $s_2 : X \to X_2$ des objets de $X/S$, et soit
$a : X_1 \to X_2$ un morphisme de $X/S$. Puisque $S$ est saturé, on a
$a \in S$ ; voir Catégories, lemme \ref{categories-lemma-what-gets-inverted}.
Par l'argument du paragraphe précédent, on peut compléter $s_1 : X \to X_1$ en
un objet
$(s_1, s'_1, s''_1) : (X, Y, Z, f, g, h) \to (X_1, Y_1, Z_1, f_1, g_1, h_1)$
de $\mathcal{I}$. En répétant l'argument, on trouve
$(a, b, c) : (X_1, Y_1, Z_1, f_1, g_1, h_1) \to (X_2, Y_2, Z_2, f_2, g_2, h_2)$
avec $a, b, c \in S$, complétant le morphisme donné $a : X_1 \to X_2$.
Alors $(a, b, c)$ est un morphisme de $\mathcal{I}$.
On en déduit que le foncteur $\mathcal{I} \to X/S$ est également surjectif sur
les flèches. Par rotation, il en va de même des foncteurs
$\mathcal{I} \to Y/S$ et $\mathcal{I} \to Z/S$.

\medskip\noindent
La catégorie $\mathcal{I}$ est non vide, car l'identité en fournit un objet.
Cela établit la condition (1) de la définition d'une catégorie filtrante ; voir
Catégories, définition \ref{categories-definition-directed}.

\medskip\noindent
Vérifions la condition (2) de Catégories, définition
\ref{categories-definition-directed}, pour la catégorie $\mathcal{I}$.
Soient donnés des objets
$(s_1, s'_1, s''_1) : (X, Y, Z, f, g, h) \to (X_1, Y_1, Z_1, f_1, g_1, h_1)$
et
$(s_2, s'_2, s''_2) : (X, Y, Z, f, g, h) \to (X_2, Y_2, Z_2, f_2, g_2, h_2)$
de $\mathcal{I}$. Nous cherchons un objet de $\mathcal{I}$ qui soit la cible
d'une flèche issue à la fois de $(X_1, Y_1, Z_1, f_1, g_1, h_1)$ et de
$(X_2, Y_2, Z_2, f_2, g_2, h_2)$. D'après Catégories, remarque
\ref{categories-remark-left-localization-morphisms-colimit}, les catégories
$X/S$, $Y/S$, $Z/S$ sont filtrantes.
On peut donc trouver $X \to X_3$ dans $X/S$ et des morphismes
$s : X_2 \to X_3$ et $a : X_1 \to X_3$. D'après ce qui précède, il existe un
morphisme
$(s, s', s'') : (X_2, Y_2, Z_2, f_2, g_2, h_2) \to
(X_3, Y_3, Z_3, f_3, g_3, h_3)$ avec $s', s'' \in S$.
Après remplacement de $(X_2, Y_2, Z_2)$ par $(X_3, Y_3, Z_3)$, on peut supposer
qu'il existe un morphisme $a : X_1 \to X_2$ dans $X/S$.
En répétant l'argument pour $Y$ et $Z$ (par rotation comme ci-dessus), on peut
supposer qu'il existe des morphismes
$a : X_1 \to X_2$ dans $X/S$,
$b : Y_1 \to Y_2$ dans $Y/S$, et
$c : Z_1 \to Z_2$ dans $Z/S$.
Ces morphismes ne définissent cependant pas nécessairement un morphisme de
triangles distingués. En revanche, les diagrammes requis commutent dans
$S^{-1}\mathcal{D}$. On voit donc, par exemple, qu'il existe un morphisme
$s'_2 : Y_2 \to Y_3$ dans $S$ tel que
$s'_2 \circ f_2 \circ a = s'_2 \circ b \circ f_1$.
Un nouveau remplacement de $(X_2, Y_2, Z_2)$ comme ci-dessus nous ramène alors
au cas où $f_2 \circ a = b \circ f_1$. Après deux autres rotations et
applications du même argument, on peut supposer que $(a, b, c)$ est un
morphisme de triangles. Cela établit la condition (2).

\medskip\noindent
Vérifions ensuite la condition (3) de Catégories, définition
\ref{categories-definition-directed}.
Supposons que
$(s_1, s_1', s_1'') : (X, Y, Z) \to (X_1, Y_1, Z_1)$ et
$(s_2, s_2', s_2'') : (X, Y, Z) \to (X_2, Y_2, Z_2)$
soient des objets de $\mathcal{I}$, et que $(a, b, c), (a', b', c')$ soient
deux morphismes entre eux. Comme $a \circ s_1 = a' \circ s_1$, il existe un
morphisme $s_3 : X_2 \to X_3$ tel que $s_3 \circ a = s_3 \circ a'$.
Par la surjectivité démontrée ci-dessus, on peut le compléter en un morphisme de
triangles
$(s_3, s_3', s_3'') : (X_2, Y_2, Z_2) \to (X_3, Y_3, Z_3)$
avec $s_3, s_3', s_3'' \in S$. Ainsi
$(s_3 \circ s_2, s_3' \circ s_2', s_3'' \circ s_2'') :
(X, Y, Z) \to (X_3, Y_3, Z_3)$ est aussi un objet de $\mathcal{I}$ ; en
composant $(a, b, c), (a', b', c')$ avec $(s_3, s_3', s_3'')$, on obtient
$a = a'$. Par rotation, on peut procéder de même pour obtenir $b = b'$ et
$c = c'$.

\medskip\noindent
Enfin, vérifions que $\mathcal{I} \to X/S$ est cofinal ; voir Catégories,
définition \ref{categories-definition-cofinal}. La première condition est
satisfaite puisque le foncteur est surjectif. Supposons donné un objet
$s : X \to X'$ de $X/S$, ainsi que deux objets
$(s_1, s'_1, s''_1) : (X, Y, Z, f, g, h) \to (X_1, Y_1, Z_1, f_1, g_1, h_1)$
et
$(s_2, s'_2, s''_2) : (X, Y, Z, f, g, h) \to (X_2, Y_2, Z_2, f_2, g_2, h_2)$
de $\mathcal{I}$, ainsi que des morphismes $t_1 : X' \to X_1$ et
$t_2 : X' \to X_2$ de $X/S$. D'après la propriété (2) de $\mathcal{I}$
démontrée ci-dessus, on peut trouver des morphismes
$(s_3, s'_3, s''_3) : (X_1, Y_1, Z_1, f_1, g_1, h_1) \to
(X_3, Y_3, Z_3, f_3, g_3, h_3)$
et
$(s_4, s'_4, s''_4) : (X_2, Y_2, Z_2, f_2, g_2, h_2) \to
(X_3, Y_3, Z_3, f_3, g_3, h_3)$ dans $\mathcal{I}$.
La démonstration serait achevée si les composées
$X' \to X_1 \to X_3$ et $X' \to X_2 \to X_3$ étaient égales (voir
l'équation affichée de Catégories, définition
\ref{categories-definition-cofinal}). Sinon, puisque $X/S$ est filtrante, on
peut choisir un morphisme $X_3 \to X_4$ dans $X/S$ tel que les composées
$X' \to X_1 \to X_3 \to X_4$ et $X' \to X_2 \to X_3 \to X_4$ soient égales.
Enfin, on complète $X_3 \to X_4$ en un morphisme
$(X_3, Y_3, Z_3) \to (X_4, Y_4, Z_4)$ de $\mathcal{I}$, puis on compose avec ce
morphisme pour obtenir le résultat.
\end{proof}

\section{Quotients de catégories triangulées}
\label{section-quotients}

\noindent
Étant données une catégorie triangulée et une sous-catégorie triangulée, on peut
construire une autre catégorie triangulée en prenant leur « quotient ».
La construction utilise une localisation. Elle est analogue au quotient d'une
catégorie abélienne par une sous-catégorie de Serre ; voir Algèbre homologique,
section \ref{homology-section-serre-subcategories}.
Avant de procéder à la construction proprement dite, examinons brièvement les
noyaux des foncteurs exacts.

\begin{definition}
\label{definition-saturated}
Soit $\mathcal{D}$ une catégorie prétriangulée. Une sous-catégorie
prétriangulée pleine $\mathcal{D}'$ de $\mathcal{D}$ est dite {\it saturée} si,
dès que $X \oplus Y$ est isomorphe à un objet de $\mathcal{D}'$, alors $X$ et
$Y$ sont tous deux isomorphes à des objets de $\mathcal{D}'$.
\end{definition}

\noindent
Une sous-catégorie triangulée saturée est parfois appelée une
{\it sous-catégorie triangulée épaisse}. Dans certaines références, cette
expression n'est employée que pour les sous-catégories triangulées strictement
pleines (et la définition est parfois formulée de manière à impliquer cette
propriété). On rencontre aussi la notion de {\it sous-catégorie triangulée
épaisse} définie par la condition suivante : pour tout diagramme commutatif
$$
\xymatrix{
& S \ar[rd] \\
X \ar[ru] \ar[rr] & & Y \ar[r] & T \ar[r] & X[1]
}
$$
dont la seconde ligne est un triangle distingué et où $S$ et $T$ sont
isomorphes à des objets de $\mathcal{D}'$, les objets $X$ et $Y$ sont également
isomorphes à des objets de $\mathcal{D}'$. Cette condition équivaut à la
saturation (c'est élémentaire ; voir \cite{Rickard-derived}), et la notion de
catégorie saturée est plus commode à manier.

\begin{lemma}
\label{lemma-triangle-functor-kernel}
Soit $F : \mathcal{D} \to \mathcal{D}'$ un foncteur exact entre catégories
prétriangulées. Soit $\mathcal{D}''$ la sous-catégorie pleine de $\mathcal{D}$
dont les objets sont
$$
\Ob(\mathcal{D}'') =
\{X \in \Ob(\mathcal{D}) \mid F(X) = 0\}
$$
Alors $\mathcal{D}''$ est une sous-catégorie prétriangulée strictement pleine
et saturée de $\mathcal{D}$. Si $\mathcal{D}$ est une catégorie triangulée,
alors $\mathcal{D}''$ est une sous-catégorie triangulée.
\end{lemma}

\begin{proof}
Il est clair que $\mathcal{D}''$ est stable par $[1]$ et $[-1]$.
Si $(X, Y, Z, f, g, h)$ est un triangle distingué de $\mathcal{D}$ et si
$F(X) = F(Y) = 0$, alors $F(Z) = 0$ également, puisque
$(F(X), F(Y), F(Z), F(f), F(g), F(h))$ est distingué.
On peut donc appliquer le lemme \ref{lemma-triangulated-subcategory} pour voir
que $\mathcal{D}''$ est une sous-catégorie prétriangulée (respectivement une
sous-catégorie triangulée si $\mathcal{D}$ est triangulée).
La saturation résulte enfin de
$F(X) \oplus F(Y) = 0 \Rightarrow F(X) = F(Y) = 0$.
\end{proof}

\begin{lemma}
\label{lemma-homological-functor-kernel}
Soit $H : \mathcal{D} \to \mathcal{A}$ un foncteur homologique d'une catégorie
prétriangulée vers une catégorie abélienne.
Soit $\mathcal{D}'$ la sous-catégorie pleine de $\mathcal{D}$ dont les objets
sont
$$
\Ob(\mathcal{D}') =
\{X \in \Ob(\mathcal{D}) \mid
H(X[n]) = 0\text{ pour tout }n \in \mathbf{Z}\}
$$
Alors $\mathcal{D}'$ est une sous-catégorie prétriangulée strictement pleine et
saturée de $\mathcal{D}$. Si $\mathcal{D}$ est une catégorie triangulée, alors
$\mathcal{D}'$ est une sous-catégorie triangulée.
\end{lemma}

\begin{proof}
Il est clair que $\mathcal{D}'$ est stable par $[1]$ et $[-1]$.
Si $(X, Y, Z, f, g, h)$ est un triangle distingué de $\mathcal{D}$ et si
$H(X[n]) = H(Y[n]) = 0$ pour tout $n$, alors $H(Z[n]) = 0$ pour tout $n$
également, d'après la longue suite exacte
(\ref{equation-long-exact-cohomology-sequence}).
On peut donc appliquer le lemme \ref{lemma-triangulated-subcategory} pour voir
que $\mathcal{D}'$ est une sous-catégorie prétriangulée (respectivement une
sous-catégorie triangulée si $\mathcal{D}$ est triangulée).
La saturation résulte de
\begin{align*}
H((X \oplus Y)[n]) = 0 & \Rightarrow H(X[n] \oplus Y[n]) = 0 \\
& \Rightarrow H(X[n]) \oplus H(Y[n]) = 0 \\
& \Rightarrow H(X[n]) = H(Y[n]) = 0
\end{align*}
pour tout $n \in \mathbf{Z}$.
\end{proof}

\begin{lemma}
\label{lemma-homological-functor-bounded}
Soit $H : \mathcal{D} \to \mathcal{A}$ un foncteur homologique d'une catégorie
prétriangulée vers une catégorie abélienne.
Soient $\mathcal{D}_H^{+}, \mathcal{D}_H^{-}, \mathcal{D}_H^b$ les
sous-catégories pleines de $\mathcal{D}$ dont les objets sont
$$
\begin{matrix}
\Ob(\mathcal{D}_H^{+}) =
\{X \in \Ob(\mathcal{D}) \mid
H(X[n]) = 0\text{ pour tout }n \ll 0\} \\
\Ob(\mathcal{D}_H^{-}) =
\{X \in \Ob(\mathcal{D}) \mid
H(X[n]) = 0\text{ pour tout }n \gg 0\} \\
\Ob(\mathcal{D}_H^b) =
\{X \in \Ob(\mathcal{D}) \mid
H(X[n]) = 0\text{ pour tout }|n| \gg 0\}
\end{matrix}
$$
Chacune est une sous-catégorie prétriangulée strictement pleine et saturée de
$\mathcal{D}$. Si $\mathcal{D}$ est une catégorie triangulée, chacune est une
sous-catégorie triangulée.
\end{lemma}

\begin{proof}
Démontrons-le pour $\mathcal{D}_H^{+}$.
Il est clair que cette sous-catégorie est stable par $[1]$ et $[-1]$.
Si $(X, Y, Z, f, g, h)$ est un triangle distingué de $\mathcal{D}$ et si
$H(X[n]) = H(Y[n]) = 0$ pour tout $n \ll 0$, alors $H(Z[n]) = 0$ pour tout
$n \ll 0$ également, d'après la longue suite exacte
(\ref{equation-long-exact-cohomology-sequence}).
On peut donc appliquer le lemme \ref{lemma-triangulated-subcategory} pour voir
que $\mathcal{D}_H^{+}$ est une sous-catégorie prétriangulée (respectivement
une sous-catégorie triangulée si $\mathcal{D}$ est triangulée).
La saturation résulte de
\begin{align*}
H((X \oplus Y)[n]) = 0 & \Rightarrow H(X[n] \oplus Y[n]) = 0 \\
& \Rightarrow H(X[n]) \oplus H(Y[n]) = 0 \\
& \Rightarrow H(X[n]) = H(Y[n]) = 0
\end{align*}
pour tout $n \in \mathbf{Z}$.
\end{proof}

\begin{definition}
\label{definition-kernel-category}
Soit $\mathcal{D}$ une catégorie (pré)triangulée.
\begin{enumerate}
\item Soit $F : \mathcal{D} \to \mathcal{D}'$ un foncteur exact.
Le {\it noyau de $F$} est la sous-catégorie (pré)triangulée strictement pleine
et saturée décrite dans le lemme \ref{lemma-triangle-functor-kernel}.
\item Soit $H : \mathcal{D} \to \mathcal{A}$ un foncteur homologique.
Le {\it noyau de $H$} est la sous-catégorie (pré)triangulée strictement pleine
et saturée décrite dans le lemme \ref{lemma-homological-functor-kernel}.
\end{enumerate}
On les note parfois $\Ker(F)$ ou $\Ker(H)$.
\end{definition}

\noindent
La démonstration du lemme suivant utilise TR4.

\begin{lemma}
\label{lemma-construct-multiplicative-system}
Soit $\mathcal{D}$ une catégorie triangulée.
Soit $\mathcal{D}' \subset \mathcal{D}$ une sous-catégorie triangulée pleine.
Posons
\begin{equation}
\label{equation-multiplicative-system}
S =
\left\{
\begin{matrix}
f \in \text{Flèches}(\mathcal{D})
\text{ tel qu'il existe un triangle distingué }\\
(X, Y, Z, f, g, h) \text{ de }\mathcal{D}\text{ avec }
Z\text{ isomorphe à un objet de }\mathcal{D}'
\end{matrix}
\right\}
\end{equation}
Alors $S$ est un système multiplicatif compatible avec la structure triangulée
de $\mathcal{D}$. Dans cette situation, les conditions suivantes sont
équivalentes :
\begin{enumerate}
\item $S$ est un système multiplicatif saturé,
\item $\mathcal{D}'$ est une sous-catégorie triangulée saturée.
\end{enumerate}
\end{lemma}

\begin{proof}
Pour démontrer la première assertion, il faut établir MS1, MS2, MS3, MS5 et
MS6.

\medskip\noindent
Démonstration de MS1. Il est clair que les identités appartiennent à $S$, car
$(X, X, 0, 1, 0, 0)$ est distingué pour tout objet $X$ de $\mathcal{D}$ et que
$0$ est un objet de $\mathcal{D}'$. Soient $f : X \to Y$ et $g : Y \to Z$ des
morphismes composables appartenant à $S$. Choisissons des triangles distingués
$(X, Y, Q_1, f, p_1, d_1)$,
$(X, Z, Q_2, g \circ f, p_2, d_2)$ et $(Y, Z, Q_3, g, p_3, d_3)$.
Par hypothèse, $Q_1$ et $Q_3$ sont isomorphes à des objets de $\mathcal{D}'$.
TR4 fournit un triangle distingué $(Q_1, Q_2, Q_3, a, b, c)$. Comme
$\mathcal{D}'$ est une sous-catégorie triangulée, $Q_2$ est isomorphe à un
objet de $\mathcal{D}'$. Ainsi $g \circ f \in S$.

\medskip\noindent
Démonstration de MS3. Soient $a : X \to Y$ un morphisme et $t : Z \to X$ un
élément de $S$ tel que $a \circ t = 0$. Pour démontrer LMS3, il suffit de
trouver $s \in S$ tel que $s \circ a = 0$ ; comparer avec la démonstration du
lemme \ref{lemma-triangle-functor-localize}. Par TR1 et TR2, choisissons un
triangle distingué $(Z, X, Q, t, g, h)$. Puisque $a \circ t = 0$, le
lemme \ref{lemma-representable-homological} fournit un morphisme
$i : Q \to Y$ tel que $i \circ g = a$. Enfin, une nouvelle application de TR1
permet de choisir un triangle $(Q, Y, W, i, s, k)$. Voici la situation :
$$
\xymatrix{
Z \ar[r]_t & X \ar[r]_g \ar[d]^1 & Q \ar[r] \ar[d]^i & Z[1] \\
& X \ar[r]_a & Y \ar[d]^s \\
& & W
}
$$
Comme $t \in S$, l'objet $Q$ est isomorphe à un objet de $\mathcal{D}'$.
Ainsi $s \in S$. Enfin, $s \circ a = s \circ i \circ g = 0$, puisque
$s \circ i = 0$ d'après le lemme \ref{lemma-composition-zero}.
On en conclut que LMS3 est satisfaite.
La démonstration de RMS3 est duale.

\medskip\noindent
Démonstration de MS5. Elle résulte de ce que les triangles distingués et
$\mathcal{D}'$ sont stables par translation.

\medskip\noindent
Démonstration de MS6. Supposons donné un diagramme commutatif
$$
\xymatrix{
X \ar[r] \ar[d]^s &
Y \ar[d]^{s'} \\
X' \ar[r] &
Y'
}
$$
avec $s, s' \in S$. La proposition \ref{proposition-9} permet de le compléter
en un diagramme à neuf carrés. Comme $s, s'$ appartiennent à $S$, les objets
$X'', Y''$ sont isomorphes à des objets de $\mathcal{D}'$. Puisque
$\mathcal{D}'$ est une sous-catégorie triangulée pleine, $Z''$ est lui aussi
isomorphe à un objet de $\mathcal{D}'$. Le morphisme $Z \to Z'$ appartient
donc à $S$. Cela démontre MS6.

\medskip\noindent
MS2 est une conséquence formelle de MS1, MS5 et MS6 ; voir le
lemme \ref{lemma-localization-conditions}.
Cela achève la démonstration de la première assertion du lemme.

\medskip\noindent
Supposons $S$ saturé. (Dans ce qui suit, nous utiliserons les rotations de
triangles distingués sans autre mention.)
Soit $X \oplus Y$ un objet isomorphe à un objet de $\mathcal{D}'$.
Considérons le morphisme $f : 0 \to X$. La composée
$0 \to X \to X \oplus Y$ appartient à $S$, car
$(0, X \oplus Y, X \oplus Y, 0, 1, 0)$ est un triangle distingué.
La composée $Y[-1] \to 0 \to X$ appartient à $S$, car
$(X, X \oplus Y, Y, (1, 0), (0, 1), 0)$ est un triangle distingué ; voir le
lemme \ref{lemma-split}.
Ainsi $0 \to X$ appartient à $S$ (puisque $S$ est saturé).
Par conséquent, $X$ est isomorphe à un objet de $\mathcal{D}'$, comme voulu.

\medskip\noindent
Supposons enfin que $\mathcal{D}'$ soit une sous-catégorie triangulée saturée.
Soient
$$
W \xrightarrow{h}
X \xrightarrow{g}
Y \xrightarrow{f} Z
$$
des morphismes composables de $\mathcal{D}$ tels que $fg, gh \in S$.
Nous allons construire progressivement un diagramme d'objets de la forme
suivante.
$$
\xymatrix{
 & &
Q_{12} \ar[rd] & &
Q_{23} \ar[rd] \\
 &
Q_1 \ar[ld]_{\! + \! 1} \ar[ru] & &
Q_2 \ar[ld]_{\! + \! 1} \ar[ll]_{\! + \! 1} \ar[ru] & &
Q_3 \ar[ld]_{\! + \! 1} \ar[ll]_{\! + \! 1} \\
W \ar[rr] & &
X \ar[lu] \ar[rr] & &
Y \ar[lu] \ar[rr] & &
Z \ar[lu]
}
$$
Choisissons d'abord des triangles distingués
$(W, X, Q_1)$, $(X, Y, Q_2)$, $(Y, Z, Q_3)$, $(W, Y, Q_{12})$ et
$(X, Z, Q_{23})$. Notons $s : Q_2 \to Q_1[1]$ la composée
$Q_2 \to X[1] \to Q_1[1]$. Notons $t : Q_3 \to Q_2[1]$ la composée
$Q_3 \to Y[1] \to Q_2[1]$.
En appliquant TR4 aux composées $W \to X \to Y$ et $X \to Y \to Z$, on
obtient des triangles distingués $(Q_1, Q_{12}, Q_2)$ et
$(Q_2, Q_{23}, Q_3)$ utilisant les morphismes $s$ et $t$.
Les objets $Q_{12}$ et $Q_{23}$ sont isomorphes à des objets de
$\mathcal{D}'$, puisque $W \to Y$ et $X \to Z$ appartiennent par hypothèse à
$S$. Ainsi $s[1]t$ appartient aussi à $S$, car $S$ est stable par composition
et par décalage. Remarquons que $s[1]t = 0$, puisque
$Y[1] \to Q_2[1] \to X[2]$ est nul ; voir le
lemme \ref{lemma-composition-zero}.
Le lemme \ref{lemma-split} montre donc que $Q_3[1] \oplus Q_1[2]$ est isomorphe
à un objet de $\mathcal{D}'$. Par l'hypothèse sur $\mathcal{D}'$, les objets
$Q_3$ et $Q_1$ sont isomorphes à des objets de $\mathcal{D}'$. Le triangle
distingué $(Q_1, Q_{12}, Q_2)$ montre alors que $Q_2$ est lui aussi isomorphe à
un objet de $\mathcal{D}'$. Enfin, le triangle distingué $(X, Y, Q_2)$ montre
que $g \in S$. (Il en résulte aussi que $h, f \in S$, mais nous n'en avons pas
besoin.)
\end{proof}

\begin{definition}
\label{definition-quotient-category}
Soit $\mathcal{D}$ une catégorie triangulée.
Soit $\mathcal{B}$ une sous-catégorie triangulée pleine.
On définit la {\it catégorie quotient $\mathcal{D}/\mathcal{B}$} par la formule
$\mathcal{D}/\mathcal{B} = S^{-1}\mathcal{D}$, où $S$ est le système
multiplicatif de $\mathcal{D}$ associé à $\mathcal{B}$ par le
lemme \ref{lemma-construct-multiplicative-system}.
Le foncteur de localisation $Q : \mathcal{D} \to \mathcal{D}/\mathcal{B}$ est
appelé dans ce cas le {\it foncteur quotient}.
\end{definition}

\noindent
Remarquons que le foncteur quotient
$Q : \mathcal{D} \to \mathcal{D}/\mathcal{B}$ est un foncteur exact entre
catégories triangulées ; voir la proposition
\ref{proposition-construct-localization}.
La propriété universelle de cette construction est la suivante.

\begin{lemma}
\label{lemma-universal-property-quotient}
\begin{slogan}
Propriété universelle du quotient de Verdier.
\end{slogan}
Soit $\mathcal{D}$ une catégorie triangulée. Soit $\mathcal{B}$ une
sous-catégorie triangulée pleine de $\mathcal{D}$. Soit
$Q : \mathcal{D} \to \mathcal{D}/\mathcal{B}$ le foncteur quotient.
\begin{enumerate}
\item Si $H : \mathcal{D} \to \mathcal{A}$ est un foncteur homologique à
valeurs dans une catégorie abélienne $\mathcal{A}$ tel que
$\mathcal{B} \subset \Ker(H)$, alors il existe une unique factorisation
$H' : \mathcal{D}/\mathcal{B} \to \mathcal{A}$ telle que $H = H' \circ Q$, et
$H'$ est lui aussi un foncteur homologique.
\item Si $F : \mathcal{D} \to \mathcal{D}'$ est un foncteur exact à valeurs
dans une catégorie prétriangulée $\mathcal{D}'$ tel que
$\mathcal{B} \subset \Ker(F)$, alors il existe une unique factorisation
$F' : \mathcal{D}/\mathcal{B} \to \mathcal{D}'$ telle que $F = F' \circ Q$, et
$F'$ est lui aussi un foncteur exact.
\end{enumerate}
\end{lemma}

\begin{proof}
Ce lemme résulte du lemme \ref{lemma-universal-property-localization}.
En effet, si $f : X \to Y$ est un morphisme de $\mathcal{D}$ tel que, pour un
certain triangle distingué $(X, Y, Z, f, g, h)$, l'objet $Z$ soit isomorphe à
un objet de $\mathcal{B}$, alors $H(f)$, respectivement $F(f)$, est un
isomorphisme sous les hypothèses de (1), respectivement de (2).
Détails omis.
\end{proof}

\noindent
Le noyau du foncteur quotient se décrit comme suit.

\begin{lemma}
\label{lemma-kernel-quotient}
Soit $\mathcal{D}$ une catégorie triangulée.
Soit $\mathcal{B}$ une sous-catégorie triangulée pleine.
Le noyau du foncteur quotient $Q : \mathcal{D} \to \mathcal{D}/\mathcal{B}$ est
la sous-catégorie strictement pleine de $\mathcal{D}$ dont les objets sont
$$
\Ob(\Ker(Q)) =
\left\{
\begin{matrix}
Z \in \Ob(\mathcal{D})
\text{ tel qu'il existe }Z' \in \Ob(\mathcal{D}) \\
\text{ tel que }Z \oplus Z'\text{ soit isomorphe à un objet de }\mathcal{B}
\end{matrix}
\right\}
$$
Autrement dit, c'est la plus petite sous-catégorie triangulée strictement
pleine et saturée de $\mathcal{D}$ contenant $\mathcal{B}$.
\end{lemma}

\begin{proof}
Remarquons d'abord que le noyau est automatiquement une sous-catégorie
triangulée strictement pleine contenant tout facteur direct de chacun de ses
objets ; voir le lemme \ref{lemma-triangle-functor-kernel}.
La description de ses objets résulte des définitions et du
lemme \ref{lemma-kernel-localization}, partie (4).
\end{proof}

\noindent
Soit $\mathcal{D}$ une catégorie triangulée.
Nous disposons maintenant de constructions qui induisent des applications
croissantes entre
\begin{enumerate}
\item l'ensemble ordonné des systèmes multiplicatifs $S$ de $\mathcal{D}$
compatibles avec la structure triangulée, et
\item l'ensemble ordonné des sous-catégories triangulées pleines
$\mathcal{B} \subset \mathcal{D}$.
\end{enumerate}
Ces constructions sont
$S \mapsto \mathcal{B}(S) = \Ker(Q : \mathcal{D} \to S^{-1}\mathcal{D})$ et
$\mathcal{B} \mapsto S(\mathcal{B})$, où $S(\mathcal{B})$ est l'ensemble
multiplicatif de (\ref{equation-multiplicative-system}), c'est-à-dire
$$
S(\mathcal{B}) =
\left\{
\begin{matrix}
f \in \text{Flèches}(\mathcal{D})
\text{ tel qu'il existe un triangle distingué }\\
(X, Y, Z, f, g, h) \text{ de }\mathcal{D}\text{ avec }
Z\text{ isomorphe à un objet de }\mathcal{B}
\end{matrix}
\right\}
$$
Ces opérations ne sont pas mutuellement inverses en général.

\begin{lemma}
\label{lemma-operations}
Soit $\mathcal{D}$ une catégorie triangulée. Les opérations décrites ci-dessus
ont les propriétés suivantes :
\begin{enumerate}
\item $S(\mathcal{B}(S))$ est la « saturation » de $S$, c'est-à-dire le plus
petit système multiplicatif saturé de $\mathcal{D}$ contenant $S$, et
\item $\mathcal{B}(S(\mathcal{B}))$ est la « saturation » de $\mathcal{B}$,
c'est-à-dire la plus petite sous-catégorie triangulée strictement pleine et
saturée de $\mathcal{D}$ contenant $\mathcal{B}$.
\end{enumerate}
En particulier, les constructions définissent des applications mutuellement
inverses entre l'ensemble ordonné des systèmes multiplicatifs saturés de
$\mathcal{D}$ compatibles avec la structure triangulée de $\mathcal{D}$ et
l'ensemble ordonné
des sous-catégories triangulées strictement pleines et saturées de
$\mathcal{D}$.
\end{lemma}

\begin{proof}
Partons d'abord d'une sous-catégorie triangulée pleine $\mathcal{B}$. Alors
$\mathcal{B}(S(\mathcal{B})) =
\Ker(Q : \mathcal{D} \to \mathcal{D}/\mathcal{B})$ ; l'assertion (2) est donc
exactement le contenu du lemme \ref{lemma-kernel-quotient}.

\medskip\noindent
Supposons ensuite que $S$ soit un système multiplicatif de $\mathcal{D}$,
compatible avec la triangulation de $\mathcal{D}$. Alors
$\mathcal{B}(S) = \Ker(Q : \mathcal{D} \to S^{-1}\mathcal{D})$.
Par conséquent (en utilisant le lemme \ref{lemma-third-object-zero} dans la
catégorie localisée),
\begin{align*}
S(\mathcal{B}(S))
& =
\left\{
\begin{matrix}
f \in \text{Flèches}(\mathcal{D})
\text{ tel qu'il existe un}\\
\text{triangle distingué }(X, Y, Z, f, g, h) \text{ de }\mathcal{D}\text{ avec }Q(Z) = 0
\end{matrix}
\right\}.
\\
& =
\{f \in \text{Flèches}(\mathcal{D}) \mid Q(f)\text{ est un isomorphisme}\} \\
& = \hat S = S'
\end{align*}
avec les notations de Catégories, lemme
\ref{categories-lemma-what-gets-inverted}.
La dernière assertion de ce lemme achève la démonstration.
\end{proof}

\begin{lemma}
\label{lemma-acyclic-general}
Soit $H : \mathcal{D} \to \mathcal{A}$ un foncteur homologique d'une catégorie
triangulée $\mathcal{D}$ vers une catégorie abélienne $\mathcal{A}$ ; voir la
définition \ref{definition-homological}.
La sous-catégorie $\Ker(H)$ de $\mathcal{D}$ est une sous-catégorie triangulée
strictement pleine et saturée de $\mathcal{D}$, et le système multiplicatif
saturé correspondant (voir le lemme \ref{lemma-operations}) est l'ensemble
$$
S = \{f \in \text{Flèches}(\mathcal{D}) \mid
H^i(f)\text{ est un isomorphisme pour tout }i \in \mathbf{Z}\}.
$$
Le foncteur $H$ se factorise par le foncteur quotient
$Q : \mathcal{D} \to \mathcal{D}/\Ker(H)$.
\end{lemma}

\begin{proof}
La catégorie $\Ker(H)$ est une sous-catégorie triangulée strictement pleine et
saturée de $\mathcal{D}$ d'après le
lemme \ref{lemma-homological-functor-kernel}.
L'ensemble $S$ est un système multiplicatif saturé compatible avec la structure
triangulée, d'après le lemme \ref{lemma-homological-functor-localize}.
Rappelons que le système multiplicatif correspondant à $\Ker(H)$ est l'ensemble
$$
\left\{
\begin{matrix}
f \in \text{Flèches}(\mathcal{D})
\text{ tel qu'il existe un triangle distingué }\\
(X, Y, Z, f, g, h)\text{ avec } H^i(Z) = 0 \text{ pour tout }i
\end{matrix}
\right\}.
$$
D'après la longue suite exacte de cohomologie
(\ref{equation-long-exact-cohomology-sequence}), il est clair que $f$ appartient
à cet ensemble si et seulement si $f$ appartient à $S$.
Enfin, la factorisation de $H$ par $Q$ résulte du
lemme \ref{lemma-universal-property-quotient}.
\end{proof}

\section{Adjoints des foncteurs exacts}
\label{section-adjoints}

\noindent
Résultats sur les foncteurs adjoints entre catégories triangulées.

\begin{lemma}
\label{lemma-adjoint-is-exact}
Soit $F : \mathcal{D} \to \mathcal{D}'$ un foncteur exact entre catégories
triangulées. Si $F$ admet un adjoint à droite
$G: \mathcal{D'} \to \mathcal{D}$, alors $G$ est lui aussi un foncteur exact.
\end{lemma}

\begin{proof}
Soit $\xi_X : F(X[1]) \to F(X)[1]$ comme dans la
définition \ref{definition-exact-functor-triangulated-categories}.
Soit $\epsilon_A : F(G(A)) \to A$ le morphisme d'adjonction.
Considérons la composée
$$
F(G(A)[1]) \xrightarrow{\xi_{G(A)}} F(G(A))[1] \xrightarrow{\epsilon_A[1]} A[1]
$$
Son morphisme adjoint est une application $G(A)[1] \to G(A[1])$, dont nous
affirmons qu'elle est un isomorphisme. Pour le voir, le lemme de Yoneda ramène
la question à montrer que l'application de $\Mor_\mathcal{D}(X[1], -)$ donne
une bijection pour tout objet $X$ de $\mathcal{D}$.
Or tout morphisme $f : X[1] \to G(A)[1]$ est de la forme $g[1]$ pour un unique
$g : X \to G(A)$, tout $g$ est adjoint à un unique $h : F(X) \to A$, et
$h[1] \circ \xi_X$ est à son tour adjoint à un unique
$i : X[1] \to G(A[1])$. La règle qui envoie $f$ sur $i$ définit une bijection
entre $\Mor_\mathcal{D}(X[1], G(A)[1])$ et
$\Mor_\mathcal{D}(X[1], G(A[1]))$.
Le fait que cette bijection soit induite par le morphisme ci-dessus résulte de
la discussion de Catégories, section \ref{categories-section-adjoint}, et du
calcul suivant :
\begin{align*}
\epsilon_A[1] \circ \xi_{G(A)} \circ F(f)
& =
\epsilon_A[1] \circ \xi_{G(A)} \circ F(g[1]) \\
& =
\epsilon_A[1] \circ F(g)[1] \circ \xi_X \\
& = 
(\epsilon_A \circ F(g))[1] \circ \xi_X \\
& =
h[1] \circ \xi_X
\end{align*}
Nous omettons certains détails. Nous allons montrer que $G$ est un foncteur
exact en utilisant les inverses des isomorphismes
$G(A)[1] \to G(A[1])$. Ces isomorphismes sont fonctoriels en $A$ (détails
omis).

\medskip\noindent
Soit $A \to B \to C \to A[1]$ un triangle distingué de $\mathcal{D}'$.
Choisissons un triangle distingué
$$
G(A) \to G(B) \to X \to G(A)[1]
$$
dans $\mathcal{D}$.
Alors $F(G(A)) \to F(G(B)) \to F(X) \to F(G(A))[1]$ est un triangle distingué
de $\mathcal{D}'$. Par TR3, on peut choisir un morphisme de triangles distingués
$$
\xymatrix{
F(G(A)) \ar[r] \ar[d]^{\epsilon_A} &
F(G(B)) \ar[r] \ar[d]^{\epsilon_B} &
F(X) \ar[r] \ar[d] &
F(G(A))[1] \ar[d]^{\epsilon_A[1]} \\
A \ar[r] & B \ar[r] & C \ar[r] & A[1]
}
$$
Rappelons que la flèche $F(X) \to F(G(A))[1]$ est la composée de
$F(X) \to F(G(A)[1])$ avec $\xi_{G(A)}$. Comme $G$ est l'adjoint à droite de
$F$, la discussion ci-dessus fournit un morphisme $X \to G(C)$ tel que le
diagramme
$$
\xymatrix{
G(A) \ar[r] \ar[d] & G(B) \ar[r] \ar[d] & X \ar[r] \ar[d] & G(A)[1] \ar[d] \\
G(A) \ar[r] & G(B) \ar[r] & G(C) \ar[r] & G(A[1])
}
$$
commute et que la flèche verticale de droite soit le morphisme construit
ci-dessus. En appliquant le foncteur homologique
$\Hom_{\mathcal{D}}(W, -)$ à un objet $W$ de $\mathcal{D}$, le $5$-lemme
montre que
$$
\Hom_{\mathcal{D}}(W, X) \to \Hom_{\mathcal{D}}(W, G(C))
$$
est une bijection. Une nouvelle application du lemme de Yoneda montre que
$X \to G(C)$ est un isomorphisme. Par conséquent,
$G(A) \to G(B) \to G(C) \to G(A)[1]$ est un triangle distingué, ce qu'il
fallait démontrer.
\end{proof}

\begin{lemma}
\label{lemma-fully-faithful-adjoint-kernel-zero}
Soient $\mathcal{D}$, $\mathcal{D}'$ des catégories triangulées.
Soient $F : \mathcal{D} \to \mathcal{D}'$ et
$G : \mathcal{D}' \to \mathcal{D}$ des foncteurs. Supposons que
\begin{enumerate}
\item $F$ et $G$ soient des foncteurs exacts,
\item $F$ soit pleinement fidèle,
\item $G$ soit un adjoint à droite de $F$, et
\item le noyau de $G$ soit nul.
\end{enumerate}
Alors $F$ est une équivalence de catégories.
\end{lemma}

\begin{proof}
Comme $F$ est pleinement fidèle, le morphisme d'adjonction
$\text{id} \to G \circ F$ est un isomorphisme (Catégories, lemme
\ref{categories-lemma-adjoint-fully-faithful}).
Soit $X$ un objet de $\mathcal{D}'$.
Choisissons un triangle distingué
$$
F(G(X)) \to X \to Y \to F(G(X))[1]
$$
dans $\mathcal{D}'$.
En appliquant $G$ et en utilisant $G(F(G(X))) = G(X)$, on obtient un triangle
distingué
$$
G(X) \to G(X) \to G(Y) \to G(X)[1]
$$
Ainsi $G(Y) = 0$. Donc $Y = 0$, et $F(G(X)) \to X$ est un isomorphisme.
\end{proof}

\section{La catégorie homotopique}
\label{section-homotopy}

\noindent
Soit $\mathcal{A}$ une catégorie additive. La catégorie homotopique
$K(\mathcal{A})$ de $\mathcal{A}$ est la catégorie des complexes de
$\mathcal{A}$ dont les morphismes sont les morphismes de complexes pris à
homotopie près. En voici la définition formelle.

\begin{definition}
\label{definition-complexes-notation}
Soit $\mathcal{A}$ une catégorie additive.
\begin{enumerate}
\item On pose $\text{Comp}(\mathcal{A}) = \text{CoCh}(\mathcal{A})$ ; c'est la
{\it catégorie des complexes (cohomologiques)}.
\item Un complexe $K^\bullet$ est dit {\it borné inférieurement} si
$K^n = 0$ pour tout $n \ll 0$.
\item Un complexe $K^\bullet$ est dit {\it borné supérieurement} si
$K^n = 0$ pour tout $n \gg 0$.
\item Un complexe $K^\bullet$ est dit {\it borné} si
$K^n = 0$ pour tout $|n| \gg 0$.
\item On note
$\text{Comp}^{+}(\mathcal{A})$, $\text{Comp}^{-}(\mathcal{A})$,
respectivement $\text{Comp}^b(\mathcal{A})$, les sous-catégories pleines de
$\text{Comp}(\mathcal{A})$ dont les objets sont les complexes bornés
inférieurement, bornés supérieurement, respectivement bornés.
\item On note $K(\mathcal{A})$ la catégorie ayant les mêmes objets que
$\text{Comp}(\mathcal{A})$, mais dont les morphismes sont les classes
d'homotopie d'applications de complexes (voir Algèbre homologique, lemme
\ref{homology-lemma-compose-homotopy-cochain}).
\item On note $K^{+}(\mathcal{A})$, $K^{-}(\mathcal{A})$,
respectivement $K^b(\mathcal{A})$, les sous-catégories pleines de
$K(\mathcal{A})$ dont les objets sont les complexes de $\mathcal{A}$ bornés
inférieurement, bornés supérieurement, respectivement bornés.
\end{enumerate}
\end{definition}

\noindent
Nous verrons que les catégories $K(\mathcal{A})$, $K^{+}(\mathcal{A})$,
$K^{-}(\mathcal{A})$ et $K^b(\mathcal{A})$ sont triangulées. Pour le démontrer,
nous développons d'abord quelques outils relatifs aux cônes et aux suites
exactes scindées.

\section{Cônes et suites scindées terme à terme}
\label{section-cones}

\noindent
Soit $\mathcal{A}$ une catégorie additive, et notons $K(\mathcal{A})$ la
catégorie des complexes de $\mathcal{A}$ dont les morphismes sont les
morphismes de complexes pris à homotopie près. Les foncteurs de décalage $[n]$
sur les complexes (voir Algèbre homologique, définition
\ref{homology-definition-shift-cochain}) induisent des foncteurs
$[n] : K(\mathcal{A}) \to K(\mathcal{A})$ tels que
$[n] \circ [m] = [n + m]$ et $[0] = \text{id}$.

\begin{definition}
\label{definition-cone}
Soit $\mathcal{A}$ une catégorie additive.
Soit $f : K^\bullet \to L^\bullet$ un morphisme de complexes de
$\mathcal{A}$. Le {\it cône} de $f$ est le complexe $C(f)^\bullet$ défini par
$C(f)^n = L^n \oplus K^{n + 1}$ et par la différentielle
$$
d_{C(f)}^n =
\left(
\begin{matrix}
d^n_L & f^{n + 1} \\
0 & -d_K^{n + 1}
\end{matrix}
\right)
$$
Il est muni de morphismes canoniques de complexes
$i : L^\bullet \to C(f)^\bullet$ et $p : C(f)^\bullet \to K^\bullet[1]$
induits par les applications évidentes $L^n \to C(f)^n \to K^{n + 1}$.
\end{definition}

\noindent
Autrement dit, $(K, L, C(f), f, i, p)$ forme un triangle :
$$
K^\bullet \to L^\bullet \to C(f)^\bullet \to K^\bullet[1]
$$
La formation de ce triangle est fonctorielle au sens suivant.

\begin{lemma}
\label{lemma-functorial-cone}
Supposons que
$$
\xymatrix{
K_1^\bullet \ar[r]_{f_1} \ar[d]_a & L_1^\bullet \ar[d]^b \\
K_2^\bullet \ar[r]^{f_2} & L_2^\bullet
}
$$
soit un diagramme de morphismes de complexes commutatif à homotopie près.
Alors il existe un morphisme $c : C(f_1)^\bullet \to C(f_2)^\bullet$ qui donne
un morphisme de triangles
$(a, b, c) : (K_1^\bullet, L_1^\bullet, C(f_1)^\bullet, f_1, i_1, p_1)
\to
(K_2^\bullet, L_2^\bullet, C(f_2)^\bullet, f_2, i_2, p_2)$
de $K(\mathcal{A})$.
\end{lemma}

\begin{proof}
Soit $h^n : K_1^n \to L_2^{n - 1}$ une famille de morphismes telle que
$b \circ f_1 - f_2 \circ a= d \circ h + h \circ d$.
Définissons $c^n$ par la matrice
$$
c^n =
\left(
\begin{matrix}
b^n & h^{n + 1} \\
0 & a^{n + 1}
\end{matrix}
\right) :
L_1^n \oplus K_1^{n + 1} \to L_2^n \oplus K_2^{n + 1}
$$
Un calcul matriciel montre que $c$ est un morphisme de complexes.
Il est immédiat que $c \circ i_1 = i_2 \circ b$, et tout aussi immédiat que
$p_2 \circ c = a \circ p_1$.
\end{proof}

\noindent
Le morphisme $c : C(f_1)^\bullet \to C(f_2)^\bullet$ construit dans la
démonstration du lemme \ref{lemma-functorial-cone} dépend en général de
l'homotopie $h$ choisie entre $f_2 \circ a$ et $b \circ f_1$.

\begin{lemma}
\label{lemma-map-from-cone}
Supposons que $f: K^\bullet \to L^\bullet$ et $g : L^\bullet \to M^\bullet$
soient des morphismes de complexes tels que $g \circ f$ soit homotope à zéro.
Alors
\begin{enumerate}
\item $g$ se factorise par un morphisme $C(f)^\bullet \to M^\bullet$, et
\item $f$ se factorise par un morphisme $K^\bullet \to C(g)^\bullet[-1]$.
\end{enumerate}
\end{lemma}

\begin{proof}
Les hypothèses signifient que le diagramme
$$
\xymatrix{
K^\bullet \ar[r]_f \ar[d] & L^\bullet \ar[d]^g \\
0 \ar[r] & M^\bullet
}
$$
commute à homotopie près.
Comme le cône de $0 \to M^\bullet$ est $M^\bullet$, l'application
$C(f)^\bullet \to C(0 \to M^\bullet) = M^\bullet$ du
lemme \ref{lemma-functorial-cone} est l'application de (1).
Le cône de $K^\bullet \to 0$ est $K^\bullet[1]$, et l'application du
lemme \ref{lemma-functorial-cone} donne un morphisme
$K^\bullet[1] \to C(g)^\bullet$. En appliquant $[-1]$, on obtient le morphisme
de (2).
\end{proof}

\noindent
Les morphismes $C(f)^\bullet \to M^\bullet$ et
$K^\bullet \to C(g)^\bullet[-1]$ construits dans la démonstration du
lemme \ref{lemma-map-from-cone} dépendent en général de l'homotopie choisie.

\begin{definition}
\label{definition-termwise-split-map}
Soit $\mathcal{A}$ une catégorie additive.
Une {\it injection scindée terme à terme
$\alpha : A^\bullet \to B^\bullet$} est un morphisme de complexes tel que
chaque $A^n \to B^n$ soit isomorphe à l'inclusion d'un facteur direct.
Une {\it surjection scindée terme à terme
$\beta : B^\bullet \to C^\bullet$} est un morphisme de complexes tel que chaque
$B^n \to C^n$ soit isomorphe à la projection sur un facteur direct.
\end{definition}

\begin{lemma}
\label{lemma-make-commute-map}
Soit $\mathcal{A}$ une catégorie additive.
Soit
$$
\xymatrix{
A^\bullet \ar[r]_f \ar[d]_a & B^\bullet \ar[d]^b \\
C^\bullet \ar[r]^g & D^\bullet
}
$$
un diagramme de morphismes de complexes commutatif à homotopie près.
Si $f$ est une injection scindée terme à terme, alors $b$ est homotope à un
morphisme qui rend le diagramme commutatif.
Si $g$ est une surjection scindée terme à terme, alors $a$ est homotope à un
morphisme qui rend le diagramme commutatif.
\end{lemma}

\begin{proof}
Soit $h^n : A^n \to D^{n - 1}$ une famille de morphismes telle que
$bf - ga = dh + hd$. Supposons que les morphismes $\pi^n : B^n \to A^n$
scindent les morphismes $f^n$.
Prenons $b' = b - dh\pi - h\pi d$.
Supposons que les morphismes $s^n : D^n \to C^n$ scindent les morphismes
$g^n : C^n \to D^n$. Prenons $a' = a + dsh + shd$.
Calculs omis.
\end{proof}

\noindent
Le lemme suivant permet de remplacer un morphisme de complexes par un morphisme
qui est, en chaque degré, l'inclusion d'un facteur direct.

\begin{lemma}
\label{lemma-make-injective}
Soit $\mathcal{A}$ une catégorie additive.
Soit $\alpha : K^\bullet \to L^\bullet$ un morphisme de complexes de
$\mathcal{A}$. Il existe une factorisation
$$
\xymatrix{
K^\bullet \ar[r]^{\tilde \alpha} \ar@/_1pc/[rr]_\alpha &
\tilde L^\bullet \ar[r]^\pi &
L^\bullet
}
$$
telle que
\begin{enumerate}
\item $\tilde \alpha$ soit une injection scindée terme à terme (voir la
définition \ref{definition-termwise-split-map}),
\item il existe une application de complexes $s : L^\bullet \to \tilde L^\bullet$
telle que $\pi \circ s = \text{id}_{L^\bullet}$ et telle que $s \circ \pi$ soit
homotope à $\text{id}_{\tilde L^\bullet}$.
\end{enumerate}
De plus, si $K^\bullet$ et $L^\bullet$ appartiennent tous deux à
$K^{+}(\mathcal{A})$, $K^{-}(\mathcal{A})$ ou $K^b(\mathcal{A})$, alors
$\tilde L^\bullet$ y appartient aussi.
\end{lemma}

\begin{proof}
Posons
$$
\tilde L^n = L^n \oplus K^n \oplus K^{n + 1}
$$
et définissons
$$
d^n_{\tilde L} =
\left(
\begin{matrix}
d^n_L & 0 & 0 \\
0 & d^n_K & \text{id}_{K^{n + 1}} \\
0 & 0 & -d^{n + 1}_K
\end{matrix}
\right)
$$
Autrement dit, $\tilde L^\bullet = L^\bullet \oplus C(1_{K^\bullet})$.
Posons en outre
$$
\tilde \alpha =
\left(
\begin{matrix}
\alpha \\
\text{id}_{K^n} \\
0
\end{matrix}
\right)
$$
qui est manifestement une injection scindée. Il est également clair que cette
matrice définit un morphisme de complexes. Définissons
$$
\pi =
\left(
\begin{matrix}
\text{id}_{L^n} &
0 &
0
\end{matrix}
\right)
$$
de sorte que $\pi \circ \tilde \alpha = \alpha$. Posons
$$
s =
\left(
\begin{matrix}
\text{id}_{L^n} \\
0 \\
0
\end{matrix}
\right)
$$
de sorte que $\pi \circ s = \text{id}_{L^\bullet}$. Enfin, soit
$h^n : \tilde L^n \to \tilde L^{n - 1}$ l'application qui envoie le facteur
$K^n$ de $\tilde L^n$ par l'identité sur le facteur $K^n$ de
$\tilde L^{n - 1}$. Il est alors immédiat (voir les calculs de la remarque
ci-dessous) que
$$
\text{id}_{\tilde L^\bullet} - s \circ \pi
=
d \circ h + h \circ d
$$
ce qui achève la démonstration du lemme.
\end{proof}

\begin{remark}
\label{remark-compute-modules}
On peut vérifier la dernière égalité affichée de la démonstration précédente en
raisonnant avec des éléments. On a
$s\pi(l, k, k^{+}) = (l, 0, 0)$.
Le morphisme du membre de gauche envoie donc
$(l, k, k^{+})$ sur $(0, k, k^{+})$.
D'autre part, $h(l, k, k^{+}) = (0, 0, k)$ et
$d(l, k, k^{+}) = (dl, dk + k^{+}, -dk^{+})$.
Ainsi $(dh + hd)(l, k, k^{+}) =
d(0, 0, k) + h(dl, dk + k^{+}, -dk^{+}) =
(0, k, -dk) + (0, 0, dk + k^{+}) = (0, k, k^{+})$, comme voulu.
\end{remark}

\begin{lemma}
\label{lemma-make-surjective}
Soit $\mathcal{A}$ une catégorie additive.
Soit $\alpha : K^\bullet \to L^\bullet$ un morphisme de complexes de
$\mathcal{A}$. Il existe une factorisation
$$
\xymatrix{
K^\bullet \ar[r]^i \ar@/_1pc/[rr]_\alpha &
\tilde K^\bullet \ar[r]^{\tilde \alpha} &
L^\bullet
}
$$
telle que
\begin{enumerate}
\item $\tilde \alpha$ soit une surjection scindée terme à terme (voir la
définition \ref{definition-termwise-split-map}),
\item il existe une application de complexes $s : \tilde K^\bullet \to K^\bullet$
telle que $s \circ i = \text{id}_{K^\bullet}$ et telle que $i \circ s$ soit
homotope à $\text{id}_{\tilde K^\bullet}$.
\end{enumerate}
De plus, si $K^\bullet$ et $L^\bullet$ appartiennent tous deux à
$K^{+}(\mathcal{A})$, $K^{-}(\mathcal{A})$ ou $K^b(\mathcal{A})$, alors
$\tilde K^\bullet$ y appartient aussi.
\end{lemma}

\begin{proof}
C'est l'énoncé dual du lemme \ref{lemma-make-injective}.
Prenons
$$
\tilde K^n = K^n \oplus L^{n - 1} \oplus L^n
$$
et définissons
$$
d^n_{\tilde K} =
\left(
\begin{matrix}
d^n_K & 0 & 0 \\
0 & - d^{n - 1}_L & \text{id}_{L^n} \\
0 & 0 & d^n_L
\end{matrix}
\right)
$$
autrement dit, $\tilde K^\bullet = K^\bullet \oplus C(1_{L^\bullet[-1]})$.
Posons en outre
$$
\tilde \alpha =
\left(
\begin{matrix}
\alpha &
0 &
\text{id}_{L^n}
\end{matrix}
\right)
$$
qui est manifestement une surjection scindée. Il est également clair que cette
matrice définit un morphisme de complexes. Définissons
$$
i =
\left(
\begin{matrix}
\text{id}_{K^n} \\
0 \\
0
\end{matrix}
\right)
$$
de sorte que $\tilde \alpha \circ i = \alpha$. Posons
$$
s =
\left(
\begin{matrix}
\text{id}_{K^n} &
0 &
0
\end{matrix}
\right)
$$
de sorte que $s \circ i = \text{id}_{K^\bullet}$. Enfin, soit
$h^n : \tilde K^n \to \tilde K^{n - 1}$ l'application qui envoie le facteur
$L^{n - 1}$ de $\tilde K^n$ par l'identité sur le facteur $L^{n - 1}$ de
$\tilde K^{n - 1}$. Il est alors immédiat que
$$
\text{id}_{\tilde K^\bullet} - i \circ s
=
d \circ h + h \circ d
$$
ce qui achève la démonstration du lemme.
\end{proof}

\begin{definition}
\label{definition-split-ses}
Soit $\mathcal{A}$ une catégorie additive.
Une {\it suite exacte de complexes de $\mathcal{A}$ scindée terme à terme} est
un complexe de complexes
$$
0 \to
A^\bullet \xrightarrow{\alpha}
B^\bullet \xrightarrow{\beta}
C^\bullet \to 0
$$
muni de décompositions en sommes directes $B^n = A^n \oplus C^n$ compatibles
avec $\alpha^n$ et $\beta^n$.
On note souvent $s^n : C^n \to B^n$ et $\pi^n : B^n \to A^n$ les applications
induites par ces décompositions en sommes directes.
D'après Algèbre homologique, lemme
\ref{homology-lemma-ses-termwise-split-cochain}, on obtient un morphisme de
complexes associé
$$
\delta : C^\bullet \longrightarrow A^\bullet[1]
$$
qui est, en degré $n$, l'application $\pi^{n + 1} \circ d_B^n \circ s^n$.
Autrement dit,
$(A^\bullet, B^\bullet, C^\bullet, \alpha, \beta, \delta)$ forme un triangle
$$
A^\bullet \to B^\bullet \to C^\bullet \to A^\bullet[1]
$$
Ce sera le {\it triangle associé à la suite exacte de complexes scindée terme
à terme}.
\end{definition}

\begin{lemma}
\label{lemma-triangle-independent-splittings}
Soit $\mathcal{A}$ une catégorie additive. Soit
$0 \to A^\bullet \to B^\bullet \to C^\bullet \to 0$ une suite exacte scindée
terme à terme comme dans la définition \ref{definition-split-ses}.
Soient $(\pi')^n$, $(s')^n$ une seconde famille de scindages.
Notons $\delta' : C^\bullet \longrightarrow A^\bullet[1]$ le morphisme associé
à cette seconde famille de scindages. Alors
$$
(1, 1, 1) :
(A^\bullet, B^\bullet, C^\bullet, \alpha, \beta, \delta)
\longrightarrow
(A^\bullet, B^\bullet, C^\bullet, \alpha, \beta, \delta')
$$
est un isomorphisme de triangles dans $K(\mathcal{A})$.
\end{lemma}

\begin{proof}
L'énoncé signifie simplement que $\delta$ et $\delta'$ sont des applications
de complexes homotopes. C'est Algèbre homologique, lemme
\ref{homology-lemma-ses-termwise-split-homotopy-cochain}.
\end{proof}

\begin{remark}
\label{remark-make-commute}
Soit $\mathcal{A}$ une catégorie additive.
Soient $0 \to A_i^\bullet \to B_i^\bullet \to C_i^\bullet \to 0$, $i = 1, 2$,
des suites exactes scindées terme à terme. Supposons que
$a : A_1^\bullet \to A_2^\bullet$,
$b : B_1^\bullet \to B_2^\bullet$ et
$c : C_1^\bullet \to C_2^\bullet$ soient des morphismes de complexes tels que
$$
\xymatrix{
A_1^\bullet \ar[d]_a \ar[r] &
B_1^\bullet \ar[r] \ar[d]_b &
C_1^\bullet \ar[d]_c \\
A_2^\bullet \ar[r] & B_2^\bullet \ar[r] & C_2^\bullet
}
$$
commute dans $K(\mathcal{A})$. En général, il n'existe {\bf pas} de morphisme
$b' : B_1^\bullet \to B_2^\bullet$ homotope à $b$ pour lequel le diagramme
ci-dessus commute dans la catégorie des complexes.
Considérons en effet Exemples, équation
(\ref{examples-equation-commutes-up-to-homotopy}). Si l'on pouvait y remplacer
l'application centrale par une application homotope rendant le diagramme
commutatif, on obtiendrait l'additivité des traces, qui est fausse.
\end{remark}

\begin{lemma}
\label{lemma-nilpotent}
Soit $\mathcal{A}$ une catégorie additive.
Soient $0 \to A_i^\bullet \to B_i^\bullet \to C_i^\bullet \to 0$,
$i = 1, 2, 3$, des suites exactes de complexes scindées terme à terme. Soient
$b : B_1^\bullet \to B_2^\bullet$ et $b' : B_2^\bullet \to B_3^\bullet$ des
morphismes de complexes tels que
$$
\vcenter{
\xymatrix{
A_1^\bullet \ar[d]_0 \ar[r] &
B_1^\bullet \ar[r] \ar[d]_b &
C_1^\bullet \ar[d]_0 \\
A_2^\bullet \ar[r] & B_2^\bullet \ar[r] & C_2^\bullet
}
}
\quad\text{et}\quad
\vcenter{
\xymatrix{
A_2^\bullet \ar[d]^0 \ar[r] &
B_2^\bullet \ar[r] \ar[d]^{b'} &
C_2^\bullet \ar[d]^0 \\
A_3^\bullet \ar[r] & B_3^\bullet \ar[r] & C_3^\bullet
}
}
$$
commutent dans $K(\mathcal{A})$. Alors $b' \circ b = 0$ dans
$K(\mathcal{A})$.
\end{lemma}

\begin{proof}
Le lemme \ref{lemma-make-commute-map} permet de remplacer $b$ et $b'$ par des
applications homotopes de sorte que le carré droit du diagramme de gauche et le
carré gauche du diagramme de droite commutent. Autrement dit,
$\Im(b^n) \subset \Im(A_2^n \to B_2^n)$ et
$\Ker((b')^n) \supset \Im(A_2^n \to B_2^n)$.
Alors $b' \circ b = 0$ comme application de complexes.
\end{proof}

\begin{lemma}
\label{lemma-third-isomorphism}
Soit $\mathcal{A}$ une catégorie additive.
Soient $f_1 : K_1^\bullet \to L_1^\bullet$ et
$f_2 : K_2^\bullet \to L_2^\bullet$ des morphismes de complexes.
Soit
$$
(a, b, c) :
(K_1^\bullet, L_1^\bullet, C(f_1)^\bullet, f_1, i_1, p_1)
\longrightarrow
(K_2^\bullet, L_2^\bullet, C(f_2)^\bullet, f_2, i_2, p_2)
$$
un morphisme quelconque de triangles de $K(\mathcal{A})$.
Si $a$ et $b$ sont des équivalences d'homotopie, alors $c$ l'est aussi.
\end{lemma}

\begin{proof}
Soit $a^{-1} : K_2^\bullet \to K_1^\bullet$ un morphisme de complexes inverse
de $a$ dans $K(\mathcal{A})$.
Soit $b^{-1} : L_2^\bullet \to L_1^\bullet$ un morphisme de complexes inverse
de $b$ dans $K(\mathcal{A})$.
Soit $c' : C(f_2)^\bullet \to C(f_1)^\bullet$ le morphisme fourni par le
lemme \ref{lemma-functorial-cone} appliqué à
$f_1 \circ a^{-1} = b^{-1} \circ f_2$. Si l'on montre que $c \circ c'$ et
$c' \circ c$ sont des isomorphismes dans $K(\mathcal{A})$, la démonstration
est achevée. Il suffit donc d'établir l'assertion suivante : étant donné un
morphisme de triangles
$(1, 1, c) : (K^\bullet, L^\bullet, C(f)^\bullet, f, i, p)$
dans $K(\mathcal{A})$, le morphisme $c$ est un isomorphisme dans
$K(\mathcal{A})$. Par hypothèse, les deux carrés du diagramme
$$
\xymatrix{
L^\bullet \ar[r] \ar[d]_1 &
C(f)^\bullet \ar[r] \ar[d]_c &
K^\bullet[1] \ar[d]_1 \\
L^\bullet \ar[r] &
C(f)^\bullet \ar[r] &
K^\bullet[1]
}
$$
commutent à homotopie près. Par construction de $C(f)^\bullet$, les lignes
forment des suites de complexes scindées terme à terme. Le
lemme \ref{lemma-nilpotent} donne donc $(c - 1)^2 = 0$ dans $K(\mathcal{A})$.
Ainsi $c$ est un isomorphisme de $K(\mathcal{A})$, d'inverse $2 - c$.
\end{proof}

\noindent
Ainsi, si $a$ et $b$ sont des équivalences d'homotopie, le morphisme de
triangles qui en résulte est un isomorphisme de triangles dans
$K(\mathcal{A})$.
Il s'avère que la famille des triangles de $K(\mathcal{A})$ définis par les
cônes et celle des triangles de $K(\mathcal{A})$ définis par les suites exactes
de complexes scindées terme à terme coïncident à isomorphisme près, au moins à
un signe près !

\begin{lemma}
\label{lemma-the-same-up-to-isomorphisms}
Soit $\mathcal{A}$ une catégorie additive.
\begin{enumerate}
\item Étant donnée une suite de complexes scindée terme à terme
$(\alpha : A^\bullet \to B^\bullet,
\beta : B^\bullet \to C^\bullet, s^n, \pi^n)$,
il existe une équivalence d'homotopie $C(\alpha)^\bullet \to C^\bullet$ telle
que le diagramme
$$
\xymatrix{
A^\bullet \ar[r] \ar[d] & B^\bullet \ar[d] \ar[r] &
C(\alpha)^\bullet \ar[r]_{-p} \ar[d] & A^\bullet[1] \ar[d] \\
A^\bullet \ar[r] & B^\bullet \ar[r] &
C^\bullet \ar[r]^\delta & A^\bullet[1]
}
$$
définisse un isomorphisme de triangles dans $K(\mathcal{A})$.
\item Étant donné un morphisme de complexes $f : K^\bullet \to L^\bullet$,
il existe un isomorphisme de triangles
$$
\xymatrix{
K^\bullet \ar[r] \ar[d] & \tilde L^\bullet \ar[d] \ar[r] &
M^\bullet \ar[r]_{\delta} \ar[d] & K^\bullet[1] \ar[d] \\
K^\bullet \ar[r] & L^\bullet \ar[r] &
C(f)^\bullet \ar[r]^{-p} & K^\bullet[1]
}
$$
où le triangle supérieur est celui qui est associé à une suite exacte scindée
terme à terme $K^\bullet \to \tilde L^\bullet \to M^\bullet$.
\end{enumerate}
\end{lemma}

\begin{proof}
Démonstration de (1). On a $C(\alpha)^n = B^n \oplus A^{n + 1}$, et l'on
définit simplement $C(\alpha)^n \to C^n$ comme la projection sur $B^n$ suivie
de $\beta^n$. C'est un morphisme de complexes, car les composées
$A^{n + 1} \to B^{n + 1} \to C^{n + 1}$ sont nulles.
Pour obtenir un inverse à homotopie près, on prend
$C^\bullet \to C(\alpha)^\bullet$ défini par
$(s^n , -\delta^n)$ en degré $n$. C'est un morphisme de complexes, car $\delta^n$ se
caractérise comme l'unique morphisme $C^n \to A^{n + 1}$ tel que
$d \circ s^n - s^{n + 1} \circ d = \alpha \circ \delta^n$ ; voir la
démonstration d'Algèbre homologique, lemme
\ref{homology-lemma-ses-termwise-split-cochain}.
La composée $C^\bullet \to C(\alpha)^\bullet \to C^\bullet$ est l'identité.
La composée $C(\alpha)^\bullet \to C^\bullet \to C(\alpha)^\bullet$ est égale
au morphisme
$$
\left(
\begin{matrix}
s^n \circ \beta^n & 0 \\
-\delta^n \circ \beta^n & 0
\end{matrix}
\right)
$$
Pour voir qu'il est homotope à l'identité, utilisons l'homotopie
$h^n : C(\alpha)^n \to C(\alpha)^{n - 1}$ définie par la matrice
$$
\left(
\begin{matrix}
0 & 0 \\
\pi^n & 0
\end{matrix}
\right) : C(\alpha)^n = B^n \oplus A^{n + 1} \to
B^{n - 1} \oplus A^n = C(\alpha)^{n - 1}
$$
On vérifie immédiatement que
$$
\left(
\begin{matrix}
1 & 0 \\
0 & 1
\end{matrix}
\right)
-
\left(
\begin{matrix}
s^n \\
-\delta^n
\end{matrix}
\right)
\left(
\begin{matrix}
\beta^n & 0
\end{matrix}
\right)
=
\left(
\begin{matrix}
d & \alpha^n \\
0 & -d
\end{matrix}
\right)
\left(
\begin{matrix}
0 & 0 \\
\pi^n & 0
\end{matrix}
\right)
+
\left(
\begin{matrix}
0 & 0 \\
\pi^{n + 1} & 0
\end{matrix}
\right)
\left(
\begin{matrix}
d & \alpha^{n + 1} \\
0 & -d
\end{matrix}
\right)
$$
Pour achever la démonstration de (1), il faut montrer que les morphismes
$-p : C(\alpha)^\bullet \to A^\bullet[1]$ (voir la
définition \ref{definition-cone}) et
$C(\alpha)^\bullet \to C^\bullet \to A^\bullet[1]$ coïncident à homotopie
près. C'est clair d'après ce qui précède. On peut en effet utiliser l'inverse à
homotopie près $(s, -\delta) : C^\bullet \to C(\alpha)^\bullet$ et vérifier à
la place que les deux applications $C^\bullet \to A^\bullet[1]$ coïncident.
Or $p \circ (s, -\delta) = -\delta$, comme voulu.

\medskip\noindent
Démonstration de (2). Soient $\tilde f : K^\bullet \to \tilde L^\bullet$,
$s : L^\bullet \to \tilde L^\bullet$ et
$\pi : \tilde L^\bullet \to L^\bullet$ comme dans le
lemme \ref{lemma-make-injective}. D'après les
lemmes \ref{lemma-functorial-cone} et \ref{lemma-third-isomorphism}, les
triangles $(K^\bullet, L^\bullet, C(f), f, i, p)$ et
$(K^\bullet, \tilde L^\bullet, C(\tilde f), \tilde f, \tilde i, \tilde p)$
sont isomorphes. On peut composer les isomorphismes de triangles. Ainsi, on
peut remplacer $L^\bullet$ par $\tilde L^\bullet$ et $f$ par $\tilde f$.
Autrement dit, on peut supposer que $f$ est une injection scindée terme à
terme. Dans ce cas, le résultat découle de la partie (1).
\end{proof}

\begin{lemma}
\label{lemma-sequence-maps-split}
Soit $\mathcal{A}$ une catégorie additive.
Soit $A_1^\bullet \to A_2^\bullet \to \ldots \to A_n^\bullet$ une suite de
morphismes composables de complexes. Il existe un diagramme commutatif
$$
\xymatrix{
A_1^\bullet \ar[r] &
A_2^\bullet \ar[r] &
\ldots \ar[r] &
A_n^\bullet \\
B_1^\bullet \ar[r] \ar[u] &
B_2^\bullet \ar[r] \ar[u] &
\ldots \ar[r] &
B_n^\bullet \ar[u]
}
$$
tel que chaque morphisme $B_i^\bullet \to B_{i + 1}^\bullet$ soit une
injection scindée terme à terme, et que chaque
$B_i^\bullet \to A_i^\bullet$ soit une équivalence d'homotopie.
De plus, si tous les $A_i^\bullet$ appartiennent à
$K^{+}(\mathcal{A})$, $K^{-}(\mathcal{A})$ ou $K^b(\mathcal{A})$, alors il en
va de même des $B_i^\bullet$.
\end{lemma}

\begin{proof}
Le cas $n = 1$ est vide.
Le lemme \ref{lemma-make-injective} est le cas $n = 2$.
Supposons le diagramme construit, à l'exception de $B_n^\bullet$.
Appliquons le lemme \ref{lemma-make-injective} à la composée
$B_{n - 1}^\bullet \to A_{n - 1}^\bullet \to A_n^\bullet$.
On obtient la factorisation
$B_{n - 1}^\bullet \to B_n^\bullet \to A_n^\bullet$ recherchée.
\end{proof}

\begin{lemma}
\label{lemma-rotate-triangle}
Soit $\mathcal{A}$ une catégorie additive. Soit
$(\alpha : A^\bullet \to B^\bullet, \beta : B^\bullet \to C^\bullet, s^n,
\pi^n)$ une suite de complexes scindée terme à terme.
Soit $(A^\bullet, B^\bullet, C^\bullet, \alpha, \beta, \delta)$ le triangle
associé. Alors le triangle
$(C^\bullet[-1], A^\bullet, B^\bullet, \delta[-1], \alpha, \beta)$
est isomorphe au triangle
$(C^\bullet[-1], A^\bullet, C(\delta[-1])^\bullet, \delta[-1], i, p)$.
\end{lemma}

\begin{proof}
Écrivons $B^n = A^n \oplus C^n$ et identifions $\alpha^n$ et $\beta^n$ aux
applications naturelles d'inclusion et de projection. Par construction de
$\delta$, on a
$$
d_B^n =
\left(
\begin{matrix}
d_A^n & \delta^n \\
0 & d_C^n
\end{matrix}
\right)
$$
D'autre part, le cône de $\delta[-1] : C^\bullet[-1] \to A^\bullet$ est donné
par $C(\delta[-1])^n = A^n \oplus C^n$, avec une différentielle identique à la
matrice ci-dessus. D'où le lemme.
\end{proof}

\begin{lemma}
\label{lemma-rotate-cone}
Soit $\mathcal{A}$ une catégorie additive.
Soit $f : K^\bullet \to L^\bullet$ un morphisme de complexes.
Le triangle $(L^\bullet, C(f)^\bullet, K^\bullet[1], i, p, f[1])$ est le
triangle associé à la suite scindée terme à terme
$$
0 \to L^\bullet \to C(f)^\bullet \to K^\bullet[1] \to 0
$$
qui provient de la définition du cône de $f$.
\end{lemma}

\begin{proof}
Cela résulte immédiatement des définitions.
\end{proof}

\section{Triangles distingués dans la catégorie d'homotopie}
\label{section-homotopy-triangulated}

\noindent
Comme nous voulons que, dans les suites exactes longues de cohomologie, nos
morphismes de bord soient donnés sans signe par les morphismes du lemme du
serpent, nous définissons comme suit les triangles distingués dans la catégorie
d'homotopie.

\begin{definition}
\label{definition-distinguished-triangle}
Soit $\mathcal{A}$ une catégorie additive.
Un triangle $(X, Y, Z, f, g, h)$ de $K(\mathcal{A})$ est appelé un
{\it triangle distingué de $K(\mathcal{A})$} s'il est isomorphe au triangle
associé à une suite exacte de complexes scindée terme à terme, voir la
définition \ref{definition-split-ses}.
Même définition pour $K^{+}(\mathcal{A})$, $K^{-}(\mathcal{A})$ et
$K^b(\mathcal{A})$.
\end{definition}

\noindent
Notons que, d'après le lemme \ref{lemma-the-same-up-to-isomorphisms}, un
triangle de la forme $(K^\bullet, L^\bullet, C(f)^\bullet, f, i, -p)$ est un
triangle distingué. Cela donne bien une catégorie triangulée, voir la
proposition \ref{proposition-homotopy-category-triangulated}. Avant de pouvoir
démontrer cette proposition, il nous faut encore un lemme afin de pouvoir
démontrer TR4.

\begin{lemma}
\label{lemma-two-split-injections}
Soit $\mathcal{A}$ une catégorie additive. Supposons que
$\alpha : A^\bullet \to B^\bullet$ et $\beta : B^\bullet \to C^\bullet$
soient des injections scindées de complexes. Alors il existe des triangles
distingués
$(A^\bullet, B^\bullet, Q_1^\bullet, \alpha, p_1, d_1)$,
$(A^\bullet, C^\bullet, Q_2^\bullet, \beta \circ \alpha, p_2, d_2)$ et
$(B^\bullet, C^\bullet, Q_3^\bullet, \beta, p_3, d_3)$
pour lesquels l'axiome TR4 est satisfait.
\end{lemma}

\begin{proof}
Notons $\pi_1^n : B^n \to A^n$ et $\pi_3^n : C^n \to B^n$ les rétractions.
Alors $A^\bullet \to C^\bullet$ est aussi une injection scindée, de
rétractions $\pi_2^n = \pi_1^n \circ \pi_3^n$. Notons $Q_1^\bullet$,
$Q_2^\bullet$ et $Q_3^\bullet$ les complexes « quotients ». Autrement dit,
$Q_1^n = \Ker(\pi_1^n)$, $Q_3^n = \Ker(\pi_3^n)$ et
$Q_2^n = \Ker(\pi_2^n)$. Notons que ces noyaux existent. On a alors
$B^n = A^n \oplus Q_1^n$ et $C^n = B^n \oplus Q_3^n$, où l'on considère
$A^n$ comme un sous-objet de $B^n$, et ainsi de suite. Il s'ensuit que
$C^n = A^n \oplus Q_1^n \oplus Q_3^n$. Notons que
$\pi_2^n = \pi_1^n \circ \pi_3^n$ est nulle tant sur $Q_1^n$ que sur $Q_3^n$.
Ainsi, $Q_2^n = Q_1^n \oplus Q_3^n$. Considérons le diagramme commutatif
$$
\begin{matrix}
0 & \to & A^\bullet & \to & B^\bullet & \to & Q_1^\bullet & \to & 0 \\
  &     & \downarrow &     & \downarrow &     & \downarrow  & \\
0 & \to & A^\bullet & \to & C^\bullet & \to & Q_2^\bullet & \to & 0 \\
  &     & \downarrow &     & \downarrow &     & \downarrow  & \\
0 & \to & B^\bullet & \to & C^\bullet & \to & Q_3^\bullet & \to & 0
\end{matrix}
$$
Les lignes de ce diagramme sont des suites exactes scindées terme à terme et
définissent donc, par définition, des triangles distingués. En outre, les
flèches descendantes du diagramme ci-dessus sont compatibles avec les
scindages choisis et définissent donc des morphismes de triangles
$$
(A^\bullet \to B^\bullet \to Q_1^\bullet \to A^\bullet[1])
\longrightarrow
(A^\bullet \to C^\bullet \to Q_2^\bullet \to A^\bullet[1])
$$
et
$$
(A^\bullet \to C^\bullet \to Q_2^\bullet \to A^\bullet[1])
\longrightarrow
(B^\bullet \to C^\bullet \to Q_3^\bullet \to B^\bullet[1]).
$$
Notons que les sections $Q_3^n \to C^n$ de la suite scindée inférieure du
diagramme fournissent une section de la suite scindée
$0 \to Q_1^\bullet \to Q_2^\bullet \to Q_3^\bullet \to 0$ après composition
avec $C^n \to Q_2^n$. On en déduit aisément que le morphisme
$\delta : Q_3^\bullet \to Q_1^\bullet[1]$ du triangle distingué correspondant
$$
(Q_1^\bullet \to Q_2^\bullet \to Q_3^\bullet \to Q_1^\bullet[1])
$$
est égal à la composée $Q_3^\bullet \to B^\bullet[1] \to Q_1^\bullet[1]$.
On obtient ainsi une configuration comme celle de la conclusion de l'axiome
TR4.
\end{proof}


\begin{proposition}
\label{proposition-homotopy-category-triangulated}
Soit $\mathcal{A}$ une catégorie additive. La catégorie $K(\mathcal{A})$ des
complexes à homotopie près, munie de ses foncteurs de translation naturels et
des triangles distingués définis ci-dessus, est une catégorie triangulée.
\end{proposition}

\begin{proof}
Démonstration de TR1. Par définition, tout triangle isomorphe à un triangle
distingué est distingué. De plus, tout triangle
$(A^\bullet, A^\bullet, 0, 1, 0, 0)$ est distingué, puisque
$0 \to A^\bullet \to A^\bullet \to 0 \to 0$ est une suite de complexes
scindée terme à terme. Enfin, pour tout morphisme de complexes
$f : K^\bullet \to L^\bullet$, le triangle
$(K, L, C(f), f, i, -p)$ est distingué d'après le
lemme \ref{lemma-the-same-up-to-isomorphisms}.

\medskip\noindent
Démonstration de TR2. Soit $(X, Y, Z, f, g, h)$ un triangle. Supposons que
$(Y, Z, X[1], g, h, -f[1])$ soit distingué. Il existe alors une suite de
complexes scindée terme à terme $A^\bullet \to B^\bullet \to C^\bullet$ telle
que le triangle associé
$(A^\bullet, B^\bullet, C^\bullet, \alpha, \beta, \delta)$ soit isomorphe à
$(Y, Z, X[1], g, h, -f[1])$. En faisant tourner en sens inverse, on voit que
$(X, Y, Z, f, g, h)$ est isomorphe à
$(C^\bullet[-1], A^\bullet, B^\bullet, -\delta[-1], \alpha, \beta)$.
Il résulte du lemme \ref{lemma-rotate-triangle} que le triangle
$(C^\bullet[-1], A^\bullet, B^\bullet, \delta[-1], \alpha, \beta)$ est
isomorphe à
$(C^\bullet[-1], A^\bullet, C(\delta[-1])^\bullet, \delta[-1], i, p)$.
En précomposant l'isomorphisme de triangles précédent par $-1$ sur $Y$, on en
déduit que $(X, Y, Z, f, g, h)$ est isomorphe à
$(C^\bullet[-1], A^\bullet, C(\delta[-1])^\bullet, \delta[-1], i, -p)$.
Il est donc distingué d'après le
lemme \ref{lemma-the-same-up-to-isomorphisms}.
Réciproquement, supposons que $(X, Y, Z, f, g, h)$ soit distingué. D'après le
lemme \ref{lemma-the-same-up-to-isomorphisms}, cela signifie qu'il est
isomorphe à un triangle de la forme
$(K^\bullet, L^\bullet, C(f), f, i, -p)$ pour un certain morphisme de
complexes $f$. Le triangle obtenu par rotation
$(Y, Z, X[1], g, h, -f[1])$ est alors isomorphe à
$(L^\bullet, C(f), K^\bullet[1], i, -p, -f[1])$, qui est isomorphe au triangle
$(L^\bullet, C(f), K^\bullet[1], i, p, f[1])$. D'après le
lemme \ref{lemma-rotate-cone}, ce triangle est distingué. Ainsi,
$(Y, Z, X[1], g, h, -f[1])$ est distingué, comme voulu.

\medskip\noindent
Démonstration de TR3. Soient
$(X, Y, Z, f, g, h)$ et $(X', Y', Z', f', g', h')$ des triangles distingués de
$K(\mathcal{A})$, et soient $a : X \to X'$ et $b : Y \to Y'$ des morphismes
tels que $f' \circ a = b \circ f$. D'après le
lemme \ref{lemma-the-same-up-to-isomorphisms}, on peut supposer que
$(X, Y, Z, f, g, h) = (X, Y, C(f), f, i, -p)$ et
$(X', Y', Z', f', g', h') = (X', Y', C(f'), f', i', -p')$.
Il suffit alors d'appliquer le lemme \ref{lemma-functorial-cone} au diagramme
commutatif donné par $f, f', a, b$.

\medskip\noindent
Démonstration de TR4. Nous savons à présent que $K(\mathcal{A})$ est une
catégorie prétriangulée. Nous pouvons donc utiliser le
lemme \ref{lemma-easier-axiom-four}. Soient
$A^\bullet \to B^\bullet$ et $B^\bullet \to C^\bullet$ des morphismes
composables de $K(\mathcal{A})$. D'après le
lemme \ref{lemma-sequence-maps-split}, on peut supposer que
$A^\bullet \to B^\bullet$ et $B^\bullet \to C^\bullet$ sont des morphismes
injectifs scindés. Dans ce cas, le résultat découle du
lemme \ref{lemma-two-split-injections}.
\end{proof}

\begin{remark}
\label{remark-boundedness-conditions-triangulated}
Soit $\mathcal{A}$ une catégorie additive. Exactement la même démonstration que
celle de la proposition \ref{proposition-homotopy-category-triangulated}
montre que les catégories
$K^{+}(\mathcal{A})$, $K^{-}(\mathcal{A})$ et $K^b(\mathcal{A})$ sont des
catégories triangulées. En effet, le cône d'un morphisme entre complexes bornés
(supérieurement, inférieurement) est borné (supérieurement, inférieurement).
Nous démontrons toutefois ci-dessous que ce sont des sous-catégories
triangulées de $K(\mathcal{A})$, ce qui fournit une autre démonstration.
\end{remark}

\begin{lemma}
\label{lemma-bounded-triangulated-subcategories}
Soit $\mathcal{A}$ une catégorie additive. Les catégories
$K^{+}(\mathcal{A})$, $K^{-}(\mathcal{A})$ et $K^b(\mathcal{A})$ sont des
sous-catégories triangulées pleines de $K(\mathcal{A})$.
\end{lemma}

\begin{proof}
Chacune des catégories mentionnées est une sous-catégorie additive pleine.
Nous utilisons le critère du lemme \ref{lemma-triangulated-subcategory} pour
montrer que ce sont des sous-catégories triangulées. Il est clair que chacune
des catégories $K^{+}(\mathcal{A})$, $K^{-}(\mathcal{A})$ et
$K^b(\mathcal{A})$ est préservée par les foncteurs de décalage $[1], [-1]$.
Enfin, supposons que $f : A^\bullet \to B^\bullet$ soit un morphisme dans
$K^{+}(\mathcal{A})$, $K^{-}(\mathcal{A})$ ou $K^b(\mathcal{A})$.
Alors $(A^\bullet, B^\bullet, C(f)^\bullet, f, i, -p)$ est un triangle
distingué de $K(\mathcal{A})$, avec
$C(f)^\bullet \in K^{+}(\mathcal{A})$, $K^{-}(\mathcal{A})$ ou
$K^b(\mathcal{A})$, comme il ressort clairement de la construction du cône.
Le lemme est donc démontré. (Une autre démonstration consiste à remarquer que
$K^\bullet \to L^\bullet$ est isomorphe, dans $K^{+}(\mathcal{A})$,
$K^{-}(\mathcal{A})$ ou $K^b(\mathcal{A})$, à une injection de complexes
scindée terme à terme, voir le lemme \ref{lemma-make-injective}; on peut alors
prendre directement le triangle distingué associé.)
\end{proof}

\begin{lemma}
\label{lemma-additive-exact-homotopy-category}
Soient $\mathcal{A}$, $\mathcal{B}$ des catégories additives. Soit
$F : \mathcal{A} \to \mathcal{B}$ un foncteur additif. Les foncteurs induits
$$
\begin{matrix}
F : K(\mathcal{A}) \longrightarrow K(\mathcal{B}) \\
F : K^{+}(\mathcal{A}) \longrightarrow K^{+}(\mathcal{B}) \\
F : K^{-}(\mathcal{A}) \longrightarrow K^{-}(\mathcal{B}) \\
F : K^b(\mathcal{A}) \longrightarrow K^b(\mathcal{B})
\end{matrix}
$$
sont des foncteurs exacts de catégories triangulées.
\end{lemma}

\begin{proof}
Supposons que $A^\bullet \to B^\bullet \to C^\bullet$ soit une suite de
complexes de $\mathcal{A}$ scindée terme à terme, de scindages
$(s^n, \pi^n)$ et de morphisme associé
$\delta : C^\bullet \to A^\bullet[1]$, voir la
définition \ref{definition-split-ses}. Alors
$F(A^\bullet) \to F(B^\bullet) \to F(C^\bullet)$ est une suite de complexes
scindée terme à terme, de scindages $(F(s^n), F(\pi^n))$ et de morphisme
associé $F(\delta) : F(C^\bullet) \to F(A^\bullet)[1]$. Ainsi, $F$ transforme
les triangles distingués en triangles distingués.
\end{proof}

\begin{lemma}
\label{lemma-improve-distinguished-triangle-homotopy}
Soit $\mathcal{A}$ une catégorie additive. Soit
$(A^\bullet, B^\bullet, C^\bullet, a, b, c)$ un triangle distingué de
$K(\mathcal{A})$. Il existe alors un triangle distingué isomorphe
$(A^\bullet, (B')^\bullet, C^\bullet, a', b', c)$ tel que
$0 \to A^n \to (B')^n \to C^n \to 0$ soit une suite exacte courte scindée
pour tout $n$.
\end{lemma}

\begin{proof}
Nous utiliserons le fait que $K(\mathcal{A})$ est une catégorie triangulée,
d'après la proposition \ref{proposition-homotopy-category-triangulated}.
Soit $W^\bullet$ le cône de $c : C^\bullet \to A^\bullet[1]$, muni de ses
morphismes $i : A^\bullet[1] \to W^\bullet$ et
$p : W^\bullet \to C^\bullet[1]$. Le triangle
$(C^\bullet, A^\bullet[1], W^\bullet, c, i, -p)$ est alors distingué d'après le
lemme \ref{lemma-the-same-up-to-isomorphisms}. Après deux rotations en sens
inverse, on voit que
$(A^\bullet, W^\bullet[-1], C^\bullet, -i[-1], p[-1], c)$ est un triangle
distingué. Par TR3, il existe un morphisme de triangles distingués
$(\text{id}, \beta, \text{id}) : (A^\bullet, B^\bullet, C^\bullet, a, b, c) \to
(A^\bullet, W^\bullet[-1], C^\bullet, -i[-1], p[-1], c)$,
qui est nécessairement un isomorphisme d'après le
lemme \ref{lemma-third-isomorphism-triangle}. Cela achève la démonstration,
puisque $0 \to A^\bullet \to W^\bullet[-1] \to C^\bullet \to 0$ est une suite
exacte courte de complexes scindée terme à terme, par la construction même des
cônes dans la section \ref{section-cones}.
\end{proof}

\begin{remark}
\label{remark-homotopy-double}
Soit $\mathcal{A}$ une catégorie additive admettant les sommes directes
dénombrables. Notons $\text{DoubleComp}(\mathcal{A})$ la catégorie des
complexes doubles de $\mathcal{A}$, voir Homologie, section
\ref{homology-section-double-complexes}. Cette catégorie permet de construire
deux catégories triangulées.
\begin{enumerate}
\item On peut considérer un objet $A^{\bullet, \bullet}$ de
$\text{DoubleComp}(\mathcal{A})$ comme un complexe de complexes de la manière
suivante :
$$
\ldots \to A^{\bullet, -1} \to A^{\bullet, 0} \to A^{\bullet, 1} \to \ldots
$$
et prendre la catégorie d'homotopie
$K_{first}(\text{DoubleComp}(\mathcal{A}))$, munie de la structure triangulée
correspondante donnée par la
proposition \ref{proposition-homotopy-category-triangulated}.
D'après Homologie, remarque
\ref{homology-remark-double-complex-complex-of-complexes-first}, le foncteur
$$
\text{Tot} :
K_{first}(\text{DoubleComp}(\mathcal{A}))
\longrightarrow
K(\mathcal{A})
$$
est un foncteur exact de catégories triangulées.
\item On peut considérer un objet $A^{\bullet, \bullet}$ de
$\text{DoubleComp}(\mathcal{A})$ comme un complexe de complexes de la manière
suivante :
$$
\ldots \to A^{-1, \bullet} \to A^{0, \bullet} \to A^{1, \bullet} \to \ldots
$$
et prendre la catégorie d'homotopie
$K_{second}(\text{DoubleComp}(\mathcal{A}))$, munie de la structure triangulée
correspondante donnée par la
proposition \ref{proposition-homotopy-category-triangulated}.
D'après Homologie, remarque
\ref{homology-remark-double-complex-complex-of-complexes-second}, le foncteur
$$
\text{Tot} :
K_{second}(\text{DoubleComp}(\mathcal{A}))
\longrightarrow
K(\mathcal{A})
$$
est un foncteur exact de catégories triangulées.
\end{enumerate}
\end{remark}

\begin{remark}
\label{remark-double-complex-as-tensor-product-of}
Soient $\mathcal{A}$, $\mathcal{B}$, $\mathcal{C}$ des catégories additives et
supposons que $\mathcal{C}$ admette les sommes directes dénombrables. Supposons
que
$$
\otimes : \mathcal{A} \times \mathcal{B} \longrightarrow \mathcal{C},
\quad
(X, Y) \longmapsto X \otimes Y
$$
soit un foncteur bilinéaire sur les morphismes. Il détermine un foncteur
$$
\text{Comp}(\mathcal{A}) \times \text{Comp}(\mathcal{B})
\longrightarrow
\text{DoubleComp}(\mathcal{C}), \quad
(X^\bullet, Y^\bullet)
\longmapsto
X^\bullet \otimes Y^\bullet
$$
Voir Homologie, exemple
\ref{homology-example-double-complex-as-tensor-product-of}.
\begin{enumerate}
\item Pour un objet fixé $X^\bullet$ de $\text{Comp}(\mathcal{A})$, le foncteur
$$
K(\mathcal{B}) \longrightarrow K(\mathcal{C}), \quad
Y^\bullet \longmapsto \text{Tot}(X^\bullet \otimes Y^\bullet)
$$
est un foncteur exact de catégories triangulées.
\item Pour un objet fixé $Y^\bullet$ de $\text{Comp}(\mathcal{B})$, le foncteur
$$
K(\mathcal{A}) \longrightarrow K(\mathcal{C}), \quad
X^\bullet \longmapsto \text{Tot}(X^\bullet \otimes Y^\bullet)
$$
est un foncteur exact de catégories triangulées.
\end{enumerate}
Cela résulte de la remarque \ref{remark-homotopy-double}, puisque les foncteurs
$\text{Comp}(\mathcal{A}) \to \text{DoubleComp}(\mathcal{C})$,
$Y^\bullet \mapsto X^\bullet \otimes Y^\bullet$ et
$\text{Comp}(\mathcal{B}) \to \text{DoubleComp}(\mathcal{C})$,
$X^\bullet \mapsto X^\bullet \otimes Y^\bullet$ sont manifestement
compatibles aux homotopies et aux suites exactes courtes scindées terme à terme
et induisent donc des foncteurs exacts de catégories triangulées
$$
K(\mathcal{B}) \to K_{first}(\text{DoubleComp}(\mathcal{C}))
\quad\text{et}\quad
K(\mathcal{A}) \to K_{second}(\text{DoubleComp}(\mathcal{C}))
$$
Observons que, pour le premier des deux, l'isomorphisme
$$
\text{Tot}(X^\bullet \otimes Y^\bullet[1]) \cong
\text{Tot}(X^\bullet \otimes Y^\bullet)[1]
$$
fait intervenir des signes (cela provient des signes choisis dans Homologie,
remarque \ref{homology-remark-shift-double-complex}).
\end{remark}


\section{Catégories dérivées}
\label{section-derived-categories}

\noindent
Dans cette section, nous construisons la catégorie dérivée d'une catégorie
abélienne $\mathcal{A}$ en inversant les quasi-isomorphismes dans
$K(\mathcal{A})$. Avant cela, rappelons que les foncteurs
$H^i : \text{Comp}(\mathcal{A}) \to \mathcal{A}$ se factorisent par
$K(\mathcal{A})$, voir Homologie, lemme
\ref{homology-lemma-map-cohomology-homotopy-cochain}. En outre, dans Homologie,
définition \ref{homology-definition-cohomology-shift}, nous avons défini des
identifications $H^i(K^\bullet[n]) = H^{i + n}(K^\bullet)$. Il est désormais
judicieux de redéfinir
$$
H^i(K^\bullet) = H^0(K^\bullet[i])
$$
afin d'éviter toute confusion et d'éventuelles erreurs de signe.

\begin{lemma}
\label{lemma-cohomology-homological}
Soit $\mathcal{A}$ une catégorie abélienne. Le foncteur
$$
H^0 : K(\mathcal{A}) \longrightarrow \mathcal{A}
$$
est homologique.
\end{lemma}

\begin{proof}
Comme $H^0$ est un foncteur, la définition de nos triangles distingués montre
qu'il suffit de prouver que, pour toute suite exacte courte de complexes
scindée terme à terme
$0 \to A^\bullet \to B^\bullet \to C^\bullet \to 0$, la suite
$H^0(A^\bullet) \to H^0(B^\bullet) \to H^0(C^\bullet)$ est exacte. Cela
résulte de Homologie, lemme
\ref{homology-lemma-long-exact-sequence-cochain}.
\end{proof}

\noindent
En particulier, ce lemme implique qu'un triangle distingué
$(X, Y, Z, f, g, h)$ de $K(\mathcal{A})$ donne naissance à une suite exacte
longue de cohomologie
\begin{equation}
\label{equation-long-exact-cohomology-sequence-D}
\xymatrix{
\ldots \ar[r] &
H^i(X) \ar[r]^{H^i(f)} &
H^i(Y) \ar[r]^{H^i(g)} &
H^i(Z) \ar[r]^{H^i(h)} &
H^{i + 1}(X) \ar[r] & \ldots
}
\end{equation}
voir (\ref{equation-long-exact-cohomology-sequence}). De plus, cette suite est
compatible avec la suite exacte longue de cohomologie associée à une suite
exacte courte de complexes. Par exemple, si
$(A^\bullet, B^\bullet, C^\bullet, \alpha, \beta, \delta)$ est le triangle
distingué associé à une suite exacte de complexes scindée terme à terme (voir
la définition \ref{definition-split-ses}), alors la suite de cohomologie
ci-dessus coïncide avec celle définie au moyen du lemme du serpent, voir
Homologie, lemme \ref{homology-lemma-long-exact-sequence-cochain}; pour la
coïncidence des suites, voir Homologie, lemme
\ref{homology-lemma-ses-termwise-split-long-cochain}.

\medskip\noindent
Rappelons qu'un complexe $K^\bullet$ est {\it acyclique} si
$H^i(K^\bullet) = 0$ pour tout $i \in \mathbf{Z}$. Rappelons aussi qu'un
morphisme de complexes $f : K^\bullet \to L^\bullet$ est un
{\it quasi-isomorphisme} si et seulement si $H^i(f)$ est un isomorphisme pour
tout $i$. Voir Homologie, définition
\ref{homology-definition-quasi-isomorphism-cochain}.

\begin{lemma}
\label{lemma-acyclic}
Soit $\mathcal{A}$ une catégorie abélienne. La sous-catégorie pleine
$\text{Ac}(\mathcal{A})$ de $K(\mathcal{A})$ formée des complexes acycliques
est une sous-catégorie triangulée strictement pleine, saturée, de
$K(\mathcal{A})$. Le système multiplicatif saturé correspondant (voir le
lemme \ref{lemma-operations}) de $K(\mathcal{A})$ est l'ensemble
$\text{Qis}(\mathcal{A})$ des quasi-isomorphismes. En particulier, le noyau du
foncteur de localisation
$Q : K(\mathcal{A}) \to \text{Qis}(\mathcal{A})^{-1}K(\mathcal{A})$ est
$\text{Ac}(\mathcal{A})$, et le foncteur $H^0$ se factorise par $Q$.
\end{lemma}

\begin{proof}
Le lemme \ref{lemma-cohomology-homological} montre que $H^0$ est un foncteur
homologique. Le présent lemme est donc un cas particulier du
lemme \ref{lemma-acyclic-general}.
\end{proof}

\begin{definition}
\label{definition-unbounded-derived-category}
Soit $\mathcal{A}$ une catégorie abélienne. Soient
$\text{Ac}(\mathcal{A})$ et $\text{Qis}(\mathcal{A})$ comme dans le
lemme \ref{lemma-acyclic}. La {\it catégorie dérivée de $\mathcal{A}$} est la
catégorie triangulée
$$
D(\mathcal{A}) =
K(\mathcal{A})/\text{Ac}(\mathcal{A}) =
\text{Qis}(\mathcal{A})^{-1} K(\mathcal{A}).
$$
Notons $H^0 : D(\mathcal{A}) \to \mathcal{A}$ l'unique foncteur dont la
composée avec le foncteur quotient redonne le foncteur $H^0$ défini ci-dessus.
À l'aide du lemme \ref{lemma-homological-functor-bounded}, introduisons les
sous-catégories triangulées strictement pleines et saturées
$D^{+}(\mathcal{A}), D^{-}(\mathcal{A}), D^b(\mathcal{A})$, dont les ensembles
d'objets sont
$$
\begin{matrix}
\Ob(D^{+}(\mathcal{A})) =
\{X \in \Ob(D(\mathcal{A})) \mid
H^n(X) = 0\text{ pour tout }n \ll 0\} \\
\Ob(D^{-}(\mathcal{A})) =
\{X \in \Ob(D(\mathcal{A})) \mid
H^n(X) = 0\text{ pour tout }n \gg 0\} \\
\Ob(D^b(\mathcal{A})) =
\{X \in \Ob(D(\mathcal{A})) \mid
H^n(X) = 0\text{ pour tout }|n| \gg 0\}
\end{matrix}
$$
La catégorie $D^b(\mathcal{A})$ est appelée la {\it catégorie dérivée bornée}
de $\mathcal{A}$.
\end{definition}

\noindent
Si $K^\bullet$ et $L^\bullet$ sont des complexes de $\mathcal{A}$, nous disons
parfois que « $K^\bullet$ est {\it quasi-isomorphe} à $L^\bullet$ » pour
indiquer que $K^\bullet$ et $L^\bullet$ sont des objets isomorphes de
$D(\mathcal{A})$.

\begin{remark}
\label{remark-existence-derived}
Dans ce chapitre, nous travaillons systématiquement avec des catégories
abéliennes « petites » (comme c'est la convention dans le projet Champs). Pour
une « grande » catégorie abélienne $\mathcal{A}$, l'existence de la catégorie
dérivée $D(\mathcal{A})$ n'est pas évidente, car il n'est pas clair que ses
morphismes forment des ensembles. En fait, ce n'est généralement pas le cas,
voir Exemples, lemme \ref{examples-lemma-big-abelian-category}. Cependant, si
$\mathcal{A}$ est une catégorie abélienne de Grothendieck et si
$K^\bullet, L^\bullet$ appartiennent à $K(\mathcal{A})$, alors, d'après
Injectifs, théorème
\ref{injectives-theorem-K-injective-embedding-grothendieck}, il existe un
quasi-isomorphisme $L^\bullet \to I^\bullet$ vers un complexe K-injectif
$I^\bullet$, et le lemme \ref{lemma-K-injective} montre que
$$
\Hom_{D(\mathcal{A})}(K^\bullet, L^\bullet) =
\Hom_{K(\mathcal{A})}(K^\bullet, I^\bullet)
$$
est un ensemble. La catégorie des modules sur un anneau, ou plus généralement
la catégorie des faisceaux de modules sur un site annelé, sont des exemples de
catégories abéliennes de Grothendieck.
\end{remark}

\noindent
Chacune des variantes
$D^{+}(\mathcal{A}), D^{-}(\mathcal{A}), D^b(\mathcal{A})$ peut être construite
comme une localisation de la catégorie d'homotopie correspondante. Cela repose
sur le lemme élémentaire suivant.

\begin{lemma}
\label{lemma-complex-cohomology-bounded}
Soit $\mathcal{A}$ une catégorie abélienne. Soit $K^\bullet$ un complexe.
\begin{enumerate}
\item Si $H^n(K^\bullet) = 0$ pour tout $n \ll 0$, il existe un
quasi-isomorphisme $K^\bullet \to L^\bullet$, avec $L^\bullet$ borné
inférieurement.
\item Si $H^n(K^\bullet) = 0$ pour tout $n \gg 0$, il existe un
quasi-isomorphisme $M^\bullet \to K^\bullet$, avec $M^\bullet$ borné
supérieurement.
\item Si $H^n(K^\bullet) = 0$ pour tout $|n| \gg 0$, il existe un diagramme
commutatif de morphismes de complexes
$$
\xymatrix{
K^\bullet \ar[r] & L^\bullet \\
M^\bullet \ar[u] \ar[r] & N^\bullet \ar[u]
}
$$
où toutes les flèches sont des quasi-isomorphismes, $L^\bullet$ est borné
inférieurement, $M^\bullet$ est borné supérieurement et $N^\bullet$ est un
complexe borné.
\end{enumerate}
\end{lemma}

\begin{proof}
Choisissons $a \ll 0 \ll b$ et posons
$M^\bullet = \tau_{\leq b}K^\bullet$,
$L^\bullet = \tau_{\geq a}K^\bullet$ et
$N^\bullet = \tau_{\leq b}L^\bullet = \tau_{\geq a}M^\bullet$.
Voir Homologie, section \ref{homology-section-truncations}, pour les foncteurs
de troncation.
\end{proof}

\noindent
Pour énoncer le lemme suivant, notons
$\text{Ac}^{+}(\mathcal{A})$, $\text{Ac}^{-}(\mathcal{A})$, resp.\
$\text{Ac}^b(\mathcal{A})$, l'intersection de
$K^{+}(\mathcal{A})$, $K^{-}(\mathcal{A})$, resp.\ $K^b(\mathcal{A})$, avec
$\text{Ac}(\mathcal{A})$. Notons
$\text{Qis}^{+}(\mathcal{A})$, $\text{Qis}^{-}(\mathcal{A})$, resp.\
$\text{Qis}^b(\mathcal{A})$, l'intersection de
$K^{+}(\mathcal{A})$, $K^{-}(\mathcal{A})$, resp.\ $K^b(\mathcal{A})$, avec
$\text{Qis}(\mathcal{A})$.

\begin{lemma}
\label{lemma-bounded-derived}
Soit $\mathcal{A}$ une catégorie abélienne. Les sous-catégories
$\text{Ac}^{+}(\mathcal{A})$, $\text{Ac}^{-}(\mathcal{A})$, resp.\
$\text{Ac}^b(\mathcal{A})$, sont des sous-catégories triangulées strictement
pleines et saturées de $K^{+}(\mathcal{A})$, $K^{-}(\mathcal{A})$, resp.\
$K^b(\mathcal{A})$. Les systèmes multiplicatifs saturés correspondants (voir
le lemme \ref{lemma-operations}) sont les ensembles
$\text{Qis}^{+}(\mathcal{A})$, $\text{Qis}^{-}(\mathcal{A})$, resp.\
$\text{Qis}^b(\mathcal{A})$.
\begin{enumerate}
\item Le noyau du foncteur $K^{+}(\mathcal{A}) \to D^{+}(\mathcal{A})$ est
$\text{Ac}^{+}(\mathcal{A})$, et cela induit une équivalence de catégories
triangulées
$$
K^{+}(\mathcal{A})/\text{Ac}^{+}(\mathcal{A}) =
\text{Qis}^{+}(\mathcal{A})^{-1}K^{+}(\mathcal{A})
\longrightarrow
D^{+}(\mathcal{A})
$$
\item Le noyau du foncteur $K^{-}(\mathcal{A}) \to D^{-}(\mathcal{A})$ est
$\text{Ac}^{-}(\mathcal{A})$, et cela induit une équivalence de catégories
triangulées
$$
K^{-}(\mathcal{A})/\text{Ac}^{-}(\mathcal{A}) =
\text{Qis}^{-}(\mathcal{A})^{-1}K^{-}(\mathcal{A})
\longrightarrow
D^{-}(\mathcal{A})
$$
\item Le noyau du foncteur $K^b(\mathcal{A}) \to D^b(\mathcal{A})$ est
$\text{Ac}^b(\mathcal{A})$, et cela induit une équivalence de catégories
triangulées
$$
K^b(\mathcal{A})/\text{Ac}^b(\mathcal{A}) =
\text{Qis}^b(\mathcal{A})^{-1}K^b(\mathcal{A})
\longrightarrow
D^b(\mathcal{A})
$$
\end{enumerate}
\end{lemma}

\begin{proof}
Les premières assertions résultent du lemme \ref{lemma-acyclic-general},
appliqué à la restriction du foncteur homologique $H^0$. L'assertion sur les
noyaux dans (1), (2), (3) est, dans chaque cas, une conséquence des
définitions. Chacun des foncteurs est essentiellement surjectif d'après le
lemme \ref{lemma-complex-cohomology-bounded}. Pour achever la démonstration,
il reste à montrer que ces foncteurs sont pleinement fidèles. Commençons par la
version bornée inférieurement.

\medskip\noindent
Supposons que $K^\bullet, L^\bullet$ soient des complexes bornés
inférieurement. Un morphisme entre eux dans $D(\mathcal{A})$ est de la forme
$s^{-1}f$ pour une paire
$f : K^\bullet \to (L')^\bullet$, $s : L^\bullet \to (L')^\bullet$, où $s$
est un quasi-isomorphisme. Il en résulte que la cohomologie de $(L')^\bullet$
est bornée inférieurement. D'après le
lemme \ref{lemma-complex-cohomology-bounded}, on peut donc choisir un
quasi-isomorphisme $s' : (L')^\bullet \to (L'')^\bullet$, avec
$(L'')^\bullet$ borné inférieurement. La paire
$(s' \circ f, s' \circ s)$ définit alors un morphisme dans
$\text{Qis}^{+}(\mathcal{A})^{-1}K^{+}(\mathcal{A})$. Le foncteur est donc
« plein ». Enfin, supposons que la paire
$f : K^\bullet \to (L')^\bullet$, $s : L^\bullet \to (L')^\bullet$ définisse
un morphisme dans $\text{Qis}^{+}(\mathcal{A})^{-1}K^{+}(\mathcal{A})$ qui est
nul dans $D(\mathcal{A})$. Cela signifie qu'il existe un quasi-isomorphisme
$s' : (L')^\bullet \to (L'')^\bullet$ tel que $s' \circ f = 0$. En appliquant
encore une fois le lemme \ref{lemma-complex-cohomology-bounded}, on obtient un
quasi-isomorphisme $s'' : (L'')^\bullet \to (L''')^\bullet$, avec
$(L''')^\bullet$ borné inférieurement. On voit donc que
$s'' \circ s' \circ f = 0$, ce qui implique que $s^{-1}f$ est nul dans
$\text{Qis}^{+}(\mathcal{A})^{-1}K^{+}(\mathcal{A})$. Cela achève la preuve
que le foncteur de (1) est une équivalence.

\medskip\noindent
La démonstration de (2) est duale de celle de (1). Pour démontrer (3), nous
pouvons utiliser le résultat de (2). Il suffit donc de montrer que le foncteur
$\text{Qis}^b(\mathcal{A})^{-1}K^b(\mathcal{A})
\to \text{Qis}^{-}(\mathcal{A})^{-1}K^{-}(\mathcal{A})$ est pleinement
fidèle. L'argument du paragraphe précédent s'applique directement, en
travaillant systématiquement avec des complexes déjà bornés supérieurement.
\end{proof}








\section{Le foncteur delta canonique}
\label{section-canonical-delta-functor}

\noindent
La catégorie dérivée devrait être le réceptacle du foncteur de cohomologie
universel. Pour énoncer le résultat, nous utilisons la notion de
$\delta$-foncteur d'une catégorie abélienne vers une catégorie triangulée, voir
la définition \ref{definition-delta-functor}.

\medskip\noindent
Considérons le foncteur $\text{Comp}(\mathcal{A}) \to K(\mathcal{A})$.
Ce foncteur n'est {\bf pas}, en général, un $\delta$-foncteur. Le moyen le plus
simple de le voir consiste à considérer une suite exacte courte non scindée
$0 \to A \to B \to C \to 0$ d'objets de $\mathcal{A}$. Comme
$\Hom_{K(\mathcal{A})}(C[0], A[1]) = 0$, tout triangle distingué provenant de
cette suite exacte serait de la forme
$(A[0], B[0], C[0], a, b, 0)$. Mais l'existence d'un tel triangle distingué
dans $K(\mathcal{A})$ implique que l'extension est scindée, ce qui est une
contradiction.

\medskip\noindent
Il se trouve que le foncteur
$\text{Comp}(\mathcal{A}) \to D(\mathcal{A})$ est un $\delta$-foncteur. Pour
le voir, nous devons définir les morphismes $\delta$ associés à une suite
exacte courte
$$
0 \to A^\bullet \xrightarrow{a} B^\bullet \xrightarrow{b} C^\bullet \to 0
$$
de complexes de la catégorie abélienne $\mathcal{A}$. Considérons le cône
$C(a)^\bullet$ du morphisme $a$. On a
$C(a)^n = B^n \oplus A^{n + 1}$, et nous définissons
$q^n : C(a)^n \to C^n$ comme la projection sur $B^n$ suivie de $b^n$. On
obtient ainsi un morphisme de complexes
$$
q : C(a)^\bullet \longrightarrow C^\bullet.
$$
Il est clair que $q \circ i = b$, où $i$ est comme dans la
définition \ref{definition-cone}. Notons que, puisque $a^\bullet$ est injectif
en chaque degré, le noyau de $q$ s'identifie au cône de
$\text{id}_{A^\bullet}$, qui est acyclique. Ainsi, $q$ est un
quasi-isomorphisme. D'après le
lemme \ref{lemma-the-same-up-to-isomorphisms}, le triangle
$$
(A, B, C(a), a, i, -p)
$$
est un triangle distingué de $K(\mathcal{A})$. Comme le foncteur de
localisation $K(\mathcal{A}) \to D(\mathcal{A})$ est exact,
$(A, B, C(a), a, i, -p)$ est un triangle distingué de $D(\mathcal{A})$.
Puisque $q$ est un quasi-isomorphisme, $q$ est un isomorphisme dans
$D(\mathcal{A})$. On en déduit donc que
$$
(A, B, C, a, b, -p \circ q^{-1})
$$
est un triangle distingué de $D(\mathcal{A})$. Cela suggère le lemme suivant.

\begin{lemma}
\label{lemma-derived-canonical-delta-functor}
Soit $\mathcal{A}$ une catégorie abélienne. Le foncteur
$\text{Comp}(\mathcal{A}) \to D(\mathcal{A})$ défini ci-dessus possède une
structure naturelle de $\delta$-foncteur, avec
$$
\delta_{A^\bullet \to B^\bullet \to C^\bullet} = - p \circ q^{-1}
$$
où $p$ et $q$ sont comme ci-dessus. La même construction munit les foncteurs
$\text{Comp}^{+}(\mathcal{A}) \to D^{+}(\mathcal{A})$,
$\text{Comp}^{-}(\mathcal{A}) \to D^{-}(\mathcal{A})$ et
$\text{Comp}^b(\mathcal{A}) \to D^b(\mathcal{A})$
de structures de $\delta$-foncteurs.
\end{lemma}

\begin{proof}
Nous avons déjà vu que ce choix produit un triangle distingué pour toute suite
exacte courte de complexes. Nous devons montrer que, pour tout diagramme
commutatif
$$
\xymatrix{
0 \ar[r] &
A^\bullet \ar[r]_a \ar[d]_f &
B^\bullet \ar[r]_b \ar[d]_g &
C^\bullet \ar[r] \ar[d]_h &
0 \\
0 \ar[r] &
(A')^\bullet \ar[r]^{a'} &
(B')^\bullet \ar[r]^{b'} &
(C')^\bullet \ar[r] &
0
}
$$
nous obtenons le diagramme commutatif voulu de la
définition \ref{definition-delta-functor} (2). D'après le
lemme \ref{lemma-functorial-cone}, la paire $(f, g)$ induit un morphisme
canonique $c : C(a)^\bullet \to C(a')^\bullet$. Un calcul simple montre que
$q' \circ c = h \circ q$ et $f[1] \circ p = p' \circ c$. Le résultat en
découle immédiatement.
\end{proof}

\begin{lemma}
\label{lemma-derived-compare-triangles-ses}
Soit $\mathcal{A}$ une catégorie abélienne. Soit
$$
\xymatrix{
0 \ar[r] &
A^\bullet \ar[r] \ar[d] &
B^\bullet \ar[r] \ar[d] &
C^\bullet \ar[r] \ar[d] &
0 \\
0 \ar[r] &
D^\bullet \ar[r] &
E^\bullet \ar[r] &
F^\bullet \ar[r] &
0
}
$$
un diagramme commutatif de morphismes de complexes tel que ses lignes soient
des suites exactes courtes de complexes et que ses flèches verticales soient
des quasi-isomorphismes. Le $\delta$-foncteur du
lemme \ref{lemma-derived-canonical-delta-functor} envoie les suites exactes
courtes
$0 \to A^\bullet \to B^\bullet \to C^\bullet \to 0$ et
$0 \to D^\bullet \to E^\bullet \to F^\bullet \to 0$ sur des triangles
distingués isomorphes.
\end{lemma}

\begin{proof}
Cela résulte immédiatement du fait que
$K(\mathcal{A}) \to D(\mathcal{A})$ transforme les quasi-isomorphismes en
isomorphismes et que les triangles distingués associés sont fonctoriels.
\end{proof}

\begin{lemma}
\label{lemma-derived-compare-triangles-split-case}
Soit $\mathcal{A}$ une catégorie abélienne. Soit
$$
\xymatrix{
0 \ar[r] &
A^\bullet \ar[r] &
B^\bullet \ar[r] &
C^\bullet \ar[r] &
0
}
$$
une suite exacte courte de complexes. Supposons que cette suite exacte courte
soit scindée terme à terme. Soit
$(A^\bullet, B^\bullet, C^\bullet, \alpha, \beta, \delta)$ le triangle
distingué de $K(\mathcal{A})$ associé à cette suite. Le $\delta$-foncteur du
lemme \ref{lemma-derived-canonical-delta-functor} envoie la suite exacte courte
$0 \to A^\bullet \to B^\bullet \to C^\bullet \to 0$ sur un triangle isomorphe
au triangle distingué
$$
(A^\bullet, B^\bullet, C^\bullet, \alpha, \beta, \delta).
$$
\end{lemma}

\begin{proof}
Cela résulte du lemme \ref{lemma-the-same-up-to-isomorphisms}.
\end{proof}

\begin{remark}
\label{remark-truncation-distinguished-triangle}
Soit $\mathcal{A}$ une catégorie abélienne. Soit $K^\bullet$ un complexe de
$\mathcal{A}$ et soit $a \in \mathbf{Z}$. Nous affirmons qu'il existe un
triangle distingué canonique
$$
\tau_{\leq a}K^\bullet \to K^\bullet \to \tau_{\geq a + 1}K^\bullet \to
(\tau_{\leq a}K^\bullet)[1]
$$
dans $D(\mathcal{A})$. Nous avons utilisé ici les foncteurs de troncation canoniques $\tau$
d'Homologie, section \ref{homology-section-truncations}. Prenons d'abord le
triangle distingué associé par notre $\delta$-foncteur
(lemme \ref{lemma-derived-canonical-delta-functor}) à la suite exacte courte de
complexes
$$
0 \to \tau_{\leq a}K^\bullet \to K^\bullet \to
K^\bullet/\tau_{\leq a}K^\bullet \to 0
$$
Ensuite, la description des groupes de cohomologie dans Homologie, section
\ref{homology-section-truncations}, montre que le morphisme
$K^\bullet \to \tau_{\geq a + 1}K^\bullet$ se factorise par un
quasi-isomorphisme
$K^\bullet/\tau_{\leq a}K^\bullet \to \tau_{\geq a + 1}K^\bullet$.
De façon analogue, on obtient les triangles distingués canoniques
$$
\tau_{\leq a}K^\bullet \to \tau_{\leq a + 1}K^\bullet \to
H^{a + 1}(K^\bullet)[-a-1] \to (\tau_{\leq a}K^\bullet)[1]
$$
et
$$
H^a(K^\bullet)[-a] \to \tau_{\geq a}K^\bullet \to
\tau_{\geq a + 1}K^\bullet \to H^a(K^\bullet)[-a + 1]
$$
\end{remark}

\begin{lemma}
\label{lemma-trick-vanishing-composition}
Soit $\mathcal{A}$ une catégorie abélienne. Soient
$$
K_0^\bullet \to K_1^\bullet \to \ldots \to K_n^\bullet
$$
des morphismes de complexes tels que
\begin{enumerate}
\item $H^i(K_0^\bullet) = 0$ pour $i > 0$,
\item $H^{-j}(K_j^\bullet) \to H^{-j}(K_{j + 1}^\bullet)$ soit nul.
\end{enumerate}
Alors la composée $K_0^\bullet \to K_n^\bullet$ se factorise par
$\tau_{\leq -n}K_n^\bullet \to K_n^\bullet$ dans $D(\mathcal{A})$.
Dualement, soient des morphismes
de complexes
$$
K_n^\bullet \to K_{n - 1}^\bullet \to \ldots \to K_0^\bullet
$$
tels que
\begin{enumerate}
\item $H^i(K_0^\bullet) = 0$ pour $i < 0$,
\item $H^j(K_{j + 1}^\bullet) \to H^j(K_j^\bullet)$ soit nul.
\end{enumerate}
Alors la composée $K_n^\bullet \to K_0^\bullet$ se factorise par
$K_n^\bullet \to \tau_{\geq n}K_n^\bullet$ dans $D(\mathcal{A})$.
\end{lemma}

\begin{proof}
Traitons le cas $n = 1$. Comme
$\tau_{\leq 0}K_0^\bullet = K_0^\bullet$ dans $D(\mathcal{A})$, on peut
remplacer $K_0^\bullet$ par $\tau_{\leq 0}K_0^\bullet$ et $K_1^\bullet$ par
$\tau_{\leq 0}K_1^\bullet$. Considérons le triangle distingué
$$
\tau_{\leq -1}K_1^\bullet \to K_1^\bullet \to
H^0(K_1^\bullet)[0] \to (\tau_{\leq -1}K_1^\bullet)[1]
$$
(remarque \ref{remark-truncation-distinguished-triangle}). La composée
$K_0^\bullet \to K_1^\bullet \to H^0(K_1^\bullet)[0]$ est nulle, car elle est
égale à
$K_0^\bullet \to H^0(K_0^\bullet)[0] \to H^0(K_1^\bullet)[0]$, qui est nulle
par hypothèse. Le fait que
$\Hom_{D(\mathcal{A})}(K_0^\bullet, -)$ soit un foncteur homologique
(lemme \ref{lemma-representable-homological}) permet de trouver la
factorisation voulue. Pour $n = 2$, nous obtenons une factorisation
$K_0^\bullet \to \tau_{\leq -1}K_1^\bullet$ par le cas $n = 1$, et nous
pouvons appliquer le cas $n = 1$ au morphisme de complexes
$\tau_{\leq -1}K_1^\bullet \to \tau_{\leq -1}K_2^\bullet$ pour obtenir une
factorisation
$\tau_{\leq -1}K_1^\bullet \to \tau_{\leq -2}K_2^\bullet$.
Le cas général se démontre exactement de la même manière.
\end{proof}




\section{Catégories dérivées filtrées}
\label{section-filtered-derived-category}

\noindent
Une référence pour cette section est \cite[I, chapitre V]{cotangent}. Soit
$\mathcal{A}$ une catégorie abélienne. Dans cette section, nous allons définir
la catégorie dérivée filtrée $DF(\mathcal{A})$ de $\mathcal{A}$. En bref, nous
la définirons comme la catégorie dérivée de la catégorie exacte des objets de
$\mathcal{A}$ munis d'une filtration finie. (Notre construction est donc un
cas particulier d'une construction plus générale de la catégorie dérivée
d'une catégorie exacte ; voir, par exemple, \cite{Buhler}, \cite{Keller}.) La
catégorie dérivée filtrée d'Illusie est la sous-catégorie pleine de la nôtre
formée des objets dont la filtration est finie. (Dans notre catégorie, la
filtration est encore finie en chaque degré, mais peut ne pas être uniformément
bornée.) La raison de notre choix est que cette construction n'est pas plus
difficile et nous permet d'appliquer la discussion aux suites spectrales du
lemme \ref{lemma-two-ss-complex-functor} ; voir aussi la
remarque \ref{remark-functorial-ss}.

\medskip\noindent
Nous emploierons les notations relatives aux objets filtrés introduites dans
Homologie, section \ref{homology-section-filtrations}. La catégorie des objets
filtrés de $\mathcal{A}$ est notée $\text{Fil}(\mathcal{A})$. Toutes les
filtrations seront, par convention, décroissantes.

\begin{definition}
\label{definition-finite-filtered}
Soit $\mathcal{A}$ une catégorie abélienne. La
{\it catégorie des objets filtrés finis de $\mathcal{A}$} est la catégorie des
objets filtrés $(A, F)$ de $\mathcal{A}$ dont la filtration $F$ est finie. Nous
la notons $\text{Fil}^f(\mathcal{A})$.
\end{definition}

\noindent
Ainsi, $\text{Fil}^f(\mathcal{A})$ est une sous-catégorie pleine de
$\text{Fil}(\mathcal{A})$. Pour tout $p \in \mathbf{Z}$, il existe un foncteur
$\text{gr}^p : \text{Fil}^f(\mathcal{A}) \to \mathcal{A}$. Il existe un
foncteur
$$
\text{gr} = \bigoplus\nolimits_{p \in \mathbf{Z}} \text{gr}^p :
\text{Fil}^f(\mathcal{A}) \to \text{Gr}(\mathcal{A})
$$
où $\text{Gr}(\mathcal{A})$ est la catégorie des objets gradués de
$\mathcal{A}$ ; voir Homologie, définition
\ref{homology-definition-graded}. Enfin, il existe un foncteur
$$
(\text{oubli de }F) : \text{Fil}^f(\mathcal{A}) \longrightarrow \mathcal{A}
$$
qui associe à l'objet filtré $(A, F)$ l'objet sous-jacent de $\mathcal{A}$.
La catégorie $\text{Fil}^f(\mathcal{A})$ est additive, mais elle n'est pas
abélienne en général ; voir Homologie, exemple
\ref{homology-example-not-abelian}.

\medskip\noindent
Comme les foncteurs $\text{gr}^p$, $\text{gr}$ et
$(\text{oubli de }F)$ sont additifs, ils induisent des foncteurs exacts de
catégories triangulées
$$
\text{gr}^p, (\text{oubli de }F) :
K(\text{Fil}^f(\mathcal{A}))
\to
K(\mathcal{A})
\quad\text{et}\quad
\text{gr} :
K(\text{Fil}^f(\mathcal{A}))
\to
K(\text{Gr}(\mathcal{A}))
$$
d'après le lemme \ref{lemma-additive-exact-homotopy-category}. Par analogie
avec le cas de la catégorie d'homotopie d'une catégorie abélienne, posons les
définitions suivantes.

\begin{definition}
\label{definition-filtered-acyclic}
Soit $\mathcal{A}$ une catégorie abélienne.
\begin{enumerate}
\item Soit $\alpha : K^\bullet \to L^\bullet$ un morphisme de
$K(\text{Fil}^f(\mathcal{A}))$. Nous disons que $\alpha$ est un
{\it quasi-isomorphisme filtré} si le morphisme $\text{gr}(\alpha)$ est un
quasi-isomorphisme.
\item Soit $K^\bullet$ un objet de $K(\text{Fil}^f(\mathcal{A}))$. Nous disons
que $K^\bullet$ est {\it filtré acyclique} si le complexe
$\text{gr}(K^\bullet)$ est acyclique.
\end{enumerate}
\end{definition}

\noindent
Notons que $\alpha : K^\bullet \to L^\bullet$ est un quasi-isomorphisme filtré
si et seulement si chaque $\text{gr}^p(\alpha)$ est un quasi-isomorphisme. De
même, un complexe $K^\bullet$ est filtré acyclique si et seulement si chaque
$\text{gr}^p(K^\bullet)$ est acyclique.

\begin{lemma}
\label{lemma-filtered-cohomology-homological}
Soit $\mathcal{A}$ une catégorie abélienne.
\begin{enumerate}
\item Le foncteur
$K(\text{Fil}^f(\mathcal{A})) \longrightarrow \text{Gr}(\mathcal{A})$,
$K^\bullet \longmapsto H^0(\text{gr}(K^\bullet))$ est homologique.
\item Le foncteur
$K(\text{Fil}^f(\mathcal{A})) \rightarrow \mathcal{A}$,
$K^\bullet \longmapsto H^0(\text{gr}^p(K^\bullet))$ est homologique.
\item Le foncteur
$K(\text{Fil}^f(\mathcal{A})) \longrightarrow \mathcal{A}$,
$K^\bullet \longmapsto H^0((\text{oubli de }F)K^\bullet)$ est homologique.
\end{enumerate}
\end{lemma}

\begin{proof}
Cela résulte du fait que $H^0 : K(\mathcal{A}) \to \mathcal{A}$ est
homologique, voir le lemme \ref{lemma-cohomology-homological}, et du fait que
les foncteurs $\text{gr}, \text{gr}^p, (\text{oubli de }F)$ sont des
foncteurs exacts de catégories triangulées. Voir le
lemme \ref{lemma-exact-compose-homological-functor}.
\end{proof}

\begin{lemma}
\label{lemma-filtered-acyclic}
Soit $\mathcal{A}$ une catégorie abélienne. La sous-catégorie pleine
$\text{FAc}(\mathcal{A})$ de $K(\text{Fil}^f(\mathcal{A}))$, formée des
complexes filtrés acycliques, est une sous-catégorie triangulée strictement
pleine et saturée de $K(\text{Fil}^f(\mathcal{A}))$. Le système multiplicatif
saturé correspondant (voir le lemme \ref{lemma-operations}) de
$K(\text{Fil}^f(\mathcal{A}))$ est l'ensemble $\text{FQis}(\mathcal{A})$ des
quasi-isomorphismes filtrés. En particulier, le noyau du foncteur de
localisation
$$
Q :
K(\text{Fil}^f(\mathcal{A}))
\longrightarrow
\text{FQis}(\mathcal{A})^{-1}K(\text{Fil}^f(\mathcal{A}))
$$
est $\text{FAc}(\mathcal{A})$, et le foncteur
$H^0 \circ \text{gr}$ se factorise par $Q$.
\end{lemma}

\begin{proof}
Le lemme \ref{lemma-filtered-cohomology-homological} montre que
$H^0 \circ \text{gr}$ est un foncteur homologique. Le présent lemme est donc un
cas particulier du lemme \ref{lemma-acyclic-general}.
\end{proof}

\begin{definition}
\label{definition-filtered-derived}
Soit $\mathcal{A}$ une catégorie abélienne. Soient
$\text{FAc}(\mathcal{A})$ et $\text{FQis}(\mathcal{A})$ comme dans le
lemme \ref{lemma-filtered-acyclic}. La
{\it catégorie dérivée filtrée de $\mathcal{A}$} est la catégorie triangulée
$$
DF(\mathcal{A}) =
K(\text{Fil}^f(\mathcal{A}))/\text{FAc}(\mathcal{A}) =
\text{FQis}(\mathcal{A})^{-1} K(\text{Fil}^f(\mathcal{A})).
$$
\end{definition}

\begin{lemma}
\label{lemma-filtered-derived-functors}
Les foncteurs $\text{gr}^p, \text{gr}, (\text{oubli de }F)$ induisent
des foncteurs exacts canoniques
$$
\text{gr}^p, (\text{oubli de }F):
DF(\mathcal{A})
\longrightarrow
D(\mathcal{A})
$$
et
$$
\text{gr}:
DF(\mathcal{A})
\longrightarrow
D(\text{Gr}(\mathcal{A}))
$$
qui commutent aux foncteurs de localisation.
\end{lemma}

\begin{proof}
Cela résulte de la propriété universelle de la localisation, voir le
lemme \ref{lemma-universal-property-localization}, pourvu que chacun des
foncteurs $\text{gr}^p, \text{gr}, (\text{oubli de }F)$ transforme un
quasi-isomorphisme filtré en un quasi-isomorphisme. C'est vrai par définition
pour les deux premiers. Pour le dernier, il faut fournir un petit argument.
Soit $f : K^\bullet \to L^\bullet$ un quasi-isomorphisme filtré dans
$K(\text{Fil}^f(\mathcal{A}))$. Choisissons un triangle distingué
$(K^\bullet, L^\bullet, M^\bullet, f, g, h)$ contenant $f$. Alors $M^\bullet$
est filtré acyclique, voir le lemme \ref{lemma-filtered-acyclic}. Le lemme
correspondant pour $K(\mathcal{A})$ montre donc qu'il suffit de prouver qu'un
complexe filtré acyclique est un complexe acyclique après oubli de la
filtration. Cela résulte d'Homologie, lemme
\ref{homology-lemma-filtered-acyclic}.
\end{proof}

\begin{definition}
\label{definition-filtered-derived-bounded}
Soit $\mathcal{A}$ une catégorie abélienne. La
{\it catégorie dérivée filtrée bornée} $DF^b(\mathcal{A})$ est la
sous-catégorie pleine de $DF(\mathcal{A})$ dont les objets sont les $X$ tels
que $\text{gr}(X) \in D^b(\mathcal{A})$. On définit de même la catégorie
dérivée filtrée bornée inférieurement $DF^{+}(\mathcal{A})$ et la catégorie
dérivée filtrée bornée supérieurement $DF^{-}(\mathcal{A})$.
\end{definition}

\begin{lemma}
\label{lemma-filtered-complex-cohomology-bounded}
Soit $\mathcal{A}$ une catégorie abélienne. Soit
$K^\bullet \in K(\text{Fil}^f(\mathcal{A}))$.
\begin{enumerate}
\item Si $H^n(\text{gr}(K^\bullet)) = 0$ pour tout $n < a$, il existe un
quasi-isomorphisme filtré $K^\bullet \to L^\bullet$ tel que $L^n = 0$ pour
tout $n < a$.
\item Si $H^n(\text{gr}(K^\bullet)) = 0$ pour tout $n > b$, il existe un
quasi-isomorphisme filtré $M^\bullet \to K^\bullet$ tel que $M^n = 0$ pour
tout $n > b$.
\item Si $H^n(\text{gr}(K^\bullet)) = 0$ pour tout $|n| \gg 0$, il existe un
diagramme commutatif de morphismes de complexes
$$
\xymatrix{
K^\bullet \ar[r] & L^\bullet \\
M^\bullet \ar[u] \ar[r] & N^\bullet \ar[u]
}
$$
où toutes les flèches sont des quasi-isomorphismes filtrés, $L^\bullet$ est
borné inférieurement, $M^\bullet$ est borné supérieurement et $N^\bullet$ est
un complexe borné.
\end{enumerate}
\end{lemma}

\begin{proof}
Supposons que $H^n(\text{gr}(K^\bullet)) = 0$ pour tout $n < a$. D'après
Homologie, lemme \ref{homology-lemma-filtered-acyclic}, la suite
$$
K^{a - 2} \xrightarrow{d^{a - 2}} K^{a - 1} \xrightarrow{d^{a - 1}} K^a
$$
est une suite exacte d'objets de $\mathcal{A}$, et les morphismes
$d^{a - 2}$ et $d^{a - 1}$ sont stricts. Ainsi,
$\Coim(d^{a - 1}) = \Im(d^{a - 1})$ dans $\text{Fil}^f(\mathcal{A})$, et le
morphisme $\text{gr}(\Im(d^{a - 1})) \to \text{gr}(K^a)$ est injectif,
d'image égale à celle de
$\text{gr}(K^{a - 1}) \to \text{gr}(K^a)$ ; voir Homologie, lemme
\ref{homology-lemma-characterize-strict}. Cela signifie que le morphisme
$K^\bullet \to \tau_{\geq a}K^\bullet$ vers le complexe tronqué
$$
\tau_{\geq a}K^\bullet =
(\ldots \to 0 \to K^a/\Im(d^{a - 1}) \to K^{a + 1} \to \ldots)
$$
est un quasi-isomorphisme filtré. Cela démontre (1). La démonstration de (2)
est duale de celle de (1). La partie (3) résulte formellement de (1) et (2).
\end{proof}

\noindent
Pour énoncer le lemme suivant, notons
$\text{FAc}^{+}(\mathcal{A})$, $\text{FAc}^{-}(\mathcal{A})$, resp.\
$\text{FAc}^b(\mathcal{A})$, l'intersection de
$K^{+}(\text{Fil}^f\mathcal{A})$, $K^{-}(\text{Fil}^f\mathcal{A})$, resp.\
$K^b(\text{Fil}^f\mathcal{A})$, avec $\text{FAc}(\mathcal{A})$. Notons
$\text{FQis}^{+}(\mathcal{A})$, $\text{FQis}^{-}(\mathcal{A})$, resp.\
$\text{FQis}^b(\mathcal{A})$, l'intersection de
$K^{+}(\text{Fil}^f\mathcal{A})$, $K^{-}(\text{Fil}^f\mathcal{A})$, resp.\
$K^b(\text{Fil}^f\mathcal{A})$, avec $\text{FQis}(\mathcal{A})$.

\begin{lemma}
\label{lemma-filtered-bounded-derived}
Soit $\mathcal{A}$ une catégorie abélienne. Les sous-catégories
$\text{FAc}^{+}(\mathcal{A})$, $\text{FAc}^{-}(\mathcal{A})$, resp.\
$\text{FAc}^b(\mathcal{A})$, sont des sous-catégories triangulées strictement
pleines et saturées de $K^{+}(\text{Fil}^f\mathcal{A})$,
$K^{-}(\text{Fil}^f\mathcal{A})$, resp.\ $K^b(\text{Fil}^f\mathcal{A})$.
Les systèmes multiplicatifs saturés correspondants (voir le
lemme \ref{lemma-operations}) sont les ensembles
$\text{FQis}^{+}(\mathcal{A})$, $\text{FQis}^{-}(\mathcal{A})$, resp.\
$\text{FQis}^b(\mathcal{A})$.
\begin{enumerate}
\item Le noyau du foncteur
$K^{+}(\text{Fil}^f\mathcal{A}) \to DF^{+}(\mathcal{A})$ est
$\text{FAc}^{+}(\mathcal{A})$, et cela induit une équivalence de catégories
triangulées
$$
K^{+}(\text{Fil}^f\mathcal{A})/\text{FAc}^{+}(\mathcal{A}) =
\text{FQis}^{+}(\mathcal{A})^{-1}K^{+}(\text{Fil}^f\mathcal{A})
\longrightarrow
DF^{+}(\mathcal{A})
$$
\item Le noyau du foncteur
$K^{-}(\text{Fil}^f\mathcal{A}) \to DF^{-}(\mathcal{A})$ est
$\text{FAc}^{-}(\mathcal{A})$, et cela induit une équivalence de catégories
triangulées
$$
K^{-}(\text{Fil}^f\mathcal{A})/\text{FAc}^{-}(\mathcal{A}) =
\text{FQis}^{-}(\mathcal{A})^{-1}K^{-}(\text{Fil}^f\mathcal{A})
\longrightarrow
DF^{-}(\mathcal{A})
$$
\item Le noyau du foncteur
$K^b(\text{Fil}^f\mathcal{A}) \to DF^b(\mathcal{A})$ est
$\text{FAc}^b(\mathcal{A})$, et cela induit une équivalence de catégories
triangulées
$$
K^b(\text{Fil}^f\mathcal{A})/\text{FAc}^b(\mathcal{A}) =
\text{FQis}^b(\mathcal{A})^{-1}K^b(\text{Fil}^f\mathcal{A})
\longrightarrow
DF^b(\mathcal{A})
$$
\end{enumerate}
\end{lemma}

\begin{proof}
Cela résulte des résultats précédents, en particulier du
lemme \ref{lemma-filtered-complex-cohomology-bounded}, exactement par les mêmes
arguments que ceux de la démonstration du lemme \ref{lemma-bounded-derived}.
\end{proof}











\section{Foncteurs dérivés en général}
\label{section-derived-functors}

\noindent
Une référence pour cette section est l'exposé XVII de Deligne dans
\cite{SGA4}. Il existe une notion très générale de foncteurs dérivés droits et
gauches lorsque l'on dispose d'un foncteur exact entre catégories triangulées
et d'un système multiplicatif dans la catégorie source, et que l'on cherche la
« bonne » extension du foncteur exact à la catégorie localisée.

\begin{situation}
\label{situation-derived-functor}
Dans cette situation, $F : \mathcal{D} \to \mathcal{D}'$ est un foncteur exact
de catégories triangulées, et $S$ est un système multiplicatif saturé de
$\mathcal{D}$ compatible avec la structure de catégorie triangulée de
$\mathcal{D}$.
\end{situation}

\noindent
Soit $X \in \Ob(\mathcal{D})$. Rappelons, d'après Catégories, remarque
\ref{categories-remark-left-localization-morphisms-colimit}, la catégorie
filtrante $X/S$ des flèches $s : X \to X'$ de $S$ de source $X$.
Dualement, dans Catégories, remarque
\ref{categories-remark-right-localization-morphisms-colimit}, nous avons défini
la catégorie cofiltrante $S/X$ des flèches $s : X' \to X$ de $S$ de but $X$.

\begin{definition}
\label{definition-right-derived-functor-defined}
Hypothèses et notations comme dans la
situation \ref{situation-derived-functor}. Soit $X \in \Ob(\mathcal{D})$.
\begin{enumerate}
\item Nous disons que le {\it foncteur dérivé droit $RF$ est défini en} $X$
si l'ind-objet
$$
(X/S) \longrightarrow \mathcal{D}', \quad
(s : X \to X') \longmapsto F(X')
$$
est essentiellement constant\footnote{Pour une discussion des conditions sous
lesquelles un ind-objet ou un pro-objet d'une catégorie est essentiellement
constant, nous renvoyons à Catégories, section
\ref{categories-section-essentially-constant}.} ; dans ce cas, la valeur $Y$
dans $\mathcal{D}'$ est appelée la {\it valeur de $RF$ en $X$}.
\item Nous disons que le {\it foncteur dérivé gauche $LF$ est défini en} $X$
si le pro-objet
$$
(S/X) \longrightarrow \mathcal{D}', \quad
(s: X' \to X) \longmapsto F(X')
$$
est essentiellement constant ; dans ce cas, la valeur $Y$ dans
$\mathcal{D}'$ est appelée la {\it valeur de $LF$ en $X$}.
\end{enumerate}
Par abus de notation, nous notons souvent ces valeurs simplement $RF(X)$ ou
$LF(X)$.
\end{definition}

\noindent
Nous verrons que la sous-catégorie pleine de $\mathcal{D}$ formée des objets
en lesquels $RF$ est défini est une sous-catégorie triangulée, et que $RF$
définit sur celle-ci un foncteur qui transforme les morphismes de $S$ en
isomorphismes.

\begin{lemma}
\label{lemma-derived-functor}
Hypothèses et notations comme dans la
situation \ref{situation-derived-functor}. Soit $f : X \to Y$ un morphisme de
$\mathcal{D}$.
\begin{enumerate}
\item Si $RF$ est défini en $X$ et en $Y$, il existe entre les valeurs un
unique morphisme $RF(f) : RF(X) \to RF(Y)$ tel que, pour tout diagramme
commutatif
$$
\xymatrix{
X \ar[d]_f \ar[r]_s & X' \ar[d]^{f'} \\
Y \ar[r]^{s'} & Y'
}
$$
avec $s, s' \in S$, le diagramme
$$
\xymatrix{
F(X) \ar[d] \ar[r] & F(X') \ar[d] \ar[r] & RF(X) \ar[d] \\
F(Y) \ar[r] & F(Y') \ar[r] & RF(Y)
}
$$
soit commutatif.
\item Si $LF$ est défini en $X$ et en $Y$, il existe entre les valeurs un
unique morphisme $LF(f) : LF(X) \to LF(Y)$ tel que, pour tout diagramme
commutatif
$$
\xymatrix{
X' \ar[d]_{f'} \ar[r]_s & X \ar[d]^f \\
Y' \ar[r]^{s'} & Y
}
$$
avec $s, s'$ dans $S$, le diagramme
$$
\xymatrix{
LF(X) \ar[d] \ar[r] & F(X') \ar[d] \ar[r] & F(X) \ar[d] \\
LF(Y) \ar[r] & F(Y') \ar[r] & F(Y)
}
$$
soit commutatif.
\end{enumerate}
\end{lemma}

\begin{proof}
La partie (1) reste valable si l'on suppose seulement que les colimites
$$
RF(X) = \colim_{s : X \to X'} F(X')
\quad\text{et}\quad
RF(Y) = \colim_{s' : Y \to Y'} F(Y')
$$
existent. En effet, se donner un morphisme $RF(X) \to RF(Y)$ entre les
colimites revient à se donner, pour chaque
$s : X \to X'$ dans $\Ob(X/S)$, un morphisme $F(X') \to RF(Y)$ compatible aux
morphismes de la catégorie $X/S$. Pour obtenir ce morphisme, choisissons un
diagramme commutatif
$$
\xymatrix{
X \ar[d]_f \ar[r]_s & X' \ar[d]^{f'} \\
Y \ar[r]^{s'} & Y'
}
$$
avec $s, s'$ dans $S$, ce qui est possible par MS2, et prenons pour
$F(X') \to RF(Y)$ la composée $F(X') \to F(Y') \to RF(Y)$. Pour voir que ce
choix ne dépend pas du diagramme ci-dessus, on utilise MS3. Les détails sont
omis. La démonstration de (2) est duale.
\end{proof}

\begin{lemma}
\label{lemma-derived-inverts}
Hypothèses et notations comme dans la
situation \ref{situation-derived-functor}. Soit $s : X \to Y$ un élément de
$S$.
\begin{enumerate}
\item $RF$ est défini en $X$ si et seulement s'il est défini en $Y$. Dans ce
cas, le morphisme $RF(s) : RF(X) \to RF(Y)$ entre les valeurs est un
isomorphisme.
\item $LF$ est défini en $X$ si et seulement s'il est défini en $Y$. Dans ce
cas, le morphisme $LF(s) : LF(X) \to LF(Y)$ entre les valeurs est un
isomorphisme.
\end{enumerate}
\end{lemma}

\begin{proof}
Omis.
\end{proof}

\begin{lemma}
\label{lemma-derived-shift}
\begin{slogan}
Les foncteurs dérivés sont compatibles aux décalages
\end{slogan}
Hypothèses et notations comme dans la
situation \ref{situation-derived-functor}. Soient $X$ un objet de
$\mathcal{D}$ et $n \in \mathbf{Z}$.
\begin{enumerate}
\item $RF$ est défini en $X$ si et seulement s'il est défini en $X[n]$. Dans
ce cas, il existe entre les valeurs un isomorphisme canonique
$RF(X)[n]= RF(X[n])$.
\item $LF$ est défini en $X$ si et seulement s'il est défini en $X[n]$. Dans
ce cas, il existe entre les valeurs un isomorphisme canonique
$LF(X)[n] \to LF(X[n])$.
\end{enumerate}
\end{lemma}

\begin{proof}
Omis.
\end{proof}

\begin{lemma}
\label{lemma-2-out-of-3-defined}
Hypothèses et notations comme dans la
situation \ref{situation-derived-functor}. Soit $(X, Y, Z, f, g, h)$ un
triangle distingué de $\mathcal{D}$. Si $RF$ est défini en deux des trois
objets $X, Y, Z$, alors il est défini en le troisième. De plus, dans ce cas,
$$
(RF(X), RF(Y), RF(Z), RF(f), RF(g), RF(h))
$$
est un triangle distingué de $\mathcal{D}'$. De même pour $LF$.
\end{lemma}

\begin{proof}
Supposons que $RF$ soit défini en $X, Y$, de valeurs $A, B$. Soit
$RF(f) : A \to B$ le morphisme induit, voir le
lemme \ref{lemma-derived-functor}. On peut choisir un triangle distingué
$(A, B, C, RF(f), b, c)$ dans $\mathcal{D}'$. Nous affirmons que $C$ est une
valeur de $RF$ en $Z$.

\medskip\noindent
Pour le voir, choisissons $s : X \to X'$ dans $S$ tel qu'il existe un
morphisme $\alpha : A \to F(X')$ comme dans Catégories, définition
\ref{categories-definition-essentially-constant-diagram}. On peut choisir par
MS2 un diagramme commutatif
$$
\xymatrix{
X \ar[d]_f \ar[r]_s & X' \ar[d]^{f'} \\
Y \ar[r]^{s'} & Y'
}
$$
avec $s' \in S$. Comme $Y/S$ est filtrante, on peut, après avoir remplacé
$s'$ par un certain $s'' : Y \to Y''$ dans $S$, supposer qu'il existe un
morphisme $\beta : B \to F(Y')$ comme dans Catégories, définition
\ref{categories-definition-essentially-constant-diagram}. Représentons la
situation :
$$
\xymatrix{
A \ar[d]_{RF(f)} \ar[r]_-\alpha &
F(X') \ar[r] \ar[d]^{F(f')} &
A \ar[d]^{RF(f)} \\
B \ar[r]^-\beta & F(Y') \ar[r] & B
}
$$
Il se peut que le carré gauche ne soit pas commutatif, mais le rectangle
extérieur et le carré droit le sont. L'hypothèse selon laquelle l'ind-objet
$\{F(Y')\}_{s' : Y \to Y'}$ est essentiellement constant signifie qu'il
existe $s'' : Y \to Y''$ dans $S$ et un morphisme $h : Y' \to Y''$ tels que
$s'' = h \circ s'$ et que $F(h)$ soit égal à la composée
$F(Y') \to B \to F(Y') \to F(Y'')$. Ainsi, après avoir remplacé $Y'$ par
$Y''$ et $\beta$ par $F(h) \circ \beta$, le diagramme est commutatif (par un
calcul direct sur les flèches).

\medskip\noindent
À l'aide de MS6, choisissons un morphisme de triangles
$$
(s, s', s'') : (X, Y, Z, f, g, h) \longrightarrow
(X', Y', Z', f', g', h')
$$
avec $s'' \in S$. Par TR3, choisissons un morphisme de triangles
$$
(\alpha, \beta, \gamma) :
(A, B, C, RF(f), b, c)
\longrightarrow
(F(X'), F(Y'), F(Z'), F(f'), F(g'), F(h'))
$$

\medskip\noindent
D'après le lemme \ref{lemma-derived-inverts}, il suffit de montrer que
$RF(Z')$ est défini et que la flèche $\gamma : C \to F(Z')$ induit un
isomorphisme $C \to RF(Z')$. En effet, on obtiendra alors un isomorphisme
$$
(A, B, C, RF(f), b, c)
\longrightarrow
(RF(X'), RF(Y'), RF(Z'), RF(f'), RF(g'), RF(h'))
$$
de triangles, et TR1 permettra de conclure que le triangle but est distingué.
Considérons la catégorie $\mathcal{I}$ du lemme \ref{lemma-limit-triangles},
formée des triangles
$$
\mathcal{I} =
\{(t, t', t'') : (X', Y', Z', f', g', h') \to
(X'', Y'', Z'', f'', g'', h'') \mid (t, t', t'') \in S\}
$$
Montrer que le système $F(Z'')$ est essentiellement constant sur la catégorie
$Z'/S$ équivaut à montrer que le système $F(Z'')$ est essentiellement constant sur
$\mathcal{I}$, puisque $\mathcal{I} \to Z'/S$ est cofinal ; voir Catégories,
lemme \ref{categories-lemma-cofinal-essentially-constant} (la cofinalité est
démontrée dans le lemme \ref{lemma-limit-triangles}). Pour tout objet $W$ de
$\mathcal{D}'$, considérons le diagramme
$$
\xymatrix{
\colim_\mathcal{I} \Mor_{\mathcal{D}'}(W, F(X'')) \ar[d] &
\Mor_{\mathcal{D}'}(W, A) \ar[l] \ar[d] \\
\colim_\mathcal{I} \Mor_{\mathcal{D}'}(W, F(Y'')) \ar[d] &
\Mor_{\mathcal{D}'}(W, B) \ar[d] \ar[l] \\
\colim_\mathcal{I} \Mor_{\mathcal{D}'}(W, F(Z'')) \ar[d] &
\Mor_{\mathcal{D}'}(W, C) \ar[d] \ar[l] \\
\colim_\mathcal{I} \Mor_{\mathcal{D}'}(W, F(X''[1])) \ar[d] &
\Mor_{\mathcal{D}'}(W, A[1]) \ar[d] \ar[l] \\
\colim_\mathcal{I} \Mor_{\mathcal{D}'}(W, F(Y''[1])) &
\Mor_{\mathcal{D}'}(W, B[1]) \ar[l]
}
$$
où les flèches horizontales sont données par composition avec
$(\alpha, \beta, \gamma)$. Comme les colimites filtrantes sont exactes
(Algèbre, lemme \ref{algebra-lemma-directed-colimit-exact}), la colonne gauche
est une suite exacte. Le lemme $5$ (Homologie, lemme
\ref{homology-lemma-five-lemma}) montre donc que le morphisme
$$
\colim_\mathcal{I} \Mor_{\mathcal{D}'}(W, F(Z''))
\longrightarrow
\Mor_{\mathcal{D}'}(W, C)
$$
est bijectif. On conclut que $F(Z'')$ est essentiellement constant sur
$\mathcal{I}$, de valeur $C$, par la partie (4) de Catégories, lemme
\ref{categories-lemma-characterize-essentially-constant-ind}.
\end{proof}

\begin{lemma}
\label{lemma-direct-sum-defined}
Hypothèses et notations comme dans la
situation \ref{situation-derived-functor}. Soient $X, Y$ des objets de
$\mathcal{D}$.
\begin{enumerate}
\item Si $RF$ est défini en $X$ et en $Y$, alors $RF$ est défini en
$X \oplus Y$.
\item Si $\mathcal{D}'$ est karoubienne et si $RF$ est défini en
$X \oplus Y$, alors $RF$ est défini à la fois en $X$ et en $Y$.
\end{enumerate}
Dans les deux cas, on a $RF(X \oplus Y) = RF(X) \oplus RF(Y)$. De même pour
$LF$.
\end{lemma}

\begin{proof}
Si $RF$ est défini en $X$ et en $Y$, le triangle distingué
$X \to X \oplus Y \to Y \to X[1]$ (lemme \ref{lemma-split}) et le
lemme \ref{lemma-2-out-of-3-defined} montrent que $RF$ est défini en
$X \oplus Y$ et que l'on a un triangle distingué
$RF(X) \to RF(X \oplus Y) \to RF(Y) \to RF(X)[1]$. En appliquant encore une
fois le lemme \ref{lemma-split}, on trouve
$RF(X \oplus Y) = RF(X) \oplus RF(Y)$. Cela démontre (1) et l'assertion
finale.

\medskip\noindent
Réciproquement, supposons que $RF$ soit défini en $X \oplus Y$ et que
$\mathcal{D}'$ soit karoubienne. Comme $S$ est un système saturé, $S$ est
l'ensemble des flèches qui deviennent inversibles sous le foncteur additif de
localisation $Q : \mathcal{D} \to S^{-1}\mathcal{D}$ ; voir Catégories, lemme
\ref{categories-lemma-what-gets-inverted}. Ainsi, pour
$s : X \to X'$ et $s' : Y \to Y'$ dans $S$, le morphisme
$s \oplus s' : X \oplus Y \to X' \oplus Y'$ appartient à $S$. On obtient de
cette manière un foncteur
$$
X/S \times Y/S \longrightarrow (X \oplus Y)/S
$$
Rappelons que les catégories $X/S, Y/S, (X \oplus Y)/S$ sont filtrantes
(Catégories, remarque
\ref{categories-remark-left-localization-morphisms-colimit}). D'après
Catégories, lemme \ref{categories-lemma-essentially-constant-over-product},
$X/S \times Y/S$ est filtrante, et
$F|_{X/S} : X/S \to \mathcal{D}'$ (resp.\
$F|_{Y/S} : Y/S \to \mathcal{D}'$) est essentiellement constant si et
seulement si
$F|_{X/S} \circ \text{pr}_1 : X/S \times Y/S \to \mathcal{D}'$ (resp.\
$F|_{Y/S} \circ \text{pr}_2 : X/S \times Y/S \to \mathcal{D}'$) est
essentiellement constant. Nous montrerons ci-dessous que le foncteur affiché
est cofinal ; ainsi, d'après Catégories, lemme
\ref{categories-lemma-cofinal-essentially-constant}, si
$F|_{(X \oplus Y)/S}$ est essentiellement constant, alors
$$
F|_{X/S} \circ \text{pr}_1 \oplus F|_{Y/S} \circ \text{pr}_2 :
X/S \times Y/S \to \mathcal{D}'
$$
est essentiellement constant. D'après Homologie, lemme
\ref{homology-lemma-direct-sum-from-product-essentially-constant}
(et c'est ici que nous utilisons que $\mathcal{D}'$ est karoubienne), le fait
que
$F|_{X/S} \circ \text{pr}_1 \oplus F|_{Y/S} \circ \text{pr}_2$ soit
essentiellement constant implique que
$F|_{X/S} \circ \text{pr}_1$ et $F|_{Y/S} \circ \text{pr}_2$ sont
essentiellement constants, ce qui prouve que $RF$ est défini en $X$ et en $Y$.

\medskip\noindent
Démontrons que le foncteur affiché est cofinal. Choisissons pour cela un
$t : X \oplus Y \to Z$ dans $S$. À l'aide de MS2, on peut trouver des
morphismes $Z \to X'$, $Z \to Y'$ et des morphismes
$s : X \to X'$, $s' : Y \to Y'$ dans $S$ tels que le diagramme
$$
\xymatrix{
X \ar[d]^s & X \oplus Y \ar[d] \ar[l] \ar[r] & Y \ar[d]_{s'} \\
X' & Z \ar[l] \ar[r] & Y'
}
$$
soit commutatif. Cela montre qu'il existe un morphisme
$Z \to X' \oplus Y'$ dans $(X \oplus Y)/S$, autrement dit que la partie (1)
de Catégories, définition \ref{categories-definition-cofinal}, est satisfaite.
Pour démontrer la partie (2), il suffit de montrer que, étant donnés
$t : X \oplus Y \to Z$ et des morphismes
$s_i \oplus s'_i : Z \to X'_i \oplus Y'_i$, $i = 1, 2$, dans
$(X \oplus Y)/S$, on peut trouver des morphismes
$a : X'_1 \to X'$, $b : X'_2 \to X'$, $c : Y'_1 \to Y'$,
$d : Y'_2 \to Y'$ dans $S$ tels que
$a \circ s_1 = b \circ s_2$ et $c \circ s'_1 = d \circ s'_2$.
Pour cela, choisissons d'abord des $X'$ et $Y'$ quelconques et des morphismes
$a, b, c, d$ dans $S$ ; c'est possible puisque $X/S$ et $Y/S$ sont
filtrantes. Les deux morphismes
$a \circ s_1, b \circ s_2 : Z \to X'$ deviennent alors égaux dans
$S^{-1}\mathcal{D}$. On peut donc trouver un morphisme $X' \to X''$ dans $S$
qui les égalise. De même, on trouve un morphisme $Y' \to Y''$ dans $S$ qui
égalise $c \circ s'_1$ et $d \circ s'_2$. En remplaçant $X'$ par $X''$ et
$Y'$ par $Y''$, on obtient $a \circ s_1 = b \circ s_2$ et
$c \circ s'_1 = d \circ s'_2$.

\medskip\noindent
La démonstration des assertions correspondantes pour $LF$ est duale.
\end{proof}

\begin{proposition}
\label{proposition-derived-functor}
Hypothèses et notations comme dans la
situation \ref{situation-derived-functor}.
\begin{enumerate}
\item La sous-catégorie pleine $\mathcal{E}$ de $\mathcal{D}$ formée des
objets en lesquels $RF$ est défini est une sous-catégorie triangulée
strictement pleine de $\mathcal{D}$.
\item On obtient un foncteur exact de catégories triangulées
$RF : \mathcal{E} \longrightarrow \mathcal{D}'$.
\item Les éléments de $S$ dont la source ou le but appartient à $\mathcal{E}$
sont des morphismes de $\mathcal{E}$.
\item Tout élément de
$S_\mathcal{E} = \text{Flèches}(\mathcal{E}) \cap S$ est envoyé par $RF$ sur
un isomorphisme.
\item L'ensemble $S_\mathcal{E}$ est un système multiplicatif saturé de
$\mathcal{E}$ compatible avec la structure triangulée.
\item Le foncteur $S_\mathcal{E}^{-1}\mathcal{E} \to S^{-1}\mathcal{D}$ est un
foncteur exact pleinement fidèle de catégories triangulées.
\item On obtient un foncteur exact
$$
RF : S_\mathcal{E}^{-1}\mathcal{E} \longrightarrow \mathcal{D}'.
$$
\item Si $\mathcal{D}'$ est karoubienne, alors $\mathcal{E}$ est une
sous-catégorie triangulée saturée de $\mathcal{D}$.
\end{enumerate}
Un résultat analogue vaut pour $LF$.
\end{proposition}

\begin{proof}
Comme $S$ est saturé, il contient tous les isomorphismes (voir la remarque qui
suit Catégories, définition
\ref{categories-definition-saturated-multiplicative-system}). Ainsi, (1)
résulte des lemmes \ref{lemma-derived-inverts},
\ref{lemma-2-out-of-3-defined} et \ref{lemma-derived-shift}. L'assertion (2)
résulte des lemmes \ref{lemma-derived-functor}, \ref{lemma-derived-shift} et
\ref{lemma-2-out-of-3-defined}. L'assertion (3) résulte du
lemme \ref{lemma-derived-inverts}. La partie (4) résulte du
lemme \ref{lemma-derived-inverts}. La partie (5) résulte des définitions et de
la partie (3). La pleine fidélité de (6) résulte de (3) et des définitions. Le
fait que $S_\mathcal{E}^{-1}\mathcal{E} \to S^{-1}\mathcal{D}$ soit exact
résulte du fait qu'un triangle de $S_\mathcal{E}^{-1}\mathcal{E}$ est
distingué si et seulement s'il est isomorphe à l'image d'un triangle distingué
de $\mathcal{E}$ ; voir la démonstration de la
proposition \ref{proposition-construct-localization}. La factorisation de
$RF : \mathcal{E} \to \mathcal{D}'$ par un foncteur exact
$S_\mathcal{E}^{-1}\mathcal{E} \to \mathcal{D}'$ résulte du
lemme \ref{lemma-universal-property-localization}. Enfin, la partie (8) résulte
du lemme \ref{lemma-direct-sum-defined}.
\end{proof}

\noindent
La proposition \ref{proposition-derived-functor} nous dit que $RF$ vit sur une
sous-catégorie triangulée strictement pleine maximale de
$S^{-1}\mathcal{D}$ et qu'il est exact sur cette catégorie triangulée. Voici le
diagramme :
$$
\xymatrix{
\mathcal{D} \ar[d]_Q \ar[rrr]_F & & & \mathcal{D}' \\
S^{-1}\mathcal{D} & &
S_\mathcal{E}^{-1}\mathcal{E}
\ar[ll]_{\text{pleinement fidèle}}^{\text{exact}} \ar[ur]_{RF}
}
$$

\begin{definition}
\label{definition-everywhere-defined}
Dans la situation \ref{situation-derived-functor}, nous disons que $F$ est
{\it dérivable à droite}, ou que $RF$ est {\it défini partout}, si $RF$ est
défini en tout objet de $\mathcal{D}$. Nous disons que $F$ est
{\it dérivable à gauche}, ou que $LF$ est {\it défini partout}, si $LF$ est
défini en tout objet de $\mathcal{D}$.
\end{definition}

\noindent
Dans ce cas, on obtient un foncteur dérivé droit (resp.\ gauche)
\begin{equation}
\label{equation-everywhere}
RF : S^{-1}\mathcal{D} \longrightarrow \mathcal{D}',
\quad\text{(resp. }
LF : S^{-1}\mathcal{D} \longrightarrow \mathcal{D}'),
\end{equation}
voir la proposition \ref{proposition-derived-functor}. Dans la plupart des
situations intéressantes, $RF \circ Q$ n'est pas égal à $F$. Il peut même
arriver que le morphisme canonique $F(X) \to RF(X)$ ne soit jamais un
isomorphisme. En pratique, cela ne se produit pas, car, en pratique, la seule
façon dont nous savons démontrer que $F$ est dérivable à droite consiste à
montrer que $RF$ peut être calculé en évaluant $F$ sur des objets judicieusement
choisis de la catégorie triangulée $\mathcal{D}$. Cela motive la définition
suivante.

\begin{definition}
\label{definition-computes}
Dans la situation \ref{situation-derived-functor} :
\begin{enumerate}
\item Un objet $X$ de $\mathcal{D}$ {\it calcule} $RF$ si $RF$ est défini en
$X$ et si le morphisme canonique $F(X) \to RF(X)$ est un isomorphisme.
\item Un objet $X$ de $\mathcal{D}$ {\it calcule} $LF$ si $LF$ est défini en
$X$ et si le morphisme canonique $LF(X) \to F(X)$ est un isomorphisme.
\end{enumerate}
\end{definition}

\begin{lemma}
\label{lemma-computes-shift}
Hypothèses et notations comme dans la
situation \ref{situation-derived-functor}. Soient $X$ un objet de
$\mathcal{D}$ et $n \in \mathbf{Z}$.
\begin{enumerate}
\item $X$ calcule $RF$ si et seulement si $X[n]$ calcule $RF$.
\item $X$ calcule $LF$ si et seulement si $X[n]$ calcule $LF$.
\end{enumerate}
\end{lemma}

\begin{proof}
Omis.
\end{proof}

\begin{lemma}
\label{lemma-2-out-of-3-computes}
Hypothèses et notations comme dans la
situation \ref{situation-derived-functor}. Soit $(X, Y, Z, f, g, h)$ un
triangle distingué de $\mathcal{D}$. Si $X, Y$ calculent $RF$, alors $Z$ le
calcule aussi. De même pour $LF$.
\end{lemma}

\begin{proof}
D'après le lemme \ref{lemma-2-out-of-3-defined}, $RF$ est défini en $Z$, et
l'application de $RF$ au triangle produit un triangle distingué. Considérons
le morphisme de triangles distingués
$$
\xymatrix{
(F(X), F(Y), F(Z), F(f), F(g), F(h)) \ar[d] \\
(RF(X), RF(Y), RF(Z), RF(f), RF(g), RF(h))
}
$$
Deux des trois morphismes sont des isomorphismes ; le troisième l'est donc
aussi.
\end{proof}

\begin{lemma}
\label{lemma-direct-sum-computes}
Hypothèses et notations comme dans la
situation \ref{situation-derived-functor}. Soient $X, Y$ des objets de
$\mathcal{D}$. Si $X \oplus Y$ calcule $RF$, alors $X$ et $Y$ calculent $RF$.
De même pour $LF$.
\end{lemma}

\begin{proof}
Si $X \oplus Y$ calcule $RF$, alors
$RF(X \oplus Y) = F(X) \oplus F(Y)$. Dans la démonstration du
lemme \ref{lemma-direct-sum-defined}, nous avons vu que le foncteur
$X/S \times Y/S \to (X \oplus Y)/S$,
$(s, s') \mapsto s \oplus s'$ est cofinal. Ainsi, d'après Catégories, lemme
\ref{categories-lemma-cofinal-essentially-constant}, et d'après la
caractérisation (4) de Catégories, lemme
\ref{categories-lemma-characterize-essentially-constant-ind}, pour tout objet
$W$ de $\mathcal{D}'$, le morphisme
$$
\Hom_{\mathcal{D}'}(F(X \oplus Y), W)
\longrightarrow
\colim_{s : X \to X', s' : Y \to Y'}
\Hom_{\mathcal{D}'}(F(X' \oplus Y'), W)
$$
est bijectif. Comme cette flèche est manifestement compatible avec les
décompositions en somme directe des deux côtés, on en conclut que le morphisme
$$
\Hom_{\mathcal{D}'}(F(X), W)
\longrightarrow
\colim_{s : X \to X'} \Hom_{\mathcal{D}'}(F(X'), W)
$$
est bijectif (un détail mineur est omis). Ainsi, d'après Catégories, lemme
\ref{categories-lemma-characterize-essentially-constant-ind}, $RF$ est défini
en $X$, de valeur $F(X)$. De même pour $Y$.
\end{proof}

\begin{lemma}
\label{lemma-existence-computes}
Hypothèses et notations comme dans la
situation \ref{situation-derived-functor}.
\begin{enumerate}
\item Si, pour tout objet $X \in \Ob(\mathcal{D})$, il existe une flèche
$s : X \to X'$ de $S$ telle que $X'$ calcule $RF$, alors $RF$ est défini
partout.
\item Si, pour tout objet $X \in \Ob(\mathcal{D})$, il existe une flèche
$s : X' \to X$ de $S$ telle que $X'$ calcule $LF$, alors $LF$ est défini
partout.
\end{enumerate}
\end{lemma}

\begin{proof}
C'est clair d'après les définitions.
\end{proof}

\begin{lemma}
\label{lemma-find-existence-computes}
Hypothèses et notations comme dans la
situation \ref{situation-derived-functor}. S'il existe un sous-ensemble
$\mathcal{I} \subset \Ob(\mathcal{D})$ tel que
\begin{enumerate}
\item pour tout $X \in \Ob(\mathcal{D})$, il existe
$s : X \to X'$ dans $S$ avec $X' \in \mathcal{I}$, et
\item pour toute flèche $s : X \to X'$ dans $S$ avec
$X, X' \in \mathcal{I}$, le morphisme $F(s) : F(X) \to F(X')$ est un
isomorphisme,
\end{enumerate}
alors $RF$ est défini partout, et tout $X \in \mathcal{I}$ calcule $RF$.
Dualement, s'il existe un sous-ensemble
$\mathcal{P} \subset \Ob(\mathcal{D})$ tel que
\begin{enumerate}
\item pour tout $X \in \Ob(\mathcal{D})$, il existe
$s : X' \to X$ dans $S$ avec $X' \in \mathcal{P}$, et
\item pour toute flèche $s : X \to X'$ dans $S$ avec
$X, X' \in \mathcal{P}$, le morphisme $F(s) : F(X) \to F(X')$ est un
isomorphisme,
\end{enumerate}
alors $LF$ est défini partout, et tout $X \in \mathcal{P}$ calcule $LF$.
\end{lemma}

\begin{proof}
Soit $X$ un objet de $\mathcal{D}$. L'hypothèse (1) implique que les flèches
$s : X \to X'$ de $S$ avec $X' \in \mathcal{I}$ sont cofinales dans la
catégorie $X/S$. L'hypothèse (2) implique que $F$ est constant sur cette
sous-catégorie cofinale. Cela implique clairement que
$F : (X/S) \to \mathcal{D}'$ est essentiellement constant, de valeur $F(X')$,
pour toute flèche $s : X \to X'$ de $S$ avec $X' \in \mathcal{I}$.
\end{proof}

\begin{lemma}
\label{lemma-compose-derived-functors-general}
Soient $\mathcal{A}, \mathcal{B}, \mathcal{C}$ des catégories triangulées.
Soient $S$, resp.\ $S'$, des systèmes multiplicatifs saturés de
$\mathcal{A}$, resp.\ $\mathcal{B}$, compatibles avec la structure triangulée.
Soient $F : \mathcal{A} \to \mathcal{B}$ et
$G : \mathcal{B} \to \mathcal{C}$ des foncteurs exacts. Notons
$F' : \mathcal{A} \to (S')^{-1}\mathcal{B}$ la composée de $F$ avec le
foncteur de localisation.
\begin{enumerate}
\item Si $RF'$, $RG$, $R(G \circ F)$ sont définis partout, il existe une
transformation canonique de foncteurs
$t : R(G \circ F) \longrightarrow RG \circ RF'$.
\item Si $LF'$, $LG$, $L(G \circ F)$ sont définis partout, il existe une
transformation canonique de foncteurs
$t : LG \circ LF' \to L(G \circ F)$.
\end{enumerate}
\end{lemma}

\begin{proof}
Dans cette démonstration, nous nous efforçons d'être précis. Considérons donc
les foncteurs dérivés comme les foncteurs
$$
RF' : S^{-1}\mathcal{A} \to (S')^{-1}\mathcal{B}, \quad
R(G \circ F) : S^{-1}\mathcal{A} \to \mathcal{C}, \quad
RG : (S')^{-1}\mathcal{B} \to \mathcal{C}.
$$
Notons $Q_A : \mathcal{A} \to S^{-1}\mathcal{A}$ et
$Q_B : \mathcal{B} \to (S')^{-1}\mathcal{B}$ les foncteurs de localisation.
Alors $F' = Q_B \circ F$. Notons que, pour tout objet $Y$ de $\mathcal{B}$,
il existe un morphisme canonique
$$
G(Y) \longrightarrow RG(Q_B(Y))
$$
autrement dit, une transformation de foncteurs
$t' : G \to RG \circ Q_B$. Soit $X$ un objet de $\mathcal{A}$. On a
\begin{align*}
R(G \circ F)(Q_A(X))
& = \colim_{s : X \to X' \in S} G(F(X')) \\
& \xrightarrow{t'} \colim_{s : X \to X' \in S} RG(Q_B(F(X'))) \\
& = \colim_{s : X \to X' \in S} RG(F'(X')) \\
& = RG(\colim_{s : X \to X' \in S} F'(X')) \\
& = RG(RF'(X)).
\end{align*}
Le système $F'(X')$ est essentiellement constant dans la catégorie
$(S')^{-1}\mathcal{B}$. On peut donc faire passer la colimite à l'intérieur du
foncteur $RG$ dans la troisième égalité du diagramme ci-dessus ; voir
Catégories, lemme \ref{categories-lemma-image-essentially-constant}, et sa
démonstration. Nous omettons la démonstration que cela définit une
transformation de foncteurs. Le cas des foncteurs dérivés gauches est analogue.
\end{proof}





\section{Foncteurs dérivés sur les catégories dérivées}
\label{section-derived-functors-classical}

\noindent
En pratique, les foncteurs dérivés apparaissent le plus souvent à partir d'un
foncteur additif entre catégories abéliennes.

\begin{situation}
\label{situation-classical}
Dans cette situation, $F : \mathcal{A} \to \mathcal{B}$ est un foncteur
additif entre catégories abéliennes. Il induit des foncteurs exacts
$$
F : K(\mathcal{A}) \to K(\mathcal{B}), \quad
K^{+}(\mathcal{A}) \to K^{+}(\mathcal{B}), \quad
K^{-}(\mathcal{A}) \to K^{-}(\mathcal{B}).
$$
Voir le lemme \ref{lemma-additive-exact-homotopy-category}. Nous notons aussi
$F$ les composées $K(\mathcal{A}) \to D(\mathcal{B})$,
$K^{+}(\mathcal{A}) \to D^{+}(\mathcal{B})$ et
$K^{-}(\mathcal{A}) \to D^-(\mathcal{B})$ de $F$ avec le foncteur de
localisation $K(\mathcal{B}) \to D(\mathcal{B})$, etc. Cette situation conduit
aux quatre foncteurs dérivés suivants.
\begin{enumerate}
\item Le foncteur dérivé droit de
$F : K(\mathcal{A}) \to D(\mathcal{B})$ relativement au système multiplicatif
$\text{Qis}(\mathcal{A})$.
\item Le foncteur dérivé droit de
$F : K^{+}(\mathcal{A}) \to D^{+}(\mathcal{B})$ relativement au système
multiplicatif $\text{Qis}^{+}(\mathcal{A})$.
\item Le foncteur dérivé gauche de
$F : K(\mathcal{A}) \to D(\mathcal{B})$ relativement au système multiplicatif
$\text{Qis}(\mathcal{A})$.
\item Le foncteur dérivé gauche de
$F : K^{-}(\mathcal{A}) \to D^{-}(\mathcal{B})$ relativement au système
multiplicatif $\text{Qis}^-(\mathcal{A})$.
\end{enumerate}
Chacun de ces cas est un exemple de la
situation \ref{situation-derived-functor}.
\end{situation}

\noindent
Le lemme suivant lève une partie de l'ambiguïté qui pourrait survenir.

\begin{lemma}
\label{lemma-irrelevant}
Dans la situation \ref{situation-classical} :
\begin{enumerate}
\item Soit $X$ un objet de $K^{+}(\mathcal{A})$. Le foncteur dérivé droit de
$K(\mathcal{A}) \to D(\mathcal{B})$ est défini en $X$ si et seulement si le
foncteur dérivé droit de
$K^{+}(\mathcal{A}) \to D^{+}(\mathcal{B})$ est défini en $X$. De plus, leurs
valeurs sont canoniquement isomorphes.
\item Soit $X$ un objet de $K^{+}(\mathcal{A})$. Alors $X$ calcule le foncteur
dérivé droit de $K(\mathcal{A}) \to D(\mathcal{B})$ si et seulement si $X$
calcule le foncteur dérivé droit de
$K^{+}(\mathcal{A}) \to D^{+}(\mathcal{B})$.
\item Soit $X$ un objet de $K^{-}(\mathcal{A})$. Le foncteur dérivé gauche de
$K(\mathcal{A}) \to D(\mathcal{B})$ est défini en $X$ si et seulement si le
foncteur dérivé gauche de
$K^{-}(\mathcal{A}) \to D^{-}(\mathcal{B})$ est défini en $X$. De plus, leurs
valeurs sont canoniquement isomorphes.
\item Soit $X$ un objet de $K^{-}(\mathcal{A})$. Alors $X$ calcule le foncteur
dérivé gauche de $K(\mathcal{A}) \to D(\mathcal{B})$ si et seulement si $X$
calcule le foncteur dérivé gauche de
$K^{-}(\mathcal{A}) \to D^{-}(\mathcal{B})$.
\end{enumerate}
\end{lemma}

\begin{proof}
Soit $X$ un objet de $K^{+}(\mathcal{A})$. Considérons un quasi-isomorphisme
$s : X \to X'$ dans $K(\mathcal{A})$. D'après le
lemme \ref{lemma-complex-cohomology-bounded}, il existe un quasi-isomorphisme
$X' \to X''$ avec $X''$ borné inférieurement. Ainsi,
$X/\text{Qis}^+(\mathcal{A})$ est cofinal dans
$X/\text{Qis}(\mathcal{A})$. L'assertion (1) est donc claire. La partie (2)
résulte directement de la partie (1). Les parties (3) et (4) sont duales des
parties (1) et (2).
\end{proof}

\noindent
À tout objet $A$ d'une catégorie abélienne $\mathcal{A}$, on associe le
complexe
$$
A[0] = ( \ldots \to 0 \to A \to 0 \to \ldots )
$$
où $A$ est placé en degré zéro. On obtient ainsi un foncteur
$\mathcal{A} \to K(\mathcal{A})$, $A \mapsto A[0]$. Disons provisoirement
qu'un foncteur partiel est un foncteur défini sur une sous-catégorie.

\begin{definition}
\label{definition-derived-functor}
Dans la situation \ref{situation-classical} :
\begin{enumerate}
\item Les {\it foncteurs dérivés droits de $F$} sont les foncteurs partiels
$RF$ associés aux cas (1) et (2) de la situation \ref{situation-classical}.
\item Les {\it foncteurs dérivés gauches de $F$} sont les foncteurs partiels
$LF$ associés aux cas (3) et (4) de la situation \ref{situation-classical}.
\item Un objet $A$ de $\mathcal{A}$ est dit {\it acyclique à droite pour $F$},
ou {\it acyclique pour $RF$}, si $A[0]$ calcule $RF$.
\item Un objet $A$ de $\mathcal{A}$ est dit {\it acyclique à gauche pour $F$},
ou {\it acyclique pour $LF$}, si $A[0]$ calcule $LF$.
\end{enumerate}
\end{definition}

\noindent
Les quelques lemmes suivants donnent des critères d'existence d'assez d'objets
acycliques.

\begin{lemma}
\label{lemma-subcategory-left-resolution}
Soit $\mathcal{A}$ une catégorie abélienne. Soit
$\mathcal{P} \subset \Ob(\mathcal{A})$ un sous-ensemble contenant $0$, tel
que tout objet de $\mathcal{A}$ soit quotient d'un élément de $\mathcal{P}$.
Soit $a \in \mathbf{Z}$.
\begin{enumerate}
\item Étant donné $K^\bullet$ tel que $K^n = 0$ pour $n > a$, il existe un
quasi-isomorphisme $P^\bullet \to K^\bullet$ tel que
$P^n \in \mathcal{P}$ et $P^n \to K^n$ soit surjectif pour tout $n$, et tel
que $P^n = 0$ pour $n > a$.
\item Étant donné $K^\bullet$ tel que $H^n(K^\bullet) = 0$ pour $n > a$, il
existe un quasi-isomorphisme $P^\bullet \to K^\bullet$ tel que
$P^n \in \mathcal{P}$ pour tout $n$ et $P^n = 0$ pour $n > a$.
\end{enumerate}
\end{lemma}

\begin{proof}
Démonstration de la partie (1). Considérons l'hypothèse de récurrence suivante
$IH_n$ : il existe des $P^j \in \mathcal{P}$, $j \geq n$, avec $P^j = 0$
pour $j > a$, des morphismes $d^j : P^j \to P^{j + 1}$ pour $j \geq n$, et
des morphismes surjectifs $\alpha^j : P^j \to K^j$ pour $j \geq n$, tels que
le diagramme
$$
\xymatrix{
& &
P^n \ar[d]^\alpha \ar[r] &
P^{n + 1} \ar[d]^\alpha \ar[r] &
P^{n + 2} \ar[d]^\alpha \ar[r] &
\ldots \\
\ldots \ar[r] &
K^{n - 1} \ar[r] &
K^n \ar[r] &
K^{n + 1} \ar[r] &
K^{n + 2} \ar[r] &
\ldots
}
$$
soit commutatif, que $d^{j + 1} \circ d^j = 0$ pour $j \geq n$, que
$\alpha$ induise des isomorphismes
$\Ker(d^j)/\Im(d^{j - 1}) \to H^j(K^\bullet)$ pour $j > n$, et que
$\alpha : \Ker(d^n) \to \Ker(d_K^n)$ soit surjectif. Choisissons alors une
surjection
$$
P^{n - 1}
\longrightarrow
K^{n - 1} \times_{K^n} \Ker(d^n) =
K^{n - 1} \times_{\Ker(d_K^n)} \Ker(d^n)
$$
avec $P^{n - 1}$ dans $\mathcal{P}$. Cela permet de prolonger le diagramme
ci-dessus en
$$
\xymatrix{
& P^{n - 1} \ar[d]^\alpha \ar[r] &
P^n \ar[d]^\alpha \ar[r] &
P^{n + 1} \ar[d]^\alpha \ar[r] &
P^{n + 2} \ar[d]^\alpha \ar[r] &
\ldots \\
\ldots \ar[r] &
K^{n - 1} \ar[r] &
K^n \ar[r] &
K^{n + 1} \ar[r] &
K^{n + 2} \ar[r] &
\ldots
}
$$
Le lecteur vérifie aisément que $IH_{n - 1}$ est satisfait avec ce choix.

\medskip\noindent
Achevons la démonstration de (1). Remarquons d'abord que $IH_n$ est vrai pour
$n = a + 1$, puisqu'on peut simplement prendre $P^j = 0$ pour $j > a$.
En procédant par récurrence descendante, on construit donc un complexe
$P^\bullet$, avec $P^n = 0$ pour $n > a$, formé d'objets de $\mathcal{P}$, et
un quasi-isomorphisme terme à terme surjectif
$\alpha : P^\bullet \to K^\bullet$, comme voulu.

\medskip\noindent
Démonstration de la partie (2). L'hypothèse implique que le morphisme
$\tau_{\leq a}K^\bullet \to K^\bullet$ (Homologie, section
\ref{homology-section-truncations}) est un quasi-isomorphisme. Appliquons la
partie (1) pour trouver $P^\bullet \to \tau_{\leq a}K^\bullet$. La composée
$P^\bullet \to K^\bullet$ est le quasi-isomorphisme recherché.
\end{proof}

\begin{lemma}
\label{lemma-subcategory-right-resolution}
Soit $\mathcal{A}$ une catégorie abélienne. Soit
$\mathcal{I} \subset \Ob(\mathcal{A})$ un sous-ensemble contenant $0$, tel
que tout objet de $\mathcal{A}$ soit un sous-objet d'un élément de
$\mathcal{I}$. Soit $a \in \mathbf{Z}$.
\begin{enumerate}
\item Étant donné $K^\bullet$ tel que $K^n = 0$ pour $n < a$, il existe un
quasi-isomorphisme $K^\bullet \to I^\bullet$ tel que
$K^n \to I^n$ soit injectif et $I^n \in \mathcal{I}$ pour tout $n$, et tel
que $I^n = 0$ pour $n < a$.
\item Étant donné $K^\bullet$ tel que $H^n(K^\bullet) = 0$ pour $n < a$, il
existe un quasi-isomorphisme $K^\bullet \to I^\bullet$ tel que
$I^n \in \mathcal{I}$ et $I^n = 0$ pour $n < a$.
\end{enumerate}
\end{lemma}

\begin{proof}
Ce lemme est dual du lemme \ref{lemma-subcategory-left-resolution}.
\end{proof}

\begin{lemma}
\label{lemma-subcategory-right-acyclics}
Dans la situation \ref{situation-classical}, soit
$\mathcal{I} \subset \Ob(\mathcal{A})$ un sous-ensemble possédant les
propriétés suivantes :
\begin{enumerate}
\item tout objet de $\mathcal{A}$ est un sous-objet d'un élément de
$\mathcal{I}$ ;
\item pour toute suite exacte courte $0 \to P \to Q \to R \to 0$ de
$\mathcal{A}$ avec $P, Q \in \mathcal{I}$, on a $R \in \mathcal{I}$ et la
suite $0 \to F(P) \to F(Q) \to F(R) \to 0$ est exacte.
\end{enumerate}
Alors tout objet de $\mathcal{I}$ est acyclique pour $RF$.
\end{lemma}

\begin{proof}
Choisissons $A \in \mathcal{I}$. Soit $A[0] \to K^\bullet$ un
quasi-isomorphisme avec $K^\bullet$ borné inférieurement. On peut trouver un
quasi-isomorphisme $K^\bullet \to I^\bullet$, avec $I^\bullet$ borné
inférieurement et chaque $I^n \in \mathcal{I}$ ; voir le
lemme \ref{lemma-subcategory-right-resolution}\footnote{D'après (1),
$\mathcal{I}$ n'est pas vide. Choisissons $P$ dans $\mathcal{I}$. La suite
exacte courte $0 \to P \to P \to 0 \to 0$ et l'hypothèse (2) montrent alors
que $0$ appartient à $\mathcal{I}$. Le lemme s'applique donc.}. Ces résolutions
sont ainsi cofinales dans la catégorie
$A[0]/\text{Qis}^{+}(\mathcal{A})$. Pour achever la démonstration, il suffit
donc de montrer que, pour tout quasi-isomorphisme
$A[0] \to I^\bullet$ avec $I^\bullet$ borné inférieurement et
$I^n \in \mathcal{I}$, le morphisme
$F(A)[0] \to F(I^\bullet)$ est un quasi-isomorphisme.
Supposons pour cela que $I^n = 0$ pour $n < n_0$. On peut évidemment supposer
que $n_0 < 0$. À partir de $n = n_0$, on démontre par récurrence que
$\Im(d^{n - 1}) = \Ker(d^n)$ et $\Im(d^{-1})$ sont des éléments de
$\mathcal{I}$, en utilisant la propriété (2) et les suites exactes
$$
0 \to \Ker(d^n) \to I^n \to \Im(d^n) \to 0.
$$
La propriété (2) garantit en outre l'exactitude du complexe
$$
0 \to F(I^{n_0}) \to F(I^{n_0 + 1}) \to \ldots \to F(I^{-1}) \to
F(\Im(d^{-1})) \to 0
$$
La suite exacte
$0 \to \Im(d^{-1}) \to I^0 \to I^0/\Im(d^{-1}) \to 0$ implique que
$I^0/\Im(d^{-1})$ appartient à $\mathcal{I}$. La suite exacte
$0 \to A \to I^0/\Im(d^{-1}) \to \Im(d^0) \to 0$ implique alors que
$\Im(d^0) = \Ker(d^1)$ appartient à $\mathcal{I}$ ; on poursuit ensuite comme
précédemment pour montrer que
$\Im(d^{n - 1}) = \Ker(d^n)$ appartient à $\mathcal{I}$ pour tout $n > 0$.
En appliquant $F$ à chacune des suites exactes courtes mentionnées ci-dessus et
en utilisant (2), on constate que
$F(A)[0] \to F(I^\bullet)$ est un quasi-isomorphisme, comme voulu.
\end{proof}

\begin{lemma}
\label{lemma-subcategory-left-acyclics}
Dans la situation \ref{situation-classical}, soit
$\mathcal{P} \subset \Ob(\mathcal{A})$ un sous-ensemble possédant les
propriétés suivantes :
\begin{enumerate}
\item tout objet de $\mathcal{A}$ est un quotient d'un élément de
$\mathcal{P}$ ;
\item pour toute suite exacte courte $0 \to P \to Q \to R \to 0$ de
$\mathcal{A}$ avec $Q, R \in \mathcal{P}$, on a $P \in \mathcal{P}$ et la
suite $0 \to F(P) \to F(Q) \to F(R) \to 0$ est exacte.
\end{enumerate}
Alors tout objet de $\mathcal{P}$ est acyclique pour $LF$.
\end{lemma}

\begin{proof}
La démonstration est duale de celle du
lemme \ref{lemma-subcategory-right-acyclics}.
\end{proof}







\section{Foncteurs dérivés supérieurs}
\label{section-higher-derived}

\noindent
Le lemme élémentaire suivant montre que les foncteurs dérivés droits « se
déplacent vers la droite ».

\begin{lemma}
\label{lemma-negative-vanishing}
Soit $F : \mathcal{A} \to \mathcal{B}$ un foncteur additif entre catégories
abéliennes. Soient $K^\bullet$ un complexe de $\mathcal{A}$ et
$a \in \mathbf{Z}$.
\begin{enumerate}
\item Si $H^i(K^\bullet) = 0$ pour tout $i < a$ et si $RF$ est défini en
$K^\bullet$, alors $H^i(RF(K^\bullet)) = 0$ pour tout $i < a$.
\item Si $RF$ est défini en $K^\bullet$ et en
$\tau_{\leq a}K^\bullet$, alors
$H^i(RF(\tau_{\leq a}K^\bullet)) = H^i(RF(K^\bullet))$ pour tout $i \leq a$.
\end{enumerate}
\end{lemma}

\begin{proof}
Supposons que $K^\bullet$ satisfasse les hypothèses de (1). Soit
$s : K^\bullet \to L^\bullet$ un quasi-isomorphisme quelconque. Alors
$K^\bullet \to \tau_{\geq a}L^\bullet$ est aussi un quasi-isomorphisme, par
l'hypothèse sur $K^\bullet$. Ainsi, dans la catégorie
$K^\bullet/\text{Qis}^{+}(\mathcal{A})$, les quasi-isomorphismes
$s : K^\bullet \to L^\bullet$ avec $L^n = 0$ pour $n < a$ sont cofinaux.
D'après Catégories, lemme
\ref{categories-lemma-cofinal-essentially-constant}, on en déduit que $RF$ est
la valeur de l'ind-objet essentiellement constant $F(L^\bullet)$ pour ces $s$.
Cela signifie que
$\text{id} : RF(K^\bullet) \to RF(K^\bullet)$ se factorise par
$F(L^\bullet)$ pour un certain complexe $L^\bullet$ tel que $L^n = 0$ pour
$n < a$. Il s'ensuit que $H^i(RF(K^\bullet)) = 0$ pour $i < a$.

\medskip\noindent
Pour démontrer (2), utilisons le triangle distingué
$$
\tau_{\leq a}K^\bullet \to K^\bullet \to \tau_{\geq a + 1}K^\bullet
\to (\tau_{\leq a}K^\bullet)[1]
$$
de la remarque \ref{remark-truncation-distinguished-triangle}. Le
lemme \ref{lemma-2-out-of-3-defined} montre que $RF$ est aussi défini en
$\tau_{\geq a + 1}K^\bullet$ et que l'on a un triangle distingué
$$
RF(\tau_{\leq a}K^\bullet) \to RF(K^\bullet) \to
RF(\tau_{\geq a + 1}K^\bullet)
\to RF(\tau_{\leq a}K^\bullet)[1]
$$
dans $D(\mathcal{B})$. La partie (1) montre que la cohomologie de
$RF(\tau_{\geq a + 1}K^\bullet)$ s'annule en degrés $< a + 1$. La suite exacte
longue de cohomologie de ce triangle distingué donne alors le résultat voulu.
\end{proof}

\begin{definition}
\label{definition-higher-derived-functors}
Soit $F : \mathcal{A} \to \mathcal{B}$ un foncteur additif entre catégories
abéliennes. Supposons que
$RF : D^{+}(\mathcal{A}) \to D^{+}(\mathcal{B})$ soit défini partout. Soit
$i \in \mathbf{Z}$. Le {\it $i$-ième foncteur dérivé droit $R^iF$ de $F$} est
le foncteur
$$
R^iF = H^i \circ RF :
\mathcal{A}
\longrightarrow
\mathcal{B}
$$
\end{definition}

\noindent
Le lemme suivant montre qu'il n'est vraiment guère pertinent de prendre le
foncteur dérivé droit si le foncteur n'est pas exact à gauche.

\begin{lemma}
\label{lemma-left-exact-higher-derived}
Soit $F : \mathcal{A} \to \mathcal{B}$ un foncteur additif entre catégories
abéliennes, et supposons que
$RF : D^{+}(\mathcal{A}) \to D^{+}(\mathcal{B})$ soit défini partout.
\begin{enumerate}
\item On a $R^iF = 0$ pour $i < 0$ ;
\item $R^0F$ est exact à gauche ;
\item le morphisme $F \to R^0F$ est un isomorphisme si et seulement si $F$ est
exact à gauche.
\end{enumerate}
\end{lemma}

\begin{proof}
Soit $A$ un objet de $\mathcal{A}$. D'après le
lemme \ref{lemma-negative-vanishing}, on a
$H^i(RF(A[0])) = 0$ pour $i < 0$. Cela démontre (1).

\medskip\noindent
Soit $0 \to A \to B \to C \to 0$ une suite exacte courte de $\mathcal{A}$.
D'après le lemme \ref{lemma-derived-canonical-delta-functor}, on obtient un
triangle distingué $(A[0], B[0], C[0], a, b, c)$ dans
$D^{+}(\mathcal{A})$. La suite exacte longue de cohomologie (et l'annulation
pour $i < 0$ démontrée ci-dessus) montre que
$0 \to R^0F(A) \to R^0F(B) \to R^0F(C)$ est exacte. Ainsi, $R^0F$ est exact à
gauche. Cela démontre aussi, bien sûr, que si $F \to R^0F$ est un
isomorphisme, alors $F$ est exact à gauche.

\medskip\noindent
Supposons $F$ exact à gauche. Rappelons que $RF(A[0])$ est la valeur du système
essentiellement constant $F(K^\bullet)$, où
$s : A[0] \to K^\bullet$ parcourt les quasi-isomorphismes. Il s'ensuit que
$R^0F(A)$ est la valeur du système essentiellement constant
$H^0(F(K^\bullet))$ lorsque $s : A[0] \to K^\bullet$ parcourt ces mêmes
quasi-isomorphismes ; voir Catégories,
lemme \ref{categories-lemma-image-essentially-constant}. Mais, si
$s : A[0] \to K^\bullet$ est un quasi-isomorphisme, alors
$A[0] \to \tau_{\geq 0}K^\bullet$ est un quasi-isomorphisme. Ainsi, dans la
catégorie $A[0]/\text{Qis}^{+}(\mathcal{A})$, les quasi-isomorphismes
$s : A[0] \to K^\bullet$ avec $K^n = 0$ pour $n < 0$ sont cofinaux. D'après
Catégories, lemme \ref{categories-lemma-cofinal-essentially-constant}, on peut
donc se restreindre à ces $s$. De plus, pour un tel $s$, la suite
$$
0 \to A \to K^0 \to K^1
$$
est exacte. Comme $F$ est exact à gauche, la suite
$0 \to F(A) \to F(K^0) \to F(K^1)$ est également exacte. Il s'ensuit que
$F(A) \to H^0(F(K^\bullet))$ est un isomorphisme et que le système est en fait
constant, de valeur $F(A)$. On conclut que $R^0F = F$, comme voulu.
\end{proof}

\begin{lemma}
\label{lemma-F-acyclic}
Soit $F : \mathcal{A} \to \mathcal{B}$ un foncteur additif entre catégories
abéliennes, et supposons que
$RF : D^{+}(\mathcal{A}) \to D^{+}(\mathcal{B})$ soit défini partout. Soit
$A$ un objet de $\mathcal{A}$.
\begin{enumerate}
\item $A$ est acyclique à droite pour $F$ si et seulement si
$F(A) \to R^0F(A)$ est un isomorphisme et $R^iF(A) = 0$ pour tout $i > 0$ ;
\item si $F$ est exact à gauche, alors $A$ est acyclique à droite pour $F$ si
et seulement si $R^iF(A) = 0$ pour tout $i > 0$.
\end{enumerate}
\end{lemma}

\begin{proof}
Si $A$ est acyclique à droite pour $F$, alors
$RF(A[0]) = F(A)[0]$ ; en particulier, $F(A) \to R^0F(A)$ est un isomorphisme
et $R^iF(A) = 0$ pour $i \not = 0$. Réciproquement, si
$F(A) \to R^0F(A)$ est un isomorphisme et $R^iF(A) = 0$ pour tout $i > 0$,
alors $F(A[0]) \to RF(A[0])$ est un quasi-isomorphisme d'après la partie (1)
du lemme \ref{lemma-left-exact-higher-derived} ; ainsi, $A$ est acyclique.
Si $F$ est exact à gauche, alors $F = R^0F$, voir le
lemme \ref{lemma-left-exact-higher-derived}.
\end{proof}

\begin{lemma}
\label{lemma-F-acyclic-ses}
Soit $F : \mathcal{A} \to \mathcal{B}$ un foncteur exact à gauche entre
catégories abéliennes, et supposons que
$RF : D^{+}(\mathcal{A}) \to D^{+}(\mathcal{B})$ soit défini partout. Soit
$0 \to A \to B \to C \to 0$ une suite exacte courte de $\mathcal{A}$.
\begin{enumerate}
\item Si $A$ et $C$ sont acycliques à droite pour $F$, alors $B$ l'est aussi.
\item Si $A$ et $B$ sont acycliques à droite pour $F$, alors $C$ l'est aussi.
\item Si $B$ et $C$ sont acycliques à droite pour $F$ et si
$F(B) \to F(C)$ est surjectif, alors $A$ est acyclique à droite pour $F$.
\end{enumerate}
Dans chacun des trois cas,
$$
0 \to F(A) \to F(B) \to F(C) \to 0
$$
est une suite exacte courte de $\mathcal{B}$.
\end{lemma}

\begin{proof}
D'après le lemme \ref{lemma-derived-canonical-delta-functor}, on obtient un
triangle distingué $(A[0], B[0], C[0], a, b, c)$ dans
$D^{+}(\mathcal{A})$. Comme $RF$ est un foncteur exact, et comme
$R^iF = 0$ pour $i < 0$ et $R^0F = F$
(lemme \ref{lemma-left-exact-higher-derived}), on obtient une suite exacte de
cohomologie
$$
0 \to F(A) \to F(B) \to F(C) \to R^1F(A) \to \ldots
$$
dans la catégorie abélienne $\mathcal{B}$. Le lemme résulte alors de la
caractérisation des objets acycliques donnée par
le lemme \ref{lemma-F-acyclic}.
\end{proof}

\begin{lemma}
\label{lemma-right-derived-delta-functor}
Soit $F : \mathcal{A} \to \mathcal{B}$ un foncteur additif entre catégories
abéliennes, et supposons que
$RF : D^{+}(\mathcal{A}) \to D^{+}(\mathcal{B})$ soit défini partout.
\begin{enumerate}
\item Les foncteurs $R^iF$, $i \geq 0$, sont munis d'une structure canonique
de $\delta$-foncteur $\mathcal{A} \to \mathcal{B}$ ; voir Homologie,
définition \ref{homology-definition-cohomological-delta-functor}.
\item Si tout objet de $\mathcal{A}$ est un sous-objet d'un objet acyclique à
droite pour $F$, alors $\{R^iF, \delta\}_{i \geq 0}$ est un
$\delta$-foncteur universel ; voir Homologie, définition
\ref{homology-definition-universal-delta-functor}.
\end{enumerate}
\end{lemma}

\begin{proof}
Le foncteur $\mathcal{A} \to \text{Comp}^{+}(\mathcal{A})$,
$A \mapsto A[0]$, est exact. Le foncteur
$\text{Comp}^{+}(\mathcal{A}) \to D^{+}(\mathcal{A})$ est un
$\delta$-foncteur, voir le
lemme \ref{lemma-derived-canonical-delta-functor}. Le foncteur
$RF : D^{+}(\mathcal{A}) \to D^{+}(\mathcal{B})$ est exact. Enfin, le foncteur
$H^0 : D^{+}(\mathcal{B}) \to \mathcal{B}$ est homologique, voir la
définition \ref{definition-unbounded-derived-category}. La structure de
$\delta$-foncteur résulte donc du
lemme \ref{lemma-compose-delta-functor-homological} et du
lemme \ref{lemma-exact-compose-delta-functor}. La partie (2) résulte
d'Homologie, lemme \ref{homology-lemma-efface-implies-universal}, et de la
description des objets acycliques donnée par le lemme \ref{lemma-F-acyclic}.
\end{proof}

\begin{lemma}[Lemme d'acyclicité de Leray]
\label{lemma-leray-acyclicity}
Soit $F : \mathcal{A} \to \mathcal{B}$ un foncteur additif entre catégories
abéliennes. Soit $A^\bullet$ un complexe borné inférieurement d'objets
acycliques à droite pour $F$, tel que $RF$ soit défini en
$A^\bullet$\footnote{C'est par exemple le cas si
$RF : D^{+}(\mathcal{A}) \to D^{+}(\mathcal{B})$ est défini partout.}.
Le morphisme canonique
$$
F(A^\bullet) \longrightarrow RF(A^\bullet)
$$
est un isomorphisme dans $D^{+}(\mathcal{B})$, autrement dit $A^\bullet$
calcule $RF$.
\end{lemma}

\begin{proof}
Soit d'abord $A^\bullet$ un complexe borné d'objets acycliques à droite pour
$F$. Nous affirmons que $RF$ est défini en $A^\bullet$ et que
$F(A^\bullet) \to RF(A^\bullet)$ est un isomorphisme dans
$D^+(\mathcal{B})$. C'est vrai pour les complexes ayant au plus un objet
acyclique à droite pour $F$ non nul, d'après la
définition \ref{definition-derived-functor}. Supposons ensuite que
$A^n = 0$ pour $n \not \in [a, b]$. Les troncations « brutales » donnent une
suite exacte courte de complexes scindée terme à terme
$$
0 \to \sigma_{\geq a + 1} A^\bullet \to A^\bullet \to
\sigma_{\leq a} A^\bullet \to 0
$$
voir Homologie, section \ref{homology-section-truncations}. On obtient donc un
triangle distingué
$(\sigma_{\geq a + 1} A^\bullet, A^\bullet,
\sigma_{\leq a} A^\bullet)$. Par hypothèse de récurrence, $RF$ est défini sur
les deux complexes extrêmes et ceux-ci calculent $RF$. Le complexe du milieu
possède alors la même propriété d'après le
lemme \ref{lemma-2-out-of-3-computes}.

\medskip\noindent
Supposons maintenant que $A^\bullet$ soit un complexe borné inférieurement
d'objets acycliques, tel que $RF$ soit défini en $A^\bullet$. Pour montrer que
$F(A^\bullet) \to RF(A^\bullet)$ est un isomorphisme dans
$D^{+}(\mathcal{B})$, il suffit de montrer que
$H^i(F(A^\bullet)) \to H^i(RF(A^\bullet))$ est un isomorphisme pour tout $i$.
Fixons $i$. Considérons la suite exacte courte de complexes scindée terme à
terme
$$
0 \to \sigma_{\geq i + 2} A^\bullet \to A^\bullet \to
\sigma_{\leq i + 1} A^\bullet \to 0.
$$
Elle induit une suite exacte courte scindée terme à terme
$$
0 \to \sigma_{\geq i + 2} F(A^\bullet) \to F(A^\bullet) \to
\sigma_{\leq i + 1} F(A^\bullet) \to 0.
$$
On obtient donc des triangles distingués
$$
(\sigma_{\geq i + 2} A^\bullet, A^\bullet,
\sigma_{\leq i + 1} A^\bullet)
\quad\text{et}\quad
(\sigma_{\geq i + 2} F(A^\bullet), F(A^\bullet),
\sigma_{\leq i + 1} F(A^\bullet))
$$
Comme $RF$ est défini en $A^\bullet$ (par hypothèse) et en
$\sigma_{\leq i + 1}A^\bullet$ (d'après le premier paragraphe), on voit que
$RF$ est défini en $\sigma_{\geq i + 2}A^\bullet$, et l'on obtient un triangle distingué
$$
(RF(\sigma_{\geq i + 2} A^\bullet), RF(A^\bullet),
RF(\sigma_{\leq i + 1} A^\bullet))
$$
Voir le lemme \ref{lemma-2-out-of-3-defined}. À l'aide de ces triangles
distingués, on obtient un morphisme de suites exactes
$$
\xymatrix{
H^i(\sigma_{\geq i + 2} F(A^\bullet)) \ar[r] \ar[d] &
H^i(F(A^\bullet)) \ar[r] \ar[d]^\alpha &
H^i(\sigma_{\leq i + 1} F(A^\bullet)) \ar[r] \ar[d]^\beta &
H^{i + 1}(\sigma_{\geq i + 2} F(A^\bullet)) \ar[d] \\
H^i(RF(\sigma_{\geq i + 2} A^\bullet)) \ar[r] &
H^i(RF(A^\bullet)) \ar[r] &
H^i(RF(\sigma_{\leq i + 1} A^\bullet)) \ar[r] &
H^{i + 1}(RF(\sigma_{\geq i + 2} A^\bullet))
}
$$
D'après les résultats du premier paragraphe, le morphisme $\beta$ est un
isomorphisme. Par inspection, les objets supérieur gauche et supérieur droit
sont nuls. Pour achever la démonstration, il suffit donc de montrer que
$H^i(RF(\sigma_{\geq i + 2} A^\bullet)) = 0$ et
$H^{i + 1}(RF(\sigma_{\geq i + 2} A^\bullet)) = 0$. Cela résulte
immédiatement du lemme \ref{lemma-negative-vanishing}.
\end{proof}

\begin{proposition}
\label{proposition-enough-acyclics}
\begin{slogan}
Un foncteur entre catégories abéliennes s'étend à la catégorie dérivée (bornée
inférieurement ou supérieurement) en résolvant par un complexe acyclique pour
ce foncteur.
\end{slogan}
Soit $F : \mathcal{A} \to \mathcal{B}$ un foncteur additif de catégories
abéliennes.
\begin{enumerate}
\item Si tout objet de $\mathcal{A}$ s'injecte dans un objet acyclique pour
$RF$, alors $RF$ est défini sur tout $K^{+}(\mathcal{A})$ et l'on obtient un
foncteur exact
$$
RF : D^{+}(\mathcal{A}) \longrightarrow D^{+}(\mathcal{B})
$$
voir (\ref{equation-everywhere}). De plus, tout complexe borné inférieurement
$A^\bullet$ dont les termes sont acycliques pour $RF$ calcule $RF$.
\item Si tout objet de $\mathcal{A}$ est quotient d'un objet acyclique pour
$LF$, alors $LF$ est défini sur tout $K^{-}(\mathcal{A})$ et l'on obtient un
foncteur exact
$$
LF : D^{-}(\mathcal{A}) \longrightarrow D^{-}(\mathcal{B})
$$
voir (\ref{equation-everywhere}). De plus, tout complexe borné supérieurement
$A^\bullet$ dont les termes sont acycliques pour $LF$ calcule $LF$.
\end{enumerate}
\end{proposition}

\begin{proof}
Supposons que tout objet de $\mathcal{A}$ s'injecte dans un objet acyclique
pour $RF$. Soit $\mathcal{I}$ l'ensemble des objets acycliques pour $RF$.
Soit $K^\bullet$ un complexe borné inférieurement de $\mathcal{A}$. D'après le
lemme \ref{lemma-subcategory-right-resolution}, il existe un
quasi-isomorphisme $\alpha : K^\bullet \to I^\bullet$, où $I^\bullet$ est
borné inférieurement et $I^n \in \mathcal{I}$. Ainsi, pour démontrer (1), il
suffit de montrer que $F(I^\bullet) \to F((I')^\bullet)$ est un
quasi-isomorphisme lorsque $s : I^\bullet \to (I')^\bullet$ est un
quasi-isomorphisme de complexes bornés inférieurement d'objets de
$\mathcal{I}$ ; voir le lemme \ref{lemma-find-existence-computes}. Notons que
le cône $C(s)^\bullet$ est un complexe acyclique borné inférieurement dont tous
les termes appartiennent à $\mathcal{I}$. Il suffit donc de montrer que, pour
tout complexe acyclique borné inférieurement $I^\bullet$ dont tous les termes
appartiennent à $\mathcal{I}$, le complexe $F(I^\bullet)$ est acyclique.

\medskip\noindent
Supposons $I^n = 0$ pour $n < n_0$. En posant $J^n = \Im(d^n)$, on décompose
$I^\bullet$ en suites exactes courtes
$0 \to J^n \to I^{n + 1} \to J^{n + 1} \to 0$ pour $n \geq n_0$.
Ces suites induisent des triangles distingués
$(J^n, I^{n + 1}, J^{n + 1})$ dans $D^+(\mathcal{A})$, d'après le
lemme \ref{lemma-derived-canonical-delta-functor}. Pour tout
$k \in \mathbf{Z}$, notons $H_k$ l'assertion suivante : pour tout $n \leq k$,
l'objet $J^n$ appartient à $\mathcal{I}$. Alors $H_k$ est trivialement vraie
pour $k < n_0$. Si $H_n$ est vraie, le
lemme \ref{lemma-2-out-of-3-computes} montre que $J^{n + 1}$ appartient à
$\mathcal{I}$, et l'on obtient $H_{n + 1}$. D'après la
proposition \ref{proposition-derived-functor}, on a un triangle distingué
$(RF(J^n), RF(I^{n + 1}), RF(J^{n + 1}))$. Comme
$J^n, I^{n + 1}, J^{n + 1}$ appartiennent à $\mathcal{I}$, la suite exacte
longue de cohomologie (\ref{equation-long-exact-cohomology-sequence-D})
associée à ce triangle distingué se réduit à une suite exacte
$$
0 \to F(J^n) \to F(I^{n + 1}) \to F(J^{n + 1}) \to 0
$$
Cela montre à son tour que $F(I^\bullet)$ est exact.

\medskip\noindent
La démonstration dans le cas de $LF$ est duale.
\end{proof}

\begin{lemma}
\label{lemma-right-derived-exact-functor}
Soit $F : \mathcal{A} \to \mathcal{B}$ un foncteur exact de catégories
abéliennes. Alors :
\begin{enumerate}
\item tout objet de $\mathcal{A}$ est acyclique à droite pour $F$ ;
\item $RF : D^{+}(\mathcal{A}) \to D^{+}(\mathcal{B})$ est défini partout ;
\item $RF : D(\mathcal{A}) \to D(\mathcal{B})$ est défini partout ;
\item tout complexe calcule $RF$, autrement dit le morphisme canonique
$F(K^\bullet) \to RF(K^\bullet)$ est un isomorphisme pour tout complexe ;
\item $R^iF = 0$ pour $i \not = 0$.
\end{enumerate}
\end{lemma}

\begin{proof}
C'est vrai parce que $F$ transforme les complexes acycliques en complexes
acycliques et les quasi-isomorphismes en quasi-isomorphismes. Les détails sont
omis.
\end{proof}









\section{Sous-catégories triangulées de la catégorie dérivée}
\label{section-triangulated-sub}

\noindent
Soit $\mathcal{A}$ une catégorie abélienne. Dans cette section, nous étudions
certaines sous-catégories triangulées strictement pleines et saturées
$\mathcal{D}' \subset D(\mathcal{A})$.

\medskip\noindent
Soit $\mathcal{B} \subset \mathcal{A}$ une sous-catégorie de Serre faible,
voir Homologie, définition \ref{homology-definition-serre-subcategory} et
lemme \ref{homology-lemma-characterize-weak-serre-subcategory}. Notons
$D_\mathcal{B}(\mathcal{A})$ la sous-catégorie pleine de $D(\mathcal{A})$ dont
les objets sont
$$
\Ob(D_\mathcal{B}(\mathcal{A}))
=
\{X \in \Ob(D(\mathcal{A})) \mid
H^n(X) \text{ est un objet de }\mathcal{B}\text{ pour tout }n\}
$$
Définissons aussi
$D^{+}_\mathcal{B}(\mathcal{A}) =
D^{+}(\mathcal{A}) \cap D_\mathcal{B}(\mathcal{A})$, et de même pour les
autres variantes bornées.

\begin{lemma}
\label{lemma-cohomology-in-serre-subcategory}
Soit $\mathcal{A}$ une catégorie abélienne. Soit
$\mathcal{B} \subset \mathcal{A}$ une sous-catégorie de Serre faible. La
catégorie $D_\mathcal{B}(\mathcal{A})$ est une sous-catégorie triangulée
strictement pleine et saturée de $D(\mathcal{A})$. Il en va de même pour les
variantes bornées.
\end{lemma}

\begin{proof}
Il est clair que $D_\mathcal{B}(\mathcal{A})$ est une sous-catégorie additive
préservée par les foncteurs de translation. Si $X \oplus Y$ appartient à
$D_\mathcal{B}(\mathcal{A})$, alors $H^n(X)$ et $H^n(Y)$ sont des noyaux de
morphismes entre objets de $\mathcal{B}$, puisque
$H^n(X \oplus Y) = H^n(X) \oplus H^n(Y)$. Ainsi, $X$ et $Y$ appartiennent
tous deux à $D_\mathcal{B}(\mathcal{A})$. D'après le
lemme \ref{lemma-triangulated-subcategory}, il suffit donc de montrer que, pour
tout triangle distingué $(X, Y, Z, f, g, h)$ tel que $X$ et $Y$ appartiennent
à $D_\mathcal{B}(\mathcal{A})$, l'objet $Z$ appartient à
$D_\mathcal{B}(\mathcal{A})$. La suite exacte longue de cohomologie
(\ref{equation-long-exact-cohomology-sequence-D}) et la définition d'une
sous-catégorie de Serre faible (voir Homologie, définition
\ref{homology-definition-serre-subcategory}) montrent que $H^n(Z)$ est un
objet de $\mathcal{B}$ pour tout $n$. Ainsi, $Z$ appartient à
$D_\mathcal{B}(\mathcal{A})$.
\end{proof}

\noindent
Nous continuons de supposer que $\mathcal{B}$ est une sous-catégorie de Serre
faible de la catégorie abélienne $\mathcal{A}$. Alors $\mathcal{B}$ est une
catégorie abélienne et le foncteur d'inclusion
$\mathcal{B} \to \mathcal{A}$ est exact. On obtient donc un foncteur dérivé
$D(\mathcal{B}) \to D(\mathcal{A})$, voir le
lemme \ref{lemma-right-derived-exact-functor}. Le foncteur
$D(\mathcal{B}) \to D(\mathcal{A})$ se factorise manifestement par un foncteur
exact canonique
\begin{equation}
\label{equation-compare}
D(\mathcal{B}) \longrightarrow D_\mathcal{B}(\mathcal{A})
\end{equation}
En effet, un complexe formé d'objets de $\mathcal{B}$ donne certainement un
objet de $D_\mathcal{B}(\mathcal{A})$ ; et, comme les triangles distingués de
$D_\mathcal{B}(\mathcal{A})$ sont exactement les triangles distingués de
$D(\mathcal{A})$ dont les sommets appartiennent à
$D_\mathcal{B}(\mathcal{A})$, le foncteur est exact puisque
$D(\mathcal{B}) \to D(\mathcal{A})$ l'est. On obtient de même les foncteurs
$D^+(\mathcal{B}) \to D^+_\mathcal{B}(\mathcal{A})$,
$D^-(\mathcal{B}) \to D^-_\mathcal{B}(\mathcal{A})$ et
$D^b(\mathcal{B}) \to D^b_\mathcal{B}(\mathcal{A})$ pour les variantes
bornées. Dans de nombreux cas, une question essentielle est de savoir si le
foncteur affiché est une équivalence.

\medskip\noindent
Supposons maintenant que $\mathcal{B}$ soit une sous-catégorie de Serre de
$\mathcal{A}$. On dispose alors du foncteur quotient
$\mathcal{A} \to \mathcal{A}/\mathcal{B}$, voir Homologie, lemme
\ref{homology-lemma-serre-subcategory-is-kernel}. Dans ce cas,
$D_\mathcal{B}(\mathcal{A})$ est le noyau du foncteur
$D(\mathcal{A}) \to D(\mathcal{A}/\mathcal{B})$. Le
lemme \ref{lemma-universal-property-quotient} fournit donc un foncteur
canonique
$$
D(\mathcal{A})/D_\mathcal{B}(\mathcal{A})
\longrightarrow
D(\mathcal{A}/\mathcal{B})
$$
Il en va de même pour les variantes bornées.

\begin{lemma}
\label{lemma-derived-of-quotient}
Soit $\mathcal{A}$ une catégorie abélienne. Soit
$\mathcal{B} \subset \mathcal{A}$ une sous-catégorie de Serre. Alors
$D(\mathcal{A}) \to D(\mathcal{A}/\mathcal{B})$ est essentiellement
surjectif.
\end{lemma}

\begin{proof}
Nous utiliserons la description de la catégorie
$\mathcal{A}/\mathcal{B}$ donnée dans la démonstration d'Homologie, lemme
\ref{homology-lemma-serre-subcategory-is-kernel}. Soit
$(X^\bullet, d^\bullet)$ un complexe de $\mathcal{A}/\mathcal{B}$. Cela
signifie que $X^i$ est un objet de $\mathcal{A}$ et que
$d^i : X^i \to X^{i + 1}$ est un morphisme de
$\mathcal{A}/\mathcal{B}$ tel que
$d^i \circ d^{i - 1} = 0$ dans $\mathcal{A}/\mathcal{B}$.

\medskip\noindent
Pour $i \geq 0$, on peut écrire $d^i = (s^i, f^i)$, où
$s^i : Y^i \to X^i$ est un morphisme de $\mathcal{A}$ dont le noyau et le
conoyau appartiennent à $\mathcal{B}$ (de façon équivalente, $s^i$ devient un
isomorphisme dans la catégorie quotient), et où
$f^i : Y^i \to X^{i + 1}$ est un morphisme de $\mathcal{A}$. Par récurrence,
nous construirons un diagramme commutatif
$$
\xymatrix{
& (X')^1 \ar@{..>}[r] & (X')^2 \ar@{..>}[r] & \ldots \\
X^0 \ar@{..>}[ru] &
X^1 \ar@{..>}[u] &
X^2 \ar@{..>}[u] &
\ldots \\
Y^0 \ar[u]_{s^0} \ar[ru]_{f^0} &
Y^1 \ar[u]_{s^1} \ar[ru]_{f^1} &
Y^2 \ar[u]_{s^2} \ar[ru]_{f^2} &
\ldots
}
$$
où les flèches verticales $X^i \to (X')^i$ deviennent des isomorphismes dans
la catégorie quotient. Posons d'abord
$(X')^1 = \Coker(Y^0 \to X^0 \oplus X^1)$ (ou plutôt, prenons le cocarré du
diagramme de flèches $s^0$ et $f^0$), ce qui fournit le premier diagramme
commutatif. Prenons ensuite
$(X')^2 = \Coker(Y^1 \to (X')^1 \oplus X^2)$, et ainsi de suite. En posant en
outre $(X')^n = X^n$ pour $n \leq 0$, on voit que le morphisme
$(X^\bullet, d^\bullet) \to ((X')^\bullet, (d')^\bullet)$ est un isomorphisme
de complexes dans $\mathcal{A}/\mathcal{B}$. On peut donc supposer que
$d^n : X^n \to X^{n + 1}$ est donné par un morphisme
$X^n \to X^{n + 1}$ de $\mathcal{A}$ pour $n \geq 0$.

\medskip\noindent
Dualement, pour $i < 0$, on peut écrire $d^i = (g^i, t^{i + 1})$, où
$t^{i + 1} : X^{i + 1} \to Z^{i + 1}$ est un isomorphisme dans la catégorie
quotient et où $g^i : X^i \to Z^{i + 1}$ est un morphisme. Par récurrence,
nous construirons un diagramme commutatif
$$
\xymatrix{
\ldots &
Z^{-2} &
Z^{-1} &
Z^0 \\
\ldots &
X^{-2} \ar[u]_{t_{-2}} \ar[ru]_{g_{-2}} &
X^{-1} \ar[u]_{t_{-1}} \ar[ru]_{g_{-1}} &
X^0 \ar[u]_{t^0} \\
\ldots &
(X')^{-2} \ar@{..>}[u] \ar@{..>}[r] &
(X')^{-1} \ar@{..>}[u] \ar@{..>}[ru]
}
$$
où les flèches verticales $(X')^i \to X^i$ deviennent des isomorphismes dans
la catégorie quotient. Prenons
$(X')^{-1} = X^{-1} \times_{Z^0} X^0$, puis
$(X')^{-2} = X^{-2} \times_{Z^{-1}} (X')^{-1}$, et ainsi de suite. En posant
en outre $(X')^n = X^n$ pour $n \geq 0$, on voit que le morphisme
$((X')^\bullet, (d')^\bullet) \to (X^\bullet, d^\bullet)$ est un isomorphisme
de complexes dans $\mathcal{A}/\mathcal{B}$. On peut donc supposer que
$d^n : X^n \to X^{n + 1}$ est donné par un morphisme
$d^n : X^n \to X^{n + 1}$ de $\mathcal{A}$ pour tout $n \in \mathbf{Z}$.

\medskip\noindent
Dans ce cas, on sait que les composées $d^n \circ d^{n - 1}$ sont nulles dans
$\mathcal{A}/\mathcal{B}$. Si, pour $n > 0$, on remplace $X^n$ par
$$
(X')^n = X^n/\sum\nolimits_{0 < k \leq n}
\Im(\Im(X^{k - 2} \to X^k) \to X^n)
$$
alors les composées $d^n \circ d^{n - 1}$ sont nulles pour $n \geq 0$.
(Comme dans le deuxième paragraphe ci-dessus, on obtient un isomorphisme de
complexes
$(X^\bullet, d^\bullet) \to ((X')^\bullet, (d')^\bullet)$.) Enfin, pour
$n < 0$, on remplace $X^n$ par
$$
(X')^n = \bigcap\nolimits_{n \leq k < 0}
(X^n \to X^k)^{-1}\Ker(X^k \to X^{k + 2})
$$
et on raisonne de la même manière pour obtenir un complexe de $\mathcal{A}$
dont l'image dans $\mathcal{A}/\mathcal{B}$ est isomorphe au complexe donné.
\end{proof}

\begin{lemma}
\label{lemma-quotient-by-serre-easy}
Soit $\mathcal{A}$ une catégorie abélienne. Soit
$\mathcal{B} \subset \mathcal{A}$ une sous-catégorie de Serre. Supposons que
le foncteur $v : \mathcal{A} \to \mathcal{A}/\mathcal{B}$ possède un adjoint
à gauche $u : \mathcal{A}/\mathcal{B} \to \mathcal{A}$ tel que
$vu \cong \text{id}$. Alors
$$
D(\mathcal{A})/D_\mathcal{B}(\mathcal{A}) = D(\mathcal{A}/\mathcal{B})
$$
et de même pour les variantes bornées.
\end{lemma}

\begin{proof}
Le foncteur $D(v) : D(\mathcal{A}) \to D(\mathcal{A}/\mathcal{B})$ est
essentiellement surjectif d'après le lemme \ref{lemma-derived-of-quotient}.
Pour un objet $X$ de $D(\mathcal{A})$, le morphisme d'adjonction
$c_X : uvX \to X$ est envoyé sur un isomorphisme dans
$D(\mathcal{A}/\mathcal{B})$, puisque $vuv \cong v$ par l'hypothèse
$vu \cong \text{id}$. Ainsi, dans un triangle distingué
$(uvX, X, Z, c_X, g, h)$, l'objet $Z$ appartient à
$D_\mathcal{B}(\mathcal{A})$, comme le montre la suite exacte longue de
cohomologie. Par conséquent, $c_X$ appartient au système multiplicatif utilisé
pour définir la catégorie quotient
$D(\mathcal{A})/D_\mathcal{B}(\mathcal{A})$. Ainsi,
$uvX \cong X$ dans $D(\mathcal{A})/D_\mathcal{B}(\mathcal{A})$.
Pour $X, Y \in \Ob(\mathcal{A})$, le morphisme
$$
\Hom_{D(\mathcal{A})/D_\mathcal{B}(\mathcal{A})}(X, Y)
\longrightarrow
\Hom_{D(\mathcal{A}/\mathcal{B})}(vX, vY)
$$
est bijectif, car $u$ fournit un inverse (d'après les remarques précédentes).
\end{proof}

\noindent
Pour certaines sous-catégories de Serre
$\mathcal{B} \subset \mathcal{A}$, on peut démontrer que le foncteur
$D(\mathcal{B}) \to D_\mathcal{B}(\mathcal{A})$ est pleinement fidèle.

\begin{lemma}
\label{lemma-fully-faithful-embedding}
Soit $\mathcal{A}$ une catégorie abélienne. Soit
$\mathcal{B} \subset \mathcal{A}$ une sous-catégorie de Serre. Supposons que,
pour toute surjection $X \to Y$ avec $X \in \Ob(\mathcal{A})$ et
$Y \in \Ob(\mathcal{B})$, il existe $X' \subset X$,
$X' \in \Ob(\mathcal{B})$, qui se surjecte sur $Y$. Alors le foncteur
$D^-(\mathcal{B}) \to D^-_\mathcal{B}(\mathcal{A})$ de
(\ref{equation-compare}) est une équivalence.
\end{lemma}

\begin{proof}
Soit $X^\bullet$ un complexe borné supérieurement de $\mathcal{A}$ tel que
$H^i(X^\bullet) \in \Ob(\mathcal{B})$ pour tout $i \in \mathbf{Z}$.
Supposons en outre donnés $B^i \subset X^i$,
$B^i \in \Ob(\mathcal{B})$, pour tout $i \in \mathbf{Z}$. Affirmation : il
existe un sous-complexe $Y^\bullet \subset X^\bullet$ tel que
\begin{enumerate}
\item $Y^\bullet \to X^\bullet$ soit un quasi-isomorphisme ;
\item $Y^i \in \Ob(\mathcal{B})$ pour tout $i \in \mathbf{Z}$ ;
\item $B^i \subset Y^i$ pour tout $i \in \mathbf{Z}$.
\end{enumerate}
Pour démontrer l'affirmation, choisissons d'abord, grâce à l'hypothèse du
lemme, des sous-objets
$C^i \subset \Ker(d^i : X^i \to X^{i + 1})$,
$C^i \in \Ob(\mathcal{B})$, se surjectant sur $H^i(X^\bullet)$. En posant
$D^i = C^i + d^{i - 1}(B^{i - 1}) + B^i$, on obtient un sous-complexe
$D^\bullet$ qui satisfait (2) et (3) et tel que
$H^i(D^\bullet) \to H^i(X^\bullet)$ soit surjectif pour tout
$i \in \mathbf{Z}$. Pour tout choix de $E^i \subset X^i$ avec
$E^i \in \Ob(\mathcal{B})$ et
$d^i(E^i) \subset D^{i + 1} + E^{i + 1}$, poser
$Y^i = D^i + E^i$ donne un sous-complexe dont les termes appartiennent à
$\mathcal{B}$ et dont la cohomologie se surjecte sur celle de $X^\bullet$. Il
est clair que, si
$d^i(E^i) = (D^{i + 1} + E^{i + 1}) \cap \Im(d^i)$, le morphisme en
cohomologie est aussi injectif. Pour $n \gg 0$, on peut prendre $E^n$ égal à
$0$. Par récurrence descendante, on peut choisir $E^i$ pour tout $i$ avec la propriété
voulue. En effet, étant donnés $E^{i + 1}, E^{i + 2}, \ldots$, choisissons
$E^i \subset X^i$ tel que
$d^i(E^i) = (D^{i + 1} + E^{i + 1}) \cap \Im(d^i)$. C'est possible par
l'hypothèse du lemme, combinée au fait que
$(D^{i + 1} + E^{i + 1}) \cap \Im(d^i)$ appartient à $\mathcal{B}$, puisque
$\mathcal{B}$ est une sous-catégorie de Serre de $\mathcal{A}$.

\medskip\noindent
L'affirmation ci-dessus implique le lemme. La surjectivité essentielle en
résulte immédiatement. Démontrons la fidélité. Supposons donné un morphisme
$f : U^\bullet \to V^\bullet$ de complexes bornés supérieurement de
$\mathcal{B}$ dont l'image dans $D(\mathcal{A})$ est nulle. Il existe alors un
quasi-isomorphisme $s : V^\bullet \to X^\bullet$ vers un complexe borné
supérieurement de $\mathcal{A}$ tel que $s \circ f$ soit homotope à zéro.
Choisissons une homotopie $h^i : U^i \to X^{i - 1}$ entre $0$ et
$s \circ f$. Appliquons l'affirmation avec
$B^i = h^{i + 1}(U^{i + 1}) + s^i(V^i)$. Le morphisme obtenu
$s' : V^\bullet \to Y^\bullet$ est encore un quasi-isomorphisme, et
$s' \circ f$ est homotope à zéro, comme le montre le fait que $h^i$ se
factorise par $Y^{i - 1}$. Cela démontre la fidélité. La plénitude se démontre
exactement de la même manière.
\end{proof}

\section{Résolutions injectives}
\label{section-injective-resolutions}

\noindent
Dans cette section, nous démontrons quelques lemmes sur l'existence de
résolutions injectives dans les catégories abéliennes ayant assez d'injectifs.

\begin{definition}
\label{definition-injective-resolution}
Soit $\mathcal{A}$ une catégorie abélienne. Soit
$A \in \Ob(\mathcal{A})$. Une {\it résolution injective de $A$} est un
complexe $I^\bullet$ muni d'un morphisme $A \to I^0$ tel que :
\begin{enumerate}
\item On a $I^n = 0$ pour $n < 0$.
\item Chaque $I^n$ est un objet injectif de $\mathcal{A}$.
\item Le morphisme $A \to I^0$ est un isomorphisme sur $\Ker(d^0)$.
\item On a $H^i(I^\bullet) = 0$ pour $i > 0$.
\end{enumerate}
Ainsi, $A[0] \to I^\bullet$ est un quasi-isomorphisme. Autrement dit, le
complexe
$$
\ldots \to 0 \to A \to I^0 \to I^1 \to \ldots
$$
est acyclique. Soit $K^\bullet$ un complexe de $\mathcal{A}$. Une
{\it résolution injective de $K^\bullet$} est un complexe $I^\bullet$ muni
d'un morphisme de complexes $\alpha : K^\bullet \to I^\bullet$ tel que
\begin{enumerate}
\item On a $I^n = 0$ pour $n \ll 0$, c'est-à-dire que $I^\bullet$ est borné
inférieurement.
\item Chaque $I^n$ est un objet injectif de $\mathcal{A}$.
\item Le morphisme $\alpha : K^\bullet \to I^\bullet$ est un
quasi-isomorphisme.
\end{enumerate}
\end{definition}

\noindent
Autrement dit, une résolution injective $K^\bullet \to I^\bullet$ donne lieu
à un diagramme
$$
\xymatrix{
\ldots \ar[r] & K^{n - 1} \ar[d] \ar[r] & K^n \ar[d] \ar[r] &
K^{n + 1} \ar[d] \ar[r] & \ldots \\
\ldots \ar[r] & I^{n - 1} \ar[r] & I^n \ar[r] & I^{n + 1} \ar[r] & \ldots
}
$$
qui induit un isomorphisme sur les objets de cohomologie en chaque degré. Une
résolution injective d'un objet $A$ de $\mathcal{A}$ est presque la même chose
qu'une résolution injective du complexe $A[0]$.

\begin{lemma}
\label{lemma-cohomology-bounded-below}
Soit $\mathcal{A}$ une catégorie abélienne. Soit $K^\bullet$ un complexe de
$\mathcal{A}$.
\begin{enumerate}
\item Si $K^\bullet$ possède une résolution injective, alors
$H^n(K^\bullet) = 0$ pour $n \ll 0$.
\item Si $H^n(K^\bullet) = 0$ pour tout $n \ll 0$, alors il existe un
quasi-isomorphisme $K^\bullet \to L^\bullet$ tel que $L^\bullet$ soit borné
inférieurement.
\end{enumerate}
\end{lemma}

\begin{proof}
La démonstration est omise. Pour la seconde assertion, prendre
$L^\bullet = \tau_{\geq n}K^\bullet$ pour un certain $n \ll 0$. Voir
Homologie, section \ref{homology-section-truncations}, pour la définition de la
troncation $\tau_{\geq n}$.
\end{proof}


\begin{lemma}
\label{lemma-injective-resolutions-exist}
Soit $\mathcal{A}$ une catégorie abélienne. Supposons que $\mathcal{A}$ ait
assez d'injectifs.
\begin{enumerate}
\item Tout objet de $\mathcal{A}$ possède une résolution injective.
\item Si $H^n(K^\bullet) = 0$ pour tout $n \ll 0$, alors $K^\bullet$ possède
une résolution injective.
\item Si $K^\bullet$ est un complexe tel que $K^n = 0$ pour $n < a$, alors il
existe une résolution injective $\alpha : K^\bullet \to I^\bullet$ telle que
$I^n = 0$ pour $n < a$ et que chaque $\alpha^n : K^n \to I^n$ soit injectif.
\end{enumerate}
\end{lemma}

\begin{proof}
Démonstration de (1). Choisissons d'abord une injection $A \to I^0$ de $A$
dans un objet injectif de $\mathcal{A}$. Choisissons ensuite une injection
$I^0/A \to I^1$ dans un objet injectif de $\mathcal{A}$. Notons $d^0$ le
morphisme induit $I^0 \to I^1$. Choisissons ensuite une injection
$I^1/\Im(d^0) \to I^2$ dans un objet injectif de $\mathcal{A}$. Notons $d^1$
le morphisme induit $I^1 \to I^2$, et ainsi de suite. D'après la partie (2) du
lemme \ref{lemma-cohomology-bounded-below}, la partie (2) résulte de la partie
(3). La partie (3) est un cas particulier du
lemme \ref{lemma-subcategory-right-resolution}.
\end{proof}

\begin{lemma}
\label{lemma-acyclic-is-zero}
Soit $\mathcal{A}$ une catégorie abélienne. Soit $K^\bullet$ un complexe
acyclique. Soit $I^\bullet$ un complexe borné inférieurement dont les termes
sont injectifs. Tout morphisme $K^\bullet \to I^\bullet$ est homotope à zéro.
\end{lemma}

\begin{proof}
Soit $\alpha : K^\bullet \to I^\bullet$ un morphisme de complexes. Remarquons
que $\alpha^j = 0$ pour $j \ll 0$, puisque $I^\bullet$ est borné
inférieurement. On peut donc trouver un $n$ tel qu'il existe des morphismes
$h^j :  K^j \to I^{j - 1}$ pour $j \leq n$ vérifiant
$\alpha^j = d^{j - 1} \circ h^j + h^{j + 1} \circ d^j$ pour $j < n$. Nous
allons montrer qu'il existe un morphisme
$h^{n + 1} : K^{n + 1} \to I^n$ tel que
$\alpha^n = d^{n - 1} \circ h^n + h^{n + 1} \circ d^n$. Remarquons que
\begin{align*}
(\alpha^n - d^{n - 1} \circ h^n) \circ d^{n - 1}
& =
\alpha^{n - 1} \circ d^{n - 1} - d^{n - 1} \circ h^n \circ d^{n - 1} \\
& =
d^{n - 1} \circ \alpha^{n - 1} - d^{n - 1} \circ h^n \circ d^{n - 1} \\
& =
d^{n - 1} \circ (d^{n - 2} \circ h^{n - 1} + h^n \circ d^{n - 1})
- d^{n - 1} \circ h^n \circ d^{n - 1} \\
& = 0
\end{align*}
Puisque $K^\bullet$ est acyclique, on a
$d^{n - 1}(K^{n - 1}) = \Ker(K^n \to K^{n + 1})$. On peut donc considérer
$\alpha^n -  d^{n - 1} \circ h^n$ comme un morphisme à valeurs dans $I^n$,
défini sur le sous-objet $\Im(K^n \to K^{n + 1})$ de $K^{n + 1}$. Comme
l'objet $I^n$ est injectif, on peut le prolonger en un morphisme
$h^{n + 1} : K^{n + 1} \to I^n$. Avec ce choix, le lecteur vérifie que l'on a
bien $\alpha^n = d^{n - 1} \circ h^n + h^{n + 1} \circ d^n$.

\medskip\noindent
Par récurrence sur $n$, on conclut que l'on peut trouver
$h  = (h^j)_{j \in \mathbf{Z}}$ qui constitue une homotopie entre $\alpha$ et
$0$, comme voulu.
\end{proof}

\begin{remark}
\label{remark-easier-proofs}
Soit $\mathcal{A}$ une catégorie abélienne. En utilisant le fait que
$K(\mathcal{A})$ est une catégorie triangulée, le
lemme \ref{lemma-acyclic-is-zero} permet d'obtenir des démonstrations de
certains des lemmes ci-dessous qui se prouvent habituellement par chasse au
diagramme. Plus précisément, supposons que
$\alpha : K^\bullet \to L^\bullet$ soit un quasi-isomorphisme de complexes.
Alors
$$
(K^\bullet, L^\bullet, C(\alpha)^\bullet, \alpha, i, -p)
$$
est un triangle distingué de $K(\mathcal{A})$
(lemme \ref{lemma-the-same-up-to-isomorphisms}), et $C(\alpha)^\bullet$ est un
complexe acyclique (lemme \ref{lemma-acyclic}). Soit ensuite $I^\bullet$ un
complexe borné inférieurement d'objets injectifs. Alors
$$
\xymatrix{
\Hom_{K(\mathcal{A})}(C(\alpha)^\bullet, I^\bullet) \ar[r] &
\Hom_{K(\mathcal{A})}(L^\bullet, I^\bullet) \ar[r] &
\Hom_{K(\mathcal{A})}(K^\bullet, I^\bullet) \ar[lld] \\
\Hom_{K(\mathcal{A})}(C(\alpha)^\bullet[-1], I^\bullet)
}
$$
est une suite exacte de groupes abéliens, voir le
lemme \ref{lemma-representable-homological}. Le
lemme \ref{lemma-acyclic-is-zero} garantit alors que les deux groupes
extrêmes sont nuls, et donc que
$\Hom_{K(\mathcal{A})}(L^\bullet, I^\bullet) =
\Hom_{K(\mathcal{A})}(K^\bullet, I^\bullet)$.
\end{remark}

\begin{lemma}
\label{lemma-morphisms-lift}
Soit $\mathcal{A}$ une catégorie abélienne. Considérons un diagramme dont les
flèches pleines sont données
$$
\xymatrix{
K^\bullet \ar[r]_\alpha \ar[d]_\gamma & L^\bullet \ar@{-->}[dl]^\beta \\
I^\bullet
}
$$
où $I^\bullet$ est borné inférieurement et formé d'objets injectifs, et où
$\alpha$ est un quasi-isomorphisme.
\begin{enumerate}
\item Il existe un morphisme de complexes $\beta$ qui rend le diagramme
commutatif à homotopie près.
\item Si $\alpha$ est injectif en chaque degré, on peut trouver un $\beta$ qui
rend le diagramme commutatif.
\end{enumerate}
\end{lemma}

\begin{proof}
La démonstration « correcte » de la partie (1) est expliquée dans la
remarque \ref{remark-easier-proofs}. Nous donnons aussi ici une démonstration
directe.

\medskip\noindent
Montrons d'abord que (2) implique (1). Soient
$\tilde \alpha : K^\bullet \to \tilde L^\bullet$, $\pi$, $s$ comme dans le
lemme \ref{lemma-make-injective}. Puisque $\tilde \alpha$ est injectif, la
partie (2) fournit un morphisme
$\tilde \beta : \tilde L^\bullet \to I^\bullet$ tel que
$\gamma = \tilde \beta \circ \tilde \alpha$. Posons
$\beta = \tilde \beta \circ s$. On a alors
$$
\beta \circ \alpha
=
\tilde \beta \circ s \circ \pi \circ \tilde \alpha
\sim
\tilde \beta \circ \tilde \alpha
=
\gamma
$$
comme voulu.

\medskip\noindent
Supposons que $\alpha : K^\bullet \to L^\bullet$ soit injectif. Supposons
$\beta$ déjà défini en tous les degrés $\leq n - 1$, de manière compatible
avec les différentielles et de sorte que
$\gamma^j = \beta^j \circ \alpha^j$ pour tout $j \leq n - 1$. Considérons le
diagramme commutatif dont les flèches pleines sont données
$$
\xymatrix{
K^{n - 1} \ar[r] \ar@/_2pc/[dd]_\gamma \ar[d]^\alpha &
K^n \ar@/^2pc/[dd]^\gamma \ar[d]^\alpha \\
L^{n - 1} \ar[r] \ar[d]^\beta &
L^n \ar@{-->}[d] \\
I^{n - 1} \ar[r] &
I^n
}
$$
La flèche en pointillés est donc prescrite sur les sous-objets
$\alpha(K^n)$ et $d^{n - 1}(L^{n - 1})$. De plus, ces deux morphismes
coïncident sur $\alpha(d^{n - 1}(K^{n - 1}))$. Par conséquent, si
\begin{equation}
\label{equation-qis}
\alpha(d^{n - 1}(K^{n - 1}))
=
\alpha(K^n) \cap d^{n - 1}(L^{n - 1})
\end{equation}
alors ces morphismes se recollent en un morphisme
$\alpha(K^n) + d^{n - 1}(L^{n - 1}) \to I^n$ et, par injectivité de $I^n$,
on peut le prolonger en un morphisme défini sur tout $L^n$ à valeurs dans
$I^n$. Une récurrence fournit alors le morphisme $\beta$ simultanément pour
tout $n$ (on peut poser $\beta^n = 0$ pour tout $n \ll 0$, puisque
$I^\bullet$ est borné inférieurement, ce qui amorce la récurrence).

\medskip\noindent
Il reste à démontrer l'égalité (\ref{equation-qis}). Le lecteur est invité à le
faire lui-même par une chasse au diagramme appropriée. Voici néanmoins notre
argument. L'inclusion
$\alpha(d^{n - 1}(K^{n - 1})) \subset \alpha(K^n) \cap
d^{n - 1}(L^{n - 1})$ est évidente. Soit $T$ un objet de $\mathcal{A}$ et soit
$x : T \to L^n$ un morphisme dont l'image est contenue dans le sous-objet
$\alpha(K^n) \cap d^{n - 1}(L^{n - 1})$. Puisque $\alpha$ est injectif, on a
$x = \alpha \circ x'$ pour un certain $x' : T \to K^n$. En outre, comme $x$
appartient à $d^{n - 1}(L^{n - 1})$, on a $d^n \circ x = 0$. En utilisant de
nouveau l'injectivité de $\alpha$, on obtient donc $d^n \circ x' = 0$. Ainsi,
$x'$ fournit un morphisme $[x'] : T \to H^n(K^\bullet)$. D'autre part, le
morphisme correspondant $[x] : T \to H^n(L^\bullet)$ induit par $x$ est nul
par hypothèse. Comme $\alpha$ est un quasi-isomorphisme, on conclut que
$[x'] = 0$. Cela signifie exactement que l'image de $x'$ est contenue dans
$d^{n - 1}(K^{n - 1})$, ce qui conclut.
\end{proof}

\begin{lemma}
\label{lemma-morphisms-equal-up-to-homotopy}
Soit $\mathcal{A}$ une catégorie abélienne. Considérons un diagramme dont les
flèches pleines sont données
$$
\xymatrix{
K^\bullet \ar[r]_\alpha \ar[d]_\gamma & L^\bullet \ar@{-->}[dl]^{\beta_i} \\
I^\bullet
}
$$
où $I^\bullet$ est borné inférieurement et formé d'objets injectifs, et où
$\alpha$ est un quasi-isomorphisme. Deux morphismes quelconques
$\beta_1, \beta_2$ qui rendent le diagramme commutatif à homotopie près sont
homotopes.
\end{lemma}

\begin{proof}
Cela résulte de la remarque \ref{remark-easier-proofs}. Nous donnons également
un argument direct.

\medskip\noindent
Soient $\tilde \alpha : K^\bullet \to \tilde L^\bullet$, $\pi$, $s$ comme
dans le lemme \ref{lemma-make-injective}. Si l'on peut montrer que
$\beta_1 \circ\pi$ est homotope à $\beta_2 \circ \pi$, alors on en déduit que
$\beta_1 \sim \beta_2$, puisque $\pi \circ s$ est l'identité. On peut donc
supposer que $\alpha^n : K^n \to L^n$ est l'inclusion d'un facteur direct pour
tout $n$. On obtient ainsi une suite exacte courte de complexes
$$
0 \to K^\bullet \to L^\bullet \to M^\bullet \to 0
$$
scindée terme à terme, telle que $M^\bullet$ soit acyclique. Choisissons des
scindages $L^n = K^n \oplus M^n$ ; on a donc
$\beta_i^n : K^n \oplus M^n \to I^n$ et $\gamma^n : K^n \to I^n$. Dans ce
cas, la condition imposée à $\beta_i$ signifie qu'il existe des morphismes
$h_i^n : K^n \to I^{n - 1}$ tels que
$$
\gamma^n - \beta_i^n|_{K^n} = d \circ h_i^n + h_i^{n + 1} \circ d
$$
On voit donc que
$$
\beta_1^n|_{K^n} - \beta_2^n|_{K^n}
=
d \circ (h_1^n - h_2^n) + (h_1^{n + 1} - h_2^{n + 1}) \circ d
$$
Considérons le morphisme $h^n : K^n \oplus M^n \to I^{n - 1}$ égal à
$h_1^n - h_2^n$ sur le premier facteur et nul sur le second. Alors
$$
\beta_1^n - \beta_2^n
-
(d \circ h^n + h^{n + 1} \circ d)
$$
est un morphisme de complexes $L^\bullet \to I^\bullet$ identiquement nul sur
le sous-complexe $K^\bullet$. Il se factorise donc en
$L^\bullet \to M^\bullet \to I^\bullet$. Le résultat du lemme découle alors
du lemme \ref{lemma-acyclic-is-zero}.
\end{proof}

\begin{lemma}
\label{lemma-morphisms-into-injective-complex}
Soit $\mathcal{A}$ une catégorie abélienne. Soit $I^\bullet$ un complexe borné
inférieurement dont les termes sont injectifs. Soit
$L^\bullet \in K(\mathcal{A})$. Alors
$$
\Mor_{K(\mathcal{A})}(L^\bullet, I^\bullet)
=
\Mor_{D(\mathcal{A})}(L^\bullet, I^\bullet).
$$
\end{lemma}

\begin{proof}
Soit $a$ un élément du membre de droite. On peut représenter
$a = \gamma\alpha^{-1}$, où $\alpha : K^\bullet \to L^\bullet$ est un
quasi-isomorphisme et $\gamma : K^\bullet \to I^\bullet$ un morphisme de
complexes. D'après le lemme \ref{lemma-morphisms-lift}, on peut trouver un
morphisme $\beta : L^\bullet \to I^\bullet$ tel que
$\beta \circ \alpha$ soit homotope à $\gamma$. Cela démontre la surjectivité
du morphisme. Soit $b$ un élément du membre de gauche dont l'image dans le
membre de droite est nulle. Alors $b$ est la classe d'homotopie d'un morphisme
$\beta : L^\bullet \to I^\bullet$ tel qu'il existe un quasi-isomorphisme
$\alpha : K^\bullet \to L^\bullet$ pour lequel $\beta  \circ \alpha$ est
homotope à zéro. Le lemme \ref{lemma-morphisms-equal-up-to-homotopy} montre
alors que $\beta$ est lui aussi homotope à zéro, c'est-à-dire que $b = 0$.
\end{proof}

\begin{lemma}
\label{lemma-injective-resolution-ses}
Soit $\mathcal{A}$ une catégorie abélienne. Supposons que $\mathcal{A}$ ait
assez d'injectifs. Pour toute suite exacte courte
$0 \to A^\bullet \to B^\bullet \to C^\bullet \to 0$ de
$\text{Comp}^{+}(\mathcal{A})$, il existe dans
$\text{Comp}^{+}(\mathcal{A})$ un diagramme commutatif
$$
\xymatrix{
0 \ar[r] &
A^\bullet \ar[r] \ar[d] &
B^\bullet \ar[r] \ar[d] &
C^\bullet \ar[r] \ar[d] &
0 \\
0 \ar[r] &
I_1^\bullet \ar[r] &
I_2^\bullet \ar[r] &
I_3^\bullet \ar[r] &
0
}
$$
où les flèches verticales sont des résolutions injectives et les lignes des
suites exactes courtes de complexes. En outre,
\begin{enumerate}
\item étant donnée une résolution injective quelconque
$A^\bullet \to I^\bullet$, on peut supposer que
$I_1^\bullet = I^\bullet$ ;
\item si $A^n = B^n = C^n = 0$ pour $n < 0$, on peut supposer que
$I_j^n = 0$ pour $n < 0$ ;
\item on peut combiner (1) et (2) si l'on a également $I^n = 0$ pour $n < 0$.
\end{enumerate}
\end{lemma}

\begin{proof}
Étape 1. Choisissons une résolution injective
$A^\bullet \to I^\bullet$ (voir le
lemme \ref{lemma-injective-resolutions-exist}) ou utilisons celle qui est
donnée. Rappelons que $\text{Comp}^{+}(\mathcal{A})$ est une catégorie
abélienne, voir Homologie, lemme
\ref{homology-lemma-cat-cochain-abelian}. On peut donc former le cocarré le
long du morphisme $A^\bullet \to I^\bullet$ et obtenir
$$
\xymatrix{
0 \ar[r] &
A^\bullet \ar[r] \ar[d] &
B^\bullet \ar[r] \ar[d] &
C^\bullet \ar[r] \ar[d] &
0 \\
0 \ar[r] &
I^\bullet \ar[r] &
E^\bullet \ar[r] &
C^\bullet \ar[r] &
0
}
$$
D'après le lemme $5$ et la dernière assertion d'Homologie, lemme
\ref{homology-lemma-long-exact-sequence-cochain}, le morphisme
$B^\bullet \to E^\bullet$ est un quasi-isomorphisme. Remarquons que la suite
exacte courte inférieure est scindée terme à terme, voir Homologie, lemme
\ref{homology-lemma-characterize-injectives}. Il suffit donc de démontrer le
lemme lorsque $0 \to A^\bullet \to B^\bullet \to C^\bullet \to 0$ est scindée
terme à terme.

\medskip\noindent
Étape 2. Choisissons des scindages. Autrement dit, écrivons
$B^n = A^n \oplus C^n$. Notons $\delta : C^\bullet \to A^\bullet[1]$ le
morphisme d'Homologie, lemme
\ref{homology-lemma-ses-termwise-split-cochain}. Choisissons des résolutions
injectives $f_1 : A^\bullet \to I_1^\bullet$ et
$f_3 : C^\bullet \to I_3^\bullet$. (Si $A^\bullet$ est un complexe
d'injectifs, prendre $I_1^\bullet = A^\bullet$. Si $A^n = C^n = 0$ pour
$n < 0$, les choisir de sorte que $I_1^n = I_3^n = 0$ pour $n < 0$, voir le
lemme \ref{lemma-injective-resolutions-exist}.) On peut supposer que $f_3$ est
injectif en chaque degré. D'après le lemme \ref{lemma-morphisms-lift}, on peut
trouver un morphisme $\delta' : I_3^\bullet \to I_1^\bullet[1]$ tel que
$\delta' \circ f_3 = f_1[1] \circ \delta$ (égalité de morphismes de complexes).
Posons $I_2^n = I_1^n \oplus I_3^n$. Définissons
$$
d_{I_2}^n =
\left(
\begin{matrix}
d_{I_1}^n & (\delta')^n \\
0 & d_{I_3}^n
\end{matrix}
\right)
$$
et définissons les morphismes $B^n \to I_2^n$ comme la somme des morphismes
$A^n \to I_1^n$ et $C^n \to I_3^n$. Tout est clair.
\end{proof}


\section{Résolutions projectives}
\label{section-projective-resolutions}

\noindent
Cette section est duale de la section \ref{section-injective-resolutions}.
Nous donnons les définitions et énonçons les résultats, sans redémontrer les
lemmes.

\begin{definition}
\label{definition-projective-resolution}
Soit $\mathcal{A}$ une catégorie abélienne. Soit
$A \in \Ob(\mathcal{A})$. Une {\it résolution projective de $A$} est un
complexe $P^\bullet$ muni d'un morphisme $P^0 \to A$ tel que :
\begin{enumerate}
\item On a $P^n = 0$ pour $n > 0$.
\item Chaque $P^n$ est un objet projectif de $\mathcal{A}$.
\item Le morphisme $P^0 \to A$ induit un isomorphisme
$\Coker(d^{-1}) \to A$.
\item On a $H^i(P^\bullet) = 0$ pour $i < 0$.
\end{enumerate}
Ainsi, $P^\bullet \to A[0]$ est un quasi-isomorphisme. Autrement dit, le
complexe
$$
\ldots \to P^{-1} \to P^0 \to A \to 0 \to \ldots
$$
est acyclique. Soit $K^\bullet$ un complexe de $\mathcal{A}$. Une
{\it résolution projective de $K^\bullet$} est un complexe $P^\bullet$ muni
d'un morphisme de complexes $\alpha : P^\bullet \to K^\bullet$ tel que
\begin{enumerate}
\item On a $P^n = 0$ pour $n \gg 0$, c'est-à-dire que $P^\bullet$ est borné
supérieurement.
\item Chaque $P^n$ est un objet projectif de $\mathcal{A}$.
\item Le morphisme $\alpha : P^\bullet \to K^\bullet$ est un
quasi-isomorphisme.
\end{enumerate}
\end{definition}

\begin{lemma}
\label{lemma-cohomology-bounded-above}
Soit $\mathcal{A}$ une catégorie abélienne. Soit $K^\bullet$ un complexe de
$\mathcal{A}$.
\begin{enumerate}
\item Si $K^\bullet$ possède une résolution projective, alors
$H^n(K^\bullet) = 0$ pour $n \gg 0$.
\item Si $H^n(K^\bullet) = 0$ pour $n \gg 0$, alors il existe un
quasi-isomorphisme $L^\bullet \to K^\bullet$ tel que $L^\bullet$ soit borné
supérieurement.
\end{enumerate}
\end{lemma}

\begin{proof}
Cela est dual du lemme \ref{lemma-cohomology-bounded-below}.
\end{proof}

\begin{lemma}
\label{lemma-projective-resolutions-exist}
Soit $\mathcal{A}$ une catégorie abélienne. Supposons que $\mathcal{A}$ ait
assez de projectifs.
\begin{enumerate}
\item Tout objet de $\mathcal{A}$ possède une résolution projective.
\item Si $H^n(K^\bullet) = 0$ pour tout $n \gg 0$, alors $K^\bullet$ possède
une résolution projective.
\item Si $K^\bullet$ est un complexe tel que $K^n = 0$ pour $n > a$, alors il
existe une résolution projective $\alpha : P^\bullet \to K^\bullet$ telle que
$P^n = 0$ pour $n > a$ et que chaque $\alpha^n : P^n \to K^n$ soit surjectif.
\end{enumerate}
\end{lemma}

\begin{proof}
Cela est dual du lemme \ref{lemma-injective-resolutions-exist}.
\end{proof}

\begin{lemma}
\label{lemma-projective-into-acyclic-is-zero}
Soit $\mathcal{A}$ une catégorie abélienne. Soit $K^\bullet$ un complexe
acyclique. Soit $P^\bullet$ un complexe borné supérieurement dont les termes
sont projectifs. Tout morphisme $P^\bullet \to K^\bullet$ est homotope à zéro.
\end{lemma}

\begin{proof}
Cela est dual du lemme \ref{lemma-acyclic-is-zero}.
\end{proof}

\begin{remark}
\label{remark-easier-projective}
Soit $\mathcal{A}$ une catégorie abélienne. Supposons que
$\alpha : K^\bullet \to L^\bullet$ soit un quasi-isomorphisme de complexes.
Soit $P^\bullet$ un complexe borné supérieurement d'objets projectifs. Alors
$$
\Hom_{K(\mathcal{A})}(P^\bullet, K^\bullet)
\longrightarrow
\Hom_{K(\mathcal{A})}(P^\bullet, L^\bullet)
$$
est un isomorphisme. Cela est dual de la remarque
\ref{remark-easier-proofs}.
\end{remark}

\begin{lemma}
\label{lemma-morphisms-lift-projective}
Soit $\mathcal{A}$ une catégorie abélienne. Considérons un diagramme dont les
flèches pleines sont données
$$
\xymatrix{
K^\bullet & L^\bullet \ar[l]^\alpha \\
P^\bullet \ar[u] \ar@{-->}[ru]_\beta
}
$$
où $P^\bullet$ est borné supérieurement et formé d'objets projectifs, et où
$\alpha$ est un quasi-isomorphisme.
\begin{enumerate}
\item Il existe un morphisme de complexes $\beta$ qui rend le diagramme
commutatif à homotopie près.
\item Si $\alpha$ est surjectif en chaque degré, on peut trouver un $\beta$
qui rend le diagramme commutatif.
\end{enumerate}
\end{lemma}

\begin{proof}
Cela est dual du lemme \ref{lemma-morphisms-lift}.
\end{proof}

\begin{lemma}
\label{lemma-morphisms-equal-up-to-homotopy-projective}
Soit $\mathcal{A}$ une catégorie abélienne. Considérons un diagramme dont les
flèches pleines sont données
$$
\xymatrix{
K^\bullet & L^\bullet \ar[l]^\alpha \\
P^\bullet \ar[u] \ar@{-->}[ru]_{\beta_i}
}
$$
où $P^\bullet$ est borné supérieurement et formé d'objets projectifs, et où
$\alpha$ est un quasi-isomorphisme. Deux morphismes quelconques
$\beta_1, \beta_2$ qui rendent le diagramme commutatif à homotopie près sont
homotopes.
\end{lemma}

\begin{proof}
Cela est dual du lemme \ref{lemma-morphisms-equal-up-to-homotopy}.
\end{proof}

\begin{lemma}
\label{lemma-morphisms-from-projective-complex}
Soit $\mathcal{A}$ une catégorie abélienne. Soit $P^\bullet$ un complexe borné
supérieurement dont les termes sont projectifs. Soit
$L^\bullet \in K(\mathcal{A})$. Alors
$$
\Mor_{K(\mathcal{A})}(P^\bullet, L^\bullet)
=
\Mor_{D(\mathcal{A})}(P^\bullet, L^\bullet).
$$
\end{lemma}

\begin{proof}
Cela est dual du lemme \ref{lemma-morphisms-into-injective-complex}.
\end{proof}

\begin{lemma}
\label{lemma-projective-resolution-ses}
Soit $\mathcal{A}$ une catégorie abélienne. Supposons que $\mathcal{A}$ ait
assez de projectifs. Pour toute suite exacte courte
$0 \to A^\bullet \to B^\bullet \to C^\bullet \to 0$ de
$\text{Comp}^{-}(\mathcal{A})$, il existe dans
$\text{Comp}^{-}(\mathcal{A})$ un diagramme commutatif
$$
\xymatrix{
0 \ar[r] &
P_1^\bullet \ar[r] \ar[d] &
P_2^\bullet \ar[r] \ar[d] &
P_3^\bullet \ar[r] \ar[d] &
0 \\
0 \ar[r] &
A^\bullet \ar[r] &
B^\bullet \ar[r] &
C^\bullet \ar[r] &
0
}
$$
où les flèches verticales sont des résolutions projectives et les lignes des
suites exactes courtes de complexes. En fait, étant donnée une résolution
projective quelconque $P^\bullet \to C^\bullet$, on peut supposer que
$P_3^\bullet = P^\bullet$.
\end{lemma}

\begin{proof}
Cela est dual du lemme \ref{lemma-injective-resolution-ses}.
\end{proof}

\begin{lemma}
\label{lemma-precise-vanishing}
Soit $\mathcal{A}$ une catégorie abélienne. Soient $P^\bullet$,
$K^\bullet$ des complexes. Soit $n \in \mathbf{Z}$. Supposons que
\begin{enumerate}
\item $P^\bullet$ soit un complexe borné dont les termes sont projectifs ;
\item $P^i = 0$ pour $i < n$ ;
\item $H^i(K^\bullet) = 0$ pour $i \geq n$.
\end{enumerate}
Alors
$\Hom_{K(\mathcal{A})}(P^\bullet, K^\bullet) =
\Hom_{D(\mathcal{A})}(P^\bullet, K^\bullet) = 0$.
\end{lemma}

\begin{proof}
La première égalité résulte du
lemme \ref{lemma-morphisms-from-projective-complex}. Remarquons qu'il existe un
triangle distingué
$$
(\tau_{\leq n - 1}K^\bullet, K^\bullet, \tau_{\geq n}K^\bullet, f, g, h)
$$
d'après la remarque \ref{remark-truncation-distinguished-triangle}. Par le
lemme \ref{lemma-representable-homological}, il suffit donc de démontrer
$\Hom_{K(\mathcal{A})}(P^\bullet, \tau_{\leq n - 1}K^\bullet) = 0$ et
$\Hom_{K(\mathcal{A})}(P^\bullet, \tau_{\geq n} K^\bullet) = 0$. La première
annulation est évidente, et la seconde résulte du
lemme \ref{lemma-projective-into-acyclic-is-zero}.
\end{proof}

\begin{lemma}
\label{lemma-lift-map}
Soit $\mathcal{A}$ une catégorie abélienne. Soient
$\beta : P^\bullet \to L^\bullet$ et $\alpha : E^\bullet \to L^\bullet$ des
morphismes de complexes. Soit $n \in \mathbf{Z}$. Supposons que
\begin{enumerate}
\item $P^\bullet$ soit un complexe borné d'objets projectifs et que
$P^i = 0$ pour $i < n$ ;
\item $H^i(\alpha)$ soit un isomorphisme pour $i > n$ et surjectif pour
$i = n$.
\end{enumerate}
Alors il existe un morphisme de complexes
$\gamma : P^\bullet \to E^\bullet$ tel que $\alpha \circ \gamma$ et $\beta$
soient homotopes.
\end{lemma}

\begin{proof}
Considérons le cône $C^\bullet = C(\alpha)^\bullet$, muni du morphisme
$i : L^\bullet \to C^\bullet$. D'après le
lemme \ref{lemma-precise-vanishing}, $i \circ \beta$ est nul. Le
lemme \ref{lemma-representable-homological} permet donc de relever $\beta$ à
$E^\bullet$.
\end{proof}
\section{Foncteurs dérivés droits et résolutions injectives}
\label{section-right-derived-functor}

\noindent
Nous pouvons maintenant utiliser les résultats précédents pour définir les
foncteurs dérivés droits d'un foncteur additif d'une catégorie abélienne ayant
assez d'injectifs vers une catégorie abélienne quelconque.

\begin{lemma}
\label{lemma-injective-acyclic}
Soit $\mathcal{A}$ une catégorie abélienne.
Soit $I \in \Ob(\mathcal{A})$ un objet injectif.
Soit $I^\bullet$ un complexe borné inférieurement d'objets injectifs de
$\mathcal{A}$.
\begin{enumerate}
\item $I^\bullet$ calcule $RF$ relativement à
$\text{Qis}^{+}(\mathcal{A})$ pour tout foncteur exact
$F : K^{+}(\mathcal{A}) \to \mathcal{D}$ à valeurs dans une catégorie
triangulée quelconque $\mathcal{D}$.
\item $I$ est acyclique à droite pour tout foncteur additif
$F : \mathcal{A} \to \mathcal{B}$ à valeurs dans une catégorie abélienne
quelconque $\mathcal{B}$.
\end{enumerate}
\end{lemma}

\begin{proof}
La partie (2) est une conséquence directe de la partie (1) et de la
définition \ref{definition-derived-functor}.
Pour prouver (1), soit $\alpha : I^\bullet \to K^\bullet$ un
quasi-isomorphisme vers un complexe. D'après le
lemme \ref{lemma-morphisms-lift}, $\alpha$ admet un inverse à gauche. La
catégorie $I^\bullet/\text{Qis}^{+}(\mathcal{A})$ est donc essentiellement
constante, de valeur $\text{id} : I^\bullet \to I^\bullet$. Par suite,
l'ind-objet
$$
I^\bullet/\text{Qis}^{+}(\mathcal{A}) \longrightarrow \mathcal{D}, \quad
(I^\bullet \to K^\bullet) \longmapsto F(K^\bullet)
$$
est lui aussi essentiellement constant, de valeur $F(I^\bullet)$. Cela prouve
(1), voir les définitions
\ref{definition-right-derived-functor-defined} et \ref{definition-computes}.
\end{proof}

\begin{lemma}
\label{lemma-enough-injectives-right-derived}
Soit $\mathcal{A}$ une catégorie abélienne ayant assez d'injectifs.
\begin{enumerate}
\item Pour tout foncteur exact $F : K^{+}(\mathcal{A}) \to \mathcal{D}$ à
valeurs dans une catégorie triangulée $\mathcal{D}$, le foncteur dérivé
droit
$$
RF : D^{+}(\mathcal{A}) \longrightarrow \mathcal{D}
$$
est défini partout.
\item Pour tout foncteur additif $F : \mathcal{A} \to \mathcal{B}$ à
valeurs dans une catégorie abélienne $\mathcal{B}$, le foncteur dérivé
droit
$$
RF : D^{+}(\mathcal{A}) \longrightarrow D^{+}(\mathcal{B})
$$
est défini partout.
\end{enumerate}
\end{lemma}

\begin{proof}
Pour la seconde assertion, combiner le
lemme \ref{lemma-injective-acyclic} et la
proposition \ref{proposition-enough-acyclics}. Pour la première, combiner le
lemme \ref{lemma-injective-resolutions-exist}, le
lemme \ref{lemma-injective-acyclic}, le
lemme \ref{lemma-existence-computes} et l'équation
(\ref{equation-everywhere}).
\end{proof}

\begin{lemma}
\label{lemma-right-derived-properties}
Soit $\mathcal{A}$ une catégorie abélienne ayant assez d'injectifs.
Soit $F : \mathcal{A} \to \mathcal{B}$ un foncteur additif.
\begin{enumerate}
\item Le foncteur $RF$ est un foncteur exact
$D^{+}(\mathcal{A}) \to D^{+}(\mathcal{B})$.
\item Le foncteur $RF$ induit un foncteur exact
$K^{+}(\mathcal{A}) \to D^{+}(\mathcal{B})$.
\item Le foncteur $RF$ induit un $\delta$-foncteur
$\text{Comp}^{+}(\mathcal{A}) \to D^{+}(\mathcal{B})$.
\item Le foncteur $RF$ induit un $\delta$-foncteur
$\mathcal{A} \to D^{+}(\mathcal{B})$.
\end{enumerate}
\end{lemma}

\begin{proof}
Ce lemme ne fait que récapituler certains des résultats obtenus jusqu'ici.
Notons que, d'après le lemme
\ref{lemma-enough-injectives-right-derived}, $RF$ est défini partout. Voici
quelques références :
\begin{enumerate}
\item Le foncteur dérivé est exact : cela se ramène au
lemme \ref{lemma-2-out-of-3-defined}.
\item Cela résulte du fait que
$K^{+}(\mathcal{A}) \to D^{+}(\mathcal{A})$ est exact et qu'un composé de
foncteurs exacts est exact.
\item Cela résulte du fait que
$\text{Comp}^{+}(\mathcal{A}) \to D^{+}(\mathcal{A})$ est un
$\delta$-foncteur, voir le
lemme \ref{lemma-derived-canonical-delta-functor}.
\item Cela résulte du fait que
$\mathcal{A} \to \text{Comp}^{+}(\mathcal{A})$ est exact et que la
précomposition d'un $\delta$-foncteur par un foncteur exact donne un
$\delta$-foncteur.
\end{enumerate}
\end{proof}

\begin{lemma}
\label{lemma-higher-derived-functors}
Soit $\mathcal{A}$ une catégorie abélienne ayant assez d'injectifs.
Soit $F : \mathcal{A} \to \mathcal{B}$ un foncteur exact à gauche.
\begin{enumerate}
\item À toute suite exacte courte
$0 \to A^\bullet \to B^\bullet \to C^\bullet \to 0$ de complexes dans
$\text{Comp}^{+}(\mathcal{A})$ est associée une suite exacte longue
$$
\ldots \to
H^i(RF(A^\bullet)) \to
H^i(RF(B^\bullet)) \to
H^i(RF(C^\bullet)) \to
H^{i + 1}(RF(A^\bullet)) \to \ldots
$$
\item Les foncteurs $R^iF : \mathcal{A} \to \mathcal{B}$ sont nuls pour
$i < 0$. De plus, $R^0F = F : \mathcal{A} \to \mathcal{B}$.
\item On a $R^iF(I) = 0$ pour $i > 0$ et $I$ injectif.
\item La suite $(R^iF, \delta)$ forme un $\delta$-foncteur universel (voir
Homologie, définition \ref{homology-definition-universal-delta-functor}) de
$\mathcal{A}$ vers $\mathcal{B}$.
\end{enumerate}
\end{lemma}

\begin{proof}
Ce lemme ne fait que récapituler certains des résultats obtenus jusqu'ici.
Notons que, d'après le lemme
\ref{lemma-enough-injectives-right-derived}, $RF$ est défini partout. Voici
quelques références :
\begin{enumerate}
\item Cela résulte de la partie (3) du
lemme \ref{lemma-right-derived-properties}, combinée à la suite exacte
longue de cohomologie (\ref{equation-long-exact-cohomology-sequence-D}) dans
$D^{+}(\mathcal{B})$.
\item C'est le lemme \ref{lemma-left-exact-higher-derived}.
\item C'est le fait que les objets injectifs sont acycliques.
\item C'est le lemme \ref{lemma-right-derived-delta-functor}.
\end{enumerate}
\end{proof}


\section{Résolutions de Cartan--Eilenberg}
\label{section-cartan-eilenberg}

\noindent
Cette section peut être développée. Les résultats peuvent être généralisés et
s'appliquer dans davantage de cas. Les résolutions ne doivent pas
nécessairement être formées d'injectifs, et la méthode fonctionne aussi dans
certaines situations non bornées.

\begin{definition}
\label{definition-cartan-eilenberg}
Soit $\mathcal{A}$ une catégorie abélienne.
Soit $K^\bullet$ un complexe borné inférieurement.
Une {\it résolution de Cartan--Eilenberg} de $K^\bullet$ est la donnée d'un
complexe double $I^{\bullet, \bullet}$ et d'un morphisme de complexes
$\epsilon : K^\bullet \to I^{\bullet, 0}$ ayant les propriétés suivantes :
\begin{enumerate}
\item Il existe $i \ll 0$ tel que $I^{p, q} = 0$ pour tout $p < i$ et tout
$q$.
\item On a $I^{p, q} = 0$ si $q < 0$.
\item Le complexe $I^{p, \bullet}$ est une résolution injective de $K^p$.
\item Le complexe $\Ker(d_1^{p, \bullet})$ est une résolution injective de
$\Ker(d_K^p)$.
\item Le complexe $\Im(d_1^{p, \bullet})$ est une résolution injective de
$\Im(d_K^p)$.
\item Le complexe $H^p_I(I^{\bullet, \bullet})$ est une résolution injective
de $H^p(K^\bullet)$.
\end{enumerate}
\end{definition}

\begin{lemma}
\label{lemma-cartan-eilenberg}
Soit $\mathcal{A}$ une catégorie abélienne ayant assez d'injectifs.
Soit $K^\bullet$ un complexe borné inférieurement.
Il existe une résolution de Cartan--Eilenberg de $K^\bullet$.
\end{lemma}

\begin{proof}
Supposons que $K^p = 0$ pour $p < n$. Décomposons $K^\bullet$ en suites
exactes courtes comme suit : posons $Z^p = \Ker(d^p)$,
$B^p = \Im(d^{p - 1})$, $H^p = Z^p/B^p$, et considérons
$$
\begin{matrix}
0 \to Z^n \to K^n \to B^{n + 1} \to 0 \\
0 \to B^{n + 1} \to Z^{n + 1} \to H^{n + 1} \to 0 \\
0 \to Z^{n + 1} \to K^{n + 1} \to B^{n + 2} \to 0 \\
0 \to B^{n + 2} \to Z^{n + 2} \to H^{n + 2} \to 0 \\
\ldots
\end{matrix}
$$
Posons $I^{p, q} = 0$ pour $p < n$. Nous choisissons par récurrence des
résolutions injectives de la manière suivante :
\begin{enumerate}
\item Choisissons une résolution injective
$Z^n \to J_Z^{n, \bullet}$ telle que $J_Z^{n, m} = 0$ pour $m < 0$, voir le
lemme \ref{lemma-injective-resolutions-exist}.
\item En utilisant le lemme \ref{lemma-injective-resolution-ses}, choisissons
des résolutions injectives $K^n \to I^{n, \bullet}$ et
$B^{n + 1} \to J_B^{n + 1, \bullet}$ telles que $I^{n, m} = 0$ pour $m < 0$
et $J_B^{n + 1, m} = 0$ pour $m < 0$, ainsi qu'une suite exacte de complexes
$0 \to J_Z^{n, \bullet} \to I^{n, \bullet} \to J_B^{n + 1, \bullet} \to 0$
compatible avec la suite exacte courte
$0 \to Z^n \to K^n \to B^{n + 1} \to 0$.
\item En utilisant le lemme \ref{lemma-injective-resolution-ses}, choisissons
des résolutions injectives $Z^{n + 1} \to J_Z^{n + 1, \bullet}$ et
$H^{n + 1} \to J_H^{n + 1, \bullet}$ telles que
$J_Z^{n + 1, m} = 0$ pour $m < 0$ et $J_H^{n + 1, m} = 0$ pour $m < 0$,
ainsi qu'une suite exacte de complexes
$0 \to J_B^{n + 1, \bullet} \to J_Z^{n + 1, \bullet}
\to J_H^{n + 1, \bullet} \to 0$ compatible avec la suite exacte courte
$0 \to B^{n + 1} \to Z^{n + 1} \to H^{n + 1} \to 0$.
\item Et ainsi de suite.
\end{enumerate}
En prenant pour applications
$d_1^\bullet : I^{p, \bullet} \to I^{p + 1, \bullet}$ les composées
$I^{p, \bullet} \to J_B^{p + 1, \bullet} \to
J_Z^{p + 1, \bullet} \to I^{p + 1, \bullet}$, tout est clair.
\end{proof}

\begin{lemma}
\label{lemma-two-ss-complex-functor}
Soit $F : \mathcal{A} \to \mathcal{B}$ un foncteur exact à gauche entre
catégories abéliennes.
Soit $K^\bullet$ un complexe borné inférieurement de $\mathcal{A}$.
Soit $I^{\bullet, \bullet}$ une résolution de Cartan--Eilenberg de
$K^\bullet$. Les suites spectrales
$({}'E_r, {}'d_r)_{r \geq 0}$ et $({}''E_r, {}''d_r)_{r \geq 0}$ associées au
complexe double $F(I^{\bullet, \bullet})$ vérifient les relations
$$
{}'E_1^{p, q} = R^qF(K^p)
\quad
\text{et}
\quad
{}''E_2^{p, q} = R^pF(H^q(K^\bullet))
$$
De plus, ces suites spectrales sont bornées, convergent vers
$H^*(RF(K^\bullet))$, et les filtrations induites correspondantes sur
$H^n(RF(K^\bullet))$ sont finies.
\end{lemma}

\begin{proof}
Nous utiliserons sans autre mention les observations suivantes :
\begin{enumerate}
\item Puisque $I^{p, \bullet}$ est une résolution injective de $K^p$, le
foncteur $RF$ est défini en $K^p[0]$, de valeur $F(I^{p, \bullet})$.
\item Puisque $H^p_I(I^{\bullet, \bullet})$ est une résolution injective de
$H^p(K^\bullet)$, le foncteur dérivé $RF$ est défini en
$H^p(K^\bullet)[0]$, de valeur $F(H^p_I(I^{\bullet, \bullet}))$.
\item D'après Homologie, lemme
\ref{homology-lemma-double-complex-gives-resolution}, le complexe total
$\text{Tot}(I^{\bullet, \bullet})$ est une résolution injective de
$K^\bullet$. Ainsi, $RF$ est défini en $K^\bullet$, de valeur
$F(\text{Tot}(I^{\bullet, \bullet}))$.
\end{enumerate}
Considérons les deux suites spectrales associées au complexe double
$L^{\bullet, \bullet} = F(I^{\bullet, \bullet})$, voir Homologie, lemme
\ref{homology-lemma-ss-double-complex}.
Elles sont toutes deux bornées, convergent vers
$H^*(\text{Tot}(L^{\bullet, \bullet}))$ et induisent des filtrations finies sur
$H^n(\text{Tot}(L^{\bullet, \bullet}))$, voir Homologie, lemme
\ref{homology-lemma-first-quadrant-ss}.
Comme
$\text{Tot}(L^{\bullet, \bullet}) =
\text{Tot}(F(I^{\bullet, \bullet})) =
F(\text{Tot}(I^{\bullet, \bullet}))$ calcule $RF(K^\bullet)$, la dernière
assertion du lemme en résulte.

\medskip\noindent
Calcul de la première suite spectrale. On a
${}'E_1^{p, q} = H^q(L^{p, \bullet})$, autrement dit
$$
{}'E_1^{p, q} = H^q(F(I^{p, \bullet})) = R^qF(K^p)
$$
comme voulu. Notons pour la suite que les applications
${}'d_1^{p, q} : {}'E_1^{p, q} \to {}'E_1^{p + 1, q}$ sont les applications
$R^qF(K^p) \to R^qF(K^{p + 1})$ induites par $K^p \to K^{p + 1}$ et par le
fait que $R^qF$ est un foncteur.

\medskip\noindent
Calcul de la seconde suite spectrale. On a
${}''E_1^{p, q} = H^q(L^{\bullet, p}) = H^q(F(I^{\bullet, p}))$.
Notons que le complexe $I^{\bullet, p}$ est borné inférieurement, qu'il est
formé d'objets injectifs, et que le noyau, l'image et la cohomologie de chacun
de ses différentiels sont en outre des objets injectifs de $\mathcal{A}$. Nous
pouvons donc scinder les différentiels : chaque différentiel est une surjection
scindée sur un facteur direct. Il en est de même après application de $F$.
Ainsi,
${}''E_1^{p, q} = F(H^q(I^{\bullet, p})) = F(H^q_I(I^{\bullet, p}))$.
Les différentielles sur ce terme sont $(-1)^q$ fois l'image par $F$ de la
différentielle du complexe $H^p_I(I^{\bullet, \bullet})$, qui est une
résolution injective de $H^p(K^\bullet)$. On en déduit la description des
termes $E_2$.
\end{proof}

\begin{remark}
\label{remark-functorial-ss}
Les suites spectrales du lemme \ref{lemma-two-ss-complex-functor} sont
fonctorielles en $K^\bullet$. Cela résulte des propriétés de
fonctorialité des résolutions de Cartan--Eilenberg. D'autre part, elles sont
toutes deux des exemples d'une suite spectrale plus générale que l'on peut
associer à un complexe filtré de $\mathcal{A}$. Sa construction entraînera la
fonctorialité. Nous y reviendrons dans la section consacrée à la catégorie
dérivée filtrée, voir la remarque \ref{remark-final-functorial}.
\end{remark}


\section{Composition de foncteurs dérivés droits}
\label{section-composition-right-derived-functors}

\noindent
Il est parfois possible de calculer le foncteur dérivé droit d'un composé.
Supposons que $\mathcal{A}, \mathcal{B}, \mathcal{C}$ soient des catégories
abéliennes. Soient $F : \mathcal{A} \to \mathcal{B}$ et
$G : \mathcal{B} \to \mathcal{C}$ des foncteurs exacts à gauche. Supposons
que les foncteurs dérivés droits
$RF : D^{+}(\mathcal{A}) \to D^{+}(\mathcal{B})$,
$RG : D^{+}(\mathcal{B}) \to D^{+}(\mathcal{C})$ et
$R(G \circ F) : D^{+}(\mathcal{A}) \to D^{+}(\mathcal{C})$ soient définis
partout. Il existe alors une transformation canonique
$$
t : R(G \circ F) \longrightarrow RG \circ RF
$$
de foncteurs de $D^{+}(\mathcal{A})$ vers $D^{+}(\mathcal{C})$, voir le
lemme \ref{lemma-compose-derived-functors-general}.
Cette transformation n'est pas toujours un isomorphisme.

\begin{lemma}
\label{lemma-compose-derived-functors}
Soient $\mathcal{A}, \mathcal{B}, \mathcal{C}$ des catégories abéliennes.
Soient $F : \mathcal{A} \to \mathcal{B}$ et
$G : \mathcal{B} \to \mathcal{C}$ des foncteurs exacts à gauche. Supposons
que $\mathcal{A}$ et $\mathcal{B}$ aient assez d'injectifs. Les assertions
suivantes sont équivalentes :
\begin{enumerate}
\item $F(I)$ est acyclique à droite pour $G$, pour tout objet injectif $I$
de $\mathcal{A}$ ;
\item l'application canonique
$$
t : R(G \circ F) \longrightarrow RG \circ RF
$$
est un isomorphisme de foncteurs de $D^{+}(\mathcal{A})$ vers
$D^{+}(\mathcal{C})$.
\end{enumerate}
\end{lemma}

\begin{proof}
Si (2) est satisfaite, alors (1) s'obtient en évaluant l'isomorphisme $t$ sur
$RF(I) = F(I)$. Réciproquement, supposons (1) satisfaite.
Soit $A^\bullet$ un complexe borné inférieurement de $\mathcal{A}$.
Choisissons une résolution injective $A^\bullet \to I^\bullet$.
L'application $t$ est donnée (voir la preuve du
lemme \ref{lemma-compose-derived-functors-general}) par les applications
$$
R(G \circ F)(A^\bullet) =
(G \circ F)(I^\bullet) =
G(F(I^\bullet)) \to
RG(F(I^\bullet)) =
RG(RF(A^\bullet))
$$
où la flèche est un isomorphisme d'après le
lemme \ref{lemma-leray-acyclicity}.
\end{proof}

\begin{lemma}[Suite spectrale de Grothendieck]
\label{lemma-grothendieck-spectral-sequence}
Sous les hypothèses du lemme \ref{lemma-compose-derived-functors}, supposons
que les conditions équivalentes (1) et (2) soient satisfaites.
Soit $X$ un objet de $D^{+}(\mathcal{A})$.
Il existe une suite spectrale $(E_r, d_r)_{r \geq 0}$ formée d'objets
bigradués $E_r$ de $\mathcal{C}$, où $d_r$ est de bidegré $(r, - r + 1)$,
et telle que
$$
E_2^{p, q} = R^pG(H^q(RF(X)))
$$
De plus, cette suite spectrale est bornée, converge vers
$H^*(R(G \circ F)(X))$ et induit une filtration finie sur chaque
$H^n(R(G \circ F)(X))$.
\end{lemma}

\noindent
Pour un objet $A$ de $\mathcal{A}$, on obtient
$E_2^{p, q} = R^pG(R^qF(A))$, qui converge vers
$R^{p + q}(G \circ F)(A)$.

\begin{proof}
Nous pouvons représenter $X$ par un complexe borné inférieurement
$A^\bullet$. Choisissons une résolution injective
$A^\bullet \to I^\bullet$. Choisissons une résolution de Cartan--Eilenberg
$F(I^\bullet) \to I^{\bullet, \bullet}$ à l'aide du
lemme \ref{lemma-cartan-eilenberg}. Appliquons la seconde suite spectrale du
lemme \ref{lemma-two-ss-complex-functor}.
\end{proof}


\section{Foncteurs de résolution}
\label{section-derived-category}

\noindent
Soit $\mathcal{A}$ une catégorie abélienne ayant assez d'injectifs.
Notons $\mathcal{I}$ la sous-catégorie additive pleine de $\mathcal{A}$ dont
les objets sont les objets injectifs de $\mathcal{A}$.
Dans ce cas, $K^{+}(\mathcal{I})$ et $D^{+}(\mathcal{A})$ sont équivalentes
(voir la proposition \ref{proposition-derived-category}).
Il est donc souvent naturel d'identifier mentalement $D^{+}(\mathcal{A})$ à
la catégorie $K^{+}(\mathcal{I})$, plus facile à appréhender dans ce cas.

\begin{proposition}
\label{proposition-derived-category}
Soit $\mathcal{A}$ une catégorie abélienne.
Supposons que $\mathcal{A}$ ait assez d'injectifs.
Notons $\mathcal{I} \subset \mathcal{A}$ la sous-catégorie additive
strictement pleine dont les objets sont les objets injectifs de
$\mathcal{A}$. Le foncteur
$$
K^{+}(\mathcal{I}) \longrightarrow D^{+}(\mathcal{A})
$$
est exact, pleinement fidèle et essentiellement surjectif, c'est-à-dire que
c'est une équivalence de catégories triangulées.
\end{proposition}

\begin{proof}
Il est clair que le foncteur est exact. Il est essentiellement surjectif
d'après le lemme \ref{lemma-injective-resolutions-exist}.
La pleine fidélité résulte du
lemme \ref{lemma-morphisms-into-injective-complex}.
\end{proof}

\noindent
La proposition \ref{proposition-derived-category} implique que nous pouvons
construire des foncteurs de résolution. On peut en fait prouver leur existence
même dans certains cas où la catégorie abélienne $\mathcal{A}$ est une
catégorie « grande », c'est-à-dire qu'elle possède une classe d'objets.

\begin{definition}
\label{definition-localization-functor}
Soit $\mathcal{A}$ une catégorie abélienne ayant assez d'injectifs.
Un {\it foncteur de résolution}\footnote{Cette terminologie n'est
vraisemblablement pas standard.} pour $\mathcal{A}$ est la donnée des objets
et morphismes suivants :
\begin{enumerate}
\item pour tout $K^\bullet \in \Ob(K^{+}(\mathcal{A}))$, un complexe
$j(K^\bullet)$ d'injectifs borné inférieurement ;
\item pour tout $K^\bullet \in \Ob(K^{+}(\mathcal{A}))$, un
quasi-isomorphisme $i_{K^\bullet} : K^\bullet \to j(K^\bullet)$.
\end{enumerate}
\end{definition}

\begin{lemma}
\label{lemma-resolution-functor}
Soit $\mathcal{A}$ une catégorie abélienne ayant assez d'injectifs.
Étant donné un foncteur de résolution $(j, i)$, il existe une unique façon
de faire de $j$ un foncteur et de $i$ un $2$-isomorphisme donnant un diagramme
$2$-commutatif
$$
\xymatrix{
K^{+}(\mathcal{A}) \ar[rd] \ar[rr]_j & & K^{+}(\mathcal{I}) \ar[ld] \\
& D^{+}(\mathcal{A})
}
$$
où $\mathcal{I}$ est la sous-catégorie additive pleine de $\mathcal{A}$
formée des objets injectifs.
\end{lemma}

\begin{proof}
Pour tout morphisme $\alpha : K^\bullet \to L^\bullet$ de
$K^{+}(\mathcal{A})$, il existe un unique morphisme
$j(\alpha) : j(K^\bullet) \to j(L^\bullet)$ dans $K^{+}(\mathcal{I})$ tel que
$$
\xymatrix{
K^\bullet \ar[r]_\alpha \ar[d]_{i_{K^\bullet}} &
L^\bullet \ar[d]^{i_{L^\bullet}} \\
j(K^\bullet) \ar[r]^{j(\alpha)} & j(L^\bullet)
}
$$
soit commutatif dans $K^{+}(\mathcal{A})$. Pour le voir, on peut utiliser les
lemmes \ref{lemma-morphisms-lift} et
\ref{lemma-morphisms-equal-up-to-homotopy}, ou le
lemme équivalent \ref{lemma-morphisms-into-injective-complex}.
L'unicité implique que $j$ est un foncteur, et la commutativité du diagramme
implique que $i$ donne un $2$-morphisme témoignant de la $2$-commutativité du
diagramme de catégories de l'énoncé.
\end{proof}

\begin{lemma}
\label{lemma-into-derived-category}
Soit $\mathcal{A}$ une catégorie abélienne.
Supposons que $\mathcal{A}$ ait assez d'injectifs.
Il existe alors un foncteur de résolution $j$, unique à un unique
isomorphisme de foncteurs près.
\end{lemma}

\begin{proof}
Considérons l'ensemble de tous les objets $K^\bullet$ de
$K^{+}(\mathcal{A})$. (Rappelons que, selon nos conventions, toute catégorie
possède un ensemble d'objets, sauf mention contraire.) D'après le
lemme \ref{lemma-injective-resolutions-exist}, tout objet admet une résolution
injective. L'axiome du choix nous permet donc de choisir, pour chaque
$K^\bullet$, une résolution injective
$i_{K^\bullet} : K^\bullet \to j(K^\bullet)$.
\end{proof}

\begin{lemma}
\label{lemma-j-is-exact}
Soit $\mathcal{A}$ une catégorie abélienne ayant assez d'injectifs.
Tout foncteur de résolution
$j : K^{+}(\mathcal{A}) \to K^{+}(\mathcal{I})$ est exact.
\end{lemma}

\begin{proof}
Notons $i_{K^\bullet} : K^\bullet \to j(K^\bullet)$ les applications
canoniques de la définition \ref{definition-localization-functor}.
Commençons par établir l'existence de l'isomorphisme fonctoriel
$j(K^\bullet[1]) \to j(K^\bullet)[1]$.
Considérons le diagramme
$$
\xymatrix{
K^\bullet[1] \ar[d]^{i_{K^\bullet[1]}} \ar@{=}[rr] & &
K^\bullet[1] \ar[d]^{i_{K^\bullet}[1]} \\
j(K^\bullet[1]) \ar@{..>}[rr]^{\xi_{K^\bullet}} & & j(K^\bullet)[1]
}
$$
D'après les lemmes \ref{lemma-morphisms-lift} et
\ref{lemma-morphisms-equal-up-to-homotopy}, il existe une unique flèche en
pointillé $\xi_{K^\bullet}$ dans $K^{+}(\mathcal{I})$ rendant le diagramme
commutatif dans $K^{+}(\mathcal{A})$.
Nous omettons la vérification que l'on obtient ainsi un isomorphisme
fonctoriel. (Indication : utiliser de nouveau le
lemme \ref{lemma-morphisms-equal-up-to-homotopy}.)

\medskip\noindent
Soit $(K^\bullet, L^\bullet, M^\bullet, f, g, h)$ un triangle distingué de
$K^{+}(\mathcal{A})$. Nous devons montrer que
$(j(K^\bullet), j(L^\bullet), j(M^\bullet), j(f), j(g),
\xi_{K^\bullet} \circ j(h))$ est un triangle distingué de
$K^{+}(\mathcal{I})$. Notons que nous avons un diagramme commutatif
$$
\xymatrix{
K^\bullet \ar[r]_f \ar[d] &
L^\bullet \ar[r]_g \ar[d] &
M^\bullet \ar[rr]_h \ar[d] & &
K^\bullet[1] \ar[d] \\
j(K^\bullet) \ar[r]^{j(f)} &
j(L^\bullet) \ar[r]^{j(g)} &
j(M^\bullet) \ar[rr]^{\xi_{K^\bullet} \circ j(h)} & &
j(K^\bullet)[1]
}
$$
dans $K^{+}(\mathcal{A})$, dont les flèches verticales sont les
quasi-isomorphismes $i_K, i_L, i_M$. Ainsi, l'image de
$(j(K^\bullet), j(L^\bullet), j(M^\bullet), j(f), j(g),
\xi_{K^\bullet} \circ j(h))$ dans $D^{+}(\mathcal{A})$ est isomorphe à un
triangle distingué, donc est elle-même un triangle distingué d'après TR1.
Le lemme \ref{lemma-exact-equivalence} montre alors que
$(j(K^\bullet), j(L^\bullet), j(M^\bullet), j(f), j(g),
\xi_{K^\bullet} \circ j(h))$ est un triangle distingué dans
$K^{+}(\mathcal{I})$.
\end{proof}

\begin{lemma}
\label{lemma-resolution-functor-quasi-inverse}
Soit $\mathcal{A}$ une catégorie abélienne ayant assez d'injectifs.
Soit $j$ un foncteur de résolution. Notons
$Q : K^{+}(\mathcal{A}) \to D^{+}(\mathcal{A})$ le foncteur naturel.
Alors $j = j' \circ Q$ pour un unique foncteur
$j' : D^{+}(\mathcal{A}) \to K^{+}(\mathcal{I})$, quasi-inverse du foncteur
canonique $K^{+}(\mathcal{I}) \to D^{+}(\mathcal{A})$.
\end{lemma}

\begin{proof}
Le foncteur $Q$ est une localisation d'après le
lemme \ref{lemma-bounded-derived}. Pour prouver l'existence de $j'$, il suffit
de montrer que tout élément de $\text{Qis}^{+}(\mathcal{A})$ est envoyé sur
un isomorphisme par le foncteur $j$, voir le
lemme \ref{lemma-universal-property-localization}.
Considérons le carré commutatif de la preuve du
lemme \ref{lemma-j-is-exact}. Dans ce carré, si $\alpha$ est un
quasi-isomorphisme, alors
$i_{L^\bullet}\circ\alpha=j(\alpha)\circ i_{K^\bullet}$ est aussi un
quasi-isomorphisme, et il en est donc de même de $j(\alpha)$. Ainsi, d'après
la proposition \ref{proposition-derived-category}, le morphisme $j(\alpha)$
est un isomorphisme dans $K^+(\mathcal{I})$. Nous omettons de vérifier que
$j'$ est quasi-inverse de
$K^{+}(\mathcal{I}) \to D^{+}(\mathcal{A})$.
\end{proof}

\begin{remark}
\label{remark-big-localization}
Supposons que $\mathcal{A}$ soit une « grande » catégorie abélienne ayant
assez d'injectifs, par exemple la catégorie des groupes abéliens. Dans ce cas,
nous devons construire notre foncteur de résolution avec un peu plus de soin,
car nous ne pouvons pas utiliser l'axiome du choix avec un quantificateur
portant sur une classe. Remarquons cependant que la preuve du lemme montre que
deux foncteurs de résolution quelconques sont canoniquement isomorphes. En
effet, étant donné des quasi-isomorphismes
$i : K^\bullet \to I^\bullet$ et $i' : K^\bullet \to J^\bullet$ d'un
complexe borné inférieurement $K^\bullet$ vers des complexes d'injectifs
bornés inférieurement, il existe un unique (!) morphisme
$a : I^\bullet \to J^\bullet$ dans $K^{+}(\mathcal{I})$ tel que
$i' = a \circ i$ comme morphismes dans $K^{+}(\mathcal{A})$. La seule
difficulté est donc l'existence, et nous verrons dans la section suivante
comment la traiter.
\end{remark}


\section{Plongements fonctoriels dans des injectifs et foncteurs de
résolution}
\label{section-functorial-injective-resolutions}

\noindent
Dans cette section, nous reprenons la construction d'un foncteur de résolution
$K^{+}(\mathcal{A}) \to K^{+}(\mathcal{I})$ lorsque la catégorie
$\mathcal{A}$ possède des plongements fonctoriels dans des objets injectifs.
Il y a deux raisons à cela : (1) la preuve est plus simple et (2) la
construction fonctionne aussi lorsque $\mathcal{A}$ est une « grande »
catégorie abélienne. Voir la remarque
\ref{remark-big-abelian-category} ci-dessous.

\medskip\noindent
Soit $\mathcal{A}$ une catégorie abélienne. Comme précédemment, notons
$\mathcal{I}$ la sous-catégorie additive pleine de $\mathcal{A}$ formée des
objets injectifs. Considérons la catégorie
$\text{InjRes}(\mathcal{A})$ des flèches
$\alpha : K^\bullet \to I^\bullet$, où $K^\bullet$ est un complexe borné
inférieurement de $\mathcal{A}$, $I^\bullet$ est un complexe borné
inférieurement d'objets injectifs de $\mathcal{A}$, et $\alpha$ est un
quasi-isomorphisme. Autrement dit, $\alpha$ est une résolution injective et
$K^\bullet$ est borné inférieurement. Il existe un foncteur évident
$$
s : \text{InjRes}(\mathcal{A}) \longrightarrow \text{Comp}^{+}(\mathcal{A})
$$
défini par $(\alpha : K^\bullet \to I^\bullet) \mapsto K^\bullet$.
Il existe aussi un foncteur
$$
t : \text{InjRes}(\mathcal{A}) \longrightarrow K^{+}(\mathcal{I})
$$
défini par $(\alpha : K^\bullet \to I^\bullet) \mapsto I^\bullet$.

\begin{lemma}
\label{lemma-functorial-injective-resolutions}
Soit $\mathcal{A}$ une catégorie abélienne.
Supposons que $\mathcal{A}$ possède des plongements fonctoriels dans des
objets injectifs, voir Homologie, définition
\ref{homology-definition-functorial-injective-embedding}.
\begin{enumerate}
\item Il existe un foncteur
$inj : \text{Comp}^{+}(\mathcal{A}) \to \text{InjRes}(\mathcal{A})$ tel que
$s \circ inj = \text{id}$.
\item Tout foncteur
$inj : \text{Comp}^{+}(\mathcal{A}) \to \text{InjRes}(\mathcal{A})$ tel que
$s \circ inj = \text{id}$ fournit un foncteur de résolution, voir la
définition \ref{definition-localization-functor}.
\end{enumerate}
\end{lemma}

\begin{proof}
Soit $A \mapsto (A \to J(A))$ un plongement fonctoriel dans des objets
injectifs, voir Homologie, définition
\ref{homology-definition-functorial-injective-embedding}.
Remarquons d'abord que nous pouvons supposer $J(0) = 0$. En effet, dans le cas
contraire, pour tout objet $A$ nous avons $0 \to A \to 0$, ce qui donne une
décomposition en somme directe
$J(A) = J(0) \oplus \Ker(J(A) \to J(0))$.
Notons que le morphisme fonctoriel $A \to J(A)$ est nécessairement à valeurs
dans le second facteur. Nous pouvons donc, au besoin, remplacer notre foncteur
par $J'(A) = \Ker(J(A) \to J(0))$.

\medskip\noindent
Soit $K^\bullet$ un complexe borné inférieurement de $\mathcal{A}$.
Disons que $K^p = 0$ si $p < B$.
Nous allons construire un complexe double $I^{\bullet, \bullet}$ d'injectifs,
ainsi qu'une application $\alpha : K^\bullet \to I^{\bullet, 0}$ telle que
$\alpha$ induise un quasi-isomorphisme de $K^\bullet$ vers le complexe total
associé à $I^{\bullet, \bullet}$.
Posons d'abord $I^{p, q} = 0$ lorsque $q < 0$.
Posons ensuite $I^{p, 0} = J(K^p)$, et prenons pour
$\alpha^p : K^p \to I^{p, 0}$ le plongement fonctoriel. Comme $J$ est un
foncteur, $I^{\bullet, 0}$ est un complexe et $\alpha$ est un morphisme de
complexes. Chaque $\alpha^p$ est injectif. De plus, $I^{p, 0} = 0$ pour
$p < B$, car $J(0) = 0$. Posons ensuite
$I^{p, 1} = J(\Coker(K^p \to I^{p, 0}))$. La fonctorialité montre de nouveau
que $I^{\bullet, 1}$ est un complexe. Nous obtenons aussi
$I^{p, 1} = 0$ pour $p < B$. Il est en outre clair que $K^p$ s'envoie
isomorphiquement sur $\Ker(I^{p, 0} \to I^{p, 1})$.
Dans une troisième étape, prenons
$I^{p, 2} = J(\Coker(I^{p, 0} \to I^{p, 1}))$, et ainsi de suite.

\medskip\noindent
Nous pouvons maintenant appliquer Homologie, lemme
\ref{homology-lemma-double-complex-gives-resolution}, pour conclure que
l'application
$$
\alpha : K^\bullet \longrightarrow \text{Tot}(I^{\bullet, \bullet})
$$
est un quasi-isomorphisme. Pour prouver que l'on obtient un foncteur $inj$, il
reste à montrer que la construction précédente est fonctorielle. Nous omettons
cette vérification.

\medskip\noindent
Supposons donné un foncteur $inj$ tel que $s \circ inj = \text{id}$.
Pour tout objet $K^\bullet$ de $\text{Comp}^{+}(\mathcal{A})$, nous pouvons
écrire
$$
inj(K^\bullet) = (i_{K^\bullet} : K^\bullet \to j(K^\bullet))
$$
Nous obtenons ainsi un foncteur de résolution au sens de la
définition \ref{definition-localization-functor}.
\end{proof}

\begin{remark}
\label{remark-match}
Supposons que $inj$ soit un foncteur tel que $s \circ inj = \text{id}$, comme
dans la partie (2) du lemme
\ref{lemma-functorial-injective-resolutions}.
Écrivons
$inj(K^\bullet) = (i_{K^\bullet} : K^\bullet \to j(K^\bullet))$, comme dans
la preuve de ce lemme.
Soit $\alpha : K^\bullet \to L^\bullet$ une application entre complexes bornés
inférieurement. Considérons l'application $inj(\alpha)$ dans la catégorie
$\text{InjRes}(\mathcal{A})$. Elle induit un diagramme commutatif
$$
\xymatrix{
K^\bullet
\ar[rr]^-{\alpha}
\ar[d]_{i_K} & &
L^\bullet \ar[d]^{i_L} \\
j(K^\bullet)
\ar[rr]^-{inj(\alpha)}
& &
j(L^\bullet)
}
$$
de morphismes de complexes.
En examinant la preuve du lemme \ref{lemma-resolution-functor}, nous voyons
donc que le foncteur
$j : K^{+}(\mathcal{A}) \to K^{+}(\mathcal{I})$ est donné par la règle
$$
j(\alpha\text{ à homotopie près}) =
inj(\alpha)\text{ à homotopie près}\in
\Hom_{K^{+}(\mathcal{I})}(j(K^\bullet), j(L^\bullet))
$$
Nous voyons ainsi que $j$ coïncide avec $t \circ inj$ dans ce cas,
c'est-à-dire que le diagramme
$$
\xymatrix{
\text{Comp}^{+}(\mathcal{A}) \ar[rr]_{t \circ inj} \ar[rd] & &
K^{+}(\mathcal{I}) \\
& K^{+}(\mathcal{A}) \ar[ru]_j
}
$$
est commutatif.
\end{remark}

\begin{remark}
\label{remark-big-abelian-category}
Soit $\textit{Mod}(\mathcal{O}_X)$ la catégorie des
$\mathcal{O}_X$-modules sur un espace annelé $(X, \mathcal{O}_X)$ (ou, plus
généralement, sur un site annelé). Nous verrons plus loin que
$\textit{Mod}(\mathcal{O}_X)$ a assez d'injectifs et possède même des
plongements fonctoriels dans des injectifs, voir Injectifs, théorème
\ref{injectives-theorem-sheaves-modules-injectives}.
Notons que la preuve du lemme \ref{lemma-into-derived-category} ne s'applique
pas à $\textit{Mod}(\mathcal{O}_X)$, mais que celle du
lemme \ref{lemma-functorial-injective-resolutions} s'applique à
$\textit{Mod}(\mathcal{O}_X)$. Nous obtenons donc
$$
j : K^{+}(\textit{Mod}(\mathcal{O}_X))
\longrightarrow
K^{+}(\mathcal{I})
$$
qui est un foncteur de résolution, où $\mathcal{I}$ est la catégorie additive
des $\mathcal{O}_X$-modules injectifs. Le même argument s'applique dans les cas
suivants :
\begin{enumerate}
\item La catégorie $\text{Mod}_R$ des $R$-modules sur un anneau $R$.
\item La catégorie $\textit{PMod}(\mathcal{O})$ des préfaisceaux de
$\mathcal{O}$-modules sur un site muni d'un préfaisceau d'anneaux.
\item La catégorie $\textit{Mod}(\mathcal{O})$ des faisceaux de
$\mathcal{O}$-modules sur un site annelé.
\end{enumerate}
\end{remark}


\section{Foncteurs dérivés droits au moyen de foncteurs de résolution}
\sectionmark{Foncteurs dérivés par résolution}
\label{section-right-derived-functor-via-resolutions}

\noindent
Le contenu du lemme suivant est que nous pouvons simplement poser
$RF(K^\bullet) = F(j(K^\bullet))$ lorsqu'un foncteur de résolution $j$ est
donné.

\begin{lemma}
\label{lemma-right-derived-functor}
Soit $\mathcal{A}$ une catégorie abélienne ayant assez d'injectifs.
Soit $F : \mathcal{A} \to \mathcal{B}$ un foncteur additif à valeurs dans une
catégorie abélienne. Soit $(j, i)$ un foncteur de résolution, voir la
définition \ref{definition-localization-functor}.
Le foncteur dérivé droit $RF$ de $F$ s'insère dans le diagramme
$2$-commutatif suivant :
$$
\xymatrix{
D^{+}(\mathcal{A}) \ar[rd]_{RF} \ar[rr]^{j'} & &
K^{+}(\mathcal{I}) \ar[ld]^F \\
& D^{+}(\mathcal{B})
}
$$
où $j'$ est le foncteur du
lemme \ref{lemma-resolution-functor-quasi-inverse}.
\end{lemma}

\begin{proof}
D'après le lemme \ref{lemma-injective-acyclic}, on a
$RF(K^\bullet) = F(j(K^\bullet))$.
\end{proof}

\begin{remark}
\label{remark-right-derived-functor}
Dans la situation du lemme \ref{lemma-right-derived-functor}, nous voyons que
nous avons en fait relevé le foncteur dérivé droit en un foncteur exact
$F \circ j' : D^{+}(\mathcal{A}) \to K^{+}(\mathcal{B})$.
Il est parfois utile de recourir à une telle factorisation.
\end{remark}


\section{Catégorie dérivée filtrée et résolutions injectives}
\label{section-filtered-derived}

\noindent
Soit $\mathcal{A}$ une catégorie abélienne. Dans cette section, nous allons
montrer que, si $\mathcal{A}$ a assez d'injectifs, il en est de même en un
certain sens de la catégorie $\text{Fil}^f(\mathcal{A})$. Cette observation
permet d'effectuer des calculs dans la catégorie dérivée filtrée de
$\mathcal{A}$.

\medskip\noindent
La catégorie $\text{Fil}^f(\mathcal{A})$ est un exemple de catégorie exacte,
voir Injectifs, remarque
\ref{injectives-remark-embed-exact-category}.
Les morphismes stricts jouent un rôle particulier, voir Homologie, définition
\ref{homology-definition-strict} ; ce sont les morphismes $f$ tels que
$\Coim(f) = \Im(f)$.
Nous dirons qu'un complexe $A \to B \to C$ dans
$\text{Fil}^f(\mathcal{A})$ est {\it exact} si la suite
$\text{gr}(A) \to \text{gr}(B) \to \text{gr}(C)$ est exacte dans
$\mathcal{A}$. Il en résulte que $A \to B$ et $B \to C$ sont des morphismes
stricts, voir Homologie, lemme
\ref{homology-lemma-filtered-acyclic}.

\begin{definition}
\label{definition-filtered-complexes-notation}
Soit $\mathcal{A}$ une catégorie abélienne.
Nous disons qu'un objet $I$ de $\text{Fil}^f(\mathcal{A})$ est
{\it injectif filtré} si chaque $\text{gr}^p(I)$ est un objet injectif de
$\mathcal{A}$.
\end{definition}

\begin{lemma}
\label{lemma-filtered-injective}
Soit $\mathcal{A}$ une catégorie abélienne.
Un objet $I$ de $\text{Fil}^f(\mathcal{A})$ est injectif filtré si et seulement
s'il existe $a \leq b$, des objets injectifs $I_n$, $a \leq n \leq b$, de
$\mathcal{A}$, et un isomorphisme
$I \cong \bigoplus_{a \leq n \leq b} I_n$ tels que
$F^pI = \bigoplus_{n \geq p} I_n$.
\end{lemma}

\begin{proof}
Cela résulte du fait que tout monomorphisme $J \to M$ de $\mathcal{A}$ est
scindé lorsque $J$ est injectif. Nous omettons les détails.
\end{proof}

\begin{lemma}
\label{lemma-split-strict-monomorphism}
Soit $\mathcal{A}$ une catégorie abélienne.
Tout monomorphisme strict $u : I \to A$ de $\text{Fil}^f(\mathcal{A})$, où
$I$ est un objet injectif filtré, est un monomorphisme scindé.
\end{lemma}

\begin{proof}
Soit $p$ le plus grand entier tel que $F^pI \not = 0$.
En particulier, $\text{gr}^p(I) = F^pI$.
Soit $I'$ l'objet de $\text{Fil}^f(\mathcal{A})$ dont l'objet sous-jacent de
$\mathcal{A}$ est $F^pI$ et dont la filtration est donnée par
$F^nI' = 0$ pour $n > p$ et $F^nI' = I' = F^pI$ pour $n \leq p$.
Notons que $I' \to I$ est aussi un monomorphisme strict.
Le fait que $u$ soit un monomorphisme strict implique que
$F^pI \to A/F^{p + 1}(A)$ est injectif, voir Homologie, lemme
\ref{homology-lemma-characterize-strict}.
Choisissons une rétraction $s : A/F^{p + 1}A \to F^pI$ dans
$\mathcal{A}$. Le morphisme induit
$s' : A \to I'$ est un morphisme strict d'objets filtrés qui scinde la
composée $I' \to I \to A$.
Nous pouvons donc écrire $A = I' \oplus \Ker(s')$ et
$I = I' \oplus \Ker(s'|_I)$. Notons que
$\Ker(s'|_I) \to \Ker(s')$ est un monomorphisme strict et que
$\Ker(s'|_I)$ est un objet injectif filtré.
Par récurrence sur la longueur de la filtration de $I$, l'application
$\Ker(s'|_I) \to \Ker(s')$ est un monomorphisme scindé. Cela conclut la
preuve.
\end{proof}

\begin{lemma}
\label{lemma-injective-property-filtered-injective}
Soit $\mathcal{A}$ une catégorie abélienne.
Soient $u : A \to B$ un monomorphisme strict de
$\text{Fil}^f(\mathcal{A})$ et $f : A \to I$ un morphisme de $A$ vers un objet
injectif filtré de $\text{Fil}^f(\mathcal{A})$.
Alors il existe un morphisme $g : B \to I$ tel que $f = g \circ u$.
\end{lemma}

\begin{proof}
Le morphisme $f' : I \to I \amalg_A B$ obtenu en poussant $f$ le long de $u$
est un monomorphisme strict, voir Homologie, lemme
\ref{homology-lemma-pushout-filtered}.
Le résultat découle donc formellement du
lemme \ref{lemma-split-strict-monomorphism}.
\end{proof}

\begin{lemma}
\label{lemma-strict-monomorphism-into-filtered-injective}
Soit $\mathcal{A}$ une catégorie abélienne ayant assez d'injectifs.
Pour tout objet $A$ de $\text{Fil}^f(\mathcal{A})$, il existe un monomorphisme
strict $A \to I$, où $I$ est un objet injectif filtré.
\end{lemma}

\begin{proof}
Choisissons $a \leq b$ tels que $\text{gr}^p(A) = 0$ sauf si
$p \in \{a, a + 1, \ldots, b\}$. Pour chaque
$n \in \{a, a + 1, \ldots, b\}$, choisissons un monomorphisme
$u_n : A/F^{n + 1}A \to I_n$, où $I_n$ est un objet injectif.
Posons $I = \bigoplus_{a \leq n \leq b} I_n$, muni de la filtration
$F^pI = \bigoplus_{n \geq p} I_n$, et prenons pour $u : A \to I$ la somme
directe des applications $u_n$.
\end{proof}

\begin{lemma}
\label{lemma-filtered-injective-right-resolution-single-object}
Soit $\mathcal{A}$ une catégorie abélienne ayant assez d'injectifs.
Pour tout objet $A$ de $\text{Fil}^f(\mathcal{A})$, il existe un
quasi-isomorphisme filtré $A[0] \to I^\bullet$, où $I^\bullet$ est un complexe
d'objets injectifs filtrés tel que $I^n = 0$ pour $n < 0$.
\end{lemma}

\begin{proof}
Choisissons d'abord un monomorphisme strict
$u_0 : A \to I^0$ de $A$ vers un objet injectif filtré, voir le
lemme \ref{lemma-strict-monomorphism-into-filtered-injective}.
Choisissons ensuite un monomorphisme strict
$u_1 : \Coker(u_0) \to I^1$ vers un objet injectif filtré de
$\text{Fil}^f(\mathcal{A})$. Notons $d^0$ l'application induite
$I^0 \to I^1$.
Choisissons ensuite un monomorphisme strict
$u_2 : \Coker(u_1) \to I^2$ vers un objet injectif filtré de
$\text{Fil}^f(\mathcal{A})$. Notons $d^1$ l'application induite
$I^1 \to I^2$, et ainsi de suite. Cette construction fonctionne parce que
chacune des suites
$$
0 \to \Coker(u_n) \to I^{n + 1} \to \Coker(u_{n + 1}) \to 0
$$
est exacte courte, c'est-à-dire qu'elle induit une suite exacte courte après
application de $\text{gr}$. Pour le voir, utiliser Homologie, lemme
\ref{homology-lemma-characterize-strict}.
\end{proof}

\begin{lemma}
\label{lemma-filtered-injective-right-resolution-map}
Soit $\mathcal{A}$ une catégorie abélienne ayant assez d'injectifs.
Soit $f : A \to B$ un morphisme de $\text{Fil}^f(\mathcal{A})$.
Étant donné des quasi-isomorphismes filtrés
$A[0] \to I^\bullet$ et $B[0] \to J^\bullet$, où
$I^\bullet, J^\bullet$ sont des complexes d'objets injectifs filtrés tels que
$I^n = J^n = 0$ pour $n < 0$, il existe un diagramme commutatif
$$
\xymatrix{
A[0] \ar[r] \ar[d] &
B[0] \ar[d] \\
I^\bullet \ar[r] &
J^\bullet
}
$$
\end{lemma}

\begin{proof}
Puisque $A[0] \to I^\bullet$ et $B[0] \to J^\bullet$ sont des
quasi-isomorphismes filtrés, nous concluons que
$a : A \to I^0$, $b : B \to J^0$ et tous les morphismes
$d_I^n$, $d_J^n$ sont stricts, voir Homologie, lemme
\ref{homology-lemma-filtered-acyclic}.
Nous allons construire par récurrence les applications $f^n$ du diagramme
commutatif suivant :
$$
\xymatrix{
A \ar[r]_a \ar[d]_f &
I^0 \ar[r] \ar[d]^{f^0} &
I^1 \ar[r] \ar[d]^{f^1} &
I^2 \ar[r] \ar[d]^{f^2} &
\ldots \\
B \ar[r]^b &
J^0 \ar[r] &
J^1 \ar[r] &
J^2 \ar[r] &
\ldots
}
$$
Comme $A \to I^0$ est un monomorphisme strict et que $J^0$ est injectif
filtré, nous pouvons trouver un morphisme $f^0 : I^0 \to J^0$ tel que
$f^0 \circ a = b \circ f$, voir le
lemme \ref{lemma-injective-property-filtered-injective}.
La composée $d_J^0 \circ b \circ f$ est nulle ; ainsi,
$d_J^0 \circ f^0 \circ a = 0$, et $d_J^0 \circ f^0$ se factorise donc par un
unique morphisme
$$
\Coker(a) = \Coim(d_I^0) = \Im(d_I^0) \longrightarrow J^1.
$$
Comme $\Im(d_I^0) \to I^1$ est un monomorphisme strict, nous pouvons prolonger
la flèche affichée en un morphisme $f^1 : I^1 \to J^1$ en appliquant de
nouveau le lemme \ref{lemma-injective-property-filtered-injective}, et ainsi de
suite.
\end{proof}


\begin{lemma}
\label{lemma-filtered-injective-right-resolution-ses}
Soit $\mathcal{A}$ une catégorie abélienne ayant assez d'injectifs.
Soit $0 \to A \to B \to C \to 0$ une suite exacte courte dans
$\text{Fil}^f(\mathcal{A})$.
Étant donné des quasi-isomorphismes filtrés
$A[0] \to I^\bullet$ et $C[0] \to J^\bullet$, où
$I^\bullet, J^\bullet$ sont des complexes d'objets injectifs filtrés tels que
$I^n = J^n = 0$ pour $n < 0$, il existe un diagramme commutatif
$$
\xymatrix{
0 \ar[r] &
A[0] \ar[r] \ar[d] &
B[0] \ar[r] \ar[d] &
C[0] \ar[r] \ar[d] &
0 \\
0 \ar[r] &
I^\bullet \ar[r] &
M^\bullet \ar[r] &
J^\bullet \ar[r] &
0
}
$$
dont la ligne inférieure est une suite de complexes scindée terme à terme.
\end{lemma}

\begin{proof}
Puisque $A[0] \to I^\bullet$ et $C[0] \to J^\bullet$ sont des
quasi-isomorphismes filtrés, nous concluons que
$a : A \to I^0$, $c : C \to J^0$ et tous les morphismes $d_I^n$, $d_J^n$
sont stricts, voir Homologie, lemme
\ref{homology-lemma-filtered-acyclic}.
Nous allons construire pas à pas les flèches orientées vers le sud-est et
vers le sud dans le diagramme commutatif suivant :
$$
\xymatrix{
B \ar[r]_\beta \ar[rd]^b &
C \ar[r]_c \ar[rd]^{\overline{b}} &
J^0 \ar[d]^{\delta^0} \ar[r] &
J^1 \ar[d]^{\delta^1} \ar[r] & \ldots \\
A \ar[u]^\alpha \ar[r]^a &
I^0 \ar[r] &
I^1 \ar[r] &
I^2 \ar[r] & \ldots
}
$$
Comme $A \to B$ est un monomorphisme strict, il existe un morphisme
$b : B \to I^0$ tel que $b \circ \alpha = a$, voir le
lemme \ref{lemma-injective-property-filtered-injective}.
Comme $A$ est le noyau du morphisme strict $I^0 \to I^1$ et que
$\beta = \Coker(\alpha)$, nous obtenons un unique morphisme
$\overline{b} : C \to I^1$ s'insérant dans le diagramme.
Comme $c$ est un monomorphisme strict et que $I^1$ est injectif filtré, il
existe $\delta^0 : J^0 \to I^1$, voir le
lemme \ref{lemma-injective-property-filtered-injective}.
Puisque $B \to C$ est un épimorphisme strict et que
$B \to I^0 \to I^1 \to I^2$ est nul, nous voyons que
$C \to I^1 \to I^2$ est nul. Ainsi, $d_I^1 \circ \delta^0$ est nul sur
$C \cong \Im(c)$. Par conséquent, $d_I^1 \circ \delta^0$ se factorise par un
unique morphisme
$$
\Coker(c) = \Coim(d_J^0) = \Im(d_J^0) \longrightarrow I^2.
$$
Comme $I^2$ est injectif filtré et que
$\Im(d_J^0) \to J^1$ est un monomorphisme strict, nous pouvons prolonger le
morphisme affiché en un morphisme $\delta^1 : J^1 \to I^2$, voir le
lemme \ref{lemma-injective-property-filtered-injective}, et ainsi de suite.
Posons $M^\bullet = I^\bullet \oplus J^\bullet$, muni de la différentielle
$$
d_M^n =
\left(
\begin{matrix}
d_I^n & (-1)^{n + 1}\delta^n \\
0 & d_J^n
\end{matrix}
\right)
$$
Enfin, l'application $B[0] \to M^\bullet$ est donnée par
$b \oplus c \circ \beta : B \to I^0 \oplus J^0$.
\end{proof}

\begin{lemma}
\label{lemma-right-resolution-by-filtered-injectives}
Soit $\mathcal{A}$ une catégorie abélienne ayant assez d'injectifs.
Pour tout $K^\bullet \in K^{+}(\text{Fil}^f(\mathcal{A}))$, il existe un
quasi-isomorphisme filtré $K^\bullet \to I^\bullet$ tel que $I^\bullet$ soit
borné inférieurement, que chaque $I^n$ soit un objet injectif filtré et que
chaque $K^n \to I^n$ soit un monomorphisme strict.
\end{lemma}

\begin{proof}
Après avoir remplacé $K^\bullet$ par un décalé, ce qui est sans conséquence
pour la preuve, nous pouvons supposer que $K^n = 0$ pour $n < 0$.
Considérons les suites exactes courtes
$$
\begin{matrix}
0 \to \Ker(d_K^0) \to K^0 \to \Coim(d_K^0) \to 0 \\
0 \to \Ker(d_K^1) \to K^1 \to \Coim(d_K^1) \to 0 \\
0 \to \Ker(d_K^2) \to K^2 \to \Coim(d_K^2) \to 0 \\
\ldots
\end{matrix}
$$
de la catégorie exacte $\text{Fil}^f(\mathcal{A})$, et les applications
$u_i : \Coim(d_K^i) \to \Ker(d_K^{i + 1})$.
Pour chaque $i \geq 0$, nous pouvons choisir des quasi-isomorphismes filtrés
$$
\begin{matrix}
\Ker(d_K^i)[0] \to I_{ker, i}^\bullet \\
\Coim(d_K^i)[0] \to I_{coim, i}^\bullet
\end{matrix}
$$
où $I_{ker, i}^n, I_{coim, i}^n$ sont injectifs filtrés et nuls pour
$n < 0$, voir le
lemme \ref{lemma-filtered-injective-right-resolution-single-object}.
D'après le lemme \ref{lemma-filtered-injective-right-resolution-map}, nous
pouvons relever $u_i$ en un morphisme de complexes
$u_i^\bullet : I_{coim, i}^\bullet \to I_{ker, i + 1}^\bullet$.
Enfin, pour chaque $i \geq 0$, nous pouvons compléter les diagrammes
$$
\xymatrix{
0 \ar[r] &
\Ker(d_K^i)[0] \ar[r] \ar[d] &
K^i[0] \ar[r] \ar[d] &
\Coim(d_K^i)[0] \ar[r] \ar[d] &
0 \\
0 \ar[r] &
I_{ker, i}^\bullet \ar[r]^{\alpha_i} &
I_i^\bullet \ar[r]^{\beta_i} &
I_{coim, i}^\bullet \ar[r] &
0
}
$$
de sorte que la suite inférieure soit une suite exacte scindée terme à
terme, voir le
lemme \ref{lemma-filtered-injective-right-resolution-ses}.
Pour $i \geq 0$, posons
$d_i : I_i^\bullet \to I_{i + 1}^\bullet$ égal à
$d_i =  \alpha_{i + 1} \circ u_i^\bullet \circ \beta_i$.
Notons que $d_i \circ d_{i - 1} = 0$, car
$\beta_i \circ \alpha_i = 0$. Nous avons donc construit un diagramme
commutatif
$$
\xymatrix{
I_0^\bullet \ar[r] &
I_1^\bullet \ar[r] &
I_2^\bullet \ar[r] & \ldots \\
K^0[0] \ar[r] \ar[u] &
K^1[0] \ar[r] \ar[u] &
K^2[0] \ar[r] \ar[u] &
\ldots
}
$$
Ici, les flèches verticales sont des quasi-isomorphismes filtrés.
La ligne supérieure est un complexe de complexes, et chacun de ces complexes
est formé d'objets injectifs filtrés sans objet non nul en degré $< 0$.
Nous obtenons donc un complexe double en posant $I^{a, b} = I_a^b$, en prenant
pour
$$
d_1^{a, b} : I^{a, b} = I_a^b \to I_{a + 1}^b = I^{a + 1, b}
$$
l'application $d_a^b$, et pour
$$
d_2^{a, b} : I^{a, b} = I_a^b \to I_a^{b + 1} = I^{a, b + 1}
$$
l'application $d_{I_a}^b$. Notons
$\text{Tot}(I^{\bullet, \bullet})$ le complexe total associé à ce complexe
double, voir Homologie, définition
\ref{homology-definition-associated-simple-complex}.
Remarquons que les applications $K^n[0] \to I_n^\bullet$ proviennent
d'applications $K^n \to I^{n, 0}$ qui donnent une application de complexes
$$
K^\bullet \longrightarrow \text{Tot}(I^{\bullet, \bullet})
$$
Nous affirmons que c'est un quasi-isomorphisme filtré.
Comme $\text{gr}(-)$ est un foncteur additif, on a
$\text{gr}(\text{Tot}(I^{\bullet, \bullet})) =
\text{Tot}(\text{gr}(I^{\bullet, \bullet}))$.
Nous pouvons donc utiliser Homologie, lemme
\ref{homology-lemma-double-complex-gives-resolution}, pour conclure que
$\text{gr}(K^\bullet) \to \text{gr}(\text{Tot}(I^{\bullet, \bullet}))$ est un
quasi-isomorphisme, comme voulu.
\end{proof}









\begin{lemma}
\label{lemma-filtered-acyclic-is-zero}
Soit $\mathcal{A}$ une catégorie abélienne.
Soient $K^\bullet, I^\bullet \in K(\text{Fil}^f(\mathcal{A}))$.
Supposons que $K^\bullet$ soit filtré acyclique et que $I^\bullet$ soit borné
inférieurement et formé d'objets injectifs filtrés.
Tout morphisme $K^\bullet \to I^\bullet$ est homotope à zéro :
$\Hom_{K(\text{Fil}^f(\mathcal{A}))}(K^\bullet, I^\bullet) = 0$.
\end{lemma}

\begin{proof}
Soit $\alpha : K^\bullet \to I^\bullet$ un morphisme de complexes.
Supposons que $\alpha^j = 0$ pour $j < n$.
Nous allons montrer qu'il existe un morphisme $h : K^{n + 1} \to I^n$ tel
que $\alpha^n = h \circ d$. Ainsi, $\alpha$ sera homotope au morphisme de
complexes $\beta$ défini par
$$
\beta^j =
\left\{
\begin{matrix}
0 & \text{si} & j \leq n \\
\alpha^{n + 1} - d \circ h & \text{si} & j = n + 1 \\
\alpha^j & \text{si} & j > n + 1
\end{matrix}
\right.
$$
Cela prouvera clairement le lemme par récurrence.
Pour montrer l'existence de $h$, notons que
$\alpha^n \circ d_K^{n - 1} = 0$, puisque $\alpha^{n - 1} = 0$.
Comme $K^\bullet$ est filtré acyclique, $d_K^{n - 1}$ et $d_K^n$ sont
stricts, et
$$
0 \to \Im(d_K^{n - 1}) \to K^n \to \Im(d_K^n) \to 0
$$
est une suite exacte de la catégorie exacte
$\text{Fil}^f(\mathcal{A})$, voir Homologie, lemme
\ref{homology-lemma-filtered-acyclic}.
Nous pouvons donc considérer $\alpha^n$ comme une application vers $I^n$
définie sur $\Im(d_K^n)$.
Comme $\Im(d_K^n) \to K^{n + 1}$ est un monomorphisme strict et que $I^n$ est
injectif filtré, nous pouvons prolonger cette application en une application
$h : K^{n + 1} \to I^n$, comme voulu, voir le
lemme \ref{lemma-injective-property-filtered-injective}.
\end{proof}

\begin{lemma}
\label{lemma-morphisms-into-filtered-injective-complex}
Soit $\mathcal{A}$ une catégorie abélienne.
Soit $I^\bullet \in K(\text{Fil}^f(\mathcal{A}))$ un complexe borné
inférieurement formé d'objets injectifs filtrés.
\begin{enumerate}
\item Soit $\alpha : K^\bullet \to L^\bullet$ dans
$K(\text{Fil}^f(\mathcal{A}))$ un quasi-isomorphisme filtré.
Alors l'application
$$
\Hom_{K(\text{Fil}^f(\mathcal{A}))}(L^\bullet, I^\bullet)
\to
\Hom_{K(\text{Fil}^f(\mathcal{A}))}(K^\bullet, I^\bullet)
$$
est bijective.
\item Pour tout $L^\bullet \in K(\text{Fil}^f(\mathcal{A}))$, on a
$$
\Hom_{K(\text{Fil}^f(\mathcal{A}))}(L^\bullet, I^\bullet)
=
\Hom_{DF(\mathcal{A})}(L^\bullet, I^\bullet).
$$
\end{enumerate}
\end{lemma}

\begin{proof}
Preuve de (1). Notons que
$$
(K^\bullet, L^\bullet, C(\alpha)^\bullet, \alpha, i, -p)
$$
est un triangle distingué dans $K(\text{Fil}^f(\mathcal{A}))$
(lemme \ref{lemma-the-same-up-to-isomorphisms}) et que
$C(\alpha)^\bullet$ est un complexe filtré acyclique
(lemme \ref{lemma-filtered-acyclic}).
Alors
$$
\xymatrix{
\Hom_{K(\text{Fil}^f(\mathcal{A}))}(C(\alpha)^\bullet, I^\bullet) \ar[r] &
\Hom_{K(\text{Fil}^f(\mathcal{A}))}(L^\bullet, I^\bullet) \ar[r] &
\Hom_{K(\text{Fil}^f(\mathcal{A}))}(K^\bullet, I^\bullet) \ar[lld] \\
\Hom_{K(\text{Fil}^f(\mathcal{A}))}(C(\alpha)^\bullet[-1], I^\bullet)
}
$$
est une suite exacte de groupes abéliens, voir le
lemme \ref{lemma-representable-homological}.
Le lemme \ref{lemma-filtered-acyclic-is-zero} garantit alors que les deux
groupes extérieurs sont nuls, et par conséquent
$\Hom_{K(\text{Fil}^f(\mathcal{A}))}(L^\bullet, I^\bullet) =
\Hom_{K(\text{Fil}^f(\mathcal{A}))}(K^\bullet, I^\bullet)$.

\medskip\noindent
Preuve de (2).
Soit $a$ un élément du membre de droite.
Nous pouvons représenter $a = \gamma\alpha^{-1}$, où
$\alpha : K^\bullet \to L^\bullet$ est un quasi-isomorphisme filtré et
$\gamma : K^\bullet \to I^\bullet$ est une application de complexes.
D'après la partie (1), il existe un morphisme
$\beta : L^\bullet \to I^\bullet$ tel que $\beta \circ \alpha$ soit homotope
à $\gamma$. Cela prouve la surjectivité de l'application.
Soit $b$ un élément du membre de gauche qui s'envoie sur zéro dans le membre
de droite. Alors $b$ est la classe d'homotopie d'un morphisme
$\beta : L^\bullet \to I^\bullet$ tel qu'il existe un quasi-isomorphisme filtré
$\alpha : K^\bullet \to L^\bullet$ pour lequel $\beta  \circ \alpha$ est
homotope à zéro. La partie (1) montre alors que $\beta$ est lui aussi
homotope à zéro, c'est-à-dire que $b = 0$.
\end{proof}

\begin{lemma}
\label{lemma-filtered-localization-functor}
Soit $\mathcal{A}$ une catégorie abélienne ayant assez d'injectifs.
Notons $\mathcal{I}^f \subset \text{Fil}^f(\mathcal{A})$ la sous-catégorie
additive strictement pleine dont les objets sont les objets injectifs filtrés.
Le foncteur canonique
$$
K^{+}(\mathcal{I}^f)
\longrightarrow
DF^{+}(\mathcal{A})
$$
est exact, pleinement fidèle et essentiellement surjectif, c'est-à-dire que
c'est une équivalence de catégories triangulées. En outre, les diagrammes
$$
\xymatrix{
K^{+}(\mathcal{I}^f) \ar[d]_{\text{gr}^p} \ar[r] &
DF^{+}(\mathcal{A}) \ar[d]_{\text{gr}^p} \\
K^{+}(\mathcal{I}) \ar[r] &
D^{+}(\mathcal{A})
}
\quad
\xymatrix{
K^{+}(\mathcal{I}^f) \ar[d]^{\text{oubli de }F} \ar[r] &
DF^{+}(\mathcal{A}) \ar[d]^{\text{oubli de }F} \\
K^{+}(\mathcal{I}) \ar[r] &
D^{+}(\mathcal{A})
}
$$
sont commutatifs, où $\mathcal{I} \subset \mathcal{A}$ est la sous-catégorie
additive strictement pleine dont les objets sont les objets injectifs.
\end{lemma}

\begin{proof}
Le foncteur $K^{+}(\mathcal{I}^f) \to DF^{+}(\mathcal{A})$ est
essentiellement surjectif d'après le
lemme \ref{lemma-right-resolution-by-filtered-injectives}.
Il est pleinement fidèle d'après le
lemme \ref{lemma-morphisms-into-filtered-injective-complex}.
Il est exact par nos définitions relatives aux triangles distingués.
La commutativité des carrés est immédiate.
\end{proof}

\begin{remark}
\label{remark-filtered-localization-big}
Nous ne pouvons inverser la flèche du lemme que si $\mathcal{A}$ est une
catégorie au sens de nos conventions, c'est-à-dire si elle possède un ensemble
d'objets. Supposons cependant donnée une grande catégorie abélienne
$\mathcal{A}$ ayant assez d'injectifs, par exemple
$\textit{Mod}(\mathcal{O}_X)$. Pour tout ensemble donné d'objets
$\{A_i\}_{i\in I}$, il existe alors une sous-catégorie abélienne
$\mathcal{A}' \subset \mathcal{A}$ qui les contient tous et qui a assez
d'injectifs, voir Ensembles, lemme
\ref{sets-lemma-abelian-injectives}.
Nous pouvons donc appliquer le lemme précédent à $\mathcal{A}'$.
Cela signifie essentiellement que, tant que nous utilisons une collection
ensembliste de diagrammes, etc., l'application du lemme ne pose jamais de
problème.
\end{remark}

\noindent
Soient $\mathcal{A}, \mathcal{B}$ des catégories abéliennes.
Soit $T : \mathcal{A} \to \mathcal{B}$ un foncteur exact à gauche.
(Nous ne pouvons pas noter ce foncteur $F$, car cette lettre entre trop en
conflit avec notre notation $F$ pour les filtrations.) Notons que $T$ induit un
foncteur additif
$$
T : \text{Fil}^f(\mathcal{A}) \to \text{Fil}^f(\mathcal{B})
$$
par la règle $T(A, F) = (T(A), F)$, où $F^pT(A) = T(F^pA)$ ; cette définition
a un sens puisque $T$ est exact à gauche. (Attention : il n'est pas
nécessaire que $\text{gr}(T(A)) = T(\text{gr}(A))$.)
Cela induit des foncteurs de catégories triangulées
\begin{equation}
\label{equation-induced-T-filtered}
T :
K^{+}(\text{Fil}^f(\mathcal{A}))
\longrightarrow
K^{+}(\text{Fil}^f(\mathcal{B}))
\end{equation}
Le foncteur dérivé droit filtré de $T$ est le foncteur dérivé droit de la
définition \ref{definition-right-derived-functor-defined} appliqué au composé
de ce foncteur exact avec le foncteur exact
$K^{+}(\text{Fil}^f(\mathcal{B})) \to DF^{+}(\mathcal{B})$, relativement au
système multiplicatif $\text{FQis}^{+}(\mathcal{A})$.
Supposons que $\mathcal{A}$ ait assez d'injectifs. Nous pouvons alors reprendre
la discussion de la section \ref{section-right-derived-functor} pour définir
les {\it foncteurs dérivés droits filtrés}
\begin{equation}
\label{equation-filtered-derived-functor}
RT : DF^{+}(\mathcal{A}) \longrightarrow DF^{+}(\mathcal{B})
\end{equation}
de notre foncteur $T$.

\medskip\noindent
Nous suivrons toutefois plutôt la méthode de la
section \ref{section-right-derived-functor-via-resolutions}, et nous verrons
qu'il est même possible de définir $RT$ lorsque $T$ est seulement additif.
Choisissons d'abord un quasi-inverse
$j' : DF^{+}(\mathcal{A}) \to K^{+}(\mathcal{I}^f)$ de l'équivalence du
lemme \ref{lemma-filtered-localization-functor}.
D'après le lemme \ref{lemma-exact-equivalence}, $j'$ est un foncteur exact de
catégories triangulées. Notons ensuite que, pour tout objet injectif filtré
$I$, il existe une décomposition non canonique
\begin{equation}
\label{equation-decompose}
I \cong \bigoplus\nolimits_{p \in \mathbf{Z}} I_p,
\quad\text{avec}\quad
F^pI = \bigoplus\nolimits_{q \geq p} I_q
\end{equation}
d'après le lemme \ref{lemma-filtered-injective}.
Par conséquent, si $T$ est un foncteur additif quelconque
$T : \mathcal{A} \to \mathcal{B}$, nous obtenons un foncteur additif
\begin{equation}
\label{equation-extend-T}
T_{ext} : \mathcal{I}^f \to \text{Fil}^f(\mathcal{B})
\end{equation}
en posant $T_{ext}(I) = \bigoplus T(I_p)$, avec
$F^pT_{ext}(I) = \bigoplus_{q \geq p} T(I_q)$. Par construction, nous avons la
propriété $\text{gr}(T_{ext}(I)) = T(\text{gr}(I))$.
Nous obtenons donc un foncteur
\begin{equation}
\label{equation-extend-T-complexes}
T_{ext} : K^{+}(\mathcal{I}^f) \to K^{+}(\text{Fil}^f(\mathcal{B}))
\end{equation}
qui commute à $\text{gr}$. Nous définissons alors le foncteur de
(\ref{equation-filtered-derived-functor}) par la composée
\begin{equation}
\label{equation-definition-filtered-derived-functor}
RT = T_{ext} \circ j'.
\end{equation}
Comme le foncteur dérivé ordinaire
$RT : D^{+}(\mathcal{A}) \to D^{+}(\mathcal{B})$ se calcule lui aussi au moyen
de résolutions injectives, voir le
lemme \ref{lemma-injective-acyclic}, la commutation de $T$ à $\text{gr}$ et
les diagrammes commutatifs du
lemme \ref{lemma-filtered-localization-functor} impliquent que
\begin{equation}
\label{equation-commute-gr}
\text{gr}^p \circ RT \cong RT \circ \text{gr}^p
\end{equation}
et
\begin{equation}
\label{equation-commute-forget}
(\text{oubli de }F) \circ RT \cong RT \circ (\text{oubli de }F)
\end{equation}
comme foncteurs $DF^{+}(\mathcal{A}) \to D^{+}(\mathcal{B})$.

\medskip\noindent
Le foncteur dérivé filtré $RT$
(\ref{equation-filtered-derived-functor}) induit des foncteurs
$$
\begin{matrix}
RT : \text{Fil}^f(\mathcal{A}) \to DF^{+}(\mathcal{B}), \\
RT : \text{Comp}^{+}(\text{Fil}^f(\mathcal{A})) \to DF^{+}(\mathcal{B}), \\
RT : KF^{+}(\mathcal{A}) \to DF^{+}(\mathcal{B}).
\end{matrix}
$$
Comme $\text{Fil}^f(\mathcal{A})$ et
$\text{Comp}^{+}(\text{Fil}^f(\mathcal{A}))$ ne sont plus abéliennes, cela
n'a pas de sens d'affirmer que la restriction de $RT$ à ces catégories est
un $\delta$-foncteur. (On peut y remédier en considérant ces catégories comme
des catégories exactes et en formulant la notion de $\delta$-foncteur d'une
catégorie exacte vers une catégorie triangulée.) En revanche, il est pertinent
d'affirmer, et il résulte de la construction, que $RT$ est un foncteur exact
sur la catégorie triangulée $KF^{+}(\mathcal{A})$.

\begin{lemma}
\label{lemma-ss-filtered-derived}
Soient $\mathcal{A}, \mathcal{B}$ des catégories abéliennes. Soit
$T : \mathcal{A} \to \mathcal{B}$ un foncteur exact à gauche.
Supposons que $\mathcal{A}$ ait assez d'injectifs.
Soit $(K^\bullet, F)$ un objet de
$\text{Comp}^{+}(\text{Fil}^f(\mathcal{A}))$.
Il existe une suite spectrale $(E_r, d_r)_{r\geq 0}$ formée d'objets
bigradués $E_r$ de $\mathcal{B}$, où $d_r$ est de bidegré $(r, - r + 1)$,
et telle que
$$
E_1^{p, q} = R^{p + q}T(\text{gr}^p(K^\bullet))
$$
De plus, cette suite spectrale est bornée, converge vers
$R^*T(K^\bullet)$ et induit une filtration finie sur chaque
$R^nT(K^\bullet)$. La construction de cette suite spectrale est fonctorielle
en l'objet $K^\bullet$ de
$\text{Comp}^{+}(\text{Fil}^f(\mathcal{A}))$, et les termes $(E_r, d_r)$ pour
$r \geq 1$ ne dépendent d'aucun choix.
\end{lemma}

\begin{proof}
Choisissons un quasi-isomorphisme filtré $K^\bullet \to I^\bullet$, où
$I^\bullet$ est un complexe borné inférieurement d'objets injectifs filtrés,
voir le lemme \ref{lemma-right-resolution-by-filtered-injectives}.
Considérons le complexe $RT(K^\bullet) = T_{ext}(I^\bullet)$, voir
(\ref{equation-definition-filtered-derived-functor}).
Nous pouvons donc considérer la suite spectrale
$(E_r, d_r)_{r \geq 0}$ associée à ce complexe filtré de $\mathcal{B}$, voir
Homologie, section \ref{homology-section-filtered-complex}.
D'après Homologie, lemme
\ref{homology-lemma-spectral-sequence-filtered-complex}, nous avons
$E_1^{p, q} = H^{p + q}(\text{gr}^p(T(I^\bullet)))$.
D'après l'équation (\ref{equation-decompose}), nous avons
$E_1^{p, q} = H^{p + q}(T(\text{gr}^p(I^\bullet)))$, et, par définition d'une
résolution injective filtrée, l'application
$\text{gr}^p(K^\bullet) \to \text{gr}^p(I^\bullet)$ est une résolution
injective. Par conséquent,
$E_1^{p, q} = R^{p + q}T(\text{gr}^p(K^\bullet))$.

\medskip\noindent
D'autre part, chaque $I^n$ possède une filtration finie, et il en est donc de
même de chaque $T(I^n)$. Nous pouvons ainsi appliquer Homologie, lemme
\ref{homology-lemma-biregular-ss-converges}, pour conclure que la suite
spectrale est bornée, converge vers
$H^*(T(I^\bullet)) = R^*T(K^\bullet)$ et induit une filtration finie sur
chaque $H^n(T(I^\bullet)) = R^nT(K^\bullet)$.

\medskip\noindent
Supposons que $K^\bullet \to L^\bullet$ soit un morphisme de
$\text{Comp}^{+}(\text{Fil}^f(\mathcal{A}))$.
Choisissons un quasi-isomorphisme filtré $L^\bullet \to J^\bullet$, où
$J^\bullet$ est un complexe borné inférieurement d'objets injectifs filtrés,
voir le lemme \ref{lemma-right-resolution-by-filtered-injectives}.
D'après les résultats précédents, par exemple le
lemme \ref{lemma-morphisms-into-filtered-injective-complex}, il existe un
diagramme
$$
\xymatrix{
K^\bullet \ar[r] \ar[d] & L^\bullet \ar[d] \\
I^\bullet \ar[r] & J^\bullet
}
$$
commutatif à homotopie près. Nous obtenons ainsi un morphisme de complexes
filtrés $T(I^\bullet) \to T(J^\bullet)$ qui donne le morphisme de suites
spectrales, voir Homologie, lemme
\ref{homology-lemma-spectral-sequence-filtered-complex-functorial}.
La dernière assertion en résulte.
\end{proof}

\begin{remark}
\label{remark-final-functorial}
Comme annoncé dans la remarque \ref{remark-functorial-ss}, examinons le lien
entre le lemme précédent et les constructions utilisant les résolutions de
Cartan--Eilenberg. Soit donc $T : \mathcal{A} \to \mathcal{B}$ un foncteur
exact à gauche entre catégories abéliennes ; supposons que $\mathcal{A}$ ait
assez d'injectifs, et soit $K^\bullet$ un complexe borné inférieurement de
$\mathcal{A}$. Nous donnons une autre construction des suites spectrales
${}'E$ et ${}''E$ du lemme
\ref{lemma-two-ss-complex-functor}.

\medskip\noindent
Première suite spectrale. Considérons sur $K^\bullet$ la filtration « bête »
définie par $F^p(K^\bullet) = \sigma_{\geq p}(K^\bullet)$, voir Homologie,
section \ref{homology-section-truncations}.
Cette filtration est bête au sens où
$d(F^p(K^\bullet)) \subset F^{p + 1}(K^\bullet)$, comparer Homologie, lemme
\ref{homology-lemma-spectral-sequence-filtered-complex-d1}.
Pour cette filtration,
$\text{gr}^p(K^\bullet) = K^p[-p]$.
D'après le lemme \ref{lemma-ss-filtered-derived}, il existe une suite spectrale
dont le terme $E_1$ est
$$
E_1^{p, q} = R^{p + q}T(K^p[-p]) = R^qT(K^p)
$$
comme dans la suite spectrale ${}'E_r$. Notons en outre que les différentielles
$E_1^{p, q} \to E_1^{p + 1, q}$ coïncident avec celles de $'{}E_1$, voir
Homologie, lemme
\ref{homology-lemma-spectral-sequence-filtered-complex-d1}, partie (2), et la
description de ${}'d_1$ dans la preuve du
lemme \ref{lemma-two-ss-complex-functor}.

\medskip\noindent
Seconde suite spectrale. Considérons sur le complexe $K^\bullet$ la filtration
définie par $F^p(K^\bullet) = \tau_{\leq -p}(K^\bullet)$, voir Homologie,
section \ref{homology-section-truncations}.
Le signe moins est nécessaire pour obtenir une filtration décroissante.
Pour cette filtration, $\text{gr}^p(K^\bullet)$ est quasi-isomorphe à
$H^{-p}(K^\bullet)[p]$.
D'après le lemme \ref{lemma-ss-filtered-derived}, il existe une suite spectrale
dont le terme $E_1$ est
$$
E_1^{p, q} = R^{p + q}T(H^{-p}(K^\bullet)[p])
= R^{2p + q}T(H^{-p}(K^\bullet)) = {}''E_2^{i, j}
$$
avec $i = 2p + q$ et $j = -p$. (Cette expression paraît peu naturelle, mais
nous aurions tout aussi bien pu développer toute la théorie des complexes
filtrés avec des filtrations croissantes ; le résultat aurait alors paru
naturel ici, mais pas pour l'autre suite.) Nous laissons au lecteur le soin de
vérifier que les différentielles coïncident.

\medskip\noindent
En fait, étant donnée une résolution de Cartan--Eilenberg
$K^\bullet \to I^{\bullet, \bullet}$, le morphisme induit
$K^\bullet \to \text{Tot}(I^{\bullet, \bullet})$ vers le complexe total
associé est une résolution injective filtrée pour chacune des deux filtrations,
si l'on munit $\text{Tot}(I^{\bullet, \bullet})$ d'une filtration appropriée.
On peut s'en servir pour identifier exactement les suites spectrales.
\end{remark}

\section{Groupes d'extensions}
\label{section-ext}

\noindent
Dans cette section, nous commençons à décrire les groupes d'extensions des
objets d'une catégorie abélienne. Nous avons d'abord la définition très
générale suivante.

\begin{definition}
\label{definition-ext}
Soit $\mathcal{A}$ une catégorie abélienne. Soit $i \in \mathbf{Z}$. Soient
$X, Y$ des objets de $D(\mathcal{A})$. Le {\it $i$-ième groupe d'extensions}
de $X$ par $Y$ est le groupe
$$
\Ext^i_\mathcal{A}(X, Y) =
\Hom_{D(\mathcal{A})}(X, Y[i]) =
\Hom_{D(\mathcal{A})}(X[-i], Y).
$$
Si $A, B \in \Ob(\mathcal{A})$, nous posons
$\Ext^i_\mathcal{A}(A, B) = \text{Ext}^i_\mathcal{A}(A[0], B[0])$.
\end{definition}

\noindent
Comme $\Hom_{D(\mathcal{A})}(X, -)$, respectivement
$\Hom_{D(\mathcal{A})}(-, Y)$, est un foncteur homologique, respectivement
cohomologique, voir le lemme \ref{lemma-representable-homological}, un triangle
distingué $(Y, Y', Y'')$, respectivement $(X, X', X'')$, donne lieu à une
suite exacte longue
$$
\ldots \to
\Ext^i_\mathcal{A}(X, Y) \to
\Ext^i_\mathcal{A}(X, Y') \to
\Ext^i_\mathcal{A}(X, Y'') \to
\Ext^{i + 1}_\mathcal{A}(X, Y) \to \ldots
$$
respectivement
$$
\ldots \to
\Ext^i_\mathcal{A}(X'', Y) \to
\Ext^i_\mathcal{A}(X', Y) \to
\Ext^i_\mathcal{A}(X, Y) \to
\Ext^{i + 1}_\mathcal{A}(X'', Y) \to \ldots
$$
Notons que, puisque $D^+(\mathcal{A})$, $D^-(\mathcal{A})$ et
$D^b(\mathcal{A})$ sont des sous-catégories pleines, nous pouvons calculer
les groupes d'extensions comme des groupes de morphismes dans ces catégories,
pourvu que $X$ et $Y$ y appartiennent.

\medskip\noindent
Lorsque la catégorie $\mathcal{A}$ a assez d'injectifs ou assez de projectifs,
nous pouvons calculer les groupes d'extensions au moyen de résolutions
injectives ou projectives. Pour éviter toute confusion, rappelons que
l'existence d'une résolution injective, respectivement projective, implique
l'annulation de l'homologie en tout degré suffisamment petit, respectivement
suffisamment grand, voir les lemmes
\ref{lemma-cohomology-bounded-below} et
\ref{lemma-cohomology-bounded-above}.

\begin{lemma}
\label{lemma-compute-ext-resolutions}
Soit $\mathcal{A}$ une catégorie abélienne.
Soient $X^\bullet, Y^\bullet \in \Ob(K(\mathcal{A}))$.
\begin{enumerate}
\item Soit $Y^\bullet \to I^\bullet$ une résolution injective
(définition \ref{definition-injective-resolution}). Alors
$$
\Ext^i_\mathcal{A}(X^\bullet, Y^\bullet) =
\Hom_{K(\mathcal{A})}(X^\bullet, I^\bullet[i]).
$$
\item Soit $P^\bullet \to X^\bullet$ une résolution projective
(définition \ref{definition-projective-resolution}). Alors
$$
\Ext^i_\mathcal{A}(X^\bullet, Y^\bullet) =
\Hom_{K(\mathcal{A})}(P^\bullet[-i], Y^\bullet).
$$
\end{enumerate}
\end{lemma}

\begin{proof}
Cela résulte immédiatement du
lemme \ref{lemma-morphisms-into-injective-complex} et du
lemme \ref{lemma-morphisms-from-projective-complex}.
\end{proof}

\noindent
Dans la suite de cette section, nous étudions les extensions d'objets de la
catégorie abélienne elle-même. Commençons par l'observation suivante.

\begin{lemma}
\label{lemma-negative-exts}
Soit $\mathcal{A}$ une catégorie abélienne.
\begin{enumerate}
\item Soient $X$, $Y$ des objets de $D(\mathcal{A})$. Étant donné
$a, b \in \mathbf{Z}$ tels que $H^i(X) = 0$ pour $i > a$ et
$H^j(Y) = 0$ pour $j < b$, on a
$\Ext^n_\mathcal{A}(X, Y) = 0$ pour $n < b - a$ et
$$
\Ext^{b - a}_\mathcal{A}(X, Y) = \Hom_\mathcal{A}(H^a(X), H^b(Y))
$$
\item Soient $A, B \in \Ob(\mathcal{A})$.
Pour $i < 0$, on a $\Ext^i_\mathcal{A}(B, A) = 0$.
On a $\Ext^0_\mathcal{A}(B, A) = \Hom_\mathcal{A}(B, A)$.
\end{enumerate}
\end{lemma}

\begin{proof}
Choisissons des complexes $X^\bullet$ et $Y^\bullet$ représentant $X$ et $Y$.
Comme $Y^\bullet \to \tau_{\geq b}Y^\bullet$ est un quasi-isomorphisme, nous
pouvons supposer que $Y^j = 0$ pour $j < b$.
Soit $L^\bullet \to X^\bullet$ un quasi-isomorphisme quelconque.
Alors $\tau_{\leq a}L^\bullet \to X^\bullet$ est un quasi-isomorphisme.
Un morphisme $X \to Y[n]$ dans $D(\mathcal{A})$ peut donc se représenter sous
la forme $fs^{-1}$, où $s : L^\bullet \to X^\bullet$ est un
quasi-isomorphisme, $f : L^\bullet \to Y^\bullet[n]$ un morphisme, et
$L^i = 0$ pour $i > a$. Notons que $f$ envoie $L^i$ dans $Y^{i + n}$.
Ainsi, $f = 0$ si $n < b - a$, car l'un au moins de $L^i$ et de
$Y^{i + n}$ est toujours nul. Si $n = b - a$, alors $f$ correspond exactement
à un morphisme $H^a(X) \to H^b(Y)$. La partie (2) est un cas particulier de
(1).
\end{proof}

\noindent
Soit $\mathcal{A}$ une catégorie abélienne.
Supposons que $0 \to A \to A' \to A'' \to 0$ soit une suite exacte courte
d'objets de $\mathcal{A}$. Alors
$0 \to A[0] \to A'[0] \to A''[0] \to 0$ donne lieu à un triangle distingué
dans $D(\mathcal{A})$ (voir le
lemme \ref{lemma-derived-canonical-delta-functor}), donc à une suite exacte
longue de groupes d'extensions
$$
0 \to \Ext^0_\mathcal{A}(B, A) \to
\Ext^0_\mathcal{A}(B, A') \to
\Ext^0_\mathcal{A}(B, A'') \to
\Ext^1_\mathcal{A}(B, A) \to \ldots
$$
De même, étant donnée une suite exacte courte
$0 \to B \to B' \to B'' \to 0$, nous obtenons une suite exacte longue de
groupes d'extensions
$$
0 \to \Ext^0_\mathcal{A}(B'', A) \to
\Ext^0_\mathcal{A}(B', A) \to
\Ext^0_\mathcal{A}(B, A) \to
\Ext^1_\mathcal{A}(B'', A) \to \ldots
$$
Nous pouvons voir ces groupes d'extensions comme une application de la
construction de la catégorie dérivée. Cette construction montre que l'on peut
définir les groupes d'extensions et construire leur suite exacte longue sans
supposer l'existence d'assez d'injectifs ou de projectifs. Une autre
construction des groupes d'extensions, due à Yoneda, évite le recours à la
catégorie dérivée, voir \cite{Yoneda}.

\begin{definition}
\label{definition-yoneda-extension}
Soit $\mathcal{A}$ une catégorie abélienne. Soient
$A, B \in \Ob(\mathcal{A})$.
Pour $i \geq 1$, une {\it extension de Yoneda de degré $i$} de $B$ par $A$
est une suite exacte
$$
E : 0 \to A \to Z_{i - 1} \to Z_{i - 2} \to \ldots \to Z_0 \to B \to 0
$$
dans $\mathcal{A}$. Nous disons que deux extensions de Yoneda $E$ et $E'$ de
même degré sont {\it équivalentes} s'il existe un diagramme commutatif
$$
\xymatrix{
0 \ar[r] & A \ar[r] & Z_{i - 1} \ar[r] & \ldots \ar[r] &
Z_0 \ar[r] & B \ar[r] & 0 \\
0 \ar[r] &
A \ar[r] \ar[u]^{\text{id}} \ar[d]_{\text{id}} &
Z''_{i - 1} \ar[r] \ar[u] \ar[d] &
\ldots \ar[r] &
Z''_0 \ar[r] \ar[u] \ar[d] &
B \ar[r] \ar[u]_{\text{id}} \ar[d]^{\text{id}} & 0 \\
0 \ar[r] & A \ar[r] & Z'_{i - 1} \ar[r] & \ldots \ar[r] &
Z'_0 \ar[r] & B \ar[r] & 0
}
$$
où la ligne du milieu est elle aussi une extension de Yoneda.
\end{definition}

\noindent
Il n'est pas immédiat que l'équivalence de la définition soit une relation
d'équivalence. Il est instructif de le prouver directement, mais cela
résultera aussi du lemme \ref{lemma-yoneda-extension} ci-dessous.

\medskip\noindent
Soit $\mathcal{A}$ une catégorie abélienne possédant des objets $A$, $B$.
Étant donnée une extension de Yoneda
$E : 0 \to A \to Z_{i - 1} \to Z_{i - 2} \to \ldots \to Z_0 \to B \to 0$,
nous définissons un élément associé
$\delta(E) \in \Ext^i_\mathcal{A}(B, A)$ comme le morphisme
$\delta(E) = fs^{-1} : B[0] \to A[i]$, où $s$ est le quasi-isomorphisme
$$
(\ldots \to 0 \to A \to Z_{i - 1} \to \ldots \to Z_0 \to 0 \to \ldots)
\longrightarrow
B[0]
$$
et $f$ le morphisme de complexes
$$
(\ldots \to 0 \to A \to Z_{i - 1} \to \ldots \to Z_0 \to 0 \to \ldots)
\longrightarrow
A[i]
$$
Nous appelons $\delta(E) = fs^{-1}$ la {\it classe} de l'extension de Yoneda.
Cette classe caractérise en fait la classe d'équivalence de l'extension de
Yoneda.

\begin{lemma}
\label{lemma-yoneda-extension}
Soit $\mathcal{A}$ une catégorie abélienne possédant des objets $A$, $B$, et
soit $i \geq 1$.
Tout élément de $\Ext^i_\mathcal{A}(B, A)$ est de la forme $\delta(E)$ pour
une extension de Yoneda de degré $i$ de $B$ par $A$.
Étant donné deux extensions de Yoneda $E$, $E'$ de même degré, $E$ est
équivalente à $E'$ si et seulement si $\delta(E) = \delta(E')$.
\end{lemma}

\begin{proof}
Soit $\xi : B[0] \to A[i]$ un élément de $\Ext^i_\mathcal{A}(B, A)$.
Nous pouvons écrire $\xi = f s^{-1}$ pour un quasi-isomorphisme
$s : L^\bullet \to B[0]$ et un morphisme $f : L^\bullet \to A[i]$.
Après avoir remplacé $L^\bullet$ par $\tau_{\leq 0}L^\bullet$, nous pouvons
supposer que $L^j = 0$ pour $j > 0$. Représentons la situation par le diagramme
$$
\xymatrix{
L^{- i - 1} \ar[r] & L^{-i} \ar[r] \ar[d] & \ldots \ar[r] &
L^0 \ar[r] & B \ar[r] & 0 \\
& A
}
$$
En posant alors $Z_{i - 1} = (L^{- i + 1} \oplus A)/L^{-i}$ et
$Z_j = L^{-j}$ pour $j = i - 2, \ldots, 0$, nous obtenons une extension de
degré $i$, notée $E$, de $B$ par $A$, dont la classe $\delta(E)$ est égale à
$\xi$.

\medskip\noindent
Il résulte immédiatement des définitions que deux extensions de Yoneda
équivalentes ont la même classe. Supposons que
$E : 0 \to A \to Z_{i - 1} \to Z_{i - 2} \to \ldots \to Z_0 \to B \to 0$ et
$E' : 0 \to A \to Z'_{i - 1} \to Z'_{i - 2} \to \ldots \to Z'_0 \to B \to 0$
soient des extensions de Yoneda de même classe.
Par construction de $D(\mathcal{A})$ comme localisation de
$K(\mathcal{A})$ par rapport à l'ensemble des quasi-isomorphismes, cela
signifie qu'il existe un complexe $L^\bullet$ et des quasi-isomorphismes
$$
t : L^\bullet \to
(\ldots \to 0 \to A \to Z_{i - 1} \to \ldots \to Z_0 \to 0 \to \ldots)
$$
et
$$
t' : L^\bullet \to
(\ldots \to 0 \to A \to Z'_{i - 1} \to \ldots \to Z'_0 \to 0 \to \ldots)
$$
tels que $s \circ t = s' \circ t'$ et $f \circ t = f' \circ t'$; voir
Catégories, section \ref{categories-section-localization}.
Soit $E''$ l'extension de degré $i$ de $B$ par $A$ construite à partir de
la paire $L^\bullet \to B[0]$ et $L^\bullet \to A[i]$ dans le premier
paragraphe de la preuve. Le lecteur constate alors aisément qu'il existe
des « morphismes » d'extensions de Yoneda de degré $i$, $E'' \to E$ et
$E'' \to E'$, comme dans la définition de l'équivalence des extensions
de Yoneda (détails omis). Cela achève la preuve.
\end{proof}

\begin{lemma}
\label{lemma-ext-1}
Soit $\mathcal{A}$ une catégorie abélienne. Soient $A$, $B$ des objets
de $\mathcal{A}$. Alors $\Ext^1_\mathcal{A}(B, A)$ est le groupe
$\Ext_\mathcal{A}(B, A)$ construit dans
Homologie, Définition \ref{homology-definition-ext-group}.
\end{lemma}

\begin{proof}
C'est le cas $i = 1$ du
Lemme \ref{lemma-yoneda-extension}.
\end{proof}

\noindent
Étant donnés une catégorie abélienne $\mathcal{A}$ et des objets $X, Y, Z$
de $D(\mathcal{A})$, il existe des lois de composition bilinéaires et
associatives
$$
\Ext^j_\mathcal{A}(Y, Z) \times \Ext^i_\mathcal{A}(X, Y)
\longrightarrow
\Ext^{i + j}_\mathcal{A}(X, Z),\quad
(\eta, \xi) \longmapsto \eta \circ \xi
$$
En effet, si $\xi : X \to Y[i]$ et $\eta : Y \to Z[j]$, nous définissons
$\eta \circ \xi$ comme la composée de $\xi$ avec
$\eta[i] : Y[i] \to Z[i + j]$. Si $A, B, C$ sont des objets de
$\mathcal{A}$ et $i, j \geq 1$, cette loi de composition
$$
\Ext^j_\mathcal{A}(B, C) \times \Ext^i_\mathcal{A}(A, B)
\longrightarrow
\Ext^{i + j}_\mathcal{A}(A, C)
$$
se décrit à l'aide d'extensions de Yoneda comme suit : la composée de
$$
0 \to A \to Z_{i - 1} \to Z_{i - 2} \to \ldots \to Z_0 \to B \to 0
$$
et de
$$
0 \to B \to Z'_{j - 1} \to Z'_{j - 2} \to \ldots \to Z'_0 \to C \to 0
$$
est l'extension de Yoneda
$$
0 \to A \to Z_{i - 1} \to Z_{i - 2} \to \ldots \to Z_0 \to 
Z'_{j - 1} \to Z'_{j - 2} \to \ldots \to Z'_0 \to C \to 0
$$

\begin{lemma}
\label{lemma-cup-ext-1-zero}
Soit $\mathcal{A}$ une catégorie abélienne. Soient
$0 \to A \to Z \to B \to 0$ et
$0 \to B \to Z' \to C \to 0$ des suites exactes courtes dans $\mathcal{A}$.
Notons $[Z] \in \Ext^1_\mathcal{A}(B, A)$ et
$[Z'] \in \Ext^1_\mathcal{A}(C, B)$ leurs classes.
Alors $[Z] \circ [Z'] \in \Ext^2_\mathcal{A}(C, A)$ est égal à $0$ si et seulement
s'il existe un diagramme commutatif
$$
\xymatrix{
&
&
0 \ar[d] &
0 \ar[d]
\\
0 \ar[r] &
A \ar[r] \ar[d]^1 &
Z \ar[r] \ar[d] &
B \ar[r] \ar[d] &
0 \\
0 \ar[r] &
A \ar[r]  &
W \ar[r] \ar[d] &
Z' \ar[r] \ar[d] &
0 \\
&
&
C \ar[r]^1  \ar[d]&
C \ar[d]\\
&
&
0 &
0 
}
$$
dont les lignes et les colonnes sont exactes dans $\mathcal{A}$.
\end{lemma}

\begin{proof}
Omis. Indications : on peut raisonner à l'aide du résultat du
Lemme \ref{lemma-yoneda-extension} et expliciter ce que signifie la nullité
d'une classe de $2$-extension. On peut aussi remarquer que, si
$[Z] \circ [Z'] \in \Ext^2_\mathcal{A}(C, A)$ est nul, alors la suite exacte
longue de cohomologie de $\Ext$ montre que l'élément
$[Z] \in \Ext^1_\mathcal{A}(B, A)$ est l'image d'un élément de
$\Ext^1_\mathcal{A}(Z', A)$.
\end{proof}

\begin{lemma}
\label{lemma-higher-ext-zero}
Soit $\mathcal{A}$ une catégorie abélienne et soit $p \geq 0$.
Si $\Ext^p_\mathcal{A}(B, A) = 0$ pour toute paire d'objets $A$, $B$
de $\mathcal{A}$, alors $\Ext^i_\mathcal{A}(B, A) = 0$ pour
$i \geq p$ et toute paire d'objets $A$, $B$ de $\mathcal{A}$.
\end{lemma}

\begin{proof}
Pour $i > p$, écrivons toute classe $\xi$ sous la forme $\delta(E)$,
où $E$ est une extension de Yoneda
$$
E : 0 \to A \to Z_{i - 1} \to Z_{i - 2} \to \ldots \to Z_0 \to B \to 0
$$
C'est possible d'après le Lemme \ref{lemma-yoneda-extension}.
Posons $C = \Ker(Z_{p - 1} \to Z_{p - 2}) = \Im(Z_p \to Z_{p - 1})$.
Alors $\delta(E)$ est la composée de $\delta(E')$ et de $\delta(E'')$, où
$$
E' : 0 \to C \to Z_{p - 1} \to \ldots \to Z_0 \to B \to 0
$$
et
$$
E'' : 0 \to A \to Z_{i - 1} \to Z_{i - 2} \to \ldots \to Z_p \to C \to 0
$$
Comme $\delta(E') \in \Ext^p_\mathcal{A}(B, C) = 0$, nous concluons.
\end{proof}

\begin{lemma}
\label{lemma-ext-2-zero-pre}
Soit $\mathcal{A}$ une catégorie abélienne. Soit $K$ un objet de
$D^b(\mathcal{A})$ tel que $\Ext^p_\mathcal{A}(H^i(K), H^j(K)) = 0$
pour tous $p \geq 2$ et $i > j$. Alors $K$ est isomorphe à la somme directe
de ses cohomologies : $K \cong \bigoplus H^i(K)[-i]$.
\end{lemma}

\begin{proof}
Choisissons $a, b$ tels que $H^i(K) = 0$ pour $i \not \in [a, b]$.
Nous allons démontrer le lemme par récurrence sur $b - a$. Si $b - a \leq 0$,
le résultat est clair. Si $b - a > 0$, considérons le triangle distingué
de troncations
$$
\tau_{\leq b - 1}K \to K \to H^b(K)[-b] \to (\tau_{\leq b - 1}K)[1]
$$
voir Remarque \ref{remark-truncation-distinguished-triangle}.
D'après le Lemme \ref{lemma-split}, si la dernière flèche est nulle, alors
$K \cong \tau_{\leq b - 1}K \oplus H^b(K)[-b]$, et nous concluons par
récurrence. En utilisant encore l'hypothèse de récurrence, nous obtenons
$$
\Hom_{D(\mathcal{A})}(H^b(K)[-b], (\tau_{\leq b - 1}K)[1]) =
\bigoplus\nolimits_{i < b} \Ext_\mathcal{A}^{b - i + 1}(H^b(K), H^i(K))
$$
Par hypothèse, cette somme directe est nulle, ce qui achève la preuve.
\end{proof}

\begin{lemma}
\label{lemma-ext-2-zero}
Soit $\mathcal{A}$ une catégorie abélienne. Supposons que
$\Ext^2_\mathcal{A}(B, A) = 0$ pour toute paire d'objets $A$, $B$
de $\mathcal{A}$.
Alors tout objet $K$ de $D^b(\mathcal{A})$ est isomorphe à la somme directe
de ses cohomologies : $K \cong \bigoplus H^i(K)[-i]$.
\end{lemma}

\begin{proof}
L'hypothèse implique que $\Ext^i_\mathcal{A}(B, A) = 0$ pour $i \geq 2$
et toute paire d'objets $A, B$ de $\mathcal{A}$, d'après le
Lemme \ref{lemma-higher-ext-zero}. Le présent lemme est donc un cas
particulier du Lemme \ref{lemma-ext-2-zero-pre}.
\end{proof}

\section{Groupes de K-théorie}
\label{section-K-groups}

\noindent
Quelques mots sur $K_0$ d'une catégorie triangulée.

\begin{definition}
\label{definition-K-zero}
Soit $\mathcal{D}$ une catégorie triangulée. Nous notons $K_0(\mathcal{D})$ le
{\it groupe de $K$-théorie d'indice zéro de $\mathcal{D}$}. C'est le groupe
abélien construit comme suit. On prend le groupe abélien libre sur les objets
de $\mathcal{D}$ et, pour tout triangle distingué $X \to Y \to Z$, on impose
la relation $[Y] - [X] - [Z] = 0$.
\end{definition}

\noindent
Remarquons que cela entraîne $[X[n]] = (-1)^n[X]$, puisque nous avons le
triangle distingué $(X, 0, X[1], 0, 0, -\text{id}[1])$.

\begin{lemma}
\label{lemma-K-bounded-derived}
Soit $\mathcal{A}$ une catégorie abélienne. Il existe alors une identification
canonique $K_0(D^b(\mathcal{A})) = K_0(\mathcal{A})$ des groupes de $K$-théorie
d'indice zéro.
\end{lemma}

\begin{proof}
Étant donné un objet $A$ de $\mathcal{A}$, notons $A[0]$ l'objet $A$
considéré comme un complexe concentré en degré $0$.
Si $0 \to A \to A' \to A'' \to 0$ est une suite exacte courte, nous obtenons
un triangle distingué $A[0] \to A'[0] \to A''[0] \to A[1]$; voir
section \ref{section-canonical-delta-functor}.
Cela montre que nous obtenons une application
$K_0(\mathcal{A}) \to K_0(D^b(\mathcal{A}))$ envoyant $[A]$ sur $[A[0]]$,
en priant le lecteur d'excuser cette notation épouvantable.

\medskip\noindent
D'autre part, étant donné un objet $X$ de $D^b(\mathcal{A})$, nous pouvons
considérer l'élément
$$
c(X) = \sum (-1)^i[H^i(X)] \in K_0(\mathcal{A})
$$
Étant donné un triangle distingué $X \to Y \to Z$, la suite exacte longue
de cohomologie (\ref{equation-long-exact-cohomology-sequence-D}) et les
relations dans $K_0(\mathcal{A})$ montrent que $c(Y) = c(X) + c(Z)$.
Ainsi, $c$ se factorise en une application
$c : K_0(D^b(\mathcal{A})) \to K_0(\mathcal{A})$.

\medskip\noindent
Nous voulons montrer que les deux applications ci-dessus sont inverses l'une
de l'autre. Il est clair que la composée $K_0(\mathcal{A}) \to
K_0(D^b(\mathcal{A})) \to K_0(\mathcal{A})$ est l'identité.
Supposons que $X^\bullet$ soit un complexe borné de $\mathcal{A}$.
L'existence des triangles distingués associés aux « troncations bêtes » (voir
Homologie, section \ref{homology-section-truncations})
$$
\sigma_{\geq n}X^\bullet \to \sigma_{\geq n - 1}X^\bullet \to
X^{n - 1}[-n + 1] \to (\sigma_{\geq n}X^\bullet)[1]
$$
et une récurrence montrent que
$$
[X^\bullet] = \sum (-1)^i[X^i[0]]
$$
dans $K_0(D^b(\mathcal{A}))$ (en priant à nouveau le lecteur d'excuser la
notation). Il s'ensuit que l'application
$K_0(\mathcal{A}) \to K_0(D^b(\mathcal{A}))$ est surjective, ce qui achève
la preuve.
\end{proof}

\begin{lemma}
\label{lemma-map-K}
Soit $F : \mathcal{D} \to \mathcal{D}'$ un foncteur exact entre catégories
triangulées. Alors $F$ induit un homomorphisme de groupes
$K_0(\mathcal{D}) \to K_0(\mathcal{D}')$.
\end{lemma}

\begin{proof}
Omis.
\end{proof}

\begin{lemma}
\label{lemma-homological-map-K}
Soit $H : \mathcal{D} \to \mathcal{A}$ un foncteur homologique d'une catégorie
triangulée vers une catégorie abélienne. Supposons que, pour tout $X$ dans
$\mathcal{D}$, seul un nombre fini des objets $H(X[i])$ soient non nuls dans
$\mathcal{A}$. Alors $H$ induit un homomorphisme de groupes
$K_0(\mathcal{D}) \to K_0(\mathcal{A})$ envoyant $[X]$ sur
$\sum (-1)^i[H(X[i])]$.
\end{lemma}

\begin{proof}
Omis.
\end{proof}

\begin{lemma}
\label{lemma-DBA-map-K}
Soit $\mathcal{B}$ une sous-catégorie de Serre faible de la catégorie
abélienne $\mathcal{A}$. Il existe un isomorphisme canonique
$$
K_0(\mathcal{B}) \longrightarrow
K_0(D^b_\mathcal{B}(\mathcal{A})),\quad
[B] \longmapsto [B[0]]
$$
L'isomorphisme inverse envoie la classe $[X]$ de $X$ sur l'élément
$\sum (-1)^i[H^i(X)]$.
\end{lemma}

\begin{proof}
Nous omettons la vérification que la règle donnée pour l'inverse définit bien
une application $K_0(D^b_\mathcal{B}(\mathcal{A})) \to K_0(\mathcal{B})$.
Il est immédiat que la composée
$K_0(\mathcal{B}) \to
K_0(D^b_\mathcal{B}(\mathcal{A})) \to
K_0(\mathcal{B})$ est l'identité. D'autre part, en utilisant les triangles
distingués de la Remarque \ref{remark-truncation-distinguished-triangle} et un
raisonnement par récurrence, le lecteur peut montrer que la flèche affichée
dans l'énoncé du lemme est surjective (détails omis). Le lemme en résulte.
\end{proof}

\begin{lemma}
\label{lemma-bilinear-map-K}
Soient $\mathcal{D}$, $\mathcal{D}'$, $\mathcal{D}''$ des catégories
triangulées. Soit
$$
\otimes : \mathcal{D} \times \mathcal{D}' \longrightarrow \mathcal{D}''
$$
un foncteur tel que, pour $X$ fixé dans $\mathcal{D}$, le foncteur
$X \otimes - : \mathcal{D}' \to \mathcal{D}''$ soit exact et que, pour $X'$
fixé dans $\mathcal{D}'$, le foncteur
$- \otimes X' : \mathcal{D} \to \mathcal{D}''$ soit exact. Alors
$\otimes$ induit une application bilinéaire
$K_0(\mathcal{D}) \times K_0(\mathcal{D}') \to K_0(\mathcal{D}'')$
qui envoie $([X], [X'])$ sur $[X \otimes X']$.
\end{lemma}

\begin{proof}
Omis.
\end{proof}

\section{Complexes non bornés}
\label{section-unbounded}

\noindent
La référence pour les résultats de cette section est \cite{Spaltenstein}.
Le lemme suivant est utile pour trouver de « bonnes » résolutions à gauche
de complexes non bornés.

\begin{lemma}
\label{lemma-special-direct-system}
Soit $\mathcal{A}$ une catégorie abélienne. Soit
$\mathcal{P} \subset \Ob(\mathcal{A})$ un sous-ensemble.
Supposons que $\mathcal{P}$ contienne $0$, soit stable par sommes directes
(finies), et que tout objet de $\mathcal{A}$ soit un quotient d'un élément
de $\mathcal{P}$. Soit $K^\bullet$ un complexe.
Il existe un diagramme commutatif
$$
\xymatrix{
P_1^\bullet \ar[d] \ar[r] & P_2^\bullet \ar[d] \ar[r] & \ldots \\
\tau_{\leq 1}K^\bullet \ar[r] & \tau_{\leq 2}K^\bullet \ar[r] & \ldots
}
$$
dans la catégorie des complexes, tel que
\begin{enumerate}
\item les flèches verticales soient des quasi-isomorphismes surjectifs
terme à terme,
\item $P_n^\bullet$ soit un complexe borné à droite dont les termes
appartiennent à $\mathcal{P}$,
\item les flèches $P_n^\bullet \to P_{n + 1}^\bullet$ soient des injections
scindées terme à terme et chaque conoyau $P^i_{n + 1}/P^i_n$ soit un élément
de $\mathcal{P}$.
\end{enumerate}
\end{lemma}

\begin{proof}
Nous allons utiliser le fait que la catégorie homotopique $K(\mathcal{A})$
est une catégorie triangulée; voir Proposition
\ref{proposition-homotopy-category-triangulated}.
D'après le Lemme \ref{lemma-subcategory-left-resolution}, nous pouvons trouver
un morphisme de complexes $P_1^\bullet \to \tau_{\leq 1}K^\bullet$,
surjectif terme à terme et qui soit un quasi-isomorphisme, tel que les termes
de $P_1^\bullet$ appartiennent à $\mathcal{P}$.
Par récurrence, il suffit, étant donnés
$P_1^\bullet, \ldots, P_n^\bullet$, de construire $P_{n + 1}^\bullet$ et les
morphismes $P_n^\bullet \to P_{n + 1}^\bullet$ et
$P_{n + 1}^\bullet \to \tau_{\leq n + 1}K^\bullet$.

\medskip\noindent
Choisissons un triangle distingué
$P_n^\bullet \to \tau_{\leq n + 1}K^\bullet \to C^\bullet \to P_n^\bullet[1]$
dans $K(\mathcal{A})$. En appliquant le
Lemme \ref{lemma-subcategory-left-resolution}, choisissons un morphisme de
complexes $Q^\bullet \to C^\bullet$ qui soit un quasi-isomorphisme et tel que
les termes de $Q^\bullet$ appartiennent à $\mathcal{P}$. D'après les axiomes
des catégories triangulées, nous pouvons insérer la composée
$Q^\bullet \to C^\bullet \to P_n^\bullet[1]$ dans un triangle distingué
$P_n^\bullet \to P_{n + 1}^\bullet \to Q^\bullet \to P_n^\bullet[1]$
de $K(\mathcal{A})$.
D'après le Lemme \ref{lemma-improve-distinguished-triangle-homotopy}, nous
pouvons supposer, et nous supposerons, que
$0 \to P_n^\bullet \to P_{n + 1}^\bullet \to Q^\bullet \to 0$
est une suite exacte courte scindée terme à terme. Il en résulte que les
termes de $P_{n + 1}^\bullet$ appartiennent à $\mathcal{P}$ et que
$P_n^\bullet \to P_{n + 1}^\bullet$ est une injection scindée terme à terme
dont les conoyaux appartiennent à $\mathcal{P}$.
Les axiomes des catégories triangulées fournissent un morphisme de triangles
distingués
$$
\xymatrix{
P_n^\bullet \ar[r] \ar[d] &
P_{n + 1}^\bullet \ar[r] \ar[d] &
Q^\bullet \ar[r] \ar[d] &
P_n^\bullet[1] \ar[d] \\
P_n^\bullet \ar[r] &
\tau_{\leq n + 1}K^\bullet \ar[r] &
C^\bullet \ar[r] &
P_n^\bullet[1]
}
$$
dans la catégorie triangulée $K(\mathcal{A})$. Choisissons un véritable
morphisme de complexes
$f : P_{n + 1}^\bullet \to \tau_{\leq n + 1}K^\bullet$.
Le carré de gauche du diagramme ci-dessus commute à homotopie près, mais,
comme $P_n^\bullet \to P_{n + 1}^\bullet$ est une injection scindée terme à
terme, nous pouvons relever l'homotopie et modifier notre choix de $f$ de
façon à rendre ce carré commutatif. Enfin, $f$ est un quasi-isomorphisme,
puisque $P_n^\bullet \to P_n^\bullet$ et $Q^\bullet \to C^\bullet$ le sont.

\medskip\noindent
À ce stade, nous avons toutes les propriétés voulues, sauf que nous ne savons
pas si le morphisme
$f : P_{n + 1}^\bullet \to \tau_{\leq n + 1}K^\bullet$ est surjectif terme à
terme. Comme nous disposons du diagramme commutatif de complexes
$$
\xymatrix{
P_n^\bullet \ar[d] \ar[r] & P_{n + 1}^\bullet \ar[d] \\
\tau_{\leq n}K^\bullet \ar[r] & \tau_{\leq n + 1}K^\bullet
}
$$
l'hypothèse de récurrence montre que $f$ est surjectif sur les termes de tous
les degrés, à l'exception possible de $n$ et $n + 1$. Choisissons un objet
$P \in \mathcal{P}$ et une surjection $q : P \to K^n$.
Considérons le morphisme
$$
g :
P^\bullet = (\ldots \to 0 \to P \xrightarrow{1} P \to 0 \to \ldots)
\longrightarrow
\tau_{\leq n + 1}K^\bullet
$$
où la première copie de $P$ est placée en degré $n$ et où les morphismes sont
donnés par $q$ en degré $n$ et par $d_K \circ q$ en degré $n + 1$.
Ce morphisme est surjectif en degré $n$, et son conoyau en degré $n + 1$ est
$H^{n + 1}(\tau_{\leq n + 1}K^\bullet)$; pour le voir, rappelons que
$\tau_{\leq n + 1}K^\bullet$ a $\Ker(d_K^{n + 1})$ en degré $n + 1$.
Cependant, puisque $f$ est un quasi-isomorphisme, nous savons que
$H^{n + 1}(f)$ est surjectif. Par conséquent, après avoir remplacé
$f : P_{n + 1}^\bullet \to \tau_{\leq n + 1}K^\bullet$ par
$f \oplus g : P_{n + 1}^\bullet \oplus P^\bullet \to \tau_{\leq n + 1}K^\bullet$,
nous avons terminé.
\end{proof}

\noindent
Dans certains cas, le lemme précédent permet de montrer qu'un foncteur dérivé
à gauche est défini partout.

\begin{proposition}
\label{proposition-left-derived-exists}
Soit $F : \mathcal{A} \to \mathcal{B}$ un foncteur exact à droite entre
catégories abéliennes. Soit $\mathcal{P} \subset \Ob(\mathcal{A})$ un
sous-ensemble. Supposons que
\begin{enumerate}
\item $\mathcal{P}$ contienne $0$, soit stable par sommes directes (finies),
et que tout objet de $\mathcal{A}$ soit un quotient d'un élément de
$\mathcal{P}$,
\item pour tout complexe acyclique borné à droite $P^\bullet$ de
$\mathcal{A}$ tel que $P^n \in \mathcal{P}$ pour tout $n$, le complexe
$F(P^\bullet)$ soit exact,
\item $\mathcal{A}$ et $\mathcal{B}$ possèdent les colimites des systèmes
indexés par $\mathbf{N}$,
\item les colimites indexées par $\mathbf{N}$ soient exactes dans
$\mathcal{A}$ et dans $\mathcal{B}$, et
\item $F$ commute aux colimites indexées par $\mathbf{N}$.
\end{enumerate}
Alors $LF$ est défini sur tout $D(\mathcal{A})$.
\end{proposition}

\begin{proof}
D'après (1) et le Lemme \ref{lemma-subcategory-left-resolution}, pour tout
complexe borné à droite $K^\bullet$, il existe un quasi-isomorphisme
$P^\bullet \to K^\bullet$, où $P^\bullet$ est borné à droite et
$P^n \in \mathcal{P}$ pour tout $n$. Supposons que
$s : P^\bullet \to (P')^\bullet$ soit un quasi-isomorphisme entre complexes
bornés à droite formés d'objets de $\mathcal{P}$. Alors
$F(P^\bullet) \to F((P')^\bullet)$ est un quasi-isomorphisme, car
$F(C(s)^\bullet)$ est acyclique par l'hypothèse (2). Cela montre déjà que
$LF$ est défini sur $D^{-}(\mathcal{A})$ et qu'un complexe borné à droite
formé d'objets de $\mathcal{P}$ calcule $LF$; voir
Lemme \ref{lemma-find-existence-computes}.

\medskip\noindent
Soit ensuite $K^\bullet$ un complexe arbitraire de $\mathcal{A}$.
Choisissons un diagramme
$$
\xymatrix{
P_1^\bullet \ar[d] \ar[r] & P_2^\bullet \ar[d] \ar[r] & \ldots \\
\tau_{\leq 1}K^\bullet \ar[r] & \tau_{\leq 2}K^\bullet \ar[r] & \ldots
}
$$
comme dans le Lemme \ref{lemma-special-direct-system}. Remarquons que le
morphisme $\colim P_n^\bullet \to K^\bullet$ est un quasi-isomorphisme,
car les colimites indexées par $\mathbf{N}$ sont exactes dans $\mathcal{A}$
et $H^i(P_n^\bullet) = H^i(K^\bullet)$ pour $n > i$. Nous affirmons que
$$
F(\colim P_n^\bullet) = \colim F(P_n^\bullet)
$$
(colimites terme à terme) est $LF(K^\bullet)$, c'est-à-dire que
$\colim P_n^\bullet$ calcule $LF$. Pour le voir, d'après le
Lemme \ref{lemma-find-existence-computes}, il suffit de démontrer
l'affirmation suivante. Supposons que
$$
P^\bullet = \colim P_n^\bullet
\xrightarrow{\ \alpha\ }
Q^\bullet = \colim Q_n^\bullet
$$
soit un quasi-isomorphisme de complexes, tel que chacun des complexes
$P_n^\bullet$, $Q_n^\bullet$ soit borné à droite, ait ses termes dans
$\mathcal{P}$, et que les morphismes
$P_n^\bullet \to \tau_{\leq n}P^\bullet$ et
$Q_n^\bullet \to \tau_{\leq n}Q^\bullet$ soient des quasi-isomorphismes.
Affirmation : $F(\alpha)$ est un quasi-isomorphisme.

\medskip\noindent
La difficulté vient de ce que nous ne supposons pas que $\alpha$ soit donné
comme colimite de morphismes entre les complexes $P_n^\bullet$ et
$Q_n^\bullet$. Cependant, pour chaque $n$, nous savons que les flèches pleines
du diagramme
$$
\xymatrix{
& R^\bullet \ar@{..>}[d] \\
P_n^\bullet \ar[d] &
L^\bullet \ar@{..>}[l] \ar@{..>}[r] &
Q_n^\bullet \ar[d] \\
\tau_{\leq n}P^\bullet \ar[rr]^{\tau_{\leq n}\alpha} & &
\tau_{\leq n}Q^\bullet
}
$$
sont des quasi-isomorphismes. Comme les quasi-isomorphismes forment un
système multiplicatif dans $K(\mathcal{A})$ (voir
Lemme \ref{lemma-acyclic}), nous pouvons trouver un quasi-isomorphisme
$L^\bullet \to P_n^\bullet$ et un morphisme de complexes
$L^\bullet \to Q_n^\bullet$ tels que le diagramme ci-dessus commute à
homotopie près. Alors $\tau_{\leq n}L^\bullet \to L^\bullet$ est un
quasi-isomorphisme. Par conséquent (d'après la première partie de la preuve),
nous pouvons trouver un complexe borné à droite $R^\bullet$, dont les termes
appartiennent à $\mathcal{P}$, et un quasi-isomorphisme
$R^\bullet \to L^\bullet$ (comme indiqué dans le diagramme). Le résultat du
premier paragraphe de la preuve montre que
$F(R^\bullet) \to F(P_n^\bullet)$ et $F(R^\bullet) \to F(Q_n^\bullet)$ sont
des quasi-isomorphismes. Nous obtenons donc un isomorphisme
$H^i(F(P_n^\bullet)) \to H^i(F(Q_n^\bullet))$ s'insérant dans le diagramme
commutatif
$$
\xymatrix{
H^i(F(P_n^\bullet)) \ar[r] \ar[d] &
H^i(F(Q_n^\bullet)) \ar[d] \\
H^i(F(P^\bullet)) \ar[r] &
H^i(F(Q^\bullet))
}
$$
Le même raisonnement montre que ces morphismes sont aussi compatibles lorsque
$n$ varie. Puisque, d'après (4) et (5), nous avons
$$
H^i(F(P^\bullet)) =
H^i(F(\colim P_n^\bullet)) =
H^i(\colim F(P_n^\bullet)) = \colim H^i(F(P_n^\bullet))
$$
et de même pour $Q^\bullet$, nous concluons que
$H^i(F(\alpha)) : H^i(F(P^\bullet)) \to H^i(F(Q^\bullet))$ est un
isomorphisme, d'où l'affirmation.
\end{proof}

\begin{lemma}
\label{lemma-special-inverse-system}
Soit $\mathcal{A}$ une catégorie abélienne. Soit
$\mathcal{I} \subset \Ob(\mathcal{A})$ un sous-ensemble.
Supposons que $\mathcal{I}$ contienne $0$, soit stable par produits (finis),
et que tout objet de $\mathcal{A}$ soit un sous-objet d'un élément de
$\mathcal{I}$. Soit $K^\bullet$ un complexe.
Il existe un diagramme commutatif
$$
\xymatrix{
\ldots \ar[r] &
\tau_{\geq -2}K^\bullet \ar[r] \ar[d] &
\tau_{\geq -1}K^\bullet \ar[d] \\
\ldots \ar[r] & I_2^\bullet \ar[r] & I_1^\bullet
}
$$
dans la catégorie des complexes, tel que
\begin{enumerate}
\item les flèches verticales soient des quasi-isomorphismes injectifs terme à
terme,
\item $I_n^\bullet$ soit un complexe borné à gauche dont les termes
appartiennent à $\mathcal{I}$,
\item les flèches $I_{n + 1}^\bullet \to I_n^\bullet$ soient des surjections
scindées terme à terme et que $\Ker(I^i_{n + 1} \to I^i_n)$ soit un élément
de $\mathcal{I}$.
\end{enumerate}
\end{lemma}

\begin{proof}
Ce lemme est dual du
Lemme \ref{lemma-special-direct-system}.
\end{proof}

\section{Dérivation de foncteurs adjoints}
\label{section-deriving-adjoints}

\noindent
Soient $F : \mathcal{D} \to \mathcal{D}'$ et $G : \mathcal{D}' \to \mathcal{D}$
des foncteurs exacts entre catégories triangulées.
Soient $S$, respectivement $S'$, un système multiplicatif de
$\mathcal{D}$, respectivement de $\mathcal{D}'$, compatible avec la structure
triangulée.
Notons $Q : \mathcal{D} \to S^{-1}\mathcal{D}$ et
$Q' : \mathcal{D}' \to (S')^{-1}\mathcal{D}'$ les foncteurs de localisation.
Dans cette situation, par abus de notation, on note souvent $RF$ le foncteur
dérivé à droite partiellement défini associé à
$Q' \circ F : \mathcal{D} \to (S')^{-1}\mathcal{D}'$ et au système
multiplicatif $S$. De même, on note $LG$ le foncteur dérivé à gauche
partiellement défini associé à
$Q \circ G : \mathcal{D}' \to S^{-1}\mathcal{D}$ et au système multiplicatif
$S'$. On a les diagrammes
$$
\vcenter{
\xymatrix{
\mathcal{D} \ar[r]_F \ar[d]_Q & \mathcal{D}' \ar[d]^{Q'} \\
S^{-1}\mathcal{D} \ar@{..>}[r]^{RF} &
(S')^{-1}\mathcal{D}'
}
}
\quad\text{et}\quad
\vcenter{
\xymatrix{
\mathcal{D}' \ar[r]_G \ar[d]_{Q'} & \mathcal{D} \ar[d]^Q \\
(S')^{-1}\mathcal{D}' \ar@{..>}[r]^{LG} &
S^{-1}\mathcal{D}
}
}
$$

\begin{lemma}
\label{lemma-pre-derived-adjoint-functors-general}
Dans la situation ci-dessus, supposons que $F$ soit adjoint à droite de $G$.
Soient $K \in \Ob(\mathcal{D})$ et $M \in \Ob(\mathcal{D}')$.
Si $RF$ est défini en $K$ et $LG$ est défini en $M$, il existe un
isomorphisme canonique
$$
\Hom_{(S')^{-1}\mathcal{D}'}(M, RF(K)) =
\Hom_{S^{-1}\mathcal{D}}(LG(M), K)
$$
Cet isomorphisme est fonctoriel en les deux variables sur les sous-catégories
triangulées de $S^{-1}\mathcal{D}$ et $(S')^{-1}\mathcal{D}'$ où $RF$ et
$LG$ sont définis.
\end{lemma}

\begin{proof}
Puisque $RF$ est défini en $K$, la règle qui associe à un
$s : K \to I$ dans $S$ l'objet $F(I)$ est essentiellement constante comme
ind-objet de $(S')^{-1}\mathcal{D}'$, de valeur $RF(K)$.
De même, la règle qui associe à un $t : P \to M$ dans $S'$ l'objet $G(P)$ est
essentiellement constante comme pro-objet de $S^{-1}\mathcal{D}$, de valeur
$LG(M)$. Ainsi,
\begin{align*}
\Hom_{(S')^{-1}\mathcal{D}'}(M, RF(K))
& =
\colim_{s : K \to I} \Hom_{(S')^{-1}\mathcal{D}'}(M, F(I)) \\
& =
\colim_{s : K \to I} \colim_{t : P \to M} \Hom_{\mathcal{D}'}(P, F(I)) \\
& =
\colim_{t : P \to M} \colim_{s : K \to I} \Hom_{\mathcal{D}'}(P, F(I)) \\
& =
\colim_{t : P \to M} \colim_{s : K \to I} \Hom_{\mathcal{D}}(G(P), I) \\
& =
\colim_{t : P \to M} \Hom_{S^{-1}\mathcal{D}}(G(P), K) \\
& =
\Hom_{S^{-1}\mathcal{D}}(LG(M), K)
\end{align*}
La première égalité résulte de
Catégories, Lemme \ref{categories-lemma-characterize-essentially-constant-ind}.
La deuxième égalité résulte de la définition des morphismes dans
$(S')^{-1}\mathcal{D}'$; voir Catégories, Remarque
\ref{categories-remark-right-localization-morphisms-colimit}.
La troisième égalité résulte de
Catégories, Lemme \ref{categories-lemma-colimits-commute}.
La quatrième égalité résulte de l'adjonction de $F$ et $G$.
La cinquième égalité résulte de la définition des morphismes dans
$S^{-1}\mathcal{D}$; voir Catégories, Remarque
\ref{categories-remark-left-localization-morphisms-colimit}.
La sixième égalité résulte de
Catégories, Lemme \ref{categories-lemma-characterize-essentially-constant-pro}.
Nous omettons la preuve de la fonctorialité.
\end{proof}

\begin{lemma}
\label{lemma-pre-derived-adjoint-functors}
Soient $F : \mathcal{A} \to \mathcal{B}$ et $G : \mathcal{B} \to \mathcal{A}$
des foncteurs entre catégories abéliennes, tels que $F$ soit adjoint à droite
de $G$. Soient $K^\bullet$ un complexe de $\mathcal{A}$ et $M^\bullet$ un
complexe de $\mathcal{B}$. Si $RF$ est défini en $K^\bullet$ et $LG$ est
défini en $M^\bullet$, il existe un isomorphisme canonique
$$
\Hom_{D(\mathcal{B})}(M^\bullet, RF(K^\bullet)) =
\Hom_{D(\mathcal{A})}(LG(M^\bullet), K^\bullet)
$$
Cet isomorphisme est fonctoriel en les deux variables sur les sous-catégories
triangulées de $D(\mathcal{A})$ et $D(\mathcal{B})$ où $RF$ et $LG$ sont
définis.
\end{lemma}

\begin{proof}
C'est un cas particulier du très général
Lemme \ref{lemma-pre-derived-adjoint-functors-general}.
\end{proof}

\noindent
Le lemme suivant illustre pourquoi il est plus commode de travailler avec des
catégories dérivées non bornées. En effet, sans les foncteurs dérivés non
bornés, il serait même impossible d'énoncer le lemme.

\begin{lemma}
\label{lemma-derived-adjoint-functors}
Soient $F : \mathcal{A} \to \mathcal{B}$ et $G : \mathcal{B} \to \mathcal{A}$
des foncteurs entre catégories abéliennes, tels que $F$ soit adjoint à droite
de $G$. Si les foncteurs dérivés
$RF : D(\mathcal{A}) \to D(\mathcal{B})$ et
$LG : D(\mathcal{B}) \to D(\mathcal{A})$ existent, alors $RF$ est adjoint à
droite de $LG$.
\end{lemma}

\begin{proof}
Cela résulte immédiatement du
Lemme \ref{lemma-pre-derived-adjoint-functors}.
\end{proof}

\section{Complexes K-injectifs}
\label{section-K-injective}

\noindent
Les types de complexes suivants permettent de calculer les foncteurs dérivés
à droite sur la catégorie dérivée non bornée.

\begin{definition}
\label{definition-K-injective}
Soit $\mathcal{A}$ une catégorie abélienne. Un complexe $I^\bullet$ est
{\it K-injectif} si, pour tout complexe acyclique $M^\bullet$, nous avons
$\Hom_{K(\mathcal{A})}(M^\bullet, I^\bullet) = 0$.
\end{definition}

\noindent
Dans la situation de la définition, nous avons en fait
$\Hom_{K(\mathcal{A})}(M^\bullet[i], I^\bullet) = 0$ pour tout $i$, puisque le
translaté d'un complexe acyclique est acyclique.

\begin{lemma}
\label{lemma-K-injective}
Soit $\mathcal{A}$ une catégorie abélienne.
Soit $I^\bullet$ un complexe. Les conditions suivantes sont équivalentes :
\begin{enumerate}
\item $I^\bullet$ est K-injectif,
\item pour tout quasi-isomorphisme $M^\bullet \to N^\bullet$, l'application
$$
\Hom_{K(\mathcal{A})}(N^\bullet, I^\bullet)
\to \Hom_{K(\mathcal{A})}(M^\bullet, I^\bullet)
$$
est bijective, et
\item pour tout complexe $N^\bullet$, l'application
$$
\Hom_{K(\mathcal{A})}(N^\bullet, I^\bullet)
\to \Hom_{D(\mathcal{A})}(N^\bullet, I^\bullet)
$$
est un isomorphisme.
\end{enumerate}
\end{lemma}

\begin{proof}
Supposons (1). Alors (2) est vérifiée, car le foncteur
$\Hom_{K(\mathcal{A})}( - , I^\bullet)$ est cohomologique et le cône d'un
quasi-isomorphisme est acyclique.

\medskip\noindent
Supposons (2). Un morphisme $N^\bullet \to I^\bullet$ dans $D(\mathcal{A})$
est de la forme $fs^{-1} : N^\bullet \to I^\bullet$, où
$s : M^\bullet \to N^\bullet$ est un quasi-isomorphisme et
$f : M^\bullet \to I^\bullet$ un morphisme. D'après (2), cela correspond à un
unique morphisme $N^\bullet \to I^\bullet$ dans $K(\mathcal{A})$;
autrement dit, (3) est vérifiée.

\medskip\noindent
Supposons (3). Si $M^\bullet$ est acyclique, alors $M^\bullet$ est isomorphe
au complexe nul dans $D(\mathcal{A})$, donc
$\Hom_{D(\mathcal{A})}(M^\bullet, I^\bullet) = 0$, et par conséquent
$\Hom_{K(\mathcal{A})}(M^\bullet, I^\bullet) = 0$ d'après (3); autrement dit,
(1) est vérifiée.
\end{proof}

\begin{lemma}
\label{lemma-triangle-K-injective}
Soit $\mathcal{A}$ une catégorie abélienne. Soit $(K, L, M, f, g, h)$ un
triangle distingué de $K(\mathcal{A})$. Si deux des complexes $K$, $L$, $M$
sont K-injectifs, alors le troisième l'est aussi.
\end{lemma}

\begin{proof}
Cela résulte de la définition, du
Lemme \ref{lemma-representable-homological} et du fait que $K(\mathcal{A})$
est une catégorie triangulée
(Proposition \ref{proposition-homotopy-category-triangulated}).
\end{proof}

\begin{lemma}
\label{lemma-bounded-below-injectives-K-injective}
Soit $\mathcal{A}$ une catégorie abélienne. Tout complexe borné à gauche
d'objets injectifs est K-injectif.
\end{lemma}

\begin{proof}
Cela résulte des Lemmes \ref{lemma-K-injective} et
\ref{lemma-morphisms-into-injective-complex}.
\end{proof}

\begin{lemma}
\label{lemma-product-K-injective}
Soit $\mathcal{A}$ une catégorie abélienne. Soit $T$ un ensemble et, pour
chaque $t \in T$, soit $I_t^\bullet$ un complexe K-injectif. Si
$I^n = \prod_t I_t^n$ existe pour tout $n$, alors $I^\bullet$ est un complexe
K-injectif. De plus, $I^\bullet$ représente le produit des objets
$I_t^\bullet$ dans $D(\mathcal{A})$.
\end{lemma}

\begin{proof}
Soit $K^\bullet$ un complexe. Remarquons que le complexe
$$
C :
\prod\nolimits_b \Hom(K^{-b}, I^{b - 1}) \to
\prod\nolimits_b \Hom(K^{-b}, I^b) \to
\prod\nolimits_b \Hom(K^{-b}, I^{b + 1})
$$
a pour cohomologie $\Hom_{K(\mathcal{A})}(K^\bullet, I^\bullet)$ au terme
central. De même, le complexe
$$
C_t :
\prod\nolimits_b \Hom(K^{-b}, I_t^{b - 1}) \to
\prod\nolimits_b \Hom(K^{-b}, I_t^b) \to
\prod\nolimits_b \Hom(K^{-b}, I_t^{b + 1})
$$
calcule $\Hom_{K(\mathcal{A})}(K^\bullet, I_t^\bullet)$.
Remarquons ensuite que
$$
C = \prod\nolimits_{t \in T} C_t
$$
comme complexes de groupes abéliens, par notre choix de $I$.
La formation des produits est un foncteur exact sur la catégorie des groupes
abéliens. Ainsi, si $K^\bullet$ est acyclique, alors
$\Hom_{K(\mathcal{A})}(K^\bullet, I_t^\bullet) = 0$, donc $C_t$ est
acyclique, donc $C$ est acyclique, et par conséquent
$\Hom_{K(\mathcal{A})}(K^\bullet, I^\bullet) = 0$.
Nous trouvons ainsi que $I^\bullet$ est K-injectif.
Ceci étant établi, nous pouvons utiliser le Lemme \ref{lemma-K-injective}
pour conclure que
$$
\Hom_{D(\mathcal{A})}(K^\bullet, I^\bullet)
=
\prod\nolimits_{t \in T} \Hom_{D(\mathcal{A})}(K^\bullet, I_t^\bullet)
$$
et $I^\bullet$ représente bien le produit dans la catégorie dérivée.
\end{proof}

\begin{lemma}
\label{lemma-K-injective-defined}
Soit $\mathcal{A}$ une catégorie abélienne.
Soit $F : K(\mathcal{A}) \to \mathcal{D}'$ un foncteur exact entre catégories
triangulées. Alors $RF$ est défini en tout complexe de $K(\mathcal{A})$ qui
est quasi-isomorphe à un complexe K-injectif. En fait, tout complexe
K-injectif calcule $RF$.
\end{lemma}

\begin{proof}
D'après le Lemme \ref{lemma-derived-inverts}, il suffit de montrer que $RF$
est défini en un complexe K-injectif, c'est-à-dire qu'il suffit de montrer
qu'un complexe K-injectif $I^\bullet$ calcule $RF$.
Considérons un quasi-isomorphisme $I^\bullet \to N^\bullet$. D'après le
Lemme \ref{lemma-K-injective}, il possède un inverse à gauche.
Ainsi, $I^\bullet \to I^\bullet$ est un objet final de
$I^\bullet/\text{Qis}(\mathcal{A})$, ce qui conclut la preuve.
\end{proof}

\begin{lemma}
\label{lemma-enough-K-injectives-implies}
Soit $\mathcal{A}$ une catégorie abélienne.
Supposons que tout complexe possède un quasi-isomorphisme vers un complexe
K-injectif. Alors tout foncteur exact
$F : K(\mathcal{A}) \to \mathcal{D}'$ entre catégories triangulées possède un
foncteur dérivé à droite
$$
RF : D(\mathcal{A}) \longrightarrow \mathcal{D}'
$$
et $RF(I^\bullet) = F(I^\bullet)$ pour les complexes K-injectifs
$I^\bullet$.
\end{lemma}

\begin{proof}
Pour le voir, appliquons le Lemme \ref{lemma-find-existence-computes} à la
collection $\mathcal{I}$ des complexes K-injectifs. Comme (1) est vérifiée par
hypothèse, il suffit de prouver que, si $I^\bullet \to J^\bullet$ est un
quasi-isomorphisme de complexes K-injectifs, alors
$F(I^\bullet) \to F(J^\bullet)$ est un isomorphisme. Cela est clair, car
$I^\bullet \to J^\bullet$ est une équivalence d'homotopie, c'est-à-dire un
isomorphisme dans $K(\mathcal{A})$, d'après le
Lemme \ref{lemma-K-injective}.
\end{proof}

\noindent
Le lemme suivant se généralise aux limites indexées par des ordinaux plus
grands.

\begin{lemma}
\label{lemma-limit-K-injectives}
\begin{slogan}
La limite d'une tour « scindée » de complexes K-injectifs est K-injective.
\end{slogan}
Soit $\mathcal{A}$ une catégorie abélienne. Soit
$$
\ldots \to I_3^\bullet \to I_2^\bullet \to I_1^\bullet
$$
un système projectif de complexes. Supposons que
\begin{enumerate}
\item chaque $I_n^\bullet$ soit $K$-injectif,
\item chaque morphisme $I_{n + 1}^m \to I_n^m$ soit une surjection scindée,
\item les limites $I^m = \lim I_n^m$ existent.
\end{enumerate}
Alors le complexe $I^\bullet$ est K-injectif.
\end{lemma}

\begin{proof}
Nous invitons vivement le lecteur à sauter la preuve de ce lemme.
Soit $M^\bullet$ un complexe acyclique. Abrégeons
$H_n(a, b) = \Hom_\mathcal{A}(M^a, I_n^b)$. Avec cette notation,
$\Hom_{K(\mathcal{A})}(M^\bullet, I^\bullet)$ est la cohomologie du complexe
$$
\prod_m \lim\limits_n H_n(m, m - 2)
\to
\prod_m \lim\limits_n H_n(m, m - 1)
\to
\prod_m \lim\limits_n H_n(m, m)
\to
\prod_m \lim\limits_n H_n(m, m + 1)
$$
au troisième terme à partir de la gauche.
Nous pouvons permuter $\prod$ et $\lim$, et chacun des complexes
$$
\prod_m H_n(m, m - 2)
\to
\prod_m H_n(m, m - 1)
\to
\prod_m H_n(m, m)
\to
\prod_m H_n(m, m + 1)
$$
est exact d'après l'hypothèse (1). D'après l'hypothèse (2), les morphismes
des systèmes
$$
\ldots \to
\prod_m H_3(m, m - 2) \to
\prod_m H_2(m, m - 2) \to
\prod_m H_1(m, m - 2)
$$
sont surjectifs. Le lemme résulte donc de
Homologie, Lemme \ref{homology-lemma-apply-Mittag-Leffler}.
\end{proof}

\noindent
Il semble qu'une combinaison des Lemmes \ref{lemma-special-inverse-system},
\ref{lemma-bounded-below-injectives-K-injective} et
\ref{lemma-limit-K-injectives} produise « assez de complexes K-injectifs »
pour toute catégorie abélienne possédant assez d'injectifs et des produits
dénombrables. En réalité, cela peut ne pas fonctionner! Voir le
Lemme \ref{lemma-difficulty-K-injectives} pour une explication.

\begin{lemma}
\label{lemma-adjoint-preserve-K-injectives}
Soient $\mathcal{A}$ et $\mathcal{B}$ des catégories abéliennes.
Soient $u : \mathcal{A} \to \mathcal{B}$ et
$v : \mathcal{B} \to \mathcal{A}$ des foncteurs additifs. Supposons que
\begin{enumerate}
\item $u$ soit adjoint à droite de $v$, et
\item $v$ soit exact.
\end{enumerate}
Alors $u$ transforme les complexes K-injectifs en complexes K-injectifs.
\end{lemma}

\begin{proof}
Soit $I^\bullet$ un complexe K-injectif de $\mathcal{A}$.
Soit $M^\bullet$ un complexe acyclique de $\mathcal{B}$.
Comme $v$ est exact, $v(M^\bullet)$ est un complexe acyclique.
Par adjonction, nous obtenons
$$
0 = \Hom_{K(\mathcal{A})}(v(M^\bullet), I^\bullet) =
\Hom_{K(\mathcal{B})}(M^\bullet, u(I^\bullet))
$$
d'où le lemme.
\end{proof}

\section{Dimension cohomologique bornée}
\label{section-bounded}

\noindent
Il existe un autre cas où le foncteur dérivé non borné existe : celui où le
foncteur est de dimension cohomologique bornée.

\begin{lemma}
\label{lemma-replace-resolution}
Soit $\mathcal{A}$ une catégorie abélienne. Soit
$d : \Ob(\mathcal{A}) \to \{0, 1, 2, \ldots, \infty\}$ une fonction.
Supposons que
\begin{enumerate}
\item tout objet de $\mathcal{A}$ soit un sous-objet d'un objet $A$ tel que
$d(A) = 0$,
\item $d(A \oplus B) \leq \max \{d(A), d(B)\}$ pour
$A, B \in \mathcal{A}$, et
\item si $0 \to A \to B \to C \to 0$ est exacte courte, alors
$d(C) \leq \max\{d(A) - 1, d(B)\}$.
\end{enumerate}
Soit $K^\bullet$ un complexe tel que $n + d(K^n)$ tende vers $-\infty$
lorsque $n \to -\infty$. Il existe alors un quasi-isomorphisme
$K^\bullet \to L^\bullet$ tel que $d(L^n) = 0$ pour tout
$n \in \mathbf{Z}$.
\end{lemma}

\begin{proof}
D'après le Lemme \ref{lemma-subcategory-right-resolution}, nous pouvons
trouver un quasi-isomorphisme
$\sigma_{\geq 0}K^\bullet \to M^\bullet$ tel que $M^n = 0$ pour $n < 0$ et
$d(M^n) = 0$ pour $n \geq 0$. Alors $K^\bullet$ est quasi-isomorphe au
complexe
$$
\ldots \to K^{-2} \to K^{-1} \to M^0 \to M^1 \to \ldots
$$
Nous pouvons donc supposer que $d(K^n) = 0$ pour $n \gg 0$.
Remarquons que ce remplacement ne détruit pas la condition
$n + d(K^n) \to -\infty$ lorsque $n \to -\infty$.

\medskip\noindent
Nous allons améliorer $K^\bullet$ par une suite (infinie) de remplacements
élémentaires. Un {\it remplacement élémentaire} est l'opération suivante.
Choisissons un indice $n$ tel que $d(K^n) > 0$. Choisissons une injection
$K^n \to M$ avec $d(M) = 0$. Posons
$M' = \Coker(K^n \to M \oplus K^{n + 1})$. Considérons le morphisme de
complexes
$$
\xymatrix{
K^\bullet : \ar[d] &
K^{n - 1} \ar[d] \ar[r] &
K^n \ar[d] \ar[r] &
K^{n + 1} \ar[d] \ar[r] &
K^{n + 2} \ar[d] \\
(K')^\bullet : &
K^{n - 1} \ar[r] &
M \ar[r] &
M' \ar[r] &
K^{n + 2}
}
$$
Il est clair que $K^\bullet \to (K')^\bullet$ est un quasi-isomorphisme.
De plus, il est clair que $d((K')^n) = 0$ et
$$
d((K')^{n + 1}) \leq \max\{d(K^n) - 1, d(M \oplus K^{n + 1})\} \leq
\max\{d(K^n) - 1, d(K^{n + 1})\}
$$
et que les autres valeurs ne changent pas.

\medskip\noindent
Pour achever la preuve, choisissons soigneusement l'ordre des remplacements
élémentaires, de façon que, pour tout entier $m$, le complexe
$\sigma_{\geq m}K^\bullet$ ne soit modifié qu'un nombre fini de fois.
Pour cela, posons
$$
\xi(K^\bullet) = \max \{n + d(K^n) \mid d(K^n) > 0\}
$$
et
$$
I = \{n \in \mathbf{Z} \mid \xi(K^\bullet) = n + d(K^n)
\text{ et }
 d(K^n) > 0\}
$$
Notre hypothèse selon laquelle $n + d(K^n)$ tend vers $-\infty$ lorsque
$n \to -\infty$, jointe au fait que $d(K^n) = 0$ pour $n >> 0$, entraîne
$\xi(K^\bullet) < +\infty$ et la finitude de $I$.
Il est clair que $\xi((K')^\bullet) \leq \xi(K^\bullet)$ pour un remplacement
élémentaire comme ci-dessus. Un remplacement élémentaire modifie le complexe
en degrés $\leq \xi(K^\bullet) + 1$. Il suffit donc de trouver une suite finie
de remplacements élémentaires qui diminue $\xi(K^\bullet)$.
Or, si nous effectuons un remplacement élémentaire en partant du plus petit
élément $n \in I$, soit nous diminuons le cardinal de $I$, soit nous
augmentons $\min I$. Comme tout élément de $I$ est
$\leq \xi(K^\bullet)$, le procédé aboutit après un nombre fini d'étapes.
\end{proof}

\begin{lemma}
\label{lemma-unbounded-right-derived}
Soit $F : \mathcal{A} \to \mathcal{B}$ un foncteur exact à gauche entre
catégories abéliennes. Supposons que
\begin{enumerate}
\item tout objet de $\mathcal{A}$ soit un sous-objet d'un objet acyclique à
droite pour $F$,
\item il existe un entier $n \geq 0$ tel que $R^nF = 0$.
\end{enumerate}
Alors
\begin{enumerate}
\item $RF : D(\mathcal{A}) \to D(\mathcal{B})$ existe,
\item tout complexe formé d'objets acycliques à droite pour $F$ calcule $RF$,
\item tout complexe est la source d'un quasi-isomorphisme vers un complexe
formé d'objets acycliques à droite pour $F$,
\item pour $E \in D(\mathcal{A})$,
\begin{enumerate}
\item $H^i(RF(\tau_{\leq a}E)) \to H^i(RF(E))$ est un isomorphisme
pour $i \leq a$,
\item $H^i(RF(E)) \to H^i(RF(\tau_{\geq b - n + 1}E))$ est un isomorphisme
pour $i \geq b$,
\item si $H^i(E) = 0$ pour $i \not \in [a, b]$, où
$-\infty \leq a \leq b \leq \infty$, alors $H^i(RF(E)) = 0$
pour $i \not \in [a, b + n - 1]$.
\end{enumerate}
\end{enumerate}
\end{lemma}

\begin{proof}
Remarquons que la première hypothèse entraîne l'existence de
$RF : D^+(\mathcal{A}) \to D^+(\mathcal{B})$; voir
Proposition \ref{proposition-enough-acyclics}.
Soit $A$ un objet de $\mathcal{A}$. Choisissons une injection $A \to A'$
avec $A'$ acyclique. La suite exacte longue de cohomologie montre alors que
$R^{n + 1}F(A) = R^nF(A'/A) = 0$. Nous concluons donc que
$R^{n + 1}F = 0$. En poursuivant de la même manière par récurrence, nous
trouvons que $R^mF = 0$ pour tout $m \geq n$.

\medskip\noindent
Nous allons utiliser le Lemme \ref{lemma-replace-resolution} avec la fonction
$d : \Ob(\mathcal{A}) \to \{0, 1, 2, \ldots \}$ donnée par
$d(A) = \max\bigl(\{0\} \cup \{i \mid R^iF(A) \not = 0\}\bigr)$.
La première hypothèse du Lemme \ref{lemma-replace-resolution} est notre
hypothèse (1). La deuxième hypothèse du
Lemme \ref{lemma-replace-resolution} résulte de l'égalité
$RF(A \oplus B) = RF(A) \oplus RF(B)$. La troisième hypothèse du
Lemme \ref{lemma-replace-resolution} résulte de la suite exacte longue de
cohomologie. Par conséquent, pour tout complexe $K^\bullet$, il existe un
quasi-isomorphisme $K^\bullet \to L^\bullet$ vers un complexe d'objets
acycliques à droite pour $F$. Cela prouve l'assertion (3).

\medskip\noindent
Nous affirmons que, si $L^\bullet \to M^\bullet$ est un quasi-isomorphisme de
complexes d'objets acycliques à droite pour $F$, alors
$F(L^\bullet) \to F(M^\bullet)$ est un quasi-isomorphisme. Une fois cette
affirmation établie, les assertions (1) et (2) du lemme résultent du
Lemme \ref{lemma-find-existence-computes}.
Pour prouver l'affirmation, choisissons un entier $i \in \mathbf{Z}$.
Considérons le triangle distingué
$$
\sigma_{\geq i - n - 1}L^\bullet \to
\sigma_{\geq i - n - 1}M^\bullet \to Q^\bullet,
$$
c'est-à-dire que $Q^\bullet$ est le cône du premier morphisme.
Remarquons que $Q^\bullet$ est borné à gauche et que $H^j(Q^\bullet)$ est nul,
sauf peut-être pour $j = i - n - 1$ ou $j = i - n - 2$.
Nous pouvons appliquer $RF$ à $Q^\bullet$.
En utilisant la deuxième suite spectrale du
Lemme \ref{lemma-two-ss-complex-functor} et l'annulation de la cohomologie
postulée en (2), nous concluons que $H^j(RF(Q^\bullet))$ est nul, sauf
peut-être pour $j \in \{i - n - 2, \ldots, i - 1\}$. Ainsi,
$RF(\sigma_{\geq i - n - 1}L^\bullet) \to RF(\sigma_{\geq i - n - 1}M^\bullet)$
induit un isomorphisme sur les objets de cohomologie en degrés $\geq i$.
D'après la Proposition \ref{proposition-enough-acyclics}, nous savons que
$RF(\sigma_{\geq i - n - 1}L^\bullet) = \sigma_{\geq i - n - 1}F(L^\bullet)$
et
$RF(\sigma_{\geq i - n - 1}M^\bullet) = \sigma_{\geq i - n - 1}F(M^\bullet)$.
Nous concluons que $F(L^\bullet) \to F(M^\bullet)$ induit un isomorphisme en
degré $i$, comme voulu.

\medskip\noindent
L'assertion (4)(a) résulte du Lemme \ref{lemma-negative-vanishing}.

\medskip\noindent
Pour l'assertion (4)(b), représentons $E$ par le complexe $L^\bullet$
d'objets acycliques à droite pour $F$. D'après (2), $RF(E)$ est représenté
par le complexe $F(L^\bullet)$, et $RF(\sigma_{\geq c}L^\bullet)$ est
représenté par $\sigma_{\geq c}F(L^\bullet)$. Considérons le triangle
distingué
$$
H^{b - n}(L^\bullet)[n - b] \to
\tau_{\geq b - n}L^\bullet \to
\tau_{\geq b - n + 1}L^\bullet
$$
de la Remarque \ref{remark-truncation-distinguished-triangle}.
L'annulation établie ci-dessus montre que
$H^i(RF(\tau_{\geq b - n}L^\bullet))$ coïncide avec
$H^i(RF(\tau_{\geq b - n + 1}L^\bullet))$ pour $i \geq b$.
Considérons la suite exacte courte de complexes
$$
0 \to
\Im(L^{b - n - 1} \to L^{b - n})[n - b] \to
\sigma_{\geq b - n}L^\bullet \to
\tau_{\geq b - n}L^\bullet \to 0
$$
En utilisant le triangle distingué associé (voir
section \ref{section-canonical-delta-functor}) et l'annulation précédente,
nous concluons que $H^i(RF(\tau_{\geq b - n}L^\bullet))$ coïncide avec
$H^i(RF(\sigma_{\geq b - n}L^\bullet))$ pour $i \geq b$.
Comme le morphisme
$RF(\sigma_{\geq b - n}L^\bullet) \to RF(L^\bullet)$ est représenté par
$\sigma_{\geq b - n}F(L^\bullet) \to F(L^\bullet)$, nous concluons que ce
dernier coïncide à son tour avec $H^i(RF(L^\bullet))$ pour $i \geq b$, comme
voulu.

\medskip\noindent
Démonstration de (4)(c). Sous l'hypothèse faite sur $E$, nous avons
$\tau_{\leq a - 1}E = 0$, et l'assertion (4)(a) donne l'annulation de
$H^i(RF(E))$ pour $i \leq a - 1$. De même, nous avons
$\tau_{\geq b + 1}E = 0$; l'assertion (4)(b) donne donc l'annulation de
$H^i(RF(E))$ pour $i \geq b + n$.
\end{proof}

\begin{lemma}
\label{lemma-unbounded-left-derived}
Soit $F : \mathcal{A} \to \mathcal{B}$ un foncteur exact à droite entre
catégories abéliennes. Si
\begin{enumerate}
\item tout objet de $\mathcal{A}$ est un quotient d'un objet acyclique à
gauche pour $F$,
\item il existe un entier $n \geq 0$ tel que $L^nF = 0$,
\end{enumerate}
alors
\begin{enumerate}
\item $LF : D(\mathcal{A}) \to D(\mathcal{B})$ existe,
\item tout complexe formé d'objets acycliques à gauche pour $F$ calcule $LF$,
\item tout complexe est la cible d'un quasi-isomorphisme provenant d'un
complexe formé d'objets acycliques à gauche pour $F$,
\item pour $E \in D(\mathcal{A})$,
\begin{enumerate}
\item $H^i(LF(\tau_{\leq a + n - 1}E)) \to H^i(LF(E))$ est un isomorphisme
pour $i \leq a$,
\item $H^i(LF(E)) \to H^i(LF(\tau_{\geq b}E))$ est un isomorphisme
pour $i \geq b$,
\item si $H^i(E) = 0$ pour $i \not \in [a, b]$, où
$-\infty \leq a \leq b \leq \infty$, alors $H^i(LF(E)) = 0$
pour $i \not \in [a - n + 1, b]$.
\end{enumerate}
\end{enumerate}
\end{lemma}

\begin{proof}
C'est le dual du Lemme \ref{lemma-unbounded-right-derived}.
\end{proof}

\section{Colimites dérivées}
\label{section-derived-colimit}

\noindent
Dans une catégorie triangulée, on dispose d'une notion de colimite dérivée.

\begin{definition}
\label{definition-derived-colimit}
Soit $\mathcal{D}$ une catégorie triangulée.
Soit $(K_n, f_n)$ un système d'objets de $\mathcal{D}$.
On dit qu'un objet $K$ est une {\it colimite dérivée}, ou une
{\it colimite homotopique}, du système $(K_n)$ si la somme directe
$\bigoplus K_n$ existe et s'il existe un triangle distingué
$$
\bigoplus K_n \to \bigoplus K_n \to K \to \bigoplus K_n[1]
$$
où le morphisme $\bigoplus K_n \to \bigoplus K_n$ est donné par
$1 - f_n$ en degré $n$. Dans ce cas, nous l'indiquons parfois par la notation
$K = \text{hocolim} K_n$.
\end{definition}

\noindent
D'après TR3, une colimite dérivée, si elle existe, est unique à isomorphisme
(non unique) près. De plus, d'après TR1, une colimite dérivée des $K_n$ existe
dès que $\bigoplus K_n$ existe. La catégorie dérivée $D(\textit{Ab})$ de la
catégorie des groupes abéliens est un exemple de catégorie triangulée dans
laquelle toutes les colimites homotopiques existent.

\medskip\noindent
Cette absence d'unicité rend difficile la détermination précise de la
colimite dérivée. Dans Algèbre complémentaire, Lemme
\ref{more-algebra-lemma-map-from-hocolim}, le lecteur trouvera une suite exacte
$$
0 \to R^1\lim \Hom(K_n, L[-1]) \to \Hom(\text{hocolim} K_n, L)
\to \lim \Hom(K_n, L) \to 0
$$
qui décrit les $\Hom$ issus d'une colimite homotopique à l'aide des $\Hom$
usuels.

\begin{remark}
\label{remark-uniqueness-derived-colimit}
Soit $\mathcal{D}$ une catégorie triangulée.
Soit $(K_n, f_n)$ un système d'objets de $\mathcal{D}$.
Nous pouvons considérer une colimite dérivée comme un objet $K$ de
$\mathcal{D}$ muni de morphismes $i_n : K_n \to K$ tels que
$i_{n + 1} \circ f_n = i_n$, et tel qu'il existe un morphisme
$c : K \to \bigoplus K_n[1]$ pour lequel
$$
\bigoplus K_n \xrightarrow{1 - f_n} \bigoplus K_n \xrightarrow{i_n}
K \xrightarrow{c} \bigoplus K_n[1]
$$
est un triangle distingué. Si $(K', i'_n, c')$ est une deuxième colimite
dérivée, il existe un isomorphisme $\varphi : K \to K'$ tel que
$\varphi \circ i_n = i'_n$ et $c' \circ \varphi = c$.
L'existence de $\varphi$ est donnée par TR3 et le fait que $\varphi$ soit un
isomorphisme résulte du Lemme \ref{lemma-third-isomorphism-triangle}.
\end{remark}

\begin{remark}
\label{remark-functoriality-derived-colimit}
Soit $\mathcal{D}$ une catégorie triangulée.
Soit $(a_n) : (K_n, f_n) \to (L_n, g_n)$ un morphisme de systèmes d'objets de
$\mathcal{D}$. Soit $(K, i_n, c)$ une colimite dérivée du premier système et
soit $(L, j_n, d)$ une colimite dérivée du second système, avec les notations
de la Remarque \ref{remark-uniqueness-derived-colimit}.
Il existe alors un morphisme $a : K \to L$ tel que
$a \circ i_n = j_n \circ a_n$ et
$d \circ a = (a_n[1]) \circ c$.
Cela résulte de TR3 appliqué aux triangles distingués qui les définissent.
\end{remark}

\begin{lemma}
\label{lemma-hocolim-subsequence}
Soit $\mathcal{D}$ une catégorie triangulée.
Soit $(K_n, f_n)$ un système d'objets de $\mathcal{D}$.
Soit $n_1 < n_2 < n_3 < \ldots$ une suite d'entiers.
Supposons que $\bigoplus K_n$ et $\bigoplus K_{n_i}$ existent.
Il existe alors un isomorphisme
$\text{hocolim} K_{n_i} \to \text{hocolim} K_n$ tel que
$$
\xymatrix{
K_{n_i} \ar[r] \ar[d]_{\text{id}} & \text{hocolim} K_{n_i} \ar[d] \\
K_{n_i} \ar[r] & \text{hocolim} K_n
}
$$
commute pour tout $i$.
\end{lemma}

\begin{proof}
Soit $g_i : K_{n_i} \to K_{n_{i + 1}}$ la composée
$f_{n_{i + 1} - 1} \circ \ldots \circ f_{n_i}$.
Construisons les diagrammes commutatifs
$$
\vcenter{
\xymatrix{
\bigoplus\nolimits_i K_{n_i} \ar[r]_{1 - g_i} \ar[d]_b &
\bigoplus\nolimits_i K_{n_i} \ar[d]^a \\
\bigoplus\nolimits_n K_n \ar[r]^{1 - f_n} &
\bigoplus\nolimits_n K_n
}
}
\quad\text{et}\quad
\vcenter{
\xymatrix{
\bigoplus\nolimits_n K_n \ar[r]_{1 - f_n} \ar[d]_d &
\bigoplus\nolimits_n K_n \ar[d]^c \\
\bigoplus\nolimits_i K_{n_i} \ar[r]^{1 - g_i} &
\bigoplus\nolimits_i K_{n_i}
}
}
$$
comme suit. Soit $a_i = a|_{K_{n_i}}$ l'inclusion de $K_{n_i}$ dans la
somme directe. Autrement dit, $a$ est l'inclusion naturelle.
Soit $b_i = b|_{K_{n_i}}$ le morphisme
$$
K_{n_i}
\xrightarrow{1,\ f_{n_i},\ f_{n_i + 1} \circ f_{n_i},
\ \ldots,\ f_{n_{i + 1} - 2} \circ \ldots \circ f_{n_i}}
K_{n_i} \oplus K_{n_i + 1} \oplus \ldots \oplus K_{n_{i + 1} - 1}
$$
Si $n_{i - 1} < j \leq n_i$, nous définissons
$c_j = c|_{K_j}$ comme le morphisme
$$
K_j \xrightarrow{f_{n_i - 1} \circ \ldots \circ f_j} K_{n_i}
$$
Nous prenons $d_j = d|_{K_j}$ nul si $j \not = n_i$ pour tout $i$, et nous
prenons pour $d_{n_i}$ l'inclusion naturelle de $K_{n_i}$ dans la somme
directe. Autrement dit, $d$ est la projection naturelle.
D'après TR3, ces diagrammes définissent des morphismes
$$
\varphi : \text{hocolim} K_{n_i} \to \text{hocolim} K_n
\quad\text{et}\quad
\psi : \text{hocolim} K_n \to \text{hocolim} K_{n_i}
$$
Comme $c \circ a$ et $d \circ b$ sont les morphismes identiques, le
Lemme \ref{lemma-third-isomorphism-triangle} montre que
$\varphi \circ \psi$ est un isomorphisme.
Dans l'autre sens, nous obtenons les morphismes $a \circ c$ et $b \circ d$.
Considérons le morphisme
$h = (h_j) : \bigoplus K_n \to \bigoplus K_n$ défini par la règle suivante :
pour $n_{i - 1} < j < n_i$, nous posons
$$
h_j : K_j
\xrightarrow{1,\ f_j,\ f_{j + 1} \circ f_j,
\ \ldots,\ f_{n_i - 1} \circ \ldots \circ f_j}
K_j \oplus \ldots \oplus K_{n_i}
$$
Le lecteur vérifie alors que
$(1 - f) \circ h = \text{id} - a \circ c$ et
$h \circ (1 - f) = \text{id} - b \circ d$. Cela signifie que
$\text{id} - \psi \circ \varphi$ est de carré nul d'après le
Lemme \ref{lemma-third-map-square-zero} (petit argument omis).
Autrement dit, $\psi \circ \varphi$ diffère de l'identité par un
endomorphisme nilpotent; c'est donc un isomorphisme. Ainsi, $\varphi$ et
$\psi$ sont des isomorphismes, comme voulu.
\end{proof}

\begin{lemma}
\label{lemma-direct-sums}
Soit $\mathcal{A}$ une catégorie abélienne.
Si $\mathcal{A}$ possède des sommes directes dénombrables exactes, alors
$D(\mathcal{A})$ possède des sommes directes dénombrables. Plus précisément,
étant donnée une famille de complexes $K_i^\bullet$ indexée par un ensemble
dénombrable $I$, la somme directe terme à terme $\bigoplus K_i^\bullet$ est la
somme directe des $K_i^\bullet$ dans $D(\mathcal{A})$.
\end{lemma}

\begin{proof}
Soit $L^\bullet$ un complexe. Supposons donnés des morphismes
$\alpha_i : K_i^\bullet \to L^\bullet$ dans $D(\mathcal{A})$.
Cela signifie qu'il existe des quasi-isomorphismes
$s_i : M_i^\bullet \to K_i^\bullet$ de complexes et des morphismes de
complexes $f_i : M_i^\bullet \to L^\bullet$ tels que
$\alpha_i = f_is_i^{-1}$. Par hypothèse, le morphisme de complexes
$$
s : \bigoplus M_i^\bullet \longrightarrow \bigoplus K_i^\bullet
$$
est un quasi-isomorphisme. En posant $f = \bigoplus f_i$, nous voyons donc que
$\alpha = fs^{-1}$ est un morphisme de $D(\mathcal{A})$ dont la composée avec
la coprojection $K_i^\bullet \to \bigoplus K_i^\bullet$ est $\alpha_i$.
Nous omettons la vérification de l'unicité de $\alpha$.
\end{proof}

\begin{lemma}
\label{lemma-compute-colimit}
Soit $\mathcal{A}$ une catégorie abélienne. Supposons que les colimites
indexées par $\mathbf{N}$ existent et soient exactes. Alors les sommes directes
dénombrables existent et sont exactes. De plus, si $(A_n, f_n)$ est un système
indexé par $\mathbf{N}$, il existe une suite exacte courte
$$
0 \to \bigoplus A_n \to \bigoplus A_n \to \colim A_n \to 0
$$
où le premier morphisme est donné en degré $n$ par $1 - f_n$.
\end{lemma}

\begin{proof}
La première assertion résulte de
$\bigoplus A_n = \colim (A_1 \oplus \ldots \oplus A_n)$.
Pour la seconde, remarquons que, pour chaque $n$, nous avons la suite exacte
courte
$$
0 \to
A_1 \oplus \ldots \oplus A_{n - 1} \to
A_1 \oplus \ldots \oplus A_n \to A_n \to 0
$$
où le premier morphisme est donné par les morphismes $1 - f_i$, et le second
est la somme des morphismes de transition.
Prenons la colimite pour obtenir la suite du lemme.
\end{proof}

\begin{lemma}
\label{lemma-colim-hocolim}
Soit $\mathcal{A}$ une catégorie abélienne. Soit $L_n^\bullet$ un système de
complexes de $\mathcal{A}$. Supposons que les colimites indexées par
$\mathbf{N}$ existent et soient exactes dans $\mathcal{A}$.
Alors la colimite terme à terme $L^\bullet = \colim L_n^\bullet$ est une
colimite homotopique du système dans $D(\mathcal{A})$.
\end{lemma}

\begin{proof}
Nous avons une suite exacte de complexes
$$
0 \to \bigoplus L_n^\bullet \to \bigoplus L_n^\bullet \to L^\bullet \to 0
$$
d'après le Lemme \ref{lemma-compute-colimit}.
Les sommes directes sont des sommes directes dans $D(\mathcal{A})$ d'après le
Lemme \ref{lemma-direct-sums}.
Le résultat découle donc de la définition des colimites dérivées donnée dans
la Définition \ref{definition-derived-colimit}, et du fait qu'une suite exacte
courte de complexes donne un triangle distingué
(Lemme \ref{lemma-derived-canonical-delta-functor}).
\end{proof}

\begin{lemma}
\label{lemma-cohomology-of-hocolim}
Soit $\mathcal{D}$ une catégorie triangulée possédant des sommes directes
dénombrables. Soit $\mathcal{A}$ une catégorie abélienne dans laquelle les
colimites indexées par $\mathbf{N}$ sont exactes.
Soit $H : \mathcal{D} \to \mathcal{A}$ un foncteur homologique commutant aux
sommes directes dénombrables.
Alors $H(\text{hocolim} K_n) = \colim H(K_n)$ pour tout système d'objets de
$\mathcal{D}$.
\end{lemma}

\begin{proof}
Écrivons $K = \text{hocolim} K_n$. Appliquons $H$ au triangle distingué qui
le définit. Nous obtenons
$$
\bigoplus H(K_n) \to \bigoplus H(K_n)
\to H(K) \to
\bigoplus H(K_n[1]) \to \bigoplus H(K_n[1])
$$
où le premier morphisme est donné par $1 - H(f_n)$ et le dernier par
$1 - H(f_n[1])$.
Le Lemme \ref{lemma-compute-colimit} montre que cela prouve le lemme.
\end{proof}

\noindent
Le lemme suivant nous dit que la formation des morphismes issus d'un objet
compact (qui sera défini plus loin) commute aux colimites dérivées.

\begin{lemma}
\label{lemma-commutes-with-countable-sums}
Soit $\mathcal{D}$ une catégorie triangulée possédant des sommes directes
dénombrables. Soit $K \in \mathcal{D}$ un objet tel que, pour toute famille
dénombrable d'objets $E_n \in \mathcal{D}$, l'application canonique
$$
\bigoplus \Hom_\mathcal{D}(K, E_n)
\longrightarrow
\Hom_\mathcal{D}(K, \bigoplus E_n)
$$
soit bijective. Alors, pour tout système $L_n$ de $\mathcal{D}$ indexé par
$\mathbf{N}$ dont la colimite dérivée $L = \text{hocolim} L_n$ existe,
l'application
$$
\colim \Hom_\mathcal{D}(K, L_n) \longrightarrow \Hom_\mathcal{D}(K, L)
$$
est bijective.
\end{lemma}

\begin{proof}
D'après le Lemme \ref{lemma-representable-homological}, c'est un cas
particulier du Lemme \ref{lemma-cohomology-of-hocolim}.
\end{proof}

\section{Limites dérivées}
\label{section-derived-limit}

\noindent
Dans une catégorie triangulée, on dispose d'une notion de limite dérivée.

\begin{definition}
\label{definition-derived-limit}
Soit $\mathcal{D}$ une catégorie triangulée.
Soit $(K_n, f_n)$ un système projectif d'objets de $\mathcal{D}$.
On dit qu'un objet $K$ est une {\it limite dérivée}, ou une
{\it limite homotopique}, du système $(K_n)$ si le produit $\prod K_n$ existe
et s'il existe un triangle distingué
$$
K \to \prod K_n \to \prod K_n \to K[1]
$$
où le morphisme $\prod K_n \to \prod K_n$ est donné par
$(k_n) \mapsto (k_n - f_{n + 1}(k_{n + 1}))$. Dans ce cas, nous l'indiquons
parfois par la notation $K = R\lim K_n$.
\end{definition}

\noindent
D'après TR3, une limite dérivée, si elle existe, est unique à isomorphisme
(non unique) près. De plus, d'après TR1, une limite dérivée $R\lim K_n$ existe
dès que $\prod K_n$ existe. La catégorie dérivée $D(\textit{Ab})$ de la
catégorie des groupes abéliens est un exemple de catégorie triangulée dans
laquelle toutes les limites dérivées existent.

\medskip\noindent
Cette absence d'unicité rend difficile la détermination précise de la limite
dérivée. Dans Algèbre complémentaire, Lemme
\ref{more-algebra-lemma-map-into-Rlim}, le lecteur trouvera une suite exacte
$$
0 \to R^1\lim \Hom(L, K_n[-1]) \to \Hom(L, R\lim K_n)
\to \lim \Hom(L, K_n) \to 0
$$
qui décrit les $\Hom$ à valeurs dans une limite dérivée à l'aide des $\Hom$
usuels.

\begin{lemma}
\label{lemma-products}
Soit $\mathcal{A}$ une catégorie abélienne dont les produits dénombrables sont
exacts. Alors
\begin{enumerate}
\item $D(\mathcal{A})$ possède des produits dénombrables,
\item les produits dénombrables $\prod K_i$ dans $D(\mathcal{A})$ s'obtiennent
en prenant les produits terme à terme de complexes quelconques représentant
les $K_i$, et
\item $H^p(\prod K_i) = \prod H^p(K_i)$.
\end{enumerate}
\end{lemma}

\begin{proof}
Soit $K_i^\bullet$ un complexe représentant $K_i$ dans $D(\mathcal{A})$.
Soit $L^\bullet$ un complexe. Supposons donnés des morphismes
$\alpha_i : L^\bullet \to K_i^\bullet$ dans $D(\mathcal{A})$.
Cela signifie qu'il existe des quasi-isomorphismes
$s_i : K_i^\bullet \to M_i^\bullet$ de complexes et des morphismes de
complexes $f_i : L^\bullet \to M_i^\bullet$ tels que
$\alpha_i = s_i^{-1}f_i$. Par hypothèse, le morphisme de complexes
$$
s : \prod K_i^\bullet \longrightarrow \prod M_i^\bullet
$$
est un quasi-isomorphisme. En posant $f = \prod f_i$, nous voyons donc que
$\alpha = s^{-1}f$ est un morphisme de $D(\mathcal{A})$ dont la composée avec
la projection $\prod K_i^\bullet \to K_i^\bullet$ est $\alpha_i$.
Nous omettons la vérification de l'unicité de $\alpha$.
\end{proof}

\noindent
Il faudrait énoncer et démontrer ici les énoncés duaux des Lemmes
\ref{lemma-compute-colimit}, \ref{lemma-colim-hocolim} et
\ref{lemma-commutes-with-countable-sums}. Nous n'en connaissons toutefois
aucune application pour le moment.

\begin{lemma}
\label{lemma-inverse-limit-bounded-below}
Soit $\mathcal{A}$ une catégorie abélienne possédant des produits dénombrables
et assez d'injectifs. Soit $(K_n)$ un système projectif d'objets de
$D^+(\mathcal{A})$. Alors $R\lim K_n$ existe.
\end{lemma}

\begin{proof}
Il suffit de montrer que $\prod K_n$ existe dans $D(\mathcal{A})$.
Pour tout $n$, nous pouvons représenter $K_n$ par un complexe borné à gauche
$I_n^\bullet$ d'objets injectifs
(Lemme \ref{lemma-injective-resolutions-exist}).
Alors $\prod K_n$ est représenté par $\prod I_n^\bullet$; voir
Lemme \ref{lemma-product-K-injective}.
\end{proof}

\begin{remark}
\label{remark-functoriality-derived-limit}
Soit $\mathcal{D}$ une catégorie triangulée. Soient $(K_n)$ et $(L_n)$ des
systèmes projectifs d'objets de $\mathcal{D}$, de limites dérivées $K$ et $L$.
Supposons donné un morphisme $a : (K_n) \to (L_n)$ de pro-objets; voir
Catégories, Exemple
\ref{categories-example-pro-morphism-inverse-systems}.
Cela signifie que des morphismes $a_n : K_{m(n)} \to L_n$ nous sont donnés.
Nous pouvons supposer que $m(1) < m(2) < m(3) < \ldots$ et que les morphismes
$a_n$ sont compatibles avec les morphismes de transition des systèmes.
Nous pouvons alors considérer les morphismes
$$
a', a'' : \prod K_n \longrightarrow \prod L_n
$$
définis par les règles $a'((k_n)) = (a_n(k_{m(n)}))$ et
$$
a''((k_n)) =
(a_n(k_{m(n)} + f(k_{m(n) + 1}) + \ldots + f(k_{m(n + 1) - 1})))
$$
où chaque occurrence de $f$ désigne un morphisme de transition approprié du
système projectif $(K_n)$. Le diagramme en traits pleins
$$
\xymatrix{
K \ar[r] \ar@{..>}[d] & \prod K_n \ar[r] \ar[d]^{a'} &
\prod K_n \ar[d]^{a''} \\
L \ar[r] & \prod L_n \ar[r] & \prod L_n
}
$$
est commutatif et, d'après TR3, nous obtenons la flèche en pointillé qui en
fait un morphisme de triangles distingués. Nous avertissons le lecteur que le
morphisme $K \to L$ n'est pas unique. Nous verrons plus loin que, si $a$ est
un pro-isomorphisme, alors $K \to L$ est un isomorphisme; voir
Algèbre complémentaire, Lemme
\ref{more-algebra-lemma-pro-isomorphism-Rlim}.
\end{remark}

\begin{remark}
\label{remark-map-into-derived-limit-truncations}
Soit $\mathcal{A}$ une catégorie abélienne. Soit $K^\bullet$ un complexe de
$\mathcal{A}$. Les complexes $\tau_{\geq -n}K^\bullet$ forment un système
projectif de complexes, et déterminent en particulier un système projectif
dans $D(\mathcal{A})$. Supposons que
$R\lim \tau_{\geq -n}K^\bullet$ existe. Les morphismes canoniques
$c_n : K^\bullet \to \tau_{\geq -n}K^\bullet$ sont alors compatibles avec les
morphismes de transition de notre système projectif. Le triangle distingué de
la Définition \ref{definition-derived-limit} et le
Lemme \ref{lemma-representable-homological} montrent qu'il existe un morphisme
$$
c : K^\bullet \longrightarrow R\lim \tau_{\geq -n}K^\bullet
$$
dans $D(\mathcal{A})$ tel que la composée de $c$ avec la projection
$R\lim \tau_{\geq -n}K^\bullet \to \tau_{\geq -m}K^\bullet$ soit égale à
$c_m$. Le morphisme $c$ peut ne pas être unique, mais nous affirmons que le
fait que $c$ soit ou non un isomorphisme est indépendant du choix de $c$ (et
de notre choix de la limite homotopique). En effet, pour
$i \in \mathbf{Z}$ et $m > -i$, la composée
$$
H^i(K^\bullet) \xrightarrow{H^i(c)} H^i(R\lim \tau_{\geq -n}K^\bullet)
\to H^i(\tau_{\geq -m}K^\bullet) = H^i(K^\bullet)
$$
est l'identité. Ainsi, $H^i(c)$ est un isomorphisme si et seulement si le
second morphisme est un isomorphisme. Cette propriété est indépendante de
$c$, ainsi que du choix de la limite homotopique (car deux choix quelconques
sont isomorphes).
\end{remark}

\begin{lemma}
\label{lemma-difficulty-K-injectives}
Soit $\mathcal{A}$ une catégorie abélienne possédant des produits dénombrables
et assez d'injectifs. Soit $K^\bullet$ un complexe. Soit $I_n^\bullet$ le
système projectif de complexes bornés à gauche d'objets injectifs produit par
le Lemme \ref{lemma-special-inverse-system}. Alors
$I^\bullet = \lim I_n^\bullet$ existe, est K-injectif, représente
$R\lim \tau_{\geq -n}K^\bullet$ dans $D(\mathcal{A})$, et les conditions
suivantes sont équivalentes :
\begin{enumerate}
\item le morphisme $K^\bullet \to I^\bullet$ (voir la preuve) est un
quasi-isomorphisme,
\item le morphisme $K^\bullet \to R\lim \tau_{\geq -n}K^\bullet$ de la
Remarque \ref{remark-map-into-derived-limit-truncations} est un isomorphisme
dans $D(\mathcal{A})$.
\end{enumerate}
\end{lemma}

\begin{proof}
L'énoncé du lemme a un sens, car
$R\lim \tau_{\geq -n}K^\bullet$ existe d'après le
Lemme \ref{lemma-inverse-limit-bounded-below}.
Chaque complexe $I_n^\bullet$ est K-injectif d'après le
Lemme \ref{lemma-bounded-below-injectives-K-injective}.
Choisissons des décompositions en somme directe
$I_{n + 1}^p = C_{n + 1}^p \oplus I_n^p$ pour tout $n \geq 1$.
Posons $C_1^p = I_1^p$. Le complexe $I^\bullet = \lim I_n^\bullet$ existe,
car nous pouvons prendre $I^p = \prod_{n \geq 1} C_n^p$.
Fixons $p \in \mathbf{Z}$. Nous affirmons qu'il existe une suite exacte
courte scindée
$$
0 \to I^p \to \prod I_n^p \to \prod I_n^p \to 0
$$
d'objets de $\mathcal{A}$. Ici, le premier morphisme est donné par les
projections $I^p \to I_n^p$, et le second par
$(x_n) \mapsto (x_n - f^p_{n + 1}(x_{n + 1}))$, où
$f^p_n : I_n^p \to I_{n - 1}^p$ sont les morphismes de transition.
Le scindage provient du morphisme
$\prod I_n^p \to \prod C_n^p = I^p$.
Nous obtenons une suite exacte courte de complexes, scindée terme à terme,
$$
0 \to I^\bullet \to \prod I_n^\bullet \to \prod I_n^\bullet \to 0
$$
et donc un triangle distingué associé dans $K(\mathcal{A})$ et dans
$D(\mathcal{A})$. D'après le Lemme \ref{lemma-product-K-injective}, les
produits sont K-injectifs et représentent les produits correspondants dans
$D(\mathcal{A})$.
Il s'ensuit que $I^\bullet$ représente $R\lim I_n^\bullet$
(Définition \ref{definition-derived-limit}).
Comme $R\lim I_n^\bullet \cong R\lim \tau_{\geq -n}K^\bullet$, puisque les
limites dérivées sont définies au niveau de la catégorie dérivée, nous voyons
que $I^\bullet$ représente $R\lim \tau_{\geq -n}K^\bullet$.
De plus, le complexe $I^\bullet$ est K-injectif d'après le
Lemme \ref{lemma-triangle-K-injective}.
Le diagramme commutatif du Lemme \ref{lemma-special-inverse-system} et
l'égalité $K^i = (\tau_{\geq -n}K^\bullet)^i$ pour $n \gg 0$ montrent que
nous obtenons un unique morphisme $\gamma : K^\bullet \to I^\bullet$ tel que
les diagrammes
$$
\xymatrix{
K^\bullet \ar[r] \ar[d]_\gamma & \tau_{\geq -n} K^\bullet \ar[d] \\
I^\bullet \ar[r] & I_n^\bullet
}
$$
commutent. Il s'ensuit que $\gamma$ est un morphisme de complexes qui
représente le morphisme
$c : K^\bullet \to R\lim \tau_{\geq -n}K^\bullet$ de la
Remarque \ref{remark-map-into-derived-limit-truncations} dans
$D(\mathcal{A})$. Autrement dit, le diagramme
$$
\xymatrix{
K^\bullet \ar[r]_-c \ar[d]_\gamma & R\lim \tau_{\geq -n} K^\bullet
\ar[d]^{\cong} \\
I^\bullet \ar[r]^-{\cong} & R\lim I_n^\bullet
}
$$
est commutatif dans $D(\mathcal{A})$. Le lemme en résulte.
\end{proof}

\begin{lemma}
\label{lemma-enough-K-injectives-Ab4-star}
Soit $\mathcal{A}$ une catégorie abélienne possédant assez d'injectifs et dont
les produits dénombrables sont exacts. Alors, pour tout complexe, il existe un
quasi-isomorphisme vers un complexe K-injectif.
\end{lemma}

\begin{proof}
D'après le Lemme \ref{lemma-difficulty-K-injectives}, il suffit de montrer que
$K \to R\lim\tau_{\geq -n}K$ est un isomorphisme pour tout $K$ dans
$D(\mathcal{A})$. Considérons le triangle distingué qui le définit :
$$
R\lim\tau_{\geq -n}K \to
\prod \tau_{\geq -n}K \to
\prod \tau_{\geq -n}K \to
(R\lim\tau_{\geq -n}K)[1]
$$
D'après le Lemme \ref{lemma-products}, nous avons
$$
H^p(\prod \tau_{\geq -n}K) =
\prod\nolimits_{n\,:\,p \geq -n} H^p(K)
$$
La suite exacte longue de cohomologie du triangle distingué affiché montre
alors immédiatement que $H^p(R\lim \tau_{\geq -n}K) = H^p(K)$.
\end{proof}

\section{Opérations sur les sous-catégories pleines}
\label{section-operate-on-full}

\noindent
Soit $\mathcal{T}$ une catégorie triangulée. Nous identifierons les
sous-catégories pleines de $\mathcal{T}$ aux sous-ensembles de
$\Ob(\mathcal{T})$. Étant données des sous-catégories pleines
$\mathcal{A}, \mathcal{B}, \ldots$, nous notons
\begin{enumerate}
\item $\mathcal{A}[a, b]$, pour $-\infty \leq a \leq b \leq \infty$, la
sous-catégorie pleine de $\mathcal{T}$ formée de tous les objets $A[-i]$ avec
$i \in [a, b] \cap \mathbf{Z}$ et $A \in \Ob(\mathcal{A})$ (attention au
signe moins!),
\item $smd(\mathcal{A})$ la sous-catégorie pleine de $\mathcal{T}$ formée de
tous les objets isomorphes à des facteurs directs d'objets de $\mathcal{A}$,
\item $add(\mathcal{A})$ la sous-catégorie pleine de $\mathcal{T}$ formée de
tous les objets isomorphes à des sommes directes finies d'objets de
$\mathcal{A}$,
\item $\mathcal{A} \star \mathcal{B}$ la sous-catégorie pleine de
$\mathcal{T}$ formée de tous les objets $X$ de $\mathcal{T}$ qui s'insèrent
dans un triangle distingué $A \to X \to B$, avec
$A \in \Ob(\mathcal{A})$ et $B \in \Ob(\mathcal{B})$,
\item $\mathcal{A}^{\star n} = \mathcal{A} \star \ldots \star \mathcal{A}$,
avec $n \geq 1$ facteurs (nous verrons ci-dessous que $\star$ est
associative),
\item $smd(add(\mathcal{A})^{\star n}) =
smd(add(\mathcal{A}) \star \ldots \star add(\mathcal{A}))$, avec
$n \geq 1$ facteurs.
\end{enumerate}
Si $E$ est un objet de $\mathcal{T}$, nous considérons parfois aussi $E$ comme
la sous-catégorie pleine de $\mathcal{T}$ dont l'unique objet est $E$.
Nous pouvons alors considérer des expressions telles que $add(E[-1, 2])$, et
ainsi de suite. Nous avertissons le lecteur que cette notation n'est pas
universellement admise.

\begin{lemma}
\label{lemma-associativity-star}
Soit $\mathcal{T}$ une catégorie triangulée.
Étant données des sous-catégories pleines $\mathcal{A}$, $\mathcal{B}$,
$\mathcal{C}$, nous avons
$(\mathcal{A} \star \mathcal{B}) \star \mathcal{C} =
\mathcal{A} \star (\mathcal{B} \star \mathcal{C})$.
\end{lemma}

\begin{proof}
Si nous avons des triangles distingués $A \to X \to B$ et
$X \to Y \to C$, l'Axiome TR4 fournit des triangles distingués
$A \to Y \to Z$ et $B \to Z \to C$.
\end{proof}

\begin{lemma}
\label{lemma-smd-star}
Soit $\mathcal{T}$ une catégorie triangulée.
Étant données des sous-catégories pleines $\mathcal{A}$, $\mathcal{B}$, nous
avons
$smd(\mathcal{A}) \star smd(\mathcal{B}) \subset
smd(\mathcal{A} \star \mathcal{B})$ et
$smd(smd(\mathcal{A}) \star smd(\mathcal{B})) =
smd(\mathcal{A} \star \mathcal{B})$.
\end{lemma}

\begin{proof}
Supposons donné un triangle distingué $A_1 \to X \to B_1$, où
$A_1 \oplus A_2 \in \Ob(\mathcal{A})$ et
$B_1 \oplus B_2 \in \Ob(\mathcal{B})$.
Nous obtenons alors un triangle distingué
$A_1 \oplus A_2 \to A_2 \oplus X \oplus B_2 \to B_1 \oplus B_2$, ce qui
prouve que $X$ appartient à $smd(\mathcal{A} \star \mathcal{B})$.
Cela prouve l'inclusion. L'égalité s'en déduit immédiatement.
\end{proof}

\begin{lemma}
\label{lemma-add-star}
Soit $\mathcal{T}$ une catégorie triangulée. Étant données des sous-catégories
pleines $\mathcal{A}$, $\mathcal{B}$, les sous-catégories pleines
$add(\mathcal{A}) \star add(\mathcal{B})$ et $smd(add(\mathcal{A}))$ sont
stables par sommes directes.
\end{lemma}

\begin{proof}
En effet, si $A \to X \to B$ et $A' \to X' \to B'$ sont des triangles
distingués, avec $A, A' \in add(\mathcal{A})$ et
$B, B' \in add(\mathcal{B})$, alors
$A \oplus A' \to X \oplus X' \to B \oplus B'$ est un triangle distingué,
avec $A \oplus A' \in add(\mathcal{A})$ et
$B \oplus B' \in add(\mathcal{B})$.
Le résultat pour $smd(add(\mathcal{A}))$ est immédiat.
\end{proof}

\begin{lemma}
\label{lemma-cone-n}
Soit $\mathcal{T}$ une catégorie triangulée. Étant donnée une sous-catégorie
pleine $\mathcal{A}$, pour $n \geq 1$ la sous-catégorie
$$
\mathcal{C}_n = smd(add(\mathcal{A})^{\star n}) =
smd(add(\mathcal{A}) \star \ldots \star add(\mathcal{A}))
$$
définie ci-dessus est une sous-catégorie strictement pleine de $\mathcal{T}$,
stable par sommes directes et facteurs directs, et
$\mathcal{C}_{n + m} = smd(\mathcal{C}_n \star \mathcal{C}_m)$ pour tous
$n, m \geq 1$.
\end{lemma}

\begin{proof}
Cela résulte immédiatement des Lemmes \ref{lemma-associativity-star},
\ref{lemma-smd-star} et \ref{lemma-add-star}.
\end{proof}

\begin{remark}
\label{remark-operations-functor}
Soit $F : \mathcal{T} \to \mathcal{T}'$ un foncteur exact entre catégories
triangulées. Étant donnée une sous-catégorie pleine $\mathcal{A}$ de
$\mathcal{T}$, nous notons $F(\mathcal{A})$ la sous-catégorie pleine de
$\mathcal{T}'$ dont les objets sont tous les objets $F(A)$, avec
$A \in \Ob(\mathcal{A})$. Nous avons
$$
F(\mathcal{A}[a, b]) = F(\mathcal{A})[a, b]
$$
$$
F(smd(\mathcal{A})) \subset smd(F(\mathcal{A})),
$$
$$
F(add(\mathcal{A})) \subset add(F(\mathcal{A})),
$$
$$
F(\mathcal{A} \star \mathcal{B}) \subset F(\mathcal{A}) \star F(\mathcal{B}),
$$
$$
F(\mathcal{A}^{\star n}) \subset F(\mathcal{A})^{\star n}.
$$
Nous omettons les vérifications immédiates.
\end{remark}

\begin{remark}
\label{remark-operations-unions}
Soit $\mathcal{T}$ une catégorie triangulée. Étant données des sous-catégories
pleines
$\mathcal{A}_1 \subset \mathcal{A}_2 \subset \mathcal{A}_3 \subset \ldots$
et $\mathcal{B}$ de $\mathcal{T}$, nous avons
$$
\left(\bigcup \mathcal{A}_i\right)[a, b] = \bigcup \mathcal{A}_i[a, b]
$$
$$
smd\left(\bigcup \mathcal{A}_i\right) = \bigcup smd(\mathcal{A}_i),
$$
$$
add\left(\bigcup \mathcal{A}_i\right) = \bigcup add(\mathcal{A}_i),
$$
$$
\left(\bigcup \mathcal{A}_i\right) \star \mathcal{B} =
\bigcup \mathcal{A}_i \star \mathcal{B},
$$
$$
\mathcal{B} \star \left(\bigcup \mathcal{A}_i\right) =
\bigcup \mathcal{B} \star \mathcal{A}_i,
$$
$$
\left(\bigcup \mathcal{A}_i\right)^{\star n} =
\bigcup \mathcal{A}_i^{\star n}.
$$
Nous omettons les vérifications immédiates.
\end{remark}

\begin{lemma}
\label{lemma-in-cone-n}
Soit $\mathcal{A}$ une catégorie abélienne. Soit
$\mathcal{D} = D(\mathcal{A})$.
Soit $\mathcal{E} \subset \Ob(\mathcal{A})$ un sous-ensemble, que nous
considérons aussi comme un sous-ensemble de $\Ob(\mathcal{D})$.
Soit $K$ un objet de $\mathcal{D}$.
\begin{enumerate}
\item Soit $b \geq a$ et supposons que $H^i(K)$ soit nul pour
$i \not \in [a, b]$ et que $H^i(K) \in \mathcal{E}$ si
$i \in [a, b]$. Alors $K$ appartient à
$smd(add(\mathcal{E}[a, b])^{\star (b - a + 1)})$.
\item Soit $b \geq a$ et supposons que $H^i(K)$ soit nul pour
$i \not \in [a, b]$ et que $H^i(K) \in smd(add(\mathcal{E}))$ si
$i \in [a, b]$. Alors $K$ appartient à
$smd(add(\mathcal{E}[a, b])^{\star (b - a + 1)})$.
\item Soit $b \geq a$ et supposons que $K$ puisse être représenté par un
complexe $K^\bullet$ tel que $K^i = 0$ pour $i \not \in [a, b]$ et
$K^i \in \mathcal{E}$ pour $i \in [a, b]$. Alors $K$ appartient à
$smd(add(\mathcal{E}[a, b])^{\star (b - a + 1)})$.
\item Soit $b \geq a$ et supposons que $K$ puisse être représenté par un
complexe $K^\bullet$ tel que $K^i = 0$ pour $i \not \in [a, b]$ et
$K^i \in smd(add(\mathcal{E}))$ pour $i \in [a, b]$. Alors $K$ appartient à
$smd(add(\mathcal{E}[a, b])^{\star (b - a + 1)})$.
\end{enumerate}
\end{lemma}

\begin{proof}
Nous utiliserons le Lemme \ref{lemma-cone-n} sans autre mention.
Nous allons démontrer (2), qui entraîne immédiatement (1). Raisonnons par
récurrence sur $b - a$. Si $b - a = 0$, alors $K$ est isomorphe à
$H^a(K)[-a]$ dans $\mathcal{D}$, et le résultat est immédiat.
Si $b - a > 0$, considérons le triangle distingué
$$
\tau_{\leq b - 1}K \to K \to H^b(K)[-b]
$$
et concluons par récurrence sur $b - a$. Nous omettons la preuve de (3) et
(4).
\end{proof}

\begin{lemma}
\label{lemma-forward-cone-n}
Soit $\mathcal{T}$ une catégorie triangulée. Soit
$H : \mathcal{T} \to \mathcal{A}$ un foncteur homologique à valeurs dans une
catégorie abélienne $\mathcal{A}$.
Soient $a \leq b$ et $\mathcal{E} \subset \Ob(\mathcal{T})$ un sous-ensemble
tel que $H^i(E) = 0$ pour $E \in \mathcal{E}$ et $i \not \in [a, b]$.
Alors, pour $X \in smd(add(\mathcal{E}[-m, m])^{\star n})$, nous avons
$H^i(X) = 0$ pour $i \not \in [a - m, b + m]$.
\end{lemma}

\begin{proof}
Omis. Agréable exercice sur les définitions.
\end{proof}

\section{Générateurs de catégories triangulées}
\label{section-generators}

\noindent
Dans cette section, nous introduisons brièvement quelques-unes des différentes
notions de générateur d'une catégorie triangulée. Notre terminologie est tirée
de \cite{BvdB} (si ce n'est que nous employons « saturée » là où les auteurs
disent « épaisse »; voir Définition \ref{definition-saturated}; notre définition
de $add(\mathcal{A})$ est également différente).

\medskip\noindent
Soit $\mathcal{D}$ une catégorie triangulée. Soit $E$ un objet de
$\mathcal{D}$. Notons $\langle E \rangle_1$ la sous-catégorie strictement
pleine de $\mathcal{D}$ formée des objets de $\mathcal{D}$ isomorphes à des
facteurs directs de sommes directes finies
$$
\bigoplus\nolimits_{i = 1, \ldots, r} E[n_i]
$$
de translatés de $E$. Il est clair qu'avec les notations de la
section \ref{section-operate-on-full}, nous avons
$$
\langle E \rangle_1 = smd(add(E[-\infty, \infty]))
$$
Pour $n > 1$, notons $\langle E \rangle_n$ la sous-catégorie pleine de
$\mathcal{D}$ formée des objets de $\mathcal{D}$ isomorphes à des facteurs
directs d'objets $X$ qui s'insèrent dans un triangle distingué
$$
A \to X \to B \to A[1]
$$
où $A$ est un objet de $\langle E \rangle_1$ et $B$ un objet de
$\langle E \rangle_{n - 1}$. Avec les notations de la
section \ref{section-operate-on-full}, nous avons
$$
\langle E \rangle_n = smd(\langle E \rangle_1 \star \langle E \rangle_{n - 1})
$$
Chacune des catégories $\langle E \rangle_n$ est une sous-catégorie additive
strictement pleine (d'après le Lemme \ref{lemma-add-star}) de $\mathcal{D}$,
stable par translations et par passage aux facteurs directs. Toutefois,
$\langle E \rangle_n$ n'est pas nécessairement stable par « formation des
cônes » ou par « extensions »; ce n'est donc pas nécessairement une
sous-catégorie triangulée. Cela sera en revanche vrai pour la sous-catégorie
$$
\langle E \rangle = \bigcup\nolimits_n \langle E \rangle_n
$$
comme le montreront les lemmes ci-dessous.

\begin{lemma}
\label{lemma-generated-by-E-explicit}
Soit $\mathcal{T}$ une catégorie triangulée. Soit $E$ un objet de
$\mathcal{T}$. Pour $n \geq 1$, nous avons
$$
\langle E \rangle_n =
smd(\langle E \rangle_1 \star \ldots \star \langle E \rangle_1) =
smd({\langle E \rangle_1}^{\star n}) =
\bigcup\nolimits_{m \geq 1} smd(add(E[-m, m])^{\star n})
$$
Pour $n, n' \geq 1$, nous avons $\langle E \rangle_{n + n'} =
smd(\langle E \rangle_n \star \langle E \rangle_{n'})$.
\end{lemma}

\begin{proof}
L'égalité de gauche dans la formule affichée résulte des Lemmes
\ref{lemma-associativity-star} et \ref{lemma-smd-star}, ainsi que d'une
récurrence. L'égalité du milieu est une question de notation.
Comme $\langle E \rangle_1 = smd(add(E[-\infty, \infty]))$ et
$E[-\infty, \infty] = \bigcup_{m \geq 1} E[-m, m]$, la
Remarque \ref{remark-operations-unions} et le Lemme \ref{lemma-smd-star}
donnent l'égalité de droite. La dernière assertion découle alors de cette
remarque et de l'assertion correspondante du Lemme \ref{lemma-cone-n}.
\end{proof}

\begin{lemma}
\label{lemma-find-smallest-containing-E}
Soit $\mathcal{D}$ une catégorie triangulée. Soit $E$ un objet de
$\mathcal{D}$. La sous-catégorie
$$
\langle E \rangle = \bigcup\nolimits_n \langle E \rangle_n
= \bigcup\nolimits_{n, m \geq 1} smd(add(E[-m, m])^{\star n})
$$
est une sous-catégorie triangulée, saturée et strictement pleine de
$\mathcal{D}$; c'est la plus petite sous-catégorie de ce type de
$\mathcal{D}$ qui contienne l'objet $E$.
\end{lemma}

\begin{proof}
L'égalité de droite résulte du
Lemme \ref{lemma-generated-by-E-explicit}.
Il est clair que $\langle E \rangle = \bigcup \langle E \rangle_n$ contient
$E$ et est stable par translations, sommes directes et facteurs directs.
Si $A \in \langle E \rangle_a$ et $B \in \langle E \rangle_b$, et si
$A \to X \to B \to A[1]$ est un triangle distingué, alors
$X \in \langle E \rangle_{a + b}$ d'après le
Lemme \ref{lemma-generated-by-E-explicit}.
Ainsi, $\bigcup \langle E \rangle_n$ est aussi stable par extensions; il
s'ensuit que c'est une sous-catégorie triangulée.

\medskip\noindent
Enfin, soit $\mathcal{D}' \subset \mathcal{D}$ une sous-catégorie triangulée,
saturée et strictement pleine de $\mathcal{D}$ qui contienne $E$. Alors
$\mathcal{D}'[-\infty, \infty] \subset \mathcal{D}'$,
$add(\mathcal{D}') \subset \mathcal{D}'$,
$smd(\mathcal{D}') \subset \mathcal{D}'$ et
$\mathcal{D}' \star \mathcal{D}' \subset \mathcal{D}'$.
Autrement dit, toutes les opérations dont nous nous sommes servis pour
construire $\langle E \rangle$ à partir de $E$ préservent $\mathcal{D}'$.
Par conséquent, $\langle E \rangle \subset \mathcal{D}'$, ce qui achève la
preuve.
\end{proof}

\begin{definition}
\label{definition-generators}
Soit $\mathcal{D}$ une catégorie triangulée. Soit $E$ un objet de
$\mathcal{D}$.
\begin{enumerate}
\item On dit que $E$ est un {\it générateur classique} de $\mathcal{D}$ si la
plus petite sous-catégorie triangulée, saturée et strictement pleine de
$\mathcal{D}$ qui contienne $E$ est égale à $\mathcal{D}$; autrement dit, si
$\langle E \rangle = \mathcal{D}$.
\item On dit que $E$ est un {\it générateur fort} de $\mathcal{D}$ si
$\langle E \rangle_n = \mathcal{D}$ pour un certain $n \geq 1$.
\item On dit que $E$ est un {\it générateur faible}, ou un
{\it générateur}, de $\mathcal{D}$ si, pour tout objet non nul $K$ de
$\mathcal{D}$, il existe un entier $n$ et un morphisme non nul $E \to K[n]$.
\end{enumerate}
\end{definition}

\noindent
Cette définition se généralise au cas d'une famille d'objets.

\begin{lemma}
\label{lemma-right-orthogonal}
Soit $\mathcal{D}$ une catégorie triangulée. Soient $E, K$ des objets de
$\mathcal{D}$. Les conditions suivantes sont équivalentes :
\begin{enumerate}
\item $\Hom(E, K[i]) = 0$ pour tout $i \in \mathbf{Z}$,
\item $\Hom(E', K) = 0$ pour tout $E' \in \langle E \rangle$.
\end{enumerate}
\end{lemma}

\begin{proof}
L'implication (2) $\Rightarrow$ (1) est immédiate. Réciproquement, supposons
(1). Alors $\Hom(X, K) = 0$ pour tout $X$ dans $\langle E \rangle_1$.
En raisonnant par récurrence sur $n$ et en utilisant le
Lemme \ref{lemma-representable-homological}, nous voyons que
$\Hom(X, K) = 0$ pour tout $X$ dans $\langle E \rangle_n$.
\end{proof}

\begin{lemma}
\label{lemma-classical-generator-generator}
Soit $\mathcal{D}$ une catégorie triangulée. Soit $E$ un objet de
$\mathcal{D}$. Si $E$ est un générateur classique de $\mathcal{D}$, alors
$E$ est un générateur.
\end{lemma}

\begin{proof}
Supposons que $E$ soit un générateur classique. Soit $K$ un objet de
$\mathcal{D}$ tel que $\Hom(E, K[i]) = 0$ pour tout $i \in \mathbf{Z}$.
D'après le Lemme \ref{lemma-right-orthogonal}, $\Hom(E', K) = 0$ pour tout
$E'$ dans $\langle E \rangle$. Or, puisque
$\mathcal{D} = \langle E \rangle$, nous concluons que
$\text{id}_K = 0$, c'est-à-dire que $K = 0$.
\end{proof}

\begin{lemma}
\label{lemma-classical-generator-strong-generator}
Soit $\mathcal{D}$ une catégorie triangulée possédant un générateur fort.
Soit $E$ un objet de $\mathcal{D}$. Si $E$ est un générateur classique de
$\mathcal{D}$, alors $E$ est un générateur fort.
\end{lemma}

\begin{proof}
Soit $E'$ un objet de $\mathcal{D}$ tel que
$\mathcal{D} = \langle E' \rangle_n$. Comme
$\mathcal{D} = \langle E \rangle$, nous avons
$E' \in \langle E \rangle_m$ pour un certain $m \geq 1$, d'après le
Lemme \ref{lemma-find-smallest-containing-E}.
Alors $\langle E' \rangle_1 \subset \langle E \rangle_m$, donc
$$
\mathcal{D} =
\langle E' \rangle_n = smd(
\langle E' \rangle_1 \star \ldots \star \langle E' \rangle_1)
\subset
smd(
\langle E \rangle_m \star \ldots \star \langle E \rangle_m)
=
\langle E \rangle_{nm}
$$
comme voulu. Nous avons utilisé ici le
Lemme \ref{lemma-generated-by-E-explicit}.
\end{proof}

\begin{remark}
\label{remark-check-on-generator}
Soit $\mathcal{D}$ une catégorie triangulée. Soit $E$ un objet de
$\mathcal{D}$. Soit $T$ une propriété des objets de $\mathcal{D}$.
Supposons que
\begin{enumerate}
\item si $K_i \in \mathcal{D}$, $i = 1, \ldots, r$, et si $T(K_i)$ pour
$i = 1, \ldots, r$, alors $T(\bigoplus K_i)$,
\item si $K \to L \to M \to K[1]$ est un triangle distingué et que $T$ est
vraie pour deux de ces objets, alors $T$ est vraie pour le troisième,
\item si $T(K \oplus L)$, alors $T(K)$ et $T(L)$, et
\item $T(E[n])$ est vraie pour tout $n$.
\end{enumerate}
Alors $T$ est vraie pour tous les objets de $\langle E \rangle$.
\end{remark}

\section{Objets compacts}
\label{section-compact}

\noindent
Voici la définition.

\begin{definition}
\label{definition-compact-object}
Soit $\mathcal{D}$ une catégorie additive possédant des sommes directes
arbitraires. Un {\it objet compact} de $\mathcal{D}$ est un objet $K$ tel que
l'application
$$
\bigoplus\nolimits_{i \in I} \Hom_{\mathcal{D}}(K, E_i)
\longrightarrow
\Hom_{\mathcal{D}}(K, \bigoplus\nolimits_{i \in I} E_i)
$$
soit bijective pour tout ensemble $I$ et tous objets
$E_i \in \Ob(\mathcal{D})$ paramétrés par $i \in I$.
\end{definition}

\noindent
Cette notion s'avère très utile en géométrie algébrique. C'est une condition
intrinsèque sur les objets qui les force à être, précisément, compacts.

\begin{lemma}
\label{lemma-compact-objects-subcategory}
Soit $\mathcal{D}$ une catégorie (pré-)triangulée possédant des sommes
directes. Alors les objets compacts de $\mathcal{D}$ sont les objets d'une
sous-catégorie (pré-)triangulée, karoubienne, saturée et strictement pleine
$\mathcal{D}_c$ de $\mathcal{D}$.
\end{lemma}

\begin{proof}
Soit $(X, Y, Z, f, g, h)$ un triangle distingué de $\mathcal{D}$ tel que
$X$ et $Y$ soient compacts. Le Lemme \ref{lemma-representable-homological} et
le lemme des cinq (Homologie, Lemme \ref{homology-lemma-five-lemma}) montrent
alors que $Z$ est lui aussi un objet compact. Il est clair que, si
$X \oplus Y$ est compact, alors $X$, $Y$ sont eux aussi des objets compacts.
Ainsi, $\mathcal{D}_c$ est une sous-catégorie triangulée saturée.
Comme $\mathcal{D}$ est karoubienne d'après le
Lemme \ref{lemma-projectors-have-images-triangulated}, nous concluons qu'il en
est de même de $\mathcal{D}_c$.
\end{proof}

\begin{lemma}
\label{lemma-write-as-colimit}
Soit $\mathcal{D}$ une catégorie triangulée possédant des sommes directes.
Soit $E_i$, $i \in I$, une famille d'objets compacts de $\mathcal{D}$ telle
que $\bigoplus E_i$ engendre $\mathcal{D}$.
Alors tout objet $X$ de $\mathcal{D}$ peut s'écrire
$$
X = \text{hocolim} X_n
$$
où $X_1$ est une somme directe de translatés des $E_i$, et où chaque morphisme
de transition s'insère dans un triangle distingué
$Y_n \to X_n \to X_{n + 1} \to Y_n[1]$, où $Y_n$ est une somme directe de
translatés des $E_i$.
\end{lemma}

\begin{proof}
Posons $X_1 = \bigoplus_{(i, m, \varphi)} E_i[m]$, où la somme directe porte
sur tous les triplets $(i, m, \varphi)$ tels que $i \in I$,
$m \in \mathbf{Z}$ et $\varphi : E_i[m] \to X$. Alors $X_1$ est muni d'un
morphisme canonique $X_1 \to X$. Étant donné $X_n \to X$, posons
$Y_n = \bigoplus_{(i, m, \varphi)} E_i[m]$, où la somme directe porte sur tous
les triplets $(i, m, \varphi)$ tels que $i \in I$, $m \in \mathbf{Z}$ et que
$\varphi : E_i[m] \to X_n$ soit un morphisme tel que
$E_i[m] \to X_n \to X$ soit nul. Choisissons un triangle distingué
$Y_n \to X_n \to X_{n + 1} \to Y_n[1]$ et prenons pour
$X_{n + 1} \to X$ un morphisme quelconque tel que
$X_n \to X_{n + 1} \to X$ soit le morphisme donné; un tel morphisme existe par
notre choix de $Y_n$.
Nous obtenons un morphisme $\text{hocolim} X_n \to X$ par construction à
partir des morphismes $X_n \to X$. Choisissons un triangle distingué
$$
C \to \text{hocolim} X_n \to X \to C[1]
$$
Soit $E_i[m] \to C$ un morphisme. Comme $E_i$ est compact, le
Lemme \ref{lemma-commutes-with-countable-sums} implique que la composée
$E_i[m] \to C \to \text{hocolim} X_n$ se factorise par $X_n$ pour un certain
$n$, disons par $E_i[m] \to X_n$. La construction de $Y_n$ montre alors que
la composée $E_i[m] \to X_n \to X_{n + 1}$ est nulle. Autrement dit, la
composée $E_i[m] \to C \to \text{hocolim} X_n$ est nulle. Cela signifie que
notre morphisme $E_i[m] \to C$ provient d'un morphisme
$E_i[m] \to X[-1]$. La construction de $X_1$ montre alors qu'un tel morphisme
se relève à $\text{hocolim} X_n$, et nous concluons que notre morphisme
$E_i[m] \to C$ est nul. L'hypothèse selon laquelle $\bigoplus E_i$ engendre
$\mathcal{D}$ entraîne que $C$ est nul, ce qui achève la preuve.
\end{proof}

\begin{lemma}
\label{lemma-factor-through}
Sous les hypothèses et avec les notations du
Lemme \ref{lemma-write-as-colimit}, si $C$ est un objet compact et
$C \to X_n$ un morphisme, alors il existe une factorisation
$C \to E \to X_n$, où $E$ est un objet de
$\langle E_{i_1} \oplus \ldots \oplus E_{i_t} \rangle$ pour certains
$i_1, \ldots, i_t \in I$.
\end{lemma}

\begin{proof}
Raisonnons par récurrence sur $n$. Le cas initial $n = 1$ est clair.
Si $n > 1$, considérons la composée $C \to X_n \to Y_{n - 1}[1]$.
Elle se factorise par un morphisme $E'[1] \to Y_{n - 1}[1]$, où $E'$ est une
somme directe finie de translatés des $E_i$. Soit $I' \subset I$ l'ensemble
fini des indices qui interviennent dans cette somme directe. Nous obtenons
ainsi
$$
\xymatrix{
E' \ar[r] \ar[d] &
C' \ar[r] \ar[d] &
C \ar[r] \ar[d] &
E'[1] \ar[d] \\
Y_{n - 1} \ar[r] &
X_{n - 1} \ar[r] &
X_n \ar[r] &
Y_{n - 1}[1]
}
$$
Par récurrence, le morphisme $C' \to X_{n - 1}$ se factorise par
$E'' \to X_{n - 1}$, avec $E''$ objet de
$\langle \bigoplus_{i \in I''} E_i \rangle$ pour un certain sous-ensemble
fini $I'' \subset I$. Choisissons un triangle distingué
$$
E' \to E'' \to E \to E'[1]
$$
Alors $E$ est un objet de
$\langle \bigoplus_{i \in I' \cup I''} E_i \rangle$.
Par construction et d'après les axiomes d'une catégorie triangulée, nous
pouvons choisir des morphismes $C \to E$ et $E \to X_n$ qui s'insèrent dans
des morphismes de triangles
$(E', C', C) \to (E', E'', E)$ et
$(E', E'', E) \to (Y_{n - 1}, X_{n - 1}, X_n)$.
La composée $C \to E \to X_n$ peut ne pas être égale au morphisme donné
$C \to X_n$, mais leurs composées vers $Y_{n - 1}$ sont égales.
Soit $C \to X_{n - 1}$ un morphisme qui relève leur différence.
Par hypothèse de récurrence, nous pouvons le factoriser par un morphisme
$E''' \to X_{n - 1}$, où $E'''$ est un objet de
$\langle \bigoplus_{i \in I'''} E_i \rangle$ pour un certain sous-ensemble
fini $I''' \subset I$. Nous obtenons donc une solution en considérant
$E \oplus E''' \to X_n$, car $E \oplus E'''$ est un objet de
$\langle \bigoplus_{i \in I' \cup I'' \cup I'''} E_i \rangle$.
\end{proof}

\begin{definition}
\label{definition-compactly-generated}
Soit $\mathcal{D}$ une catégorie triangulée possédant des sommes directes
arbitraires. On dit que $\mathcal{D}$ est {\it compactement engendrée} s'il
existe un ensemble d'objets compacts $E_i$, $i \in I$, tel que
$\bigoplus E_i$ engendre $\mathcal{D}$.
\end{definition}

\noindent
La proposition suivante précise la relation entre générateurs classiques et
générateurs faibles.

\begin{proposition}
\label{proposition-generator-versus-classical-generator}
Soit $\mathcal{D}$ une catégorie triangulée possédant des sommes directes.
Soit $E$ un objet compact de $\mathcal{D}$.
Les conditions suivantes sont équivalentes :
\begin{enumerate}
\item $E$ est un générateur classique de $\mathcal{D}_c$ et $\mathcal{D}$ est
compactement engendrée, et
\item $E$ est un générateur de $\mathcal{D}$.
\end{enumerate}
\end{proposition}

\begin{proof}
Si $E$ est un générateur classique de $\mathcal{D}_c$, alors
$\mathcal{D}_c = \langle E \rangle$. L'hypothèse que $\mathcal{D}$ est
compactement engendrée et le Lemme \ref{lemma-right-orthogonal} entraînent
formellement que $E$ est un générateur de $\mathcal{D}$.

\medskip\noindent
La réciproque est plus intéressante. Supposons que $E$ soit un générateur de
$\mathcal{D}$. Soit $X$ un objet compact de $\mathcal{D}$.
Appliquons le Lemme \ref{lemma-write-as-colimit} avec $I = \{1\}$ et
$E_1 = E$ pour écrire
$$
X = \text{hocolim} X_n
$$
comme dans le lemme. Puisque $X$ est compact, nous trouvons que
$X \to \text{hocolim} X_n$ se factorise par $X_n$ pour un certain $n$
(Lemme \ref{lemma-commutes-with-countable-sums}).
Ainsi, $X$ est un facteur direct de $X_n$.
D'après le Lemme \ref{lemma-factor-through}, $X$ est un objet de
$\langle E \rangle$, ce qui prouve la proposition.
\end{proof}

\section{Représentabilité de Brown}
\label{section-brown}

\noindent
La référence pour les résultats de cette section est
\cite{Neeman-Grothendieck}.

\begin{lemma}
\label{lemma-brown}
\begin{reference}
\cite[Theorem 3.1]{Neeman-Grothendieck}.
\end{reference}
Soit $\mathcal{D}$ une catégorie triangulée possédant des sommes directes et
compactement engendrée. Soit $H : \mathcal{D} \to \textit{Ab}$ un foncteur
cohomologique contravariant qui transforme les sommes directes en produits.
Alors $H$ est représentable.
\end{lemma}

\begin{proof}
Soit $E_i$, $i \in I$, un ensemble d'objets compacts tel que
$\bigoplus_{i \in I} E_i$ engendre $\mathcal{D}$. Nous pouvons supposer, et
nous supposerons, que l'ensemble d'objets $\{E_i\}$ est stable par
translations. Considérons les paires $(i, a)$, où $i \in I$ et
$a \in H(E_i)$, et posons
$$
X_1 = \bigoplus\nolimits_{(i, a)} E_i
$$
Comme $H(X_1) = \prod_{(i, a)} H(E_i)$, nous voyons que
$(a)_{(i, a)}$ définit un élément $a_1 \in H(X_1)$.
Posons $H_1 = \Hom_\mathcal{D}(- , X_1)$.
D'après le lemme de Yoneda
(Catégories, Lemme \ref{categories-lemma-yoneda}), l'élément $a_1$ définit une
transformation naturelle $H_1 \to H$.

\medskip\noindent
Nous allons construire par récurrence des $X_n$ et des transformations
$a_n : H_n \to H$, où $H_n = \Hom_\mathcal{D}(-, X_n)$.
Plus précisément, appliquons le procédé précédent au foncteur
$\Ker(H_n \to H)$ pour obtenir un objet
$$
K_{n + 1} = \bigoplus\nolimits_{(i, k),\ k \in \Ker(H_n(E_i) \to H(E_i))} E_i
$$
et une transformation
$\Hom_\mathcal{D}(-, K_{n + 1}) \to \Ker(H_n \to H)$.
D'après le lemme de Yoneda, la composée
$\Hom_\mathcal{D}(-, K_{n + 1}) \to H_n$ donne un morphisme
$K_{n + 1} \to X_n$. Choisissons un triangle distingué
$$
K_{n + 1} \to X_n \to X_{n + 1} \to K_{n + 1}[1]
$$
dans $\mathcal{D}$. Par construction, l'élément $a_n \in H(X_n)$ a une image
nulle dans $H(K_{n + 1})$. Comme $H$ est cohomologique, nous pouvons le relever
en un élément $a_{n + 1} \in H(X_{n + 1})$.

\medskip\noindent
Nous affirmons que $X = \text{hocolim} X_n$ représente $H$.
En appliquant $H$ au triangle distingué qui définit cette colimite,
$$
\bigoplus X_n \to
\bigoplus X_n \to X \to
\bigoplus X_n[1]
$$
nous obtenons une suite exacte
$$
\prod H(X_n) \leftarrow 
\prod H(X_n) \leftarrow 
H(X)
$$
Il existe donc un élément $a \in H(X)$ qui s'envoie sur $(a_n)$ dans
$\prod H(X_n)$. Nous obtenons ainsi une transformation naturelle
$\Hom_\mathcal{D}(- , X) \to H$ telle que le diagramme
$$
\Hom_\mathcal{D}(-, X_1) \to
\Hom_\mathcal{D}(-, X_2) \to
\Hom_\mathcal{D}(-, X_3) \to \ldots \to
\Hom_\mathcal{D}(-, X) \to H
$$
commute. Pour chaque $i$, l'application
$\Hom_\mathcal{D}(E_i, X) \to H(E_i)$ est surjective, par construction de
$X_1$. D'autre part, par construction de $X_n \to X_{n + 1}$, le noyau de
$\Hom_\mathcal{D}(E_i, X_n) \to H(E_i)$ est annulé par l'application
$\Hom_\mathcal{D}(E_i, X_n) \to \Hom_\mathcal{D}(E_i, X_{n + 1})$.
Comme
$$
\Hom_\mathcal{D}(E_i, X) = \colim \Hom_\mathcal{D}(E_i, X_n)
$$
d'après le Lemme \ref{lemma-commutes-with-countable-sums}, nous voyons que
$\Hom_\mathcal{D}(E_i, X) \to H(E_i)$ est injective.

\medskip\noindent
Pour achever la preuve, considérons la sous-catégorie
$$
\mathcal{D}' =
\{Y \in \Ob(\mathcal{D}) \mid \Hom_\mathcal{D}(Y[n], X) \to H(Y[n])
\text{ est un isomorphisme pour tout }n\}
$$
Comme $\Hom_\mathcal{D}(-, X) \to H$ est une transformation entre foncteurs
cohomologiques, la sous-catégorie $\mathcal{D}'$ est une sous-catégorie
triangulée, saturée et strictement pleine de $\mathcal{D}$ (détails omis;
voir la preuve du Lemme \ref{lemma-homological-functor-kernel}).
De plus, comme $H$ et $\Hom_\mathcal{D}(-, X)$ transforment tous deux les
sommes directes en produits, les sommes directes d'objets de $\mathcal{D}'$
appartiennent à $\mathcal{D}'$. Ainsi, les colimites dérivées d'objets de
$\mathcal{D}'$ appartiennent à $\mathcal{D}'$.
Puisque $\{E_i\}$ est stable par translations, nous voyons que $E_i$ est un
objet de $\mathcal{D}'$ pour tout $i$. Le
Lemme \ref{lemma-write-as-colimit} entraîne que
$\mathcal{D}' = \mathcal{D}$, ce qui achève la preuve.
\end{proof}

\begin{proposition}
\label{proposition-brown}
\begin{reference}
\cite[Theorem 4.1]{Neeman-Grothendieck}.
\end{reference}
Soit $\mathcal{D}$ une catégorie triangulée possédant des sommes directes et
compactement engendrée. Soit $F : \mathcal{D} \to \mathcal{D}'$ un foncteur
exact entre catégories triangulées qui transforme les sommes directes en
sommes directes. Alors $F$ possède un adjoint à droite exact.
\end{proposition}

\begin{proof}
Pour un objet $Y$ de $\mathcal{D}'$, considérons le foncteur contravariant
$$
\mathcal{D} \to \textit{Ab},\quad W \mapsto \Hom_{\mathcal{D}'}(F(W), Y)
$$
C'est un foncteur cohomologique puisque $F$ est exact, et il transforme les
sommes directes en produits puisque $F$ transforme les sommes directes en
sommes directes. Le Lemme \ref{lemma-brown} fournit donc un objet $X$ de
$\mathcal{D}$ tel que
$\Hom_\mathcal{D}(W, X) = \Hom_{\mathcal{D}'}(F(W), Y)$.
L'existence de l'adjoint résulte de
Catégories, Lemme \ref{categories-lemma-adjoint-exists}.
Son exactitude résulte du Lemme \ref{lemma-adjoint-is-exact}.
\end{proof}

\section{Représentabilité de Brown, bis}
\label{section-brown-bis}

\noindent
Dans cette section, nous exposons une version de la représentabilité de Brown
pour les catégories triangulées qui possèdent un ensemble convenable de
générateurs; pour d'autres versions, voir \cite{Franke}, \cite{Neeman} et
\cite{Krause}.

\begin{lemma}
\label{lemma-brown-bis}
\begin{reference}
Version faible de \cite[Theorem A]{Krause}
\end{reference}
Soit $\mathcal{D}$ une catégorie triangulée possédant des sommes directes.
Supposons donné un ensemble $\mathcal{E}$ d'objets de $\mathcal{D}$ tel que
\begin{enumerate}
\item si $X$ est un objet non nul de $\mathcal{D}$, alors il existe un
$E \in \mathcal{E}$ et un morphisme non nul $E \to X$, et
\item étant donnés des objets $X_n$, $n \in \mathbf{N}$, de $\mathcal{D}$,
un objet $E \in \mathcal{E}$ et un morphisme
$\alpha : E \to \bigoplus X_n$, il existe des $E_n \in \mathcal{E}$, des
$\beta_n : E_n \to X_n$ et un morphisme
$\gamma : E \to \bigoplus E_n$ tels que
$\alpha = (\bigoplus \beta_n) \circ \gamma$.
\end{enumerate}
Soit $H : \mathcal{D} \to \textit{Ab}$ un foncteur cohomologique contravariant
qui transforme les sommes directes en produits. Alors $H$ est représentable.
\end{lemma}

\begin{proof}
Cette preuve est très semblable à celle du Lemme \ref{lemma-brown}.
Nous pouvons remplacer $\mathcal{E}$ par
$\bigcup_{i \in \mathbf{Z}} \mathcal{E}[i]$ et supposer que $\mathcal{E}$ est
stable par translations.
Considérons les paires $(E, a)$, où $E \in \mathcal{E}$ et $a \in H(E)$, et
posons
$$
X_1 = \bigoplus\nolimits_{(E, a)} E
$$
Comme $H(X_1) = \prod_{(E, a)} H(E)$, nous voyons que $(a)_{(E, a)}$ définit
un élément $a_1 \in H(X_1)$. Posons
$H_1 = \Hom_\mathcal{D}(- , X_1)$.
D'après le lemme de Yoneda
(Catégories, Lemme \ref{categories-lemma-yoneda}), l'élément $a_1$ définit une
transformation naturelle $H_1 \to H$.

\medskip\noindent
Nous allons construire par récurrence des $X_n$ et des transformations
$a_n : H_n \to H$, où $H_n = \Hom_\mathcal{D}(-, X_n)$.
Plus précisément, appliquons le procédé précédent au foncteur
$\Ker(H_n \to H)$ pour obtenir un objet
$$
K_{n + 1} = \bigoplus\nolimits_{(E, k),\ k \in \Ker(H_n(E) \to H(E))} E
$$
et une transformation
$\Hom_\mathcal{D}(-, K_{n + 1}) \to \Ker(H_n \to H)$.
D'après le lemme de Yoneda, la composée
$\Hom_\mathcal{D}(-, K_{n + 1}) \to H_n$ donne un morphisme
$K_{n + 1} \to X_n$. Choisissons un triangle distingué
$$
K_{n + 1} \to X_n \to X_{n + 1} \to K_{n + 1}[1]
$$
dans $\mathcal{D}$. Par construction, l'élément $a_n \in H(X_n)$ a une image
nulle dans $H(K_{n + 1})$. Comme $H$ est cohomologique, nous pouvons le relever
en un élément $a_{n + 1} \in H(X_{n + 1})$.

\medskip\noindent
Posons $X = \text{hocolim} X_n$. En appliquant $H$ au triangle distingué qui
le définit,
$$
\bigoplus X_n \to
\bigoplus X_n \to X \to
\bigoplus X_n[1]
$$
nous obtenons une suite exacte
$$
\prod H(X_n) \leftarrow
\prod H(X_n) \leftarrow
H(X)
$$
Il existe donc un élément $a \in H(X)$ qui s'envoie sur $(a_n)$ dans
$\prod H(X_n)$. Nous obtenons ainsi une transformation naturelle
$\Hom_\mathcal{D}(- , X) \to H$ telle que le diagramme
$$
\Hom_\mathcal{D}(-, X_1) \to
\Hom_\mathcal{D}(-, X_2) \to
\Hom_\mathcal{D}(-, X_3) \to \ldots \to
\Hom_\mathcal{D}(-, X) \to H
$$
commute. Nous affirmons que $\Hom_\mathcal{D}(-, X) \to H(-)$ est un
isomorphisme.

\medskip\noindent
Soit $E \in \mathcal{E}$. Montrons que
$$
\Hom_\mathcal{D}(E, \bigoplus X_n) \to  \Hom_\mathcal{D}(E, \bigoplus X_n)
$$
est injective. Soit donc $\alpha : E \to \bigoplus X_n$. L'hypothèse (2)
fournit une factorisation
$\alpha = (\bigoplus \beta_n) \circ \gamma$.
Comme $E_n \to X_n \to X_{n + 1}$ est nul par construction, la composée
$\bigoplus E_n \to \bigoplus X_n \to \bigoplus X_n$ est égale à
$\bigoplus \beta_n$. Par conséquent, la composée
$E \to \bigoplus X_n \to \bigoplus X_n$ est elle aussi égale à $\alpha$.
Cela prouve l'injectivité annoncée, et donc également l'injectivité de
$$
\Hom_\mathcal{D}(E, \bigoplus X_n[1]) \to \Hom_\mathcal{D}(E, \bigoplus X_n[1])
$$
Il s'ensuit que nous avons une suite exacte
$$
\Hom_\mathcal{D}(E, \bigoplus X_n) \to
\Hom_\mathcal{D}(E, \bigoplus X_n) \to
\Hom_\mathcal{D}(E, X) \to 0
$$
pour tout $E \in \mathcal{E}$.

\medskip\noindent
Soit $E \in \mathcal{E}$ et soit $f : E \to X$ un morphisme.
D'après le paragraphe précédent, nous pouvons choisir un
$\alpha : E \to \bigoplus X_n$ qui relève $f$. L'hypothèse (2) fournit alors
une factorisation $\alpha = (\bigoplus \beta_n) \circ \gamma$.
Pour chaque $n$, il existe un morphisme $\delta_n : E_n \to X_1$ tel que
$\delta_n$ et $\beta_n$ aient la même image dans $H(E_n)$. Alors les composées
$$
E_n \to X_n \to X_{n + 1}
\quad\text{et}\quad
E_n \to X_1 \to X_{n + 1}
$$
sont égales par construction de $X_n \to X_{n + 1}$. Il s'ensuit que
$$
\bigoplus E_n \to \bigoplus X_n \to X
\quad\text{et}\quad
\bigoplus E_n \to \bigoplus X_1 \to X
$$
sont elles aussi égales. En remarquant que $\bigoplus X_1 \to X$ se factorise
en $\bigoplus X_1 \to X_1 \to X$, nous concluons que
$$
\Hom_\mathcal{D}(E, X_1) \to \Hom_\mathcal{D}(E, X)
$$
est surjective. Comme, par construction, l'application
$\Hom_\mathcal{D}(E, X_1) \to H(E)$ est surjective et que le noyau de cette
application est annulé par
$\Hom_\mathcal{D}(E, X_1) \to \Hom_\mathcal{D}(E, X)$, nous concluons que
$\Hom_\mathcal{D}(E, X) \to H(E)$ est une bijection pour tout
$E \in \mathcal{E}$.

\medskip\noindent
Pour achever la preuve, considérons la sous-catégorie
$$
\mathcal{D}' =
\{Y \in \Ob(\mathcal{D}) \mid \Hom_\mathcal{D}(Y[n], X) \to H(Y[n])
\text{ est un isomorphisme pour tout }n\}
$$
Comme $\Hom_\mathcal{D}(-, X) \to H$ est une transformation entre foncteurs
cohomologiques, la sous-catégorie $\mathcal{D}'$ est une sous-catégorie
triangulée, saturée et strictement pleine de $\mathcal{D}$ (détails omis;
voir la preuve du Lemme \ref{lemma-homological-functor-kernel}).
De plus, comme $H$ et $\Hom_\mathcal{D}(-, X)$ transforment tous deux les
sommes directes en produits, les sommes directes d'objets de $\mathcal{D}'$
appartiennent à $\mathcal{D}'$. Ainsi, les colimites dérivées d'objets de
$\mathcal{D}'$ appartiennent à $\mathcal{D}'$.
Puisque $\mathcal{E}$ est stable par translations, le résultat du paragraphe
précédent donne $\mathcal{E} \subset \Ob(\mathcal{D}')$.
Pour achever la preuve, il reste à montrer que
$\mathcal{D}' = \mathcal{D}$.

\medskip\noindent
Soit $Y$ un objet de $\mathcal{D}$ et posons
$H(-) = \Hom_\mathcal{D}(-, Y)$. Alors $H$ est un foncteur cohomologique qui
transforme les sommes directes en produits.
La construction de la première partie de la preuve fournit un morphisme
$\text{hocolim} X_n = X \to Y$ tel que
$\Hom_\mathcal{D}(E, X) \to \Hom_\mathcal{D}(E, Y)$ soit bijectif pour tout
$E \in \mathcal{E}$. L'hypothèse (1) nous dit alors que $X \to Y$ est un
isomorphisme! D'autre part, par construction, $X_1, X_2, \ldots$ appartiennent
à $\mathcal{D}'$, et il en est donc de même de $X$. Ainsi,
$Y \in \mathcal{D}'$, ce qui achève la preuve.
\end{proof}

\begin{proposition}
\label{proposition-brown-bis}
Soit $\mathcal{D}$ une catégorie triangulée possédant des sommes directes.
Supposons qu'il existe un ensemble $\mathcal{E}$ d'objets de $\mathcal{D}$
satisfaisant aux conditions (1) et (2) du Lemme \ref{lemma-brown-bis}.
Soit $F : \mathcal{D} \to \mathcal{D}'$ un foncteur exact entre catégories
triangulées qui transforme les sommes directes en sommes directes.
Alors $F$ possède un adjoint à droite exact.
\end{proposition}

\begin{proof}
Pour un objet $Y$ de $\mathcal{D}'$, considérons le foncteur contravariant
$$
\mathcal{D} \to \textit{Ab},\quad W \mapsto \Hom_{\mathcal{D}'}(F(W), Y)
$$
C'est un foncteur cohomologique puisque $F$ est exact, et il transforme les
sommes directes en produits puisque $F$ transforme les sommes directes en
sommes directes. Le Lemme \ref{lemma-brown-bis} fournit donc un objet $X$ de
$\mathcal{D}$ tel que
$\Hom_\mathcal{D}(W, X) = \Hom_{\mathcal{D}'}(F(W), Y)$.
L'existence de l'adjoint résulte de
Catégories, Lemme \ref{categories-lemma-adjoint-exists}.
Son exactitude résulte du Lemme \ref{lemma-adjoint-is-exact}.
\end{proof}

\section{Sous-catégories admissibles}
\label{section-admissible}

\noindent
La référence pour cette section est \cite[section 1]{Bondal-Kapranov}.

\begin{definition}
\label{definition-orthogonal}
Soit $\mathcal{D}$ une catégorie additive. Soit
$\mathcal{A} \subset \mathcal{D}$ une sous-catégorie pleine.
L'{\it orthogonal à droite} $\mathcal{A}^\perp$ de $\mathcal{A}$ est la
sous-catégorie pleine formée des objets $X$ de $\mathcal{D}$ tels que
$\Hom(A, X) = 0$ pour tout $A \in \Ob(\mathcal{A})$.
L'{\it orthogonal à gauche} ${}^\perp\mathcal{A}$ de $\mathcal{A}$ est la
sous-catégorie pleine formée des objets $X$ de $\mathcal{D}$ tels que
$\Hom(X, A) = 0$ pour tout $A \in \Ob(\mathcal{A})$.
\end{definition}

\begin{lemma}
\label{lemma-pre-prepare-adjoint}
Soit $\mathcal{D}$ une catégorie triangulée.
Soit $\mathcal{A} \subset \mathcal{D}$ une sous-catégorie pleine stable par
toutes les translations. Considérons un triangle distingué
$$
X \to Y \to Z \to X[1]
$$
de $\mathcal{D}$. Les conditions suivantes sont équivalentes :
\begin{enumerate}
\item $Z$ appartient à $\mathcal{A}^\perp$, et
\item $\Hom(A, X) = \Hom(A, Y)$ pour tout $A \in \Ob(\mathcal{A})$.
\end{enumerate}
\end{lemma}

\begin{proof}
D'après le Lemme \ref{lemma-representable-homological}, le foncteur
$\Hom(A, -)$ est homologique; nous obtenons donc une suite exacte longue comme
dans (\ref{equation-long-exact-cohomology-sequence}).
Supposons (1) et soit $A \in \Ob(\mathcal{A})$.
Considérons alors la suite exacte
$$
\Hom(A[1], Z) \to \Hom(A, X) \to \Hom(A, Y) \to \Hom(A, Z)
$$
Comme $A[1] \in \Ob(\mathcal{A})$, le premier et le dernier groupe sont nuls.
Nous obtenons donc (2). Supposons (2) et soit
$A \in \Ob(\mathcal{A})$. Considérons alors la suite exacte
$$
\Hom(A, X) \to \Hom(A, Y) \to \Hom(A, Z) \to \Hom(A[-1], X) \to \Hom(A[-1], Y)
$$
et concluons que $\Hom(A, Z) = 0$, comme voulu.
\end{proof}

\begin{lemma}
\label{lemma-pre-prepare-adjoint-dual}
Soit $\mathcal{D}$ une catégorie triangulée.
Soit $\mathcal{B} \subset \mathcal{D}$ une sous-catégorie pleine stable par
toutes les translations. Considérons un triangle distingué
$$
X \to Y \to Z \to X[1]
$$
de $\mathcal{D}$. Les conditions suivantes sont équivalentes :
\begin{enumerate}
\item $X$ appartient à ${}^\perp\mathcal{B}$, et
\item $\Hom(Y, B) = \Hom(Z, B)$ pour tout $B \in \Ob(\mathcal{B})$.
\end{enumerate}
\end{lemma}

\begin{proof}
C'est le dual du Lemme \ref{lemma-pre-prepare-adjoint}.
\end{proof}

\begin{lemma}
\label{lemma-orthogonal-triangulated}
Soit $\mathcal{D}$ une catégorie triangulée. Soit
$\mathcal{A} \subset \mathcal{D}$ une sous-catégorie pleine stable par toutes
les translations. Alors l'orthogonal à droite $\mathcal{A}^\perp$ et
l'orthogonal à gauche ${}^\perp\mathcal{A}$ de $\mathcal{A}$ sont des
sous-catégories triangulées, saturées\footnote{Définition
\ref{definition-saturated}.} et strictement pleines de $\mathcal{D}$.
\end{lemma}

\begin{proof}
Il résulte immédiatement des définitions que les orthogonaux sont stables par
translations, sommes directes et facteurs directs.
Considérons un triangle distingué
$$
X \to Y \to Z \to X[1]
$$
de $\mathcal{D}$. D'après le Lemme \ref{lemma-triangulated-subcategory}, il
suffit de montrer que, si $X$ et $Y$ appartiennent à $\mathcal{A}^\perp$,
alors $Z$ appartient à $\mathcal{A}^\perp$. Cela résulte immédiatement du
Lemme \ref{lemma-pre-prepare-adjoint}.
\end{proof}

\begin{lemma}
\label{lemma-prepare-adjoint}
Soit $\mathcal{D}$ une catégorie triangulée. Soit $\mathcal{A}$ une
sous-catégorie triangulée pleine de $\mathcal{D}$. Pour un objet $X$ de
$\mathcal{D}$, considérons la propriété $P(X)$ : il existe un triangle
distingué $A \to X \to B \to A[1]$ dans $\mathcal{D}$, avec $A$ dans
$\mathcal{A}$ et $B$ dans $\mathcal{A}^\perp$.
\begin{enumerate}
\item Si $X_1 \to X_2 \to X_3 \to X_1[1]$ est un triangle distingué et que
$P$ est vraie pour deux de ces trois objets, alors elle est vraie pour le
troisième.
\item Si $P$ est vraie pour $X_1$ et $X_2$, alors elle est vraie pour
$X_1 \oplus X_2$.
\end{enumerate}
\end{lemma}

\begin{proof}
Soit $X_1 \to X_2 \to X_3 \to X_1[1]$ un triangle distingué et supposons que
$P$ soit vraie pour $X_1$ et $X_2$. Choisissons des triangles distingués
$$
A_1 \to X_1 \to B_1 \to A_1[1]
\quad\text{et}\quad
A_2 \to X_2 \to B_2 \to A_2[1]
$$
comme dans la condition $P$. Puisque
$\Hom(A_1, A_2) = \Hom(A_1, X_2)$ d'après le
Lemme \ref{lemma-pre-prepare-adjoint}, il existe un unique morphisme
$A_1 \to A_2$ tel que le diagramme
$$
\xymatrix{
A_1 \ar[d] \ar[r] & X_1 \ar[d] \\
A_2 \ar[r] & X_2
}
$$
commute. Prolongeons-le en un diagramme
$$
\xymatrix{
A_1 \ar[r] \ar[d] & X_1 \ar[r] \ar[d] & Q_1 \ar[r] \ar[d] & A_1[1] \ar[d] \\
A_2 \ar[r] \ar[d] & X_2 \ar[r] \ar[d] & Q_2 \ar[r] \ar[d] & A_2[1] \ar[d] \\
A_3 \ar[r] \ar[d] & X_3 \ar[r] \ar[d] & Q_3 \ar[r] \ar[d] & A_3[1] \ar[d] \\
A_1[1] \ar[r] & X_1[1] \ar[r] & Q_1[1] \ar[r] & A_1[2]
}
$$
comme dans la Proposition \ref{proposition-9}. D'après TR3, nous avons
$Q_1 \cong B_1$ et $Q_2 \cong B_2$, donc
$Q_1, Q_2 \in \Ob(\mathcal{A}^\perp)$.
Comme $Q_1 \to Q_2 \to Q_3 \to Q_1[1]$ est un triangle distingué, nous avons
$Q_3 \in \Ob(\mathcal{A}^\perp)$ d'après le
Lemme \ref{lemma-orthogonal-triangulated}.
Puisque $\mathcal{A}$ est une sous-catégorie triangulée pleine, nous voyons
que $A_3$ est isomorphe à un objet de $\mathcal{A}$.
Ainsi, $X_3$ satisfait $P$. Les autres cas de (1) s'en déduisent par
translation. L'assertion (2) est un cas particulier de (1), par le
Lemme \ref{lemma-split}.
\end{proof}

\begin{lemma}
\label{lemma-prepare-adjoint-dual}
Soit $\mathcal{D}$ une catégorie triangulée. Soit $\mathcal{B}$ une
sous-catégorie triangulée pleine de $\mathcal{D}$. Pour un objet $X$ de
$\mathcal{D}$, considérons la propriété $P(X)$ : il existe un triangle
distingué $A \to X \to B \to A[1]$ dans $\mathcal{D}$, avec $B$ dans
$\mathcal{B}$ et $A$ dans ${}^\perp\mathcal{B}$.
\begin{enumerate}
\item Si $X_1 \to X_2 \to X_3 \to X_1[1]$ est un triangle distingué et que
$P$ est vraie pour deux de ces trois objets, alors elle est vraie pour le
troisième.
\item Si $P$ est vraie pour $X_1$ et $X_2$, alors elle est vraie pour
$X_1 \oplus X_2$.
\end{enumerate}
\end{lemma}

\begin{proof}
C'est le dual du Lemme \ref{lemma-prepare-adjoint}.
\end{proof}

\begin{lemma}
\label{lemma-right-adjoint}
Soit $\mathcal{D}$ une catégorie triangulée. Soit
$\mathcal{A} \subset \mathcal{D}$ une sous-catégorie triangulée pleine.
Les conditions suivantes sont équivalentes :
\begin{enumerate}
\item le foncteur d'inclusion $\mathcal{A} \to \mathcal{D}$ possède un
adjoint à droite, et
\item pour tout $X$ dans $\mathcal{D}$, il existe un triangle distingué
$$
A \to X \to B \to A[1]
$$
dans $\mathcal{D}$, avec $A \in \Ob(\mathcal{A})$ et
$B \in \Ob(\mathcal{A}^\perp)$.
\end{enumerate}
Si ces conditions sont vérifiées, alors $\mathcal{A}$ est saturée
(Définition \ref{definition-saturated}); et si $\mathcal{A}$ est strictement
pleine dans $\mathcal{D}$, alors
$\mathcal{A} = {}^\perp(\mathcal{A}^\perp)$.
\end{lemma}

\begin{proof}
Supposons (1) et notons $v : \mathcal{D} \to \mathcal{A}$ l'adjoint à droite.
Soit $X \in \Ob(\mathcal{D})$. Posons $A = v(X)$. Nous pouvons prolonger le
morphisme d'adjonction $A \to X$ en un triangle distingué
$A \to X \to B \to A[1]$. Puisque
$$
\Hom_\mathcal{A}(A', A) =
\Hom_\mathcal{A}(A', v(X)) =
\Hom_\mathcal{D}(A', X)
$$
pour $A' \in \Ob(\mathcal{A})$, nous concluons que
$B \in \Ob(\mathcal{A}^\perp)$ d'après le
Lemme \ref{lemma-pre-prepare-adjoint}.

\medskip\noindent
Supposons (2). Construisons explicitement l'adjoint $v$.
Soit $X \in \Ob(\mathcal{D})$. Choisissons $A \to X \to B \to A[1]$ comme
en (2). Posons $v(X) = A$. Soit $f : X \to Y$ un morphisme de $\mathcal{D}$.
Choisissons $A' \to Y \to B' \to A'[1]$ comme en (2). Puisque
$\Hom(A, A') = \Hom(A, Y)$ d'après le
Lemme \ref{lemma-pre-prepare-adjoint}, il existe un unique morphisme
$f' : A \to A'$ tel que le diagramme
$$
\xymatrix{
A \ar[d]_{f'} \ar[r] & X \ar[d]^f \\
A' \ar[r] & Y
}
$$
commute. Nous pouvons donc poser $v(f) = f'$ et obtenir ainsi un foncteur.
Pour voir que $v$ est adjoint au foncteur d'inclusion, appliquons de nouveau le
Lemme \ref{lemma-pre-prepare-adjoint}.

\medskip\noindent
Démontrons la dernière assertion. Pour prouver que $\mathcal{A}$ est saturée,
nous pouvons remplacer $\mathcal{A}$ par la sous-catégorie strictement pleine
ayant les mêmes classes d'isomorphie que $\mathcal{A}$; détails omis.
Supposons $\mathcal{A}$
strictement pleine. Si nous montrons que
$\mathcal{A} = {}^\perp(\mathcal{A}^\perp)$, alors $\mathcal{A}$ sera saturée
d'après le Lemme \ref{lemma-orthogonal-triangulated}.
Puisque l'inclusion
$\mathcal{A} \subset {}^\perp(\mathcal{A}^\perp)$ est claire, il suffit de
démontrer l'inclusion opposée.
Soit $X$ un objet de ${}^\perp(\mathcal{A}^\perp)$.
Choisissons un triangle distingué $A \to X \to B \to A[1]$ comme en (2).
Comme $\Hom(X, B) = 0$ par hypothèse, nous avons
$A \cong X \oplus B[-1]$ d'après le Lemme \ref{lemma-split}.
Puisque $\Hom(A, B[-1]) = 0$, car $B \in \mathcal{A}^\perp$, cela entraîne
$B[-1] = 0$ et $A \cong X$, comme voulu.
\end{proof}

\begin{lemma}
\label{lemma-left-adjoint}
Soit $\mathcal{D}$ une catégorie triangulée. Soit
$\mathcal{B} \subset \mathcal{D}$ une sous-catégorie triangulée pleine.
Les conditions suivantes sont équivalentes :
\begin{enumerate}
\item le foncteur d'inclusion $\mathcal{B} \to \mathcal{D}$ possède un
adjoint à gauche, et
\item pour tout $X$ dans $\mathcal{D}$, il existe un triangle distingué
$$
A \to X \to B \to A[1]
$$
dans $\mathcal{D}$, avec $B \in \Ob(\mathcal{B})$ et
$A \in \Ob({}^\perp\mathcal{B})$.
\end{enumerate}
Si ces conditions sont vérifiées, alors $\mathcal{B}$ est saturée
(Définition \ref{definition-saturated}); et si $\mathcal{B}$ est strictement
pleine dans $\mathcal{D}$, alors
$\mathcal{B} = ({}^\perp\mathcal{B})^\perp$.
\end{lemma}

\begin{proof}
C'est le dual du Lemme \ref{lemma-right-adjoint}.
\end{proof}

\begin{definition}
\label{definition-admissible}
Soit $\mathcal{D}$ une catégorie triangulée. Une sous-catégorie
{\it admissible à droite} de $\mathcal{D}$ est une sous-catégorie triangulée
strictement pleine qui satisfait aux conditions équivalentes du
Lemme \ref{lemma-right-adjoint}.
Une sous-catégorie {\it admissible à gauche} de $\mathcal{D}$ est une
sous-catégorie triangulée strictement pleine qui satisfait aux conditions
équivalentes du Lemme \ref{lemma-left-adjoint}.
Une sous-catégorie {\it admissible des deux côtés} est admissible à droite et
à gauche.
\end{definition}

\noindent
Soit $\mathcal{A}$ une sous-catégorie admissible à droite de la catégorie
triangulée $\mathcal{D}$. Remarquons que, pour $X \in \mathcal{D}$, le
triangle distingué
$$
A \to X \to B \to A[1]
$$
avec $A \in \mathcal{A}$ et $B \in \mathcal{A}^\perp$ est canonique au sens
suivant : pour tout autre triangle distingué
$A' \to X \to B' \to A'[1]$, avec $A' \in \mathcal{A}$ et
$B' \in \mathcal{A}^\perp$, il existe un isomorphisme
$(\alpha, \text{id}_X, \beta) : (A, X, B) \to (A', X, B')$ de triangles.
La proposition suivante résume ce qui précède.

\begin{proposition}
\label{proposition-summarize-admissible}
Soit $\mathcal{D}$ une catégorie triangulée. Soient
$\mathcal{A} \subset \mathcal{D}$ et $\mathcal{B} \subset \mathcal{D}$ des
sous-catégories. Les conditions suivantes sont équivalentes :
\begin{enumerate}
\item $\mathcal{A}$ est admissible à droite et
$\mathcal{B} = \mathcal{A}^\perp$,
\item $\mathcal{B}$ est admissible à gauche et
$\mathcal{A} = {}^\perp\mathcal{B}$,
\item $\Hom(A, B) = 0$ pour tous $A \in \mathcal{A}$ et
$B \in \mathcal{B}$, et, pour tout $X$ dans $\mathcal{D}$, il existe un
triangle distingué $A \to X \to B \to A[1]$ dans $\mathcal{D}$, avec
$A \in \mathcal{A}$ et $B \in \mathcal{B}$.
\end{enumerate}
Si ces conditions sont vérifiées, alors
$\mathcal{A} \to \mathcal{D}/\mathcal{B}$ et
$\mathcal{B} \to \mathcal{D}/\mathcal{A}$ sont des équivalences de catégories
triangulées; l'adjoint à droite du foncteur d'inclusion
$\mathcal{A} \to \mathcal{D}$ est
$\mathcal{D} \to \mathcal{D}/\mathcal{B} \to \mathcal{A}$, et l'adjoint à
gauche du foncteur d'inclusion $\mathcal{B} \to \mathcal{D}$ est
$\mathcal{D} \to \mathcal{D}/\mathcal{A} \to \mathcal{B}$.
\end{proposition}

\begin{proof}
L'équivalence entre (1), (2) et (3) résulte immédiatement des Lemmes
\ref{lemma-right-adjoint} et \ref{lemma-left-adjoint} (un petit détail est
omis). Notons $v : \mathcal{D} \to \mathcal{A}$ l'adjoint à droite du foncteur
d'inclusion $i : \mathcal{A} \to \mathcal{D}$.
Il est immédiat que $\Ker(v) = \mathcal{A}^\perp = \mathcal{B}$.
La propriété universelle du quotient permet donc de factoriser $v$ en un
foncteur $\overline{v} : \mathcal{D}/\mathcal{B} \to \mathcal{A}$.
Comme $v \circ i = \text{id}_\mathcal{A}$ d'après
Catégories, Lemme \ref{categories-lemma-adjoint-fully-faithful}, nous voyons
que $\overline{v}$ est un quasi-inverse à gauche de
$\overline{i} : \mathcal{A} \to \mathcal{D}/\mathcal{B}$.
Nous affirmons que la composée $\overline{i} \circ \overline{v}$ est elle
aussi isomorphe à $\text{id}_{\mathcal{D}/\mathcal{B}}$.
En effet, supposons que $X$ s'insère dans un triangle distingué
$A \to X \to B \to A[1]$ comme en (3). Alors $v(X) = A$, comme nous l'avons vu
dans la preuve du Lemme \ref{lemma-right-adjoint}.
En considérant $X$ comme un objet de $\mathcal{D}/\mathcal{B}$, nous avons
$\overline{i}(\overline{v}(X)) = A$, et il existe un isomorphisme fonctoriel
$\overline{i}(\overline{v}(X)) = A \to X$ dans
$\mathcal{D}/\mathcal{B}$. Nous trouvons donc bien que
$\overline{v} : \mathcal{D}/\mathcal{B} \to \mathcal{A}$ est une équivalence.
L'assertion selon laquelle $\mathcal{B} \to \mathcal{D}/\mathcal{A}$ est une
équivalence et l'adjoint à gauche du foncteur d'inclusion
$\mathcal{B} \to \mathcal{D}$ est
$\mathcal{D} \to \mathcal{D}/\mathcal{A} \to \mathcal{B}$ est duale de ce
que nous venons de démontrer.
\end{proof}

\section{Systèmes de Postnikov}
\label{section-postnikov}

\noindent
Une référence pour cette section est \cite{Orlov-K3}. Soit $\mathcal{D}$
une catégorie triangulée. Soit
$$
X_n \to X_{n - 1} \to \ldots \to X_0
$$
un complexe de $\mathcal{D}$. Dans cette section, nous étudions le problème
de construire une « totalisation » de ce complexe.

\begin{definition}
\label{definition-postnikov-system}
Soit $\mathcal{D}$ une catégorie triangulée. Soit
$$
X_n \to X_{n - 1} \to \ldots \to X_0
$$
un complexe de $\mathcal{D}$. Un {\it système de Postnikov} est défini
par récurrence comme suit.
\begin{enumerate}
\item Si $n = 0$, c'est un isomorphisme $Y_0 \to X_0$.
\item Si $n = 1$, c'est le choix d'un isomorphisme $Y_0 \to X_0$ et
le choix d'un triangle distingué
$$
Y_1 \to X_1 \to Y_0 \to Y_1[1]
$$
tel que le composé de $X_1 \to Y_0$ avec $Y_0 \to X_0$ soit le morphisme
donné $X_1 \to X_0$.
\item Si $n > 1$, c'est le choix d'un système de Postnikov
pour $X_{n - 1} \to \ldots \to X_0$ et le choix d'un triangle
distingué
$$
Y_n \to X_n \to Y_{n - 1} \to Y_n[1]
$$
tel que le composé du morphisme $X_n \to Y_{n - 1}$ avec
$Y_{n - 1} \to X_{n - 1}$ soit le morphisme donné
$X_n \to X_{n - 1}$.
\end{enumerate}
Étant donné un morphisme
\begin{equation}
\label{equation-map-complexes}
\vcenter{
\xymatrix{
X_n \ar[r] \ar[d] &
X_{n - 1} \ar[r] \ar[d] &
\ldots \ar[r] &
X_0 \ar[d] \\
X'_n \ar[r] &
X'_{n - 1} \ar[r] &
\ldots \ar[r] &
X'_0
}
}
\end{equation}
entre des complexes de même longueur dans $\mathcal{D}$,
on dispose d'une notion évidente de {\it morphisme de systèmes de Postnikov}.
\end{definition}

\noindent
Voici un exemple fondamental.

\begin{example}
\label{example-key-postnikov}
Soit $\mathcal{A}$ une catégorie abélienne. Soit
$\ldots \to A_2 \to A_1 \to A_0$ un complexe de chaînes de $\mathcal{A}$.
Nous pouvons alors considérer les objets
$$
X_n = A_n
\quad\text{et}\quad
Y_n = (A_n \to A_{n - 1} \to \ldots \to A_0)[-n]
$$
de $D(\mathcal{A})$. Munis des morphismes canoniques évidents
$Y_n \to X_n$ et
$Y_0 \to Y_1[1] \to Y_2[2] \to \ldots$, les triangles distingués
$Y_n \to X_n \to Y_{n - 1} \to Y_n[1]$ définissent un système de Postnikov,
au sens de la Définition \ref{definition-postnikov-system}, pour
$\ldots \to X_2 \to X_1 \to X_0$. Nous utilisons ici l'extension évidente
de la notion de système de Postnikov à un complexe infini de
$D(\mathcal{A})$. Enfin, si les limites inductives indexées par
$\mathbf{N}$ existent et sont exactes dans $\mathcal{A}$, alors
$$
\text{hocolim} Y_n[n] = (\ldots \to A_2 \to A_1 \to A_0 \to 0 \to \ldots)
$$
dans $D(\mathcal{A})$. Cela résulte immédiatement du
Lemme \ref{lemma-colim-hocolim}.
\end{example}

\noindent
Étant donnés un complexe $X_n \to X_{n - 1} \to \ldots \to X_0$ et un
système de Postnikov comme dans la Définition
\ref{definition-postnikov-system}, nous pouvons considérer les morphismes
$$
Y_0 \to Y_1[1] \to \ldots \to Y_n[n]
$$
Ces morphismes s'insèrent dans certains triangles distingués et sont
compatibles avec les morphismes donnés entre les $X_i$. Voici le diagramme
correspondant lorsque $n = 3$ :
$$
\xymatrix{
Y_0 \ar[rr] & &
Y_1[1] \ar[dl] \ar[rr] & &
Y_2[2] \ar[dl] \ar[rr] & &
Y_3[3] \ar[dl] \\
& X_1[1] \ar[lu]_{+1} & &
X_2[2] \ar[ll]_{+1} \ar[lu]_{+1} & &
X_3[3] \ar[ll]_{+1} \ar[lu]_{+1}
}
$$
Nous invitons le lecteur à voir $Y_n[n]$ comme construit à partir de
$X_0, X_1[1], \ldots, X_n[n]$ ; par exemple, si les morphismes
$X_i \to X_{i - 1}$ sont nuls, nous pouvons prendre
$Y_n[n] = \bigoplus_{i = 0, \ldots, n} X_i[i]$.
Les systèmes de Postnikov n'existent pas toujours.
Voici un lemme élémentaire pour les petites valeurs de $n$.

\begin{lemma}
\label{lemma-postnikov-system-small-cases}
Soit $\mathcal{D}$ une catégorie triangulée. Considérons les systèmes de
Postnikov associés à des complexes de longueur $n$.
\begin{enumerate}
\item Pour $n = 0$, les systèmes de Postnikov existent toujours et tout
morphisme de complexes (\ref{equation-map-complexes}) se prolonge de manière
unique en un morphisme de systèmes de Postnikov.
\item Pour $n = 1$, les systèmes de Postnikov existent toujours et tout
morphisme de complexes (\ref{equation-map-complexes}) se prolonge en un
morphisme, non nécessairement unique, de systèmes de Postnikov.
\item Pour $n = 2$, les systèmes de Postnikov existent toujours, mais les
morphismes de complexes (\ref{equation-map-complexes}) ne se prolongent pas,
en général, en morphismes de systèmes de Postnikov.
\item Pour $n > 2$, les systèmes de Postnikov n'existent pas toujours.
\end{enumerate}
\end{lemma}

\begin{proof}
Le cas $n = 0$ est immédiat puisque les isomorphismes sont inversibles.
Le cas $n = 1$ résulte immédiatement de TR1 (existence des triangles) et de
TR3 (prolongement des morphismes aux triangles).
Pour le cas $n = 2$, raisonnons comme suit.
Posons $Y_0 = X_0$. D'après le cas $n = 1$, nous pouvons choisir un système
de Postnikov
$$
Y_1 \to X_1 \to Y_0 \to Y_1[1]
$$
Comme le composé $X_2 \to X_1 \to X_0$ est nul, le morphisme
$X_2 \to X_1$ se factorise, de manière non unique, sous la forme
$X_2 \to Y_1 \to X_1$, d'après le
Lemme \ref{lemma-representable-homological}.
Il suffit alors d'insérer le morphisme $X_2 \to Y_1$ dans un triangle
distingué
$$
Y_2 \to X_2 \to Y_1 \to Y_2[1]
$$
pour obtenir le système de Postnikov lorsque $n = 2$.
Pour $n > 2$, nous ne pouvons pas raisonner de la même manière, car nous ne
savons pas si le composé $X_n \to X_{n - 1} \to Y_{n - 1}$ est nul dans
$\mathcal{D}$.
\end{proof}

\begin{lemma}
\label{lemma-maps-postnikov-systems-vanishing}
Soit $\mathcal{D}$ une catégorie triangulée. Étant donné un morphisme
(\ref{equation-map-complexes}), considérons la condition
\begin{equation}
\label{equation-P}
\Hom(X_i[i - j - 1], X'_j) = 0 \text{ pour }i > j + 1
\end{equation}
Alors :
\begin{enumerate}
\item si nous disposons d'un système de Postnikov pour
$X'_n \to X'_{n - 1} \to \ldots \to X'_0$, la propriété
(\ref{equation-P}) implique que
$$
\Hom(X_i[i - j - 1], Y'_j) = 0 \text{ pour }i > j + 1
$$
\item si nous disposons de systèmes de Postnikov pour les deux complexes et
si (\ref{equation-P}) est satisfaite, alors le morphisme se prolonge en un
morphisme, non nécessairement unique, de systèmes de Postnikov.
\end{enumerate}
\end{lemma}

\begin{proof}
Démontrons d'abord (1) par récurrence sur $j$. Pour le cas initial $j = 0$,
il n'y a rien à démontrer puisque $Y'_0 \to X'_0$ est un isomorphisme.
Supposons le résultat démontré pour $j - 1$. Considérons le triangle distingué
$$
Y'_j \to X'_j \to Y'_{j - 1} \to Y'_j[1]
$$
La suite exacte longue du Lemme \ref{lemma-representable-homological}
fournit une suite exacte
$$
\Hom(X_i[i - j - 1], Y'_{j - 1}[-1]) \to
\Hom(X_i[i - j - 1], Y'_j) \to
\Hom(X_i[i - j - 1], X'_j)
$$
L'hypothèse de récurrence et (\ref{equation-P}) montrent que les groupes
extrêmes sont nuls, d'où le résultat.

\medskip\noindent
Démonstration de (2). Pour $n = 1$, l'existence des morphismes a été établie
dans le Lemme \ref{lemma-postnikov-system-small-cases}.
Pour $n > 1$, nous pouvons supposer, par récurrence, que le morphisme de
systèmes de Postnikov de longueur $n - 1$ est donné. Le problème est que nous
ne savons pas si le diagramme
$$
\xymatrix{
X_n \ar[r] \ar[d] & Y_{n - 1} \ar[d] \\
X'_n \ar[r] & Y'_{n - 1}
}
$$
est commutatif. Notons $\alpha : X_n \to Y'_{n - 1}$ la différence.
Nous savons toutefois que le composé de $\alpha$ avec
$Y'_{n - 1} \to X'_{n - 1}$ est nul (par définition d'un morphisme de
systèmes de Postnikov de longueur $n - 1$).
Le triangle distingué
$Y'_{n - 1} \to X'_{n - 1} \to Y'_{n - 2} \to Y'_{n - 1}[1]$
montre alors que $\alpha$ est le composé de
$Y'_{n - 2}[-1] \to Y'_{n - 1}$ avec un morphisme
$\alpha' : X_n \to Y'_{n - 2}[-1]$. La propriété
(\ref{equation-P}) assure, par la partie (1) du lemme, que $\alpha'$ est nul.
Ainsi, $\alpha$ est nul.
Pour achever la démonstration de l'existence, la commutativité permet de
choisir la flèche pointillée qui complète le diagramme
$$
\xymatrix{
Y_{n - 1}[-1] \ar[d] \ar[r] &
Y_n \ar[r] \ar@{..>}[d] &
X_n \ar[r] \ar[d] &
Y_{n - 1} \ar[d] \\
Y'_{n - 1}[-1] \ar[r] &
Y'_n \ar[r] &
X'_n \ar[r] &
Y'_{n - 1}
}
$$
en vertu de TR3.
\end{proof}

\begin{lemma}
\label{lemma-uniqueness-maps-postnikov-systems}
Soit $\mathcal{D}$ une catégorie triangulée. Étant donné un morphisme
(\ref{equation-map-complexes}), supposons choisis des systèmes de Postnikov
pour les deux complexes. Si l'une des conditions suivantes est satisfaite :
\begin{enumerate}
\item $\Hom(X_i[i], Y'_n[n]) = 0$ pour $i = 1, \ldots, n$, ou
\item $\Hom(Y_n[n], X'_{n - i}[n - i]) = 0$ pour $i = 1, \ldots, n$, ou
\item $\Hom(X_{j - i}[-i + 1], X'_j) = 0$ et
$\Hom(X_j, X'_{j - i}[-i]) = 0$ pour $j \geq i > 0$,
\end{enumerate}
alors il existe au plus un morphisme entre ces systèmes de Postnikov.
\end{lemma}

\begin{proof}
Démonstration de (1). Considérons le diagramme suivant :
$$
\xymatrix{
Y_0 \ar[r] \ar[d] &
Y_1[1] \ar[r] \ar[ld] &
Y_2[2] \ar[r] \ar[lld] &
\ldots \ar[r] &
Y_n[n] \ar[lllld] \\
Y'_n[n]
}
$$
Les flèches sont les composés du morphisme $Y_n[n] \to Y'_n[n]$ avec les
morphismes $Y_i[i] \to Y_n[n]$. La flèche $Y_0 \to Y'_n[n]$ est déterminée,
car elle est le composé
$Y_0 = X_0 \to X'_0 = Y'_0 \to Y'_n[n]$.
Comme nous avons le triangle distingué $Y_0 \to Y_1[1] \to X_1[1]$,
l'égalité $\Hom(X_1[1], Y'_n[n]) = 0$ assure l'unicité de la deuxième
flèche verticale. De même, le triangle distingué
$Y_1[1] \to Y_2[2] \to X_2[2]$ montre que l'égalité
$\Hom(X_2[2], Y'_n[n]) = 0$ assure l'unicité de la troisième flèche verticale,
et ainsi de suite.

\medskip\noindent
Démonstration de (2). Le composé $Y_n[n] \to Y'_n[n] \to X'_n[n]$ est
égal au composé $Y_n[n] \to X_n[n] \to X'_n[n]$, et il est donc unique.
Le triangle distingué
$Y'_{n - 1}[n - 1] \to Y'_n[n] \to X'_n[n]$ montre alors qu'il suffit de
prouver que $\Hom(Y_n[n], Y'_{n - 1}[n - 1]) = 0$. À l'aide des triangles
distingués
$$
Y'_{n - i - 1}[n - i - 1] \to Y'_{n - i}[n - i] \to X'_{n - i}[n - i]
$$
nous déduisons cette annulation de notre hypothèse. Nous omettons les détails
élémentaires.

\medskip\noindent
Démonstration de (3). En reprenant la démonstration du
Lemme \ref{lemma-maps-postnikov-systems-vanishing} et en raisonnant par
récurrence sur $n$, il suffit de montrer l'unicité de la flèche pointillée
dans le morphisme de triangles
$$
\xymatrix{
Y_{n - 1}[-1] \ar[d] \ar[r] &
Y_n \ar[r] \ar@{..>}[d] &
X_n \ar[r] \ar[d] &
Y_{n - 1} \ar[d] \\
Y'_{n - 1}[-1] \ar[r] &
Y'_n \ar[r] &
X'_n \ar[r] &
Y'_{n - 1}
}
$$
D'après la partie (5) du Lemme \ref{lemma-uniqueness-third-arrow}, il suffit
de montrer que $\Hom(Y_{n - 1}, X'_n) = 0$ et
$\Hom(X_n, Y'_{n - 1}[-1]) = 0$.
Pour la première annulation, nous utilisons les triangles distingués
$Y_{n - i - 1}[-i] \to Y_{n - i}[-(i - 1)] \to X_{n - i}[-(i - 1)]$
pour $i > 0$ et raisonnons par récurrence sur $i$ ; l'annulation supposée de
$\Hom(X_{n - i}[-i + 1], X'_n)$ suffit.
Pour la seconde, nous utilisons de même les triangles distingués
$Y'_{n - i - 1}[-i - 1] \to Y'_{n - i}[-i] \to X'_{n - i}[-i]$
pour voir que l'annulation supposée de
$\Hom(X_n, X'_{n - i}[-i])$ suffit également.
\end{proof}

\begin{lemma}
\label{lemma-existence-postnikov-system}
Soit $\mathcal{D}$ une catégorie triangulée.
Soit $X_n \to X_{n - 1} \to \ldots \to X_0$ un complexe de
$\mathcal{D}$. Si
$$
\Hom(X_i[i - j - 2], X_j) = 0 \text{ pour }i > j + 2
$$
alors il existe un système de Postnikov. Si
$$
\Hom(X_i[i - j - 1], X_j) = 0 \text{ pour }i > j + 1
$$
alors deux systèmes de Postnikov quelconques sont isomorphes.
\end{lemma}

\begin{proof}
Nous raisonnons par récurrence sur $n$. Les cas $n = 0, 1, 2$ résultent du
Lemme \ref{lemma-postnikov-system-small-cases}.
Supposons $n > 2$.
Supposons donné un système de Postnikov pour le complexe
$X_{n - 1} \to X_{n - 2} \to \ldots \to X_0$.
Le seul obstacle à son prolongement en un système de Postnikov de longueur
$n$ est la nécessité de trouver un morphisme $X_n \to Y_{n - 1}$ tel que le
composé $X_n \to Y_{n - 1} \to X_{n - 1}$ soit égal au morphisme donné
$X_n \to X_{n - 1}$. Considérons le triangle distingué
$$
Y_{n - 1} \to X_{n - 1} \to Y_{n - 2} \to Y_{n - 1}[1]
$$
et la suite exacte longue qui lui est associée par le foncteur
$\Hom(X_n, -)$ (voir le Lemme \ref{lemma-representable-homological}).
Il suffit alors de montrer que le composé
$X_n \to X_{n - 1} \to Y_{n - 2}$ est nul.
Comme nous savons que $X_n \to X_{n - 1} \to X_{n - 2}$ est nul, nous
pouvons appliquer le triangle distingué
$$
Y_{n - 2} \to X_{n - 2} \to Y_{n - 3} \to Y_{n - 2}[1]
$$
pour voir qu'il suffit que $\Hom(X_n, Y_{n - 3}[-1]) = 0$.
En raisonnant exactement comme dans la partie (1) de la démonstration du
Lemme \ref{lemma-maps-postnikov-systems-vanishing}, le lecteur vérifiera
aisément que cette annulation résulte de la condition énoncée dans le lemme.

\medskip\noindent
L'assertion sur les isomorphismes résulte de l'existence d'un morphisme entre
les systèmes de Postnikov prolongeant l'identité du complexe, établie dans la
partie (2) du Lemme \ref{lemma-maps-postnikov-systems-vanishing}, et du
Lemme \ref{lemma-third-isomorphism-triangle}, qui montre que tous les
morphismes ainsi obtenus sont des isomorphismes.
\end{proof}

\section{Systèmes essentiellement constants}
\label{section-essentially-constant}

\noindent
Quelques lemmes préliminaires sur les systèmes essentiellement constants
dans les catégories triangulées.

\begin{lemma}
\label{lemma-essentially-constant}
Soit $\mathcal{D}$ une catégorie triangulée. Soit $(A_i)$ un système
projectif de $\mathcal{D}$. Alors $(A_i)$ est essentiellement constant (voir
Catégories, Définition
\ref{categories-definition-essentially-constant-diagram}) si et seulement
s'il existe un $i$ tel que, pour tout $j \geq i$, on ait une décomposition en
somme directe $A_j = A \oplus Z_j$ vérifiant les propriétés suivantes :
(a) les morphismes $A_{j'} \to A_j$ sont compatibles avec les décompositions
en somme directe et induisent l'identité sur $A$ ; (b) pour tout $j \geq i$,
il existe un $j' \geq j$ tel que $Z_{j'} \to Z_j$ soit nul.
\end{lemma}

\begin{proof}
Supposons $(A_i)$ essentiellement constant de valeur $A$. Alors
$A = \lim A_i$, et il existe un $i$ et un morphisme $A_i \to A$ tels que
(1) le composé $A \to A_i \to A$ soit l'identité de $A$ et (2), pour tout
$j \geq i$, il existe un $j' \geq j$ tel que $A_{j'} \to A_j$ se factorise
sous la forme $A_{j'} \to A_i \to A \to A_j$. D'après (1), pour
$j \geq i$, les morphismes $A \to A_j$ et $A_j \to A_i \to A$ ont pour
composé l'identité de $A$. Il s'ensuit que $A_j \to A$ possède un noyau
$Z_j$ et que le morphisme $A \oplus Z_j \to A_j$ est un isomorphisme ; voir
les Lemmes \ref{lemma-when-split} et \ref{lemma-split}.
Ces décompositions en somme directe satisfont manifestement (a).
D'après (2), pour tout $j$, il existe un $j' \geq j$ tel que
$Z_{j'} \to Z_j$ soit nul ; ainsi, (b) est satisfaite. Nous omettons la
démonstration de la réciproque.
\end{proof}

\begin{lemma}
\label{lemma-essentially-constant-2-out-of-3}
Soit $\mathcal{D}$ une catégorie triangulée. Soit
$$
A_n \to B_n \to C_n \to A_n[1]
$$
un système projectif de triangles distingués de $\mathcal{D}$.
Si $(A_n)$ et $(C_n)$ sont essentiellement constants, alors $(B_n)$ est
essentiellement constant et leurs valeurs s'insèrent dans un triangle
distingué $A \to B \to C \to A[1]$ tel que, pour un certain $n \geq 1$,
il existe un morphisme
$$
\xymatrix{
A_n \ar[d] \ar[r] &
B_n \ar[d] \ar[r] &
C_n \ar[d] \ar[r] &
A_n[1] \ar[d] \\
A \ar[r] &
B \ar[r] &
C \ar[r] &
A[1]
}
$$
de triangles distingués qui induit un isomorphisme
$\lim_{n' \geq n} A_{n'} \to A$, et de même pour $B$ et $C$.
\end{lemma}

\begin{proof}
Après renumérotation, nous pouvons supposer que
$A_n = A \oplus A'_n$ et $C_n = C \oplus C'_n$, où les systèmes projectifs
$(A'_n)$ et $(C'_n)$ sont essentiellement nuls ; voir le
Lemme \ref{lemma-essentially-constant}.
En particulier, le morphisme
$$
C \oplus C'_n \to (A \oplus A'_n)[1]
$$
envoie le facteur direct $C$ dans le facteur direct $A[1]$, pour tout $n$,
par un morphisme $\delta : C \to A[1]$ indépendant de $n$. Choisissons un
triangle distingué
$$
A \to B \to C \xrightarrow{\delta} A[1]
$$
Choisissons ensuite un morphisme de triangles distingués
$$
(A_1 \to B_1 \to C_1 \to A_1[1]) \to
(A \to B \to C \to A[1])
$$
ce qui est possible d'après TR3. Pour tout objet $D$ de $\mathcal{D}$, il
induit un diagramme commutatif
$$
\xymatrix{
\ldots \ar[r] &
\Hom_\mathcal{D}(C, D) \ar[r] \ar[d] &
\Hom_\mathcal{D}(B, D) \ar[r] \ar[d] &
\Hom_\mathcal{D}(A, D) \ar[r] \ar[d] &
\ldots \\
\ldots \ar[r] &
\colim \Hom_\mathcal{D}(C_n, D) \ar[r] &
\colim \Hom_\mathcal{D}(B_n, D) \ar[r] &
\colim \Hom_\mathcal{D}(A_n, D) \ar[r] &
\ldots
}
$$
Les flèches verticales de gauche et de droite sont des isomorphismes, de même
que celles placées immédiatement à leur gauche et à leur droite. Le lemme des
cinq montre donc que la flèche médiane est un isomorphisme. Il s'ensuit que
$(B_n)$ est isomorphe au système projectif constant de valeur $B$, d'après la
discussion de Catégories, Remarque
\ref{categories-remark-pro-category-copresheaves}.
Comme ceci équivaut à dire que $(B_n)$ est essentiellement constant de valeur
$B$, d'après Catégories, Remarque
\ref{categories-remark-pro-category}, la démonstration est achevée.
\end{proof}

\begin{lemma}
\label{lemma-essentially-constant-cohomology}
Soit $\mathcal{A}$ une catégorie abélienne. Soit $A_n$ un système projectif
d'objets de $D(\mathcal{A})$. Supposons que :
\begin{enumerate}
\item il existe des entiers $a \leq b$ tels que $H^i(A_n) = 0$
pour $i \not \in [a, b]$ ;
\item les systèmes projectifs $H^i(A_n)$ de $\mathcal{A}$ sont
essentiellement constants pour tout $i \in \mathbf{Z}$.
\end{enumerate}
Alors $A_n$ est un système essentiellement constant de $D(\mathcal{A})$ dont
la valeur $A$ est telle que $H^i(A)$ soit la valeur du système constant
$H^i(A_n)$ pour tout $i \in \mathbf{Z}$.
\end{lemma}

\begin{proof}
D'après la Remarque \ref{remark-truncation-distinguished-triangle}, nous
obtenons un système projectif de triangles distingués
$$
\tau_{\leq a}A_n \to A_n \to \tau_{\geq a + 1}A_n \to (\tau_{\leq a}A_n)[1]
$$
Bien entendu, $\tau_{\leq a}A_n = H^a(A_n)[-a]$ dans $D(\mathcal{A})$.
Par hypothèse, ces objets forment donc un système essentiellement constant.
Une récurrence sur $b - a$ montre que le système projectif
$\tau_{\geq a + 1}A_n$ est essentiellement constant, disons de valeur $A'$.
Le Lemme \ref{lemma-essentially-constant-2-out-of-3} montre alors que $A_n$
est un système essentiellement constant. Nous omettons la démonstration de
l'assertion sur les cohomologies (indication : utiliser la dernière partie du
Lemme \ref{lemma-essentially-constant-2-out-of-3}).
\end{proof}

\begin{lemma}
\label{lemma-pro-isomorphism}
Soit $\mathcal{D}$ une catégorie triangulée. Soit
$$
A_n \to B_n \to C_n \to A_n[1]
$$
un système projectif de triangles distingués. Si le système $C_n$ est
pro-nul (essentiellement constant de valeur $0$), alors les morphismes
$A_n \to B_n$ déterminent un pro-isomorphisme entre le pro-objet $(A_n)$ et
le pro-objet $(B_n)$.
\end{lemma}

\begin{proof}
Pour tout objet $X$ de $\mathcal{D}$, considérons la suite exacte
$$
\colim \Hom(C_n, X) \to
\colim \Hom(B_n, X) \to
\colim \Hom(A_n, X) \to
\colim \Hom(C_n[-1], X) \to
$$
L'exactitude résulte du Lemme \ref{lemma-representable-homological}, combiné
avec Algèbre, Lemme \ref{algebra-lemma-directed-colimit-exact}.
Par hypothèse, le premier et le dernier terme sont nuls. Par conséquent, le
morphisme
$\colim \Hom(B_n, X) \to \colim \Hom(A_n, X)$ est un isomorphisme pour tout
$X$. Le lemme résulte de ce fait et de Catégories, Remarque
\ref{categories-remark-pro-category-copresheaves}.
\end{proof}

\begin{lemma}
\label{lemma-pro-isomorphism-bis}
Soit $\mathcal{A}$ une catégorie abélienne. Soit
$$
A_n \to B_n
$$
un système projectif de morphismes de $D(\mathcal{A})$. Supposons que :
\begin{enumerate}
\item il existe des entiers $a \leq b$ tels que $H^i(A_n) = 0$
et $H^i(B_n) = 0$ pour $i \not \in [a, b]$ ;
\item le système projectif de morphismes
$H^i(A_n) \to H^i(B_n)$ de $\mathcal{A}$ définit un isomorphisme de
pro-objets de $\mathcal{A}$ pour tout $i \in \mathbf{Z}$.
\end{enumerate}
Alors les morphismes $A_n \to B_n$ déterminent un pro-isomorphisme entre le
pro-objet $(A_n)$ et le pro-objet $(B_n)$.
\end{lemma}

\begin{proof}
Nous pouvons, par récurrence, prolonger les morphismes $A_n \to B_n$ en un
système projectif de triangles distingués
$A_n \to B_n \to C_n \to A_n[1]$ à l'aide de l'axiome TR3.
D'après le Lemme \ref{lemma-pro-isomorphism}, il suffit de montrer que $C_n$
est pro-nul. D'après le Lemme
\ref{lemma-essentially-constant-cohomology}, il suffit de montrer que
$H^p(C_n)$ est pro-nul pour tout $p$.
Cela résulte de l'hypothèse (2) et des suites exactes longues
$$
H^p(A_n) \xrightarrow{\alpha_n} H^p(B_n)
\xrightarrow{\beta_n}
H^p(C_n) \xrightarrow{\delta_n} H^{p + 1}(A_n)
\xrightarrow{\epsilon_n}
H^{p + 1}(B_n)
$$
En effet, pour tout $n$, nous pouvons trouver un $m > n$ tel que
$\Im(\beta_m)$ s'envoie sur zéro dans $H^p(C_n)$ : il suffit de choisir $m$
de telle sorte que $H^p(B_m) \to H^p(B_n)$ se factorise par
$\alpha_n : H^p(A_n) \to H^p(B_n)$. Pour une raison analogue, nous pouvons
ensuite choisir $k > m$ tel que $\Im(\delta_k)$ s'envoie sur zéro dans
$H^{p + 1}(A_m)$. Alors $H^p(C_k) \to H^p(C_n)$ est nul, car
$H^p(C_k) \to H^p(C_m)$ prend ses valeurs dans $\Ker(\delta_m)$ et
$H^p(C_m) \to H^p(C_n)$ annule
$\Ker(\delta_m) = \Im(\beta_m)$.
\end{proof}








\begin{multicols}{2}[\section{Autres chapitres}]
\noindent
Préliminaires
\begin{enumerate}
\item \hyperref[introduction-section-phantom]{Introduction}
\item \hyperref[conventions-section-phantom]{Conventions}
\item \hyperref[sets-section-phantom]{Théorie des ensembles}
\item \hyperref[categories-section-phantom]{Catégories}
\item \hyperref[topology-section-phantom]{Topologie}
\item \hyperref[sheaves-section-phantom]{Faisceaux sur les espaces}
\item \hyperref[sites-section-phantom]{Sites et faisceaux}
\item \hyperref[stacks-section-phantom]{Champs}
\item \hyperref[fields-section-phantom]{Corps}
\item \hyperref[algebra-section-phantom]{Algèbre commutative}
\item \hyperref[brauer-section-phantom]{Groupes de Brauer}
\item \hyperref[homology-section-phantom]{Algèbre homologique}
\item \hyperref[derived-section-phantom]{Catégories dérivées}
\item \hyperref[simplicial-section-phantom]{Méthodes simpliciales}
\item \hyperref[more-algebra-section-phantom]{Compléments d'algèbre}
\item \hyperref[smoothing-section-phantom]{Lissage des morphismes d'anneaux}
\item \hyperref[modules-section-phantom]{Faisceaux de modules}
\item \hyperref[sites-modules-section-phantom]{Modules sur les sites}
\item \hyperref[injectives-section-phantom]{Objets injectifs}
\item \hyperref[cohomology-section-phantom]{Cohomologie des faisceaux}
\item \hyperref[sites-cohomology-section-phantom]{Cohomologie sur les sites}
\item \hyperref[dga-section-phantom]{Algèbre différentielle graduée}
\item \hyperref[dpa-section-phantom]{Algèbre à puissances divisées}
\item \hyperref[sdga-section-phantom]{Faisceaux différentiels gradués}
\item \hyperref[hypercovering-section-phantom]{Hyperrecouvrements}
\end{enumerate}
Schémas
\begin{enumerate}
\setcounter{enumi}{25}
\item \hyperref[schemes-section-phantom]{Schémas}
\item \hyperref[constructions-section-phantom]{Constructions de schémas}
\item \hyperref[properties-section-phantom]{Propriétés des schémas}
\item \hyperref[morphisms-section-phantom]{Morphismes de schémas}
\item \hyperref[coherent-section-phantom]{Cohomologie des schémas}
\item \hyperref[divisors-section-phantom]{Diviseurs}
\item \hyperref[limits-section-phantom]{Limites de schémas}
\item \hyperref[varieties-section-phantom]{Variétés}
\item \hyperref[topologies-section-phantom]{Topologies sur les schémas}
\item \hyperref[descent-section-phantom]{Descente}
\item \hyperref[perfect-section-phantom]{Catégories dérivées de schémas}
\item \hyperref[more-morphisms-section-phantom]{Compléments sur les morphismes}
\item \hyperref[flat-section-phantom]{Compléments sur la platitude}
\item \hyperref[groupoids-section-phantom]{Schémas en groupoïdes}
\item \hyperref[more-groupoids-section-phantom]{Compléments sur les schémas en groupoïdes}
\item \hyperref[etale-section-phantom]{Morphismes étales de schémas}
\end{enumerate}
Thèmes de théorie des schémas
\begin{enumerate}
\setcounter{enumi}{41}
\item \hyperref[chow-section-phantom]{Homologie de Chow}
\item \hyperref[intersection-section-phantom]{Théorie de l'intersection}
\item \hyperref[pic-section-phantom]{Schémas de Picard des courbes}
\item \hyperref[weil-section-phantom]{Théories cohomologiques de Weil}
\item \hyperref[adequate-section-phantom]{Modules adéquats}
\item \hyperref[dualizing-section-phantom]{Complexes dualisants}
\item \hyperref[duality-section-phantom]{Dualité pour les schémas}
\item \hyperref[discriminant-section-phantom]{Discriminants et différentes}
\item \hyperref[derham-section-phantom]{Cohomologie de de Rham}
\item \hyperref[local-cohomology-section-phantom]{Cohomologie locale}
\item \hyperref[algebraization-section-phantom]{Géométrie algébrique et formelle}
\item \hyperref[curves-section-phantom]{Courbes algébriques}
\item \hyperref[resolve-section-phantom]{Résolution des surfaces}
\item \hyperref[models-section-phantom]{Réduction semi-stable}
\item \hyperref[functors-section-phantom]{Foncteurs et morphismes}
\item \hyperref[equiv-section-phantom]{Catégories dérivées de variétés}
\item \hyperref[pione-section-phantom]{Groupes fondamentaux de schémas}
\item \hyperref[etale-cohomology-section-phantom]{Cohomologie étale}
\item \hyperref[crystalline-section-phantom]{Cohomologie cristalline}
\item \hyperref[proetale-section-phantom]{Cohomologie pro-étale}
\item \hyperref[relative-cycles-section-phantom]{Cycles relatifs}
\item \hyperref[more-etale-section-phantom]{Compléments de cohomologie étale}
\item \hyperref[trace-section-phantom]{La formule des traces}
\end{enumerate}
Espaces algébriques
\begin{enumerate}
\setcounter{enumi}{64}
\item \hyperref[spaces-section-phantom]{Espaces algébriques}
\item \hyperref[spaces-properties-section-phantom]{Propriétés des espaces algébriques}
\item \hyperref[spaces-morphisms-section-phantom]{Morphismes d'espaces algébriques}
\item \hyperref[decent-spaces-section-phantom]{Espaces algébriques décents}
\item \hyperref[spaces-cohomology-section-phantom]{Cohomologie des espaces algébriques}
\item \hyperref[spaces-limits-section-phantom]{Limites d'espaces algébriques}
\item \hyperref[spaces-divisors-section-phantom]{Diviseurs sur les espaces algébriques}
\item \hyperref[spaces-over-fields-section-phantom]{Espaces algébriques sur des corps}
\item \hyperref[spaces-topologies-section-phantom]{Topologies sur les espaces algébriques}
\item \hyperref[spaces-descent-section-phantom]{Descente et espaces algébriques}
\item \hyperref[spaces-perfect-section-phantom]{Catégories dérivées d'espaces}
\item \hyperref[spaces-more-morphisms-section-phantom]{Compléments sur les morphismes d'espaces}
\item \hyperref[spaces-flat-section-phantom]{Platitude sur les espaces algébriques}
\item \hyperref[spaces-groupoids-section-phantom]{Groupoïdes dans les espaces algébriques}
\item \hyperref[spaces-more-groupoids-section-phantom]{Compléments sur les groupoïdes dans les espaces}
\item \hyperref[bootstrap-section-phantom]{Amorçage}
\item \hyperref[spaces-pushouts-section-phantom]{Sommes amalgamées d'espaces algébriques}
\end{enumerate}
Thèmes de géométrie
\begin{enumerate}
\setcounter{enumi}{81}
\item \hyperref[spaces-chow-section-phantom]{Groupes de Chow des espaces}
\item \hyperref[groupoids-quotients-section-phantom]{Quotients de groupoïdes}
\item \hyperref[spaces-more-cohomology-section-phantom]{Compléments sur la cohomologie des espaces}
\item \hyperref[spaces-simplicial-section-phantom]{Espaces simpliciaux}
\item \hyperref[spaces-duality-section-phantom]{Dualité pour les espaces}
\item \hyperref[formal-spaces-section-phantom]{Espaces algébriques formels}
\item \hyperref[restricted-section-phantom]{Algébrisation des espaces formels}
\item \hyperref[spaces-resolve-section-phantom]{Résolution des surfaces revisitée}
\end{enumerate}
Théorie des déformations
\begin{enumerate}
\setcounter{enumi}{89}
\item \hyperref[formal-defos-section-phantom]{Théorie formelle des déformations}
\item \hyperref[defos-section-phantom]{Théorie des déformations}
\item \hyperref[cotangent-section-phantom]{Le complexe cotangent}
\item \hyperref[examples-defos-section-phantom]{Problèmes de déformation}
\end{enumerate}
Champs algébriques
\begin{enumerate}
\setcounter{enumi}{93}
\item \hyperref[algebraic-section-phantom]{Champs algébriques}
\item \hyperref[examples-stacks-section-phantom]{Exemples de champs}
\item \hyperref[stacks-sheaves-section-phantom]{Faisceaux sur les champs algébriques}
\item \hyperref[criteria-section-phantom]{Critères de représentabilité}
\item \hyperref[artin-section-phantom]{Axiomes d'Artin}
\item \hyperref[quot-section-phantom]{Espaces de Quot et de Hilbert}
\item \hyperref[stacks-properties-section-phantom]{Propriétés des champs algébriques}
\item \hyperref[stacks-morphisms-section-phantom]{Morphismes de champs algébriques}
\item \hyperref[stacks-limits-section-phantom]{Limites de champs algébriques}
\item \hyperref[stacks-cohomology-section-phantom]{Cohomologie des champs algébriques}
\item \hyperref[stacks-perfect-section-phantom]{Catégories dérivées de champs}
\item \hyperref[stacks-introduction-section-phantom]{Introduction aux champs algébriques}
\item \hyperref[stacks-more-morphisms-section-phantom]{Compléments sur les morphismes de champs}
\item \hyperref[stacks-geometry-section-phantom]{La géométrie des champs}
\end{enumerate}
Thèmes de théorie des espaces de modules
\begin{enumerate}
\setcounter{enumi}{107}
\item \hyperref[moduli-section-phantom]{Champs de modules}
\item \hyperref[moduli-curves-section-phantom]{Espaces de modules de courbes}
\end{enumerate}
Divers
\begin{enumerate}
\setcounter{enumi}{109}
\item \hyperref[examples-section-phantom]{Exemples}
\item \hyperref[exercises-section-phantom]{Exercices}
\item \hyperref[guide-section-phantom]{Guide bibliographique}
\item \hyperref[desirables-section-phantom]{Desiderata}
\item \hyperref[coding-section-phantom]{Style de programmation}
\item \hyperref[obsolete-section-phantom]{Obsolète}
\item \hyperref[fdl-section-phantom]{Licence de documentation libre GNU}
\item \hyperref[index-section-phantom]{Index généré automatiquement}
\end{enumerate}
\end{multicols}


\bibliography{my}
\bibliographystyle{amsalpha}

\end{document}
